%============================================================
%  Bd. 19 - Reelle Zahlen
%============================================================

\documentclass[main.tex]{subfiles}

\ifSubfilesClassLoaded{
  \BandSetupStandalone{B19}
}{
}

\title{Bd. 19 - Reelle Zahlen}
\author{Martin Kunze}
\date{}
\setcounter{file}{19}

\hypersetup{
  pdftitle={Bd. 19 - Reelle Zahlen},
  pdfauthor={Martin Kunze}
}

\begin{document}

\title{Bd. 19 - Reelle Zahlen}
\author{Martin Kunze}
\date{}
\hypersetup{pageanchor=false}
\maketitle
\hypersetup{pageanchor=true}
\setcounter{tocdepth}{3}
\tableofcontents
\setcounter{file}{19}
\renewcommand{\FormulaBandID}{19}

% Lokale Notation dieses Bandes. Die Definitionen stehen absichtlich innerhalb
% des Dokumentteils, damit sie sowohl im Gesamtband als auch im Standalone-Build
% zur Verfuegung stehen.
\providecommand{\RealCut}[1]{\operatorname{Cut}_{\mathbb Q}\!\left(#1\right)}
\providecommand{\RealHat}[1]{\widehat{#1}}
\providecommand{\RealEmb}{\iota_{\mathbb R}}
\providecommand{\RealLe}{\mathrel{\leq_{\mathbb R}}}
\providecommand{\RealLt}{\mathrel{<_{\mathbb R}}}
\providecommand{\RealZero}{0_{\mathbb R}}
\providecommand{\RealOne}{1_{\mathbb R}}
\providecommand{\RealAdd}{\mathbin{\oplus}}
\providecommand{\RealMul}{\mathbin{\odot}}
\providecommand{\RealPosMul}{\mathbin{\odot_{+}}}
\providecommand{\RealNeg}[1]{\mathop{\ominus}\!#1}
\providecommand{\RealPosInv}[1]{\operatorname{inv}_{+}\!\left(#1\right)}
\providecommand{\RealInv}[1]{\operatorname{inv}_{\mathbb R}\!\left(#1\right)}
\providecommand{\RealAbs}[1]{\left\lvert #1\right\rvert_{\mathbb R}}
\providecommand{\RealDist}[2]{d_{\mathbb R}\!\left(#1,#2\right)}
\providecommand{\CompleteOrderedField}[1]{\mathsf{COF}\!\left(#1\right)}
\providecommand{\IntegerPrimeEmbeddingLaws}[1]{%
  \mathsf{ZEmb}_{K}\!\left(#1\right)}
\providecommand{\RationalPrimeEmbeddingLaws}[1]{%
  \mathsf{QEmb}_{K}\!\left(#1\right)}

\chapter{Einleitung}

Band~18 konstruiert die rationalen Zahlen mit ihrer Addition,
Multiplikation und totalen Ordnung; die maßgeblichen Grundlagen sind in
\FormulaRefAuto{RationalOrderedFieldAxioms} zusammengefasst. In diesem Band
werden daraus die reellen Zahlen als Dedekindsche Schnitte gewonnen. Diese
Konstruktion verwendet nur Mengen, rationale Arithmetik und Ordnung;
insbesondere setzt sie weder Folgen noch einen Grenzwertbegriff voraus.

Ein reeller Wert wird durch die Menge aller rationalen Zahlen beschrieben,
die echt unter ihm liegen. Die Inklusion solcher Mengen liefert die Ordnung.
Das Supremum einer nichtleeren, nach oben beschränkten Familie ist ihre
Vereinigung. Damit ist die Vollständigkeit bereits auf der Ebene der
Konstruktion sichtbar. Folgen und ihre Grenzwerte können deshalb erst in
einem späteren Band ohne begrifflichen Zirkel eingeführt werden.

\chapter{Dedekindsche Schnitte}

\section{Schnittprädikat und Trägermenge}

\FormulaDefDeltaK[Dedekindscher Schnitt der rationalen Zahlen]{%
  \RealCut{A}%
}{DedekindCutPredicate}{%
  \DeltaRow{Mengen}{A\dsep p\dsep q\dsep r}
  \DeltaRow{Relationsoperatoren}{<}
  \DeltaPrem{Geordneter rationaler Körper}{\mathbb Q}
  \DeltaRow{\textbf{Axiome}}{}
  \DeltaPrem{Rationale Teilmenge}{A\subseteq\mathbb Q}
  \DeltaPrem{Nichtleerheit}{A\neq\varnothing}
  \DeltaPrem{Echtheit}{A\neq\mathbb Q}
  \DeltaRow{Nach unten abgeschlossen}{%
    p\in A\dsep q\in\mathbb Q\dsep q<p\vdash q\in A}
  \DeltaRow{Kein größtes Element}{p\in A\vdash\exists r\in A\,(p<r)}
  \DeltaRow{\textbf{Neues Symbol}}{}
  \DeltaPrem{Schnittprädikat}{\RealCut{A}}
}

\FormulaAxiomDeltaK[Rationale Trägermenge]{%
  \RealCut{A}\vdash A\subseteq\mathbb Q%
}{DedekindCutSubsetAxiom}{%
  \DeltaRow{Mengen}{A\dsep p\dsep q\dsep r}
}

\FormulaAxiomDeltaK[Nichtleerheit]{%
  \RealCut{A}\vdash A\neq\varnothing%
}{DedekindCutNonemptyAxiom}{%
  \DeltaRow{Mengen}{A\dsep p\dsep q\dsep r}
}

\FormulaAxiomDeltaK[Echter Schnitt]{%
  \RealCut{A}\vdash A\neq\mathbb Q%
}{DedekindCutProperAxiom}{%
  \DeltaRow{Mengen}{A\dsep p\dsep q\dsep r}
}

\FormulaAxiomDeltaK[Abwärtsabgeschlossenheit]{%
  \RealCut{A}\vdash \forall p\in A\,\forall q\in\mathbb Q\,(q<p\rightarrow q\in A)%
}{DedekindCutDownwardAxiom}{%
  \DeltaRow{Mengen}{A\dsep p\dsep q\dsep r}
}

\FormulaAxiomDeltaK[Kein größtes Element]{%
  \RealCut{A}\vdash \forall p\in A\,\exists r\in A\,(p<r)%
}{DedekindCutNoMaximumAxiom}{%
  \DeltaRow{Mengen}{A\dsep p\dsep q\dsep r}
}

\FormulaThmDeltaK[Schnittkriterium]{%
  \begin{aligned}[t]
  \RealCut{A}\eqvdash{}&
    A\subseteq\mathbb Q\land A\neq\varnothing\land A\neq\mathbb Q
    \land{}\\[-2pt]
  &\forall p\in A\,\forall q\in\mathbb Q\,
      \bigl(q<p\rightarrow q\in A\bigr)
    \land{}\\[-2pt]
  &\forall p\in A\,\exists r\in A\,(p<r).
  \end{aligned}%
}{DedekindCutCharacterization}{%
  \DeltaRow{Mengen}{A\dsep p\dsep q\dsep r}
  \DeltaRow{Relationsoperatoren}{<}
}
\begin{tabproofsplitwide}
  \proofpartwide{Hinrichtung}
    \proofstepwidestar[1]{\RealCut{A}}{\rA}
    \proofstepwidestar[1]{A\subseteq\mathbb Q}{%
      \FormulaRefAuto{DedekindCutSubsetAxiom}[ax]{1}}
    \proofstepwidestar[1]{A\neq\varnothing}{%
      \FormulaRefAuto{DedekindCutNonemptyAxiom}[ax]{1}}
    \proofstepwidestar[1]{A\neq\mathbb Q}{%
      \FormulaRefAuto{DedekindCutProperAxiom}[ax]{1}}
    \proofstepwidestar[1]{\forall p\in A\,\forall q\in\mathbb Q\,(q<p\rightarrow q\in A)}{%
      \FormulaRefAuto{DedekindCutDownwardAxiom}[ax]{1}}
    \proofstepwidestar[1]{\forall p\in A\,\exists r\in A\,(p<r)}{%
      \FormulaRefAuto{DedekindCutNoMaximumAxiom}[ax]{1}}
    \proofstepwidestar[1]{%
      \begin{aligned}[t]
      &A\subseteq\mathbb Q\land A\neq\varnothing\land A\neq\mathbb Q\land{}\\[-2pt]
      &\forall p\in A\,\forall q\in\mathbb Q\,(q<p\rightarrow q\in A)\land{}\\[-2pt]
      &\forall p\in A\,\exists r\in A\,(p<r)
    \end{aligned}}{\rAIStar{2,3,4,5,6}}
  \closeproofpartwide

  \proofpartwide{Rückrichtung}
    \proofstepwidestar[1]{%
      \begin{aligned}[t]
      &A\subseteq\mathbb Q\land A\neq\varnothing\land A\neq\mathbb Q\land{}\\[-2pt]
      &\forall p\in A\,\forall q\in\mathbb Q\,(q<p\rightarrow q\in A)\land{}\\[-2pt]
      &\forall p\in A\,\exists r\in A\,(p<r)
    \end{aligned}}{\rA}
    \proofstepwidestar[1]{A\subseteq\mathbb Q}{\rAEa{1}}
    \proofstepwidestar[1]{A\neq\varnothing}{\rAEa{\rAEb{1}}}
    \proofstepwidestar[1]{A\neq\mathbb Q}{\rAEa{\rAEb{\rAEb{1}}}}
    \proofstepwidestar[1]{\forall p\in A\,\forall q\in\mathbb Q\,(q<p\rightarrow q\in A)}{\rAEa{\rAEb{\rAEb{\rAEb{1}}}}}
    \proofstepwidestar[1]{\forall p\in A\,\exists r\in A\,(p<r)}{\rAEb{\rAEb{\rAEb{\rAEb{1}}}}}
    \proofstepwidestar[1]{\RealCut{A}}{%
      \FormulaRefAuto{DedekindCutPredicate}[def]{2,3,4,5,6}}
  \closeproofpartwide
\end{tabproofsplitwide}

Neben der rationalen Typisierung sind die vier Schnittbedingungen wesentlich. Die
Abwärtsabgeschlossenheit macht einen Schnitt zu einem rationalen
Anfangsabschnitt; die letzte Bedingung entfernt die ansonsten mögliche
Mehrdeutigkeit am Rand eines rationalen Anfangsabschnitts.

\FormulaDefDeltaK[Die Menge der reellen Zahlen]{%
  \mathbb R\coloneqq
  \bigl\{A\in\powerset(\mathbb Q)\bigm|\RealCut{A}\bigr\}%
}{RealCutCarrier}{%
  \DeltaRow{Mengen}{A\dsep\mathbb R}%
  \DeltaRow{Prädikate}{\RealCut{\,\cdot\,}}%
  \DeltaRow{\textbf{Neues Symbol}}{}%
  \DeltaRow{Reelle Zahlen}{\mathbb R}%
}

\FormulaThmDeltaK[Elementcharakterisierung der reellen Zahlen]{%
  A\in\mathbb R\eqvdash\RealCut{A}%
}{RealCutMembership}{%
  \DeltaRow{Mengen}{A}%
}
\begin{tabproofsplitwide}
  \proofpartwideindR[Hinrichtung]{%
    A\in\mathbb R\rightarrow\RealCut A%
  }{A\in\mathbb R\rightarrow\RealCut A}[RealCutMembershipForward]
    \proofstepwidestar[1]{A\in\mathbb R}{\rA}
    \proofstepwidestar[]{%
      \mathbb R=\{X\in\powerset(\mathbb Q)\mid\RealCut X\}}{%
      \FormulaRefAuto{RealCutCarrier}[def]}
    \proofstepwidestar[1]{%
      A\in\powerset(\mathbb Q)\land\RealCut A}{%
      \rIE{1,2}}
    \proofstepwidestar[1]{\RealCut A}{\rAEb{3}}
  \proofstepwidestar[]{A\in\mathbb R\rightarrow\RealCut A}{\rRI{1,4}}
  \closeproofpartwideindR

  \proofpartwideindR[Rückrichtung]{%
    \RealCut A\rightarrow A\in\mathbb R%
  }{\RealCut A\rightarrow A\in\mathbb R}[RealCutMembershipBackward]
    \proofstepwidestar[1]{\RealCut A}{\rA}
    \proofstepwidestar[1]{A\subseteq\mathbb Q}{%
      \FormulaRefAuto{DedekindCutCharacterization}[thm]{1}}
    \proofstepwidestar[1]{A\in\powerset(\mathbb Q)}{%
      \FormulaRefAuto{x\in\powerset(A)\eqvdash x\subseteq A}{2}}
    \proofstepwidestar[1]{%
      A\in\powerset(\mathbb Q)\land\RealCut A}{\rAI{3,1}}
    \proofstepwidestar[1]{A\in\mathbb R}{%
      \FormulaRefAuto{RealCutCarrier}[def]{4}}
  \proofstepwidestar[]{\RealCut A\rightarrow A\in\mathbb R}{\rRI{1,5}}
  \closeproofpartwideindR

  \proofpartwide{Zusammenfassung}
    \proofstepwidestar[]{A\in\mathbb R\rightarrow\RealCut A}{%
      \FormulaRefAuto{RealCutMembershipForward}[thm]}
    \proofstepwidestar[]{\RealCut A\rightarrow A\in\mathbb R}{%
      \FormulaRefAuto{RealCutMembershipBackward}[thm]}
    \proofstepwidestar[]{A\in\mathbb R\leftrightarrow\RealCut A}{%
      \rLRI{1,2}}
  \closeproofpartwide
\end{tabproofsplitwide}

Diese Charakterisierung erlaubt es im Folgenden, eine reelle Zahl und den sie
darstellenden Schnitt ohne zusätzliche Klassenklammern miteinander zu
identifizieren.

\section{Hauptschnitte und rationale Einbettung}

\FormulaDefDeltaK[Hauptschnitt einer rationalen Zahl]{%
  q\in\mathbb Q\vdash
  \RealHat q\coloneqq\{p\in\mathbb Q\mid p<q\}%
}{RealPrincipalCut}{%
  \DeltaRow{Mengen}{p\dsep q}%
  \DeltaRow{Relationsoperatoren}{<}%
  \DeltaRow{\textbf{Neues Symbol}}{}%
  \DeltaRow{Hauptschnitt}{\RealHat q}%
}

\FormulaThmDeltaK[Hauptschnitte sind Dedekindsche Schnitte]{%
  q\in\mathbb Q\vdash\RealCut{\RealHat q}%
}{RealPrincipalCutIsCut}{%
  \DeltaRow{Mengen}{p\dsep q\dsep r}%
}
\begin{tabproofsplitwide}
\proofpartwide{Nichtleerheit}
\proofstepwidestarbreakkeep[1]{q\in\mathbb Q}{\rA}
    \proofstepwidestarbreakkeep[]{\IntOne\in\PositiveIntegers}{%
      \FormulaRefAuto{IntegerOnePositive}[thm]}
    \proofstepwidestarbreakkeep[]{\IntZero<\IntOne}{%
      \FormulaRefAuto{PositiveIntegersDef}[def]{2}}
    \proofstepwidestarbreakkeep[]{%
      \IntZero\leq\IntOne\land\IntZero\neq\IntOne}{%
      \FormulaRefAuto{IntegerStrictOrderDef}[def]{3}}
    \proofstepwidestarbreakkeep[]{\IntZero\in\mathbb Z}{%
      \FormulaRefAuto{IntZeroCarrier}[thm]}
    \proofstepwidestarbreakkeep[]{\IntOne\in\mathbb Z}{%
      \FormulaRefAuto{IntOneCarrier}[thm]}
    \proofstepwidestarbreakkeep[]{%
      \RatEmbed(\IntZero)\leq\RatEmbed(\IntOne)}{%
      \FormulaRefAuto{IntegerRationalEmbeddingOrder}[thm]{4,5,6}}
    \proofstepwidestarbreakkeep[]{%
      \RatEmbed(\IntZero)\neq\RatEmbed(\IntOne)}{%
      \FormulaRefAuto{IntegerRationalEmbeddingInjective}[thm]{4,5,6}}
    \proofstepwidestar[]{\RatOne=\RatEmbed(\IntOne)}{%
      \FormulaRefAuto{RationalOneDef}[def]}
    \proofstepwidestar[]{\RatZero=\RatEmbed(\IntZero)}{%
      \FormulaRefAuto{RationalZeroDef}[def]}
    \proofstepwidestarbreakkeep[]{%
      \RatZero=\RatEmbed(\IntZero)\land
      \RatOne=\RatEmbed(\IntOne)}{%
      \rAI{10,
        9}}
    \proofstepwidestarbreakkeep[]{\RatZero<\RatOne}{%
      \FormulaRefAuto{RationalStrictOrderDef}[def]{7,8,11}}
    \proofstepwide[1,12]{q-\RatOne}{<}{q}{%
      \FormulaRefAuto{RationalOrderTranslation}[thm]{1,12}}
    \proofstepwidestarbreakkeep[1,12]{%
      \exists p\in\mathbb Q\,(q-\RatOne<p\land p<q)}{%
      \FormulaRefAuto{RationalOrderDense}[thm]{13}}
    \proofstepwidestarbreakkeep[1,12]{p\in\mathbb Q\land p<q}{\rA}
    \proofstepwidestarbreakkeep[1,12]{p\in\RealHat q}{%
      \FormulaRefAuto{RealPrincipalCut}[def]{15}}
    \proofstepwidestarbreakkeep[1,12]{\RealHat q\neq\varnothing}{%
      \FormulaRefAuto{x\in A\vdash A\neq\varnothing}{16}}

\closeproofpartwide

\proofpartwide{Echtheit}
  \setcounter{proofstepnr}{17}% Fortlaufende Schritte dieses Beweises.
\proofstepwidestar[1]{\TotOrd{\mathbb Q,\leq}}{%
      \FormulaRefAuto{RationalOrderTotal}[thm]{1}}
\proofstepwidestarbreakkeep[1]{q\notin\RealHat q}{%
      \FormulaRefAuto{RealPrincipalCut}[def]{%
        18}}
    \proofstepwidestarbreakkeep[1]{\RealHat q\neq\mathbb Q}{%
      \FormulaRefAuto{x\in B, x\not\in A\vdash A\neq B}{1,19}}

\closeproofpartwide

\proofpartwideindR[Nichtleerheit und Echtheit]{%
    q\in\mathbb Q\vdash
    \RealHat q\neq\varnothing\land\RealHat q\neq\mathbb Q
  }{%
    q\in\mathbb Q\vdash
    \RealHat q\neq\varnothing\land\RealHat q\neq\mathbb Q
  }[RealPrincipalCutNonemptyProper]
  \setcounter{proofstepnr}{20}% Fortlaufende Schritte dieses Beweises.
\proofstepwidestarbreakkeep[1]{%
      \RealHat q\neq\varnothing\land\RealHat q\neq\mathbb Q}{%
      \rAI{17,20}}
    \proofstepwidestarbreakkeep[1]{%
      \RealHat q\neq\varnothing\land\RealHat q\neq\mathbb Q}{%
      \rEE{14,15,21}}

\closeproofpartwideindR


  \proofpartwideindR[Abwärtsabschluss]{%
    \begin{aligned}[t]
      &q\in\mathbb Q\dsep p\in\RealHat q\dsep s\in\mathbb Q\dsep s<p\\[-2pt]
      &\quad\vdash s\in\RealHat q
    \end{aligned}
  }{%
    q\in\mathbb Q\dsep p\in\RealHat q\dsep s\in\mathbb Q\dsep s<p
    \vdash s\in\RealHat q
  }[RealPrincipalCutDownwardClosed]
    \proofstepwidestar[1]{q\in\mathbb Q}{\rA}
    \proofstepwidestar[2]{p\in\RealHat q}{\rA}
    \proofstepwidestar[3]{s\in\mathbb Q}{\rA}
    \proofstepwidestar[4]{s<p}{\rA}
    \proofstepwide[1,2]{p}{<}{q}{%
      \FormulaRefAuto{RealPrincipalCut}[def]{1,2}}
    \proofstepwide[1,2,4]{s}{<}{q}{%
      \FormulaRefAuto{RationalOrderTotal}[thm]{4,5}}
    \proofstepwidestar[1,2,3,4]{s\in\RealHat q}{%
      \FormulaRefAuto{RealPrincipalCut}[def]{3,6}}
  \closeproofpartwideindR

  \proofpartwideindR[Kein größtes Element]{%
    q\in\mathbb Q\dsep p\in\RealHat q
    \vdash\exists r\in\RealHat q\,(p<r)
  }{%
    q\in\mathbb Q\dsep p\in\RealHat q
    \vdash\exists r\in\RealHat q\,(p<r)
  }[RealPrincipalCutNoMaximum]
    \proofstepwidestar[1]{q\in\mathbb Q}{\rA}
    \proofstepwidestar[2]{p\in\RealHat q}{\rA}
    \proofstepwide[1,2]{p}{<}{q}{%
      \FormulaRefAuto{RealPrincipalCut}[def]{1,2}}
    \proofstepwidestar[1,2]{%
      \exists r\in\mathbb Q\,(p<r\land r<q)}{%
      \FormulaRefAuto{RationalOrderDense}[thm]{3}}
    \proofstepwidestar[1,2]{r\in\mathbb Q\land p<r\land r<q}{\rA}
    \proofstepwidestar[1,2]{r\in\RealHat q}{%
      \FormulaRefAuto{RealPrincipalCut}[def]{5}}
    \proofstepwidestar[1,2]{\exists r\in\RealHat q\,(p<r)}{%
      \rEI{5,6}}
    \proofstepwidestar[1,2]{\exists r\in\RealHat q\,(p<r)}{%
      \rEE{4,5,7}}
  \closeproofpartwideindR

  \proofpartwide{Schluss}
    \proofstepwidestar[1]{q\in\mathbb Q}{\rA}
    \proofstepwidestar[1]{\RealHat q\subseteq\mathbb Q}{%
      \FormulaRefAuto{RealPrincipalCut}[def]{1}}
    \proofstepwidestar[1]{%
      \RealHat q\neq\varnothing\land\RealHat q\neq\mathbb Q}{%
      \FormulaRefAuto{RealPrincipalCutNonemptyProper}[thm]{1}}
    \proofstepwidestar[1]{%
      \forall p\in\RealHat q\,\forall s\in\mathbb Q\,
      (s<p\rightarrow s\in\RealHat q)}{%
      \FormulaRefAuto{RealPrincipalCutDownwardClosed}[thm]{1}}
    \proofstepwidestar[1]{%
      \forall p\in\RealHat q\,\exists r\in\RealHat q\,(p<r)}{%
      \FormulaRefAuto{RealPrincipalCutNoMaximum}[thm]{1}}
    \proofstepwidestar[1]{\RealCut{\RealHat q}}{%
      \FormulaRefAuto{DedekindCutCharacterization}[thm]{2,3,4,5}}
  \closeproofpartwide
\end{tabproofsplitwide}

\FormulaDefDeltaK[Kanonische Einbettung der rationalen Zahlen]{%
  \begin{aligned}[t]
    &\RealEmb\colon\mathbb Q\to\mathbb R,\\[-2pt]
    &\forall q\in\mathbb Q\,
      \bigl(\RealEmb(q)\coloneqq\RealHat q\bigr)
  \end{aligned}%
}{RationalToRealEmbedding}{%
  \DeltaRow{Mengen}{p\dsep q}%
  \DeltaRow{Funktionen}{\RealEmb}%
  \DeltaRow{\textbf{Neues Symbol}}{}%
  \DeltaRow{Rationale Einbettung}{\RealEmb}%
}
\begin{tabproof}
  \proofstep{1}{q\in\mathbb Q}{\rA}
  \proofstep{1}{\RealCut{\RealHat q}}{%
    \FormulaRefAuto{RealPrincipalCutIsCut}[thm]{1}}
  \proofstep{1}{\RealHat q\in\mathbb R}{%
    \FormulaRefAuto{A\in\mathbb R\eqvdash\RealCut{A}}[thm]{2}}
  \proofstep{}{\forall q\in\mathbb Q\,\bigl(\RealHat q\in\mathbb R\bigr)}{%
    \rUI{\rRI{1,3}}}
\end{tabproof}

\FormulaThmDeltaK[Die rationale Einbettung ist eine Funktion]{%
  \RealEmb\colon\mathbb Q\to\mathbb R%
}{RationalToRealEmbeddingFunction}{%
  \DeltaRow{Funktionen}{\RealEmb}%
}
\begin{tabproof}
  \proofstep{}{\RealEmb\colon\mathbb Q\to\mathbb R}{%
    \FormulaRefAuto{RationalToRealEmbedding}[def]}
\end{tabproof}

\chapter{Ordnung und Vollständigkeit}

\section{Inklusionsordnung}

\FormulaDefDeltaK[Ordnung der reellen Zahlen]{%
  A,B\in\mathbb R\vdash
  A\RealLe B\coloneqq A\subseteq B%
}{RealCutOrder}{%
  \DeltaRow{Mengen}{A\dsep B}%
  \DeltaRow{Relationsoperatoren}{\RealLe}%
  \DeltaRow{\textbf{Neues Symbol}}{}%
  \DeltaRow{Reelle Ordnung}{\RealLe}%
}

\FormulaDefDeltaK[Strikte Ordnung der reellen Zahlen]{%
  A,B\in\mathbb R\vdash
  A\RealLt B\coloneqq A\subsetneq B%
}{RealCutStrictOrder}{%
  \DeltaRow{Mengen}{A\dsep B}%
  \DeltaRow{Relationsoperatoren}{\RealLe\dsep\RealLt}%
  \DeltaRow{\textbf{Neues Symbol}}{}%
  \DeltaRow{Strikte reelle Ordnung}{\RealLt}%
}

\FormulaThmDeltaK[Charakterisierung der reellen Vergleichsrelationen]{%
  A,B\in\mathbb R\vdash
  \bigl(A\RealLe B\leftrightarrow A\subseteq B\bigr)
  \land
  \bigl(A\RealLt B\leftrightarrow A\subsetneq B\bigr)%
}{RealCutOrderCharacterization}{%
  \DeltaRow{Mengen}{A\dsep B}%
  \DeltaRow{Relationsoperatoren}{\RealLe\dsep\RealLt}%
}
\begin{tabproofwide}
  \proofstepwidestar[1]{A,B\in\mathbb R}{\rA}
  \proofstepwide[1]{A\RealLe B}{\leftrightarrow}{A\subseteq B}{%
    \FormulaRefAuto{RealCutOrder}[def]{1}}
  \proofstepwide[1]{A\RealLt B}{\leftrightarrow}{A\subsetneq B}{%
    \FormulaRefAuto{RealCutStrictOrder}[def]{1}}
  \proofstepwidestar[1]{%
    (A\RealLe B\leftrightarrow A\subseteq B)\land
    (A\RealLt B\leftrightarrow A\subsetneq B)}{\rAI{2,3}}
\end{tabproofwide}

\FormulaThmDeltaK[Vergleichbarkeit zweier Dedekindscher Schnitte]{%
  A,B\in\mathbb R\vdash A\subseteq B\lor B\subseteq A%
}{RealCutInclusionComparable}{%
  \DeltaRow{Mengen}{A\dsep B\dsep a\dsep b}%
}
\begin{tabproofsplitwide}
  \proofpartwideindR[Ausschluss zweier Gegenzeugen]{%
    \begin{aligned}[t]
      &A,B\in\mathbb R\dsep a\in A\setminus B\dsep b\in B\setminus A\\[-2pt]
      &\quad\vdash\bot
    \end{aligned}
  }{%
    A,B\in\mathbb R\dsep a\in A\setminus B\dsep b\in B\setminus A
    \vdash\bot
  }[RealCutOppositeWitnessesImpossible]
    \proofstepwidestar[1]{A,B\in\mathbb R}{\rA}
    \proofstepwidestar[2]{a\in A\setminus B}{\rA}
    \proofstepwidestar[3]{b\in B\setminus A}{\rA}
    \proofstepwidestarbreakkeep[1]{\RealCut A\land\RealCut B}{%
      \FormulaRefAuto{RealCutMembership}[thm]{1}}
    \proofstepwidestar[2]{a\in A\land a\notin B}{%
      \FormulaRefAuto{x\in A\setminus B\eqvdash x\in A\land x\notin B}{2}}
    \proofstepwidestar[3]{b\in B\land b\notin A}{%
      \FormulaRefAuto{x\in A\setminus B\eqvdash x\in A\land x\notin B}{3}}
    \proofstepwidestar[1,2,3]{a<b\lor a=b\lor b<a}{%
      \FormulaRefAuto{RationalOrderTotal}[thm]{1,5,6}}
    \proofstepwidestar[8]{a<b}{\rA}
    \proofstepwidestar[1,2,3,8]{a\in B}{%
      \FormulaRefAuto{DedekindCutCharacterization}[thm]{4,6,8}}
    \proofstepwidestar[1,2,3,8]{\bot}{\rBI{5,9}}
    \proofstepwidestar[11]{a=b}{\rA}
    \proofstepwidestar[1,2,3,11]{a\in B}{\rIE{11,6}}
    \proofstepwidestar[1,2,3,11]{\bot}{\rBI{5,12}}
    \proofstepwidestar[14]{b<a}{\rA}
    \proofstepwidestar[1,2,3,14]{b\in A}{%
      \FormulaRefAuto{DedekindCutCharacterization}[thm]{4,5,14}}
    \proofstepwidestar[1,2,3,14]{\bot}{\rBI{6,15}}
    \proofstepwidestar[1,2,3]{\bot}{\rOE{7,8,10,11,13,14,16}}
  \closeproofpartwideindR

  \proofpartwide{Schluss}
    \proofstepwidestar[1]{A,B\in\mathbb R}{\rA}
    \proofstepwidestar[2]{\neg(A\subseteq B\lor B\subseteq A)}{\rA}
    \proofstepwidestar[2]{A\nsubseteq B\land B\nsubseteq A}{%
      \FormulaRefAuto{\neg(P\lor Q)\eqvdash\neg P\land\neg Q}{2}}
    \proofstepwidestar[1,2]{\exists a\in A\setminus B}{%
      \FormulaRefAuto{NonSubsetWitness}[thm]{3}}
    \proofstepwidestar[1,2]{\exists b\in B\setminus A}{%
      \FormulaRefAuto{NonSubsetWitness}[thm]{3}}
    \proofstepwidestar[1,2]{a\in A\setminus B}{\rA}
    \proofstepwidestar[1,2]{b\in B\setminus A}{\rA}
    \proofstepwidestar[1,2]{\bot}{%
      \FormulaRefAuto{RealCutOppositeWitnessesImpossible}[thm]{1,6,7}}
    \proofstepwidestar[1,2]{\bot}{\rEE{5,7,8}}
    \proofstepwidestar[1,2]{\bot}{\rEE{4,6,9}}
    \proofstepwidestar[1]{A\subseteq B\lor B\subseteq A}{\rCI{2,10}}
  \closeproofpartwide
\end{tabproofsplitwide}

\FormulaThmDeltaK[Die Inklusionsordnung der Schnitte ist total]{%
  \TotOrd{\mathbb R,\RealLe}%
}{RealCutTotalOrder}{%
  \DeltaRow{Mengen}{A\dsep B}%
  \DeltaRow{Relationsoperatoren}{\RealLe}%
}
\begin{tabproofsplitwide}
  \proofpartwideindR[Partielle Ordnung]{%
    \PartOrd{\mathbb R,\RealLe}%
  }{%
    \PartOrd{\mathbb R,\RealLe}%
  }[RealCutInclusionPartialOrder]
    \proofstepwidestarbreakkeep[1]{A\in\mathbb R}{\rA}
    \proofstepwidestarbreakkeep[1]{A\subseteq A}{%
      \FormulaRefAuto{A\subseteq A}}
    \proofstepwidestarbreakkeep[1]{A\RealLe A}{%
      \FormulaRefAuto{RealCutOrder}[def]{1,2}}
    \proofstepwidestarbreakkeep[]{%
      \forall A\in\mathbb R\,(A\RealLe A)}{\rUI{\rRI{1,3}}}
    \proofstepwidestarbreakkeep[5]{A,B,C\in\mathbb R}{\rA}
    \proofstepwidestarbreakkeep[6]{A\RealLe B\land B\RealLe C}{\rA}
    \proofstepwidestarbreakkeep[5,6]{A\subseteq B\land B\subseteq C}{%
      \FormulaRefAuto{RealCutOrderCharacterization}[thm]{5,6}}
    \proofstepwidestarbreakkeep[5,6]{A\subseteq C}{%
      \FormulaRefAuto{A\subseteq B\dsep B\subseteq C\vdash A\subseteq C}{7}}
    \proofstepwidestarbreakkeep[5,6]{A\RealLe C}{%
      \FormulaRefAuto{RealCutOrder}[def]{5,8}}
    \proofstepwidestarbreakkeep[]{%
      \forall A,B,C\in\mathbb R\,
      ((A\RealLe B\land B\RealLe C)\rightarrow A\RealLe C)}{%
      \rUI{\rRI{5,\rRI{6,9}}}}
    \proofstepwidestarbreakkeep[11]{A,B\in\mathbb R}{\rA}
    \proofstepwidestarbreakkeep[12]{A\RealLe B\land B\RealLe A}{\rA}
    \proofstepwidestarbreakkeep[11,12]{A\subseteq B\land B\subseteq A}{%
      \FormulaRefAuto{RealCutOrderCharacterization}[thm]{11,12}}
    \proofstepwidestarbreakkeep[11,12]{A=B}{%
      \FormulaRefAuto{A\subseteq B\dsep B\subseteq A\vdash A=B}{13}}
    \proofstepwidestarbreakkeep[]{%
      \forall A,B\in\mathbb R\,
      ((A\RealLe B\land B\RealLe A)\rightarrow A=B)}{%
      \rUI{\rRI{11,\rRI{12,14}}}}
    \proofstepwidestarbreakkeep[]{\PartOrd{\mathbb R,\RealLe}}{%
      \FormulaRefAuto{Partielle Ordnung}[def]{4,10,15}}
  \closeproofpartwideindR

  \proofpartwide{Vergleichbarkeit und Schluss}
    \proofstepwidestar[]{\PartOrd{\mathbb R,\RealLe}}{%
      \FormulaRefAuto{RealCutInclusionPartialOrder}[thm]}
    \proofstepwidestar[2]{A,B\in\mathbb R}{\rA}
    \proofstepwidestar[2]{A\subseteq B\lor B\subseteq A}{%
      \FormulaRefAuto{RealCutInclusionComparable}[thm]{2}}
    \proofstepwidestar[2]{A\RealLe B\lor B\RealLe A}{%
      \FormulaRefAuto{RealCutOrderCharacterization}[thm]{2,3}}
    \proofstepwidestar[]{%
      \forall A,B\in\mathbb R\,(A\RealLe B\lor B\RealLe A)}{%
      \rUI{\rRI{2,4}}}
    \proofstepwidestar[]{\TotOrd{\mathbb R,\RealLe}}{%
      \FormulaRefAuto{Totale Ordnung}[def]{1,5}}
  \closeproofpartwide
\end{tabproofsplitwide}

\FormulaThmDeltaK[Hauptschnitte bewahren und spiegeln die Ordnung]{%
  p,q\in\mathbb Q\vdash
  \bigl(p\leq q\leftrightarrow\RealEmb(p)\RealLe\RealEmb(q)\bigr)
  \land
  \bigl(p<q\leftrightarrow\RealEmb(p)\RealLt\RealEmb(q)\bigr)%
}{RationalToRealEmbeddingOrder}{%
  \DeltaRow{Mengen}{p\dsep q}%
  \DeltaRow{Funktionen}{\RealEmb}%
  \DeltaRow{Relationsoperatoren}{\RealLe\dsep\RealLt}%
}
\begin{tabproofsplitwide}
  \proofpartwideindR[Nichtstrikte Hinrichtung]{%
    p,q\in\mathbb Q\dsep p\leq q
    \vdash\RealHat p\subseteq\RealHat q
  }{%
    p,q\in\mathbb Q\dsep p\leq q
    \vdash\RealHat p\subseteq\RealHat q
  }[RealPrincipalCutOrderForward]
    \proofstepwidestar[1]{p,q\in\mathbb Q}{\rA}
    \proofstepwidestar[2]{p\leq q}{\rA}
    \proofstepwidestar[3]{x\in\RealHat p}{\rA}
    \proofstepwide[1,3]{x}{<}{p}{%
      \FormulaRefAuto{RealPrincipalCut}[def]{1,3}}
    \proofstepwide[1,2,3]{x}{<}{q}{%
      \FormulaRefAuto{RationalOrderTotal}[thm]{2,4}}
    \proofstepwidestar[1,2,3]{x\in\RealHat q}{%
      \FormulaRefAuto{RealPrincipalCut}[def]{1,5}}
    \proofstepwidestar[1,2]{\RealHat p\subseteq\RealHat q}{%
      \FormulaRefAuto{SubsetIntroductionRule}[thm]{[3]\vdots 6}}
  \closeproofpartwideindR

  \proofpartwideindR[Nichtstrikte Rückrichtung]{%
    p,q\in\mathbb Q\dsep\RealHat p\subseteq\RealHat q
    \vdash p\leq q
  }{%
    p,q\in\mathbb Q\dsep\RealHat p\subseteq\RealHat q
    \vdash p\leq q
  }[RealPrincipalCutOrderBackward]
    \proofstepwidestar[1]{p,q\in\mathbb Q}{\rA}
    \proofstepwidestar[2]{\RealHat p\subseteq\RealHat q}{\rA}
    \proofstepwidestar[3]{q<p}{\rA}
    \proofstepwidestar[1,3]{%
      \exists r\in\mathbb Q\,(q<r\land r<p)}{%
      \FormulaRefAuto{RationalOrderDense}[thm]{1,3}}
    \proofstepwidestar[1,3]{r\in\mathbb Q\land q<r\land r<p}{\rA}
    \proofstepwidestar[1,3]{r\in\RealHat p}{%
      \FormulaRefAuto{RealPrincipalCut}[def]{5}}
    \proofstepwidestar[1,2,3]{r\in\RealHat q}{%
      \FormulaRefAuto{A\subseteq B\dsep x\in A\vdash x\in B}{2,6}}
    \proofstepwidestar[1,3]{r\notin\RealHat q}{%
      \FormulaRefAuto{RealPrincipalCut}[def]{5}}
    \proofstepwidestar[1,2,3]{\bot}{\rBI{7,8}}
    \proofstepwidestar[1,2,3]{\bot}{\rEE{4,5,9}}
    \proofstepwidestar[1,2]{\neg(q<p)}{\rCI{3,10}}
    \proofstepwidestar[1,2]{p\leq q}{%
      \FormulaRefAuto{RationalOrderTotal}[thm]{1,11}}
  \closeproofpartwideindR

  \proofpartwideindR[Strikte Hinrichtung]{%
    p,q\in\mathbb Q\dsep p<q
    \vdash\RealHat p\subsetneq\RealHat q
  }{%
    p,q\in\mathbb Q\dsep p<q
    \vdash\RealHat p\subsetneq\RealHat q
  }[RealPrincipalCutStrictOrderForward]
    \proofstepwidestar[1]{p,q\in\mathbb Q}{\rA}
    \proofstepwidestar[2]{p<q}{\rA}
    \proofstepwidestar[1,2]{p\leq q}{%
      \FormulaRefAuto{RationalStrictOrderDef}[def]{1,2}}
    \proofstepwidestar[1,2]{\RealHat p\subseteq\RealHat q}{%
      \FormulaRefAuto{RealPrincipalCutOrderForward}[thm]{1,3}}
    \proofstepwidestar[1,2]{%
      \exists r\in\mathbb Q\,(p<r\land r<q)}{%
      \FormulaRefAuto{RationalOrderDense}[thm]{1,2}}
    \proofstepwidestar[6]{r\in\mathbb Q\land p<r\land r<q}{\rA}
    \proofstepwidestar[1,2,6]{r\in\RealHat q}{%
      \FormulaRefAuto{RealPrincipalCut}[def]{6}}
    \proofstepwidestar[1,2,6]{r\notin\RealHat p}{%
      \FormulaRefAuto{RealPrincipalCut}[def]{6}}
    \proofstepwidestar[1,2,6]{\RealHat p\neq\RealHat q}{%
      \FormulaRefAuto{x\in B, x\not\in A\vdash A\neq B}{7,8}}
    \proofstepwidestar[1,2,6]{\RealHat p\subsetneq\RealHat q}{%
      \FormulaRefAuto{A\subset B\coloneq A\subseteq B\land A\neq B}{4,9}}
    \proofstepwidestar[1,2]{\RealHat p\subsetneq\RealHat q}{%
      \rEE{5,6,10}}
  \closeproofpartwideindR

  \proofpartwideindR[Strikte Rückrichtung]{%
    p,q\in\mathbb Q\dsep\RealHat p\subsetneq\RealHat q
    \vdash p<q
  }{%
    p,q\in\mathbb Q\dsep\RealHat p\subsetneq\RealHat q
    \vdash p<q
  }[RealPrincipalCutStrictOrderBackward]
    \proofstepwidestar[1]{p,q\in\mathbb Q}{\rA}
    \proofstepwidestar[2]{\RealHat p\subsetneq\RealHat q}{\rA}
    \proofstepwidestar[2]{%
      \RealHat p\subseteq\RealHat q\land\RealHat p\neq\RealHat q}{%
      \FormulaRefAuto{A\subset B\coloneq A\subseteq B\land A\neq B}{2}}
    \proofstepwidestar[1,2]{p\leq q}{%
      \FormulaRefAuto{RealPrincipalCutOrderBackward}[thm]{1,3}}
    \proofstepwidestar[5]{p=q}{\rA}
    \proofstepwidestar[1,5]{\RealHat p=\RealHat q}{\rIE{5}}
    \proofstepwidestar[1,2,5]{\bot}{\rBI{3,6}}
    \proofstepwidestar[1,2]{p\neq q}{\rCI{5,7}}
    \proofstepwidestar[1,2]{p<q}{%
      \FormulaRefAuto{RationalStrictOrderDef}[def]{1,4,8}}
  \closeproofpartwideindR

  \proofpartwide{Einbettung und Zusammenfassung}
    \proofstepwidestarbreakkeep[1]{p,q\in\mathbb Q}{\rA}
    \proofstepwidestar[1]{p,q\in\mathbb Q\dsep p\leq q\vdash\RealHat p\subseteq\RealHat q}{%
      \FormulaRefAuto{RealPrincipalCutOrderForward}[thm]}
    \proofstepwidestar[1]{p,q\in\mathbb Q\dsep\RealHat p\subseteq\RealHat q\vdash p\leq q}{%
      \FormulaRefAuto{RealPrincipalCutOrderBackward}[thm]}
    \proofstepwide[1]{p\leq q}{\leftrightarrow}{%
      \RealHat p\subseteq\RealHat q}{%
      \rLRI{2,
        3}}
    \proofstepwidestar[1]{p,q\in\mathbb Q\dsep p<q\vdash\RealHat p\subsetneq\RealHat q}{%
      \FormulaRefAuto{RealPrincipalCutStrictOrderForward}[thm]}
    \proofstepwidestar[1]{p,q\in\mathbb Q\dsep\RealHat p\subsetneq\RealHat q\vdash p<q}{%
      \FormulaRefAuto{RealPrincipalCutStrictOrderBackward}[thm]}
    \proofstepwide[1]{p<q}{\leftrightarrow}{%
      \RealHat p\subsetneq\RealHat q}{%
      \rLRI{5,
        6}}
    \proofstepwidestarbreakkeep[1]{%
      \RealEmb(p)=\RealHat p\land\RealEmb(q)=\RealHat q}{%
      \FormulaRefAuto{RationalToRealEmbedding}[def]{1}}
    \proofstepwide[1]{p\leq q}{\leftrightarrow}{%
      \RealEmb(p)\RealLe\RealEmb(q)}{%
      \FormulaRefAuto{RealCutOrderCharacterization}[thm]{4,8}}
    \proofstepwide[1]{p<q}{\leftrightarrow}{%
      \RealEmb(p)\RealLt\RealEmb(q)}{%
      \FormulaRefAuto{RealCutOrderCharacterization}[thm]{7,8}}
    \proofstepwidestarbreakkeep[1]{%
      \begin{aligned}[t]
        &(p\leq q\leftrightarrow\RealEmb(p)\RealLe\RealEmb(q))\land{}\\[-2pt]
        &(p<q\leftrightarrow\RealEmb(p)\RealLt\RealEmb(q))
      \end{aligned}}{%
      \rAI{9,10}}
  \closeproofpartwide
\end{tabproofsplitwide}

\FormulaThmDeltaK[Injektivität der rationalen Einbettung]{%
  \RealEmb\colon\mathbb Q\inj\mathbb R%
}{RationalToRealEmbeddingInjective}{%
  \DeltaRow{Funktionen}{\RealEmb}%
}
\begin{tabproofsplitwide}
  \proofpartwideindR[Punktweise Injektivität]{%
    p,q\in\mathbb Q\dsep\RealEmb(p)=\RealEmb(q)\vdash p=q
  }{%
    p,q\in\mathbb Q\dsep\RealEmb(p)=\RealEmb(q)\vdash p=q
  }[RationalToRealEmbeddingPointInjective]
    \proofstepwidestar[1]{p,q\in\mathbb Q}{\rA}
    \proofstepwidestar[2]{\RealEmb(p)=\RealEmb(q)}{\rA}
    \proofstepwidestar[1,2]{%
      \RealEmb(p)\RealLe\RealEmb(q)\land
      \RealEmb(q)\RealLe\RealEmb(p)}{%
      \FormulaRefAuto{RealCutOrderCharacterization}[thm]{1,2}}
    \proofstepwidestar[1,2]{p\leq q\land q\leq p}{%
      \FormulaRefAuto{RationalToRealEmbeddingOrder}[thm]{1,3}}
    \proofstepwidestar[1,2]{p=q}{%
      \FormulaRefAuto{RationalOrderTotal}[thm]{1,4}}
  \closeproofpartwideindR

  \proofpartwide{Funktionsschluss}
    \proofstepwidestar[]{\RealEmb\colon\mathbb Q\to\mathbb R}{%
      \FormulaRefAuto{RationalToRealEmbeddingFunction}[thm]}
    \proofstepwidestar[]{%
      \forall p,q\in\mathbb Q\,
      (\RealEmb(p)=\RealEmb(q)\rightarrow p=q)}{%
      \FormulaRefAuto{RationalToRealEmbeddingPointInjective}[thm]}
    \proofstepwidestar[]{\RealEmb\colon\mathbb Q\inj\mathbb R}{%
      \FormulaRefAuto{Injektive Funktion}[def]{1,2}}
  \closeproofpartwide
\end{tabproofsplitwide}

\section{Kanonische Teilmengenkopien}

Die rationale Einbettung transportiert die in \(\mathbb Q\) bereits
geschachtelten natürlichen und ganzzahligen Kopien in den reellen Träger.
Damit entsteht innerhalb von \(\mathbb R\) eine wörtliche Inklusionskette.

\FormulaDefDeltaK[Kanonische rationale Kopie in den reellen Zahlen]{%
  \RatCopyInReal\coloneqq\RealEmb[\mathbb Q]%
}{RationalCopyInReals}{%
  \DeltaRow{Mengen}{\RatCopyInReal}%
  \DeltaRow{Funktion}{\RealEmb}%
  \DeltaRow{Neue Symbole}{\RatCopyInReal}%
}

\FormulaDefDeltaK[Kanonische ganzzahlige Kopie in den reellen Zahlen]{%
  \IntCopyInReal\coloneqq\RealEmb[\IntCopyInRat]%
}{IntegerCopyInReals}{%
  \DeltaRow{Mengen}{\IntCopyInRat\dsep\IntCopyInReal}%
  \DeltaRow{Funktion}{\RealEmb}%
  \DeltaRow{Neue Symbole}{\IntCopyInReal}%
}

\FormulaDefDeltaK[Kanonische natürliche Kopie in den reellen Zahlen]{%
  \NatCopyInReal\coloneqq\RealEmb[\NatCopyInRat]%
}{NaturalCopyInReals}{%
  \DeltaRow{Mengen}{\NatCopyInRat\dsep\NatCopyInReal}%
  \DeltaRow{Funktion}{\RealEmb}%
  \DeltaRow{Neue Symbole}{\NatCopyInReal}%
}

\FormulaThmDeltaK[Geschachtelter Zahlbereichsturm in den reellen Zahlen]{%
  \begin{aligned}[t]
    &\bigl(\NatCopyInReal\subseteq\IntCopyInReal\land
      \IntCopyInReal\subseteq\RatCopyInReal\land
      \RatCopyInReal\subseteq\mathbb R\bigr)\land{}\\[-2pt]
    &\bigl(\NatCopyInReal\subseteq\mathbb R\land
      \IntCopyInReal\subseteq\mathbb R\bigr)
  \end{aligned}%
}{EmbeddedNumberTowerInReals}{%
  \DeltaRow{Mengen}{\NatCopyInRat\dsep\IntCopyInRat\dsep
    \NatCopyInReal\dsep\IntCopyInReal\dsep\RatCopyInReal}%
  \DeltaRow{Funktion}{\RealEmb}%
}
\begin{tabproofsplitwide}
\proofpartwide{Einbettungen und Bildinklusionen}
\proofstepwidestar[]{\RealEmb\colon\mathbb Q\to\mathbb R}{%
    \FormulaRefAuto{RationalToRealEmbeddingFunction}[thm]}
  \proofstepwidestar[]{%
    \bigl(\NatCopyInRat\subseteq\IntCopyInRat\land
    \IntCopyInRat\subseteq\mathbb Q\bigr)
    \land\NatCopyInRat\subseteq\mathbb Q}{%
    \FormulaRefAuto{EmbeddedNaturalIntegerRationalTower}[thm]}
  \proofstepwidestar[]{\NatCopyInRat\subseteq\IntCopyInRat}{%
    \rAEa{\rAEa{2}}}
  \proofstepwidestar[]{\IntCopyInRat\subseteq\mathbb Q}{%
    \rAEb{\rAEa{2}}}
  \proofstepwidestar[]{\mathbb Q\subseteq\mathbb Q}{%
    \FormulaRefAuto{A\subseteq A}}
  \proofstepwidestar[]{%
    \RealEmb[\NatCopyInRat]\subseteq\RealEmb[\IntCopyInRat]}{%
    \FormulaRefAuto{F\colon M\to N\dsep A\subseteq B\dsep
      B\subseteq M\vdash F[A]\subseteq F[B]}[thm]{1,3,4}}
  \proofstepwidestar[]{%
    \RealEmb[\IntCopyInRat]\subseteq\RealEmb[\mathbb Q]}{%
    \FormulaRefAuto{F\colon M\to N\dsep A\subseteq B\dsep
      B\subseteq M\vdash F[A]\subseteq F[B]}[thm]{1,4,5}}
  \proofstepwidestar[]{\RealEmb[\mathbb Q]\subseteq\mathbb R}{%
    \FormulaRefAuto{F\colon A\to B\vdash F[A]\subseteq B}[thm]{1}}
  \proofstepwidestar[]{\NatCopyInReal=\RealEmb[\NatCopyInRat]}{%
    \FormulaRefAuto{NaturalCopyInReals}[def]}
  \proofstepwidestar[]{\IntCopyInReal=\RealEmb[\IntCopyInRat]}{%
    \FormulaRefAuto{IntegerCopyInReals}[def]}
  \proofstepwidestar[]{\RatCopyInReal=\RealEmb[\mathbb Q]}{%
    \FormulaRefAuto{RationalCopyInReals}[def]}

\closeproofpartwide

\proofpartwide{Natürliche Kopie in der ganzzahligen Kopie}
  \setcounter{proofstepnr}{11}% Fortlaufende Schritte dieses Beweises.
\proofstepwidestar[]{\NatCopyInReal\subseteq\IntCopyInReal}{\rIE{9,10,6}}

\closeproofpartwide

\proofpartwide{Ganzzahlige Kopie in der rationalen Kopie}
  \setcounter{proofstepnr}{12}% Fortlaufende Schritte dieses Beweises.
\proofstepwidestar[]{\IntCopyInReal\subseteq\RatCopyInReal}{\rIE{10,11,7}}

\closeproofpartwide

\proofpartwide{Rationale Kopie im reellen Träger}
  \setcounter{proofstepnr}{13}% Fortlaufende Schritte dieses Beweises.
\proofstepwidestar[]{\RatCopyInReal\subseteq\mathbb R}{\rIE{11,8}}

\closeproofpartwide

\proofpartwide{Übrige Trägerinklusionen}
  \setcounter{proofstepnr}{14}% Fortlaufende Schritte dieses Beweises.
\proofstepwidestar[]{\IntCopyInReal\subseteq\mathbb R}{%
    \FormulaRefAuto{A\subseteq B\dsep B\subseteq C\vdash A\subseteq C}{13,14}}
  \proofstepwidestar[]{\NatCopyInReal\subseteq\mathbb R}{%
    \FormulaRefAuto{A\subseteq B\dsep B\subseteq C\vdash A\subseteq C}{12,15}}

\closeproofpartwide

\proofpartwide{Zusammenführung}
  \setcounter{proofstepnr}{16}% Fortlaufende Schritte dieses Beweises.
\proofstepwidestar[]{%
    \begin{aligned}[t]
      &\bigl(\NatCopyInReal\subseteq\IntCopyInReal\land
        \IntCopyInReal\subseteq\RatCopyInReal\land
        \RatCopyInReal\subseteq\mathbb R\bigr)\land{}\\[-2pt]
      &\bigl(\NatCopyInReal\subseteq\mathbb R\land
        \IntCopyInReal\subseteq\mathbb R\bigr)
    \end{aligned}}{\rAI{\rAI{12,13,14},\rAI{16,15}}}

\closeproofpartwide
\end{tabproofsplitwide}

\FormulaThmDeltaK[Bijektive Identifikation der natürlichen Kopie]{%
  \RealEmb\restriction_{\NatCopyInRat}\colon
  \NatCopyInRat\bij\NatCopyInReal%
}{NaturalCopyRealBijection}{%
  \DeltaRow{Mengen}{\NatCopyInRat\dsep\NatCopyInReal}%
  \DeltaRow{Funktion}{\RealEmb}%
}
\begin{tabproofwide}
  \proofstepwidestar[]{\RealEmb\colon\mathbb Q\inj\mathbb R}{%
    \FormulaRefAuto{RationalToRealEmbeddingInjective}[thm]}
  \proofstepwidestar[]{( \NatCopyInRat \subseteq \IntCopyInRat \land \IntCopyInRat \subseteq \mathbb Q ) \land \NatCopyInRat \subseteq \mathbb Q}{%
      \FormulaRefAuto{EmbeddedNaturalIntegerRationalTower}[thm]}
  \proofstepwidestar[]{\NatCopyInRat\subseteq\mathbb Q}{%
    \rAEb{2}}
  \proofstepwidestar[]{%
    \RealEmb\restriction_{\NatCopyInRat}\colon
    \NatCopyInRat\inj\mathbb R}{%
    \FormulaRefAuto{F\colon A\inj B\dsep C\subseteq A
      \vdash F\restriction_C\colon C\inj B}[thm]{1,3}}
  \proofstepwidestar[]{%
    \RealEmb\restriction_{\NatCopyInRat}\colon\NatCopyInRat\bij
    (\RealEmb\restriction_{\NatCopyInRat})[\NatCopyInRat]}{%
    \FormulaRefAuto{F\colon A\inj B\vdash F\colon A\bij F[A]}[thm]{4}}
  \proofstepwidestar[]{\RealEmb \colon \mathbb Q \to \mathbb R}{%
      \FormulaRefAuto{Injektive Funktion}[def]{1}}
  \proofstepwidestar[]{%
    (\RealEmb\restriction_{\NatCopyInRat})[\NatCopyInRat]
    =\RealEmb[\NatCopyInRat]}{%
    \FormulaRefAuto{C\subseteq A\dsep F\colon A\to B
      \vdash (F\restriction_C)[C]=F[C]}[thm]{%
        3,6}}
  \proofstepwidestar[]{\NatCopyInReal=\RealEmb[\NatCopyInRat]}{%
    \FormulaRefAuto{NaturalCopyInReals}[def]}
  \proofstepwidestar[]{%
    \RealEmb\restriction_{\NatCopyInRat}\colon
    \NatCopyInRat\bij\NatCopyInReal}{\rIE{7,8,5}}
\end{tabproofwide}

\FormulaThmDeltaK[Bijektive Identifikation der ganzzahligen Kopie]{%
  \RealEmb\restriction_{\IntCopyInRat}\colon
  \IntCopyInRat\bij\IntCopyInReal%
}{IntegerCopyRealBijection}{%
  \DeltaRow{Mengen}{\IntCopyInRat\dsep\IntCopyInReal}%
  \DeltaRow{Funktion}{\RealEmb}%
}
\begin{tabproofwide}
  \proofstepwidestar[]{\RealEmb\colon\mathbb Q\inj\mathbb R}{%
    \FormulaRefAuto{RationalToRealEmbeddingInjective}[thm]}
  \proofstepwidestar[]{( \NatCopyInRat \subseteq \IntCopyInRat \land \IntCopyInRat \subseteq \mathbb Q ) \land \NatCopyInRat \subseteq \mathbb Q}{%
      \FormulaRefAuto{EmbeddedNaturalIntegerRationalTower}[thm]}
  \proofstepwidestar[]{\IntCopyInRat\subseteq\mathbb Q}{%
    \rAEb{\rAEa{2}}}
  \proofstepwidestar[]{%
    \RealEmb\restriction_{\IntCopyInRat}\colon
    \IntCopyInRat\inj\mathbb R}{%
    \FormulaRefAuto{F\colon A\inj B\dsep C\subseteq A
      \vdash F\restriction_C\colon C\inj B}[thm]{1,3}}
  \proofstepwidestar[]{%
    \RealEmb\restriction_{\IntCopyInRat}\colon\IntCopyInRat\bij
    (\RealEmb\restriction_{\IntCopyInRat})[\IntCopyInRat]}{%
    \FormulaRefAuto{F\colon A\inj B\vdash F\colon A\bij F[A]}[thm]{4}}
  \proofstepwidestar[]{\RealEmb \colon \mathbb Q \to \mathbb R}{%
      \FormulaRefAuto{Injektive Funktion}[def]{1}}
  \proofstepwidestar[]{%
    (\RealEmb\restriction_{\IntCopyInRat})[\IntCopyInRat]
    =\RealEmb[\IntCopyInRat]}{%
    \FormulaRefAuto{C\subseteq A\dsep F\colon A\to B
      \vdash (F\restriction_C)[C]=F[C]}[thm]{%
        3,6}}
  \proofstepwidestar[]{\IntCopyInReal=\RealEmb[\IntCopyInRat]}{%
    \FormulaRefAuto{IntegerCopyInReals}[def]}
  \proofstepwidestar[]{%
    \RealEmb\restriction_{\IntCopyInRat}\colon
    \IntCopyInRat\bij\IntCopyInReal}{\rIE{7,8,5}}
\end{tabproofwide}

\FormulaThmDeltaK[Bijektive Identifikation der rationalen Kopie]{%
  \RealEmb\colon\mathbb Q\bij\RatCopyInReal%
}{RationalCopyRealBijection}{%
  \DeltaRow{Mengen}{\RatCopyInReal}%
  \DeltaRow{Funktion}{\RealEmb}%
}
\begin{tabproof}
  \proofstep{}{\RealEmb\colon\mathbb Q\inj\mathbb R}{%
    \FormulaRefAuto{RationalToRealEmbeddingInjective}[thm]}
  \proofstep{}{\RealEmb\colon\mathbb Q\bij\RealEmb[\mathbb Q]}{%
    \FormulaRefAuto{F\colon A\inj B\vdash F\colon A\bij F[A]}[thm]{1}}
  \proofstep{}{\RatCopyInReal=\RealEmb[\mathbb Q]}{%
    \FormulaRefAuto{RationalCopyInReals}[def]}
  \proofstep{}{\RealEmb\colon\mathbb Q\bij\RatCopyInReal}{\rIE{3,2}}
\end{tabproof}

\section{Suprema als Vereinigungen}

\FormulaThmDeltaK[Vereinigung einer beschränkten Schnittfamilie]{%
  \begin{aligned}[t]
    &\varnothing\neq\mathcal A\subseteq\mathbb R\dsep
      U\in\mathbb R\dsep
      \forall A\in\mathcal A\,(A\RealLe U)\\[-2pt]
    &\hspace{34mm}\vdash\RealCut{\bigcup\mathcal A}.
  \end{aligned}%
}{RealBoundedCutUnionIsCut}{%
  \DeltaRow{Mengen}{\mathcal A\dsep A\dsep U}%
  \DeltaRow{Relationsoperatoren}{\RealLe}%
}
\begin{tabproofsplitwide}
\proofpartwide{Nichtleerheit}
\proofstepwidestarbreakkeep[1]{\varnothing\neq\mathcal A\subseteq\mathbb R}{\rA}
    \proofstepwidestarbreakkeep[2]{U\in\mathbb R}{\rA}
    \proofstepwidestarbreakkeep[3]{\forall A\in\mathcal A\,(A\RealLe U)}{\rA}
    \proofstepwidestarbreakkeep[1]{\exists A_0\in\mathcal A}{%
      \FormulaRefAuto{A\neq\varnothing\eqvdash\exists x\,(x\in A)}{1}}
    \proofstepwidestarbreakkeep[1]{A_0\in\mathcal A}{\rA}
    \proofstepwidestarbreakkeep[1,5]{A_0\in\mathbb R}{%
      \FormulaRefAuto{A\subseteq B\dsep x\in A\vdash x\in B}{1,5}}
    \proofstepwidestarbreakkeep[1,5]{\RealCut{A_0}}{%
      \FormulaRefAuto{RealCutMembership}[thm]{6}}
    \proofstepwidestarbreakkeep[1,5]{\exists a_0\in A_0}{%
      \FormulaRefAuto{DedekindCutCharacterization}[thm]{7}}
    \proofstepwidestarbreakkeep[1,5]{a_0\in A_0}{\rA}
    \proofstepwidestarbreakkeep[1,5]{a_0\in\bigcup\mathcal A}{%
      \FormulaRefAuto{FamilyMemberSubsetUnion}[thm]{5,9}}
    \proofstepwidestarbreakkeep[1,5]{\bigcup\mathcal A\neq\varnothing}{%
      \FormulaRefAuto{x\in A\vdash A\neq\varnothing}{10}}
    \proofstepwidestarbreakkeep[1]{\bigcup\mathcal A\neq\varnothing}{%
      \rEE{8,9,11}}
    \proofstepwidestarbreakkeep[1]{\bigcup\mathcal A\neq\varnothing}{%
      \rEE{4,5,12}}

\closeproofpartwide

\proofpartwide{Trägerinklusion}
  \setcounter{proofstepnr}{13}% Fortlaufende Schritte dieses Beweises.
\proofstepwidestarbreakkeep[14]{x\in\bigcup\mathcal A}{\rA}
    \proofstepwidestarbreakkeep[14]{\exists A\in\mathcal A\,(x\in A)}{%
      \FormulaRefAuto{x\in\bigcup A\eqvdash\exists B\in A\,(x\in B)}{14}}
    \proofstepwidestarbreakkeep[14]{A\in\mathcal A\land x\in A}{\rA}
    \proofstepwidestarbreakkeep[1,14]{A\in\mathbb R}{%
      \FormulaRefAuto{A\subseteq B\dsep x\in A\vdash x\in B}{1,16}}
    \proofstepwidestarbreakkeep[1,14]{A\subseteq\mathbb Q}{%
      \FormulaRefAuto{RealCutMembership}[thm]{17}}
    \proofstepwidestarbreakkeep[1,14]{x\in\mathbb Q}{%
      \FormulaRefAuto{A\subseteq B\dsep x\in A\vdash x\in B}{18,16}}
    \proofstepwidestarbreakkeep[1,14]{x\in\mathbb Q}{\rEE{15,16,19}}
    \proofstepwidestarbreakkeep[1]{\bigcup\mathcal A\subseteq\mathbb Q}{%
      \FormulaRefAuto{SubsetIntroductionRule}[thm]{[14]\vdots 20}}

\closeproofpartwide

\proofpartwide{Echtheit}
  \setcounter{proofstepnr}{21}% Fortlaufende Schritte dieses Beweises.
\proofstepwidestarbreakkeep[22]{x\in\bigcup\mathcal A}{\rA}
    \proofstepwidestarbreakkeep[22]{\exists A\in\mathcal A\,(x\in A)}{%
      \FormulaRefAuto{x\in\bigcup A\eqvdash\exists B\in A\,(x\in B)}{22}}
    \proofstepwidestarbreakkeep[22]{A\in\mathcal A\land x\in A}{\rA}
    \proofstepwidestarbreakkeep[3,22]{A\RealLe U}{\rUE{3,24}}
    \proofstepwidestarbreakkeep[3,22]{A\subseteq U}{%
      \FormulaRefAuto{RealCutOrderCharacterization}[thm]{2,24,25}}
    \proofstepwidestarbreakkeep[3,22]{x\in U}{%
      \FormulaRefAuto{A\subseteq B\dsep x\in A\vdash x\in B}{26,24}}
    \proofstepwidestarbreakkeep[3,22]{x\in U}{\rEE{23,24,27}}
    \proofstepwidestarbreakkeep[3]{\bigcup\mathcal A\subseteq U}{%
      \FormulaRefAuto{SubsetIntroductionRule}[thm]{[22]\vdots 28}}
    \proofstepwidestarbreakkeep[2]{\RealCut U}{%
      \FormulaRefAuto{RealCutMembership}[thm]{2}}
    \proofstepwidestarbreakkeep[2]{U\neq\mathbb Q}{%
      \FormulaRefAuto{DedekindCutCharacterization}[thm]{30}}
    \proofstepwidestarbreakkeep[32]{\bigcup\mathcal A=\mathbb Q}{\rA}
    \proofstepwidestarbreakkeep[2,3,32]{\mathbb Q\subseteq U}{\rIE{32,29}}
    \proofstepwidestarbreakkeep[2]{U\subseteq\mathbb Q}{%
      \FormulaRefAuto{DedekindCutCharacterization}[thm]{30}}
    \proofstepwidestarbreakkeep[2,3,32]{U=\mathbb Q}{%
      \FormulaRefAuto{A\subseteq B\dsep B\subseteq A\vdash A=B}{34,33}}
    \proofstepwidestarbreakkeep[2,3,32]{\bot}{\rBI{31,35}}
    \proofstepwidestarbreakkeep[2,3]{\bigcup\mathcal A\neq\mathbb Q}{%
      \rCI{32,36}}

\closeproofpartwide

\proofpartwideindR[Träger, Nichtleerheit und Echtheit]{%
    \begin{aligned}[t]
      &\varnothing\neq\mathcal A\subseteq\mathbb R\dsep U\in\mathbb R\dsep
        \forall A\in\mathcal A\,(A\RealLe U)\\[-2pt]
      &\quad\vdash
        \bigcup\mathcal A\subseteq\mathbb Q\land
        \bigcup\mathcal A\neq\varnothing\land
        \bigcup\mathcal A\neq\mathbb Q
    \end{aligned}
  }{%
    \varnothing\neq\mathcal A\subseteq\mathbb R\dsep U\in\mathbb R\dsep
    \forall A\in\mathcal A\,(A\RealLe U)
    \vdash
    \bigcup\mathcal A\subseteq\mathbb Q\land
    \bigcup\mathcal A\neq\varnothing\land
    \bigcup\mathcal A\neq\mathbb Q
  }[RealBoundedCutUnionCarrierProper]
  \setcounter{proofstepnr}{37}% Fortlaufende Schritte dieses Beweises.
\proofstepwidestarbreakkeep[1,2,3]{%
      \bigcup\mathcal A\subseteq\mathbb Q\land
      \bigcup\mathcal A\neq\varnothing\land
      \bigcup\mathcal A\neq\mathbb Q}{%
      \rAI{21,\rAI{13,37}}}

\closeproofpartwideindR


  \proofpartwideindR[Abwärtsabschluss]{%
    \begin{aligned}[t]
      &\mathcal A\subseteq\mathbb R\dsep p\in\bigcup\mathcal A\dsep
        q\in\mathbb Q\dsep q<p\\[-2pt]
      &\quad\vdash q\in\bigcup\mathcal A
    \end{aligned}
  }{%
    \mathcal A\subseteq\mathbb R\dsep p\in\bigcup\mathcal A\dsep
    q\in\mathbb Q\dsep q<p\vdash q\in\bigcup\mathcal A
  }[RealCutUnionDownwardClosed]
    \proofstepwidestar[1]{\mathcal A\subseteq\mathbb R}{\rA}
    \proofstepwidestar[2]{p\in\bigcup\mathcal A}{\rA}
    \proofstepwidestar[3]{q\in\mathbb Q}{\rA}
    \proofstepwidestar[4]{q<p}{\rA}
    \proofstepwidestar[2]{\exists A\in\mathcal A\,(p\in A)}{%
      \FormulaRefAuto{x\in\bigcup A\eqvdash\exists B\in A\,(x\in B)}{2}}
    \proofstepwidestar[2]{A\in\mathcal A\land p\in A}{\rA}
    \proofstepwidestar[1,2]{A\in\mathbb R}{%
      \FormulaRefAuto{A\subseteq B\dsep x\in A\vdash x\in B}{1,6}}
    \proofstepwidestar[1,2]{\RealCut A}{%
      \FormulaRefAuto{RealCutMembership}[thm]{7}}
    \proofstepwidestar[1,2,3,4]{q\in A}{%
      \FormulaRefAuto{DedekindCutCharacterization}[thm]{8,6,3,4}}
    \proofstepwidestar[1,2,3,4]{q\in\bigcup\mathcal A}{%
      \FormulaRefAuto{FamilyMemberSubsetUnion}[thm]{6,9}}
    \proofstepwidestar[1,2,3,4]{q\in\bigcup\mathcal A}{%
      \rEE{5,6,10}}
  \closeproofpartwideindR

  \proofpartwideindR[Kein größtes Element]{%
    \mathcal A\subseteq\mathbb R\dsep p\in\bigcup\mathcal A
    \vdash\exists r\in\bigcup\mathcal A\,(p<r)
  }{%
    \mathcal A\subseteq\mathbb R\dsep p\in\bigcup\mathcal A
    \vdash\exists r\in\bigcup\mathcal A\,(p<r)
  }[RealCutUnionNoMaximum]
    \proofstepwidestar[1]{\mathcal A\subseteq\mathbb R}{\rA}
    \proofstepwidestar[2]{p\in\bigcup\mathcal A}{\rA}
    \proofstepwidestar[2]{\exists A\in\mathcal A\,(p\in A)}{%
      \FormulaRefAuto{x\in\bigcup A\eqvdash\exists B\in A\,(x\in B)}{2}}
    \proofstepwidestar[2]{A\in\mathcal A\land p\in A}{\rA}
    \proofstepwidestar[1,2]{A\in\mathbb R}{%
      \FormulaRefAuto{A\subseteq B\dsep x\in A\vdash x\in B}{1,4}}
    \proofstepwidestar[1,2]{\RealCut A}{%
      \FormulaRefAuto{RealCutMembership}[thm]{5}}
    \proofstepwidestar[1,2]{\exists r\in A\,(p<r)}{%
      \FormulaRefAuto{DedekindCutCharacterization}[thm]{6,4}}
    \proofstepwidestar[1,2]{r\in A\land p<r}{\rA}
    \proofstepwidestar[1,2]{r\in\bigcup\mathcal A}{%
      \FormulaRefAuto{FamilyMemberSubsetUnion}[thm]{4,8}}
    \proofstepwidestar[1,2]{\exists r\in\bigcup\mathcal A\,(p<r)}{%
      \rEI{8,9}}
    \proofstepwidestar[1,2]{\exists r\in\bigcup\mathcal A\,(p<r)}{%
      \rEE{7,8,10}}
    \proofstepwidestar[1,2]{\exists r\in\bigcup\mathcal A\,(p<r)}{%
      \rEE{3,4,11}}
  \closeproofpartwideindR

  \proofpartwide{Schluss}
    \proofstepwidestar[1]{\varnothing\neq\mathcal A\subseteq\mathbb R}{\rA}
    \proofstepwidestar[2]{U\in\mathbb R}{\rA}
    \proofstepwidestar[3]{\forall A\in\mathcal A\,(A\RealLe U)}{\rA}
    \proofstepwidestar[1,2,3]{%
      \bigcup\mathcal A\subseteq\mathbb Q\land
      \bigcup\mathcal A\neq\varnothing\land
      \bigcup\mathcal A\neq\mathbb Q}{%
      \FormulaRefAuto{RealBoundedCutUnionCarrierProper}[thm]{1,2,3}}
    \proofstepwidestar[1]{%
      \forall p\in\bigcup\mathcal A\,\forall q\in\mathbb Q\,
      (q<p\rightarrow q\in\bigcup\mathcal A)}{%
      \FormulaRefAuto{RealCutUnionDownwardClosed}[thm]{1}}
    \proofstepwidestar[1]{%
      \forall p\in\bigcup\mathcal A\,
      \exists r\in\bigcup\mathcal A\,(p<r)}{%
      \FormulaRefAuto{RealCutUnionNoMaximum}[thm]{1}}
    \proofstepwidestar[1,2,3]{\RealCut{\bigcup\mathcal A}}{%
      \FormulaRefAuto{DedekindCutCharacterization}[thm]{4,5,6}}
  \closeproofpartwide
\end{tabproofsplitwide}

\FormulaThmDeltaK[Vollständigkeit der reellen Ordnung]{%
  \begin{aligned}[t]
    &\varnothing\neq\mathcal A\subseteq\mathbb R\dsep
      \exists U\in\mathbb R\,
        \forall A\in\mathcal A\,(A\RealLe U)\\[-2pt]
    &\hspace{19mm}\vdash
      \exists S\in\mathbb R\,
      \left(
        \begin{aligned}[t]
        &\forall A\in\mathcal A\,(A\RealLe S)\land{}\\[-2pt]
        &\forall V\in\mathbb R\,
          \bigl(\forall A\in\mathcal A\,(A\RealLe V)
          \rightarrow S\RealLe V\bigr)
        \end{aligned}
      \right).
  \end{aligned}%
}{RealLeastUpperBoundProperty}{%
  \DeltaRow{Mengen}{\mathcal A\dsep A\dsep S\dsep U\dsep V}%
  \DeltaRow{Relationsoperatoren}{\RealLe}%
}
\begin{tabproofsplitwide}
  \proofpartwideindR[Vereinigung ist obere Schranke]{%
    \varnothing\neq\mathcal A\subseteq\mathbb R\dsep
    S=\bigcup\mathcal A
    \vdash\forall A\in\mathcal A\,(A\RealLe S)
  }{%
    \varnothing\neq\mathcal A\subseteq\mathbb R\dsep
    S=\bigcup\mathcal A
    \vdash\forall A\in\mathcal A\,(A\RealLe S)
  }[RealCutUnionUpperBound]
    \proofstepwidestar[1]{\varnothing\neq\mathcal A\subseteq\mathbb R}{\rA}
    \proofstepwidestar[2]{S=\bigcup\mathcal A}{\rA}
    \proofstepwidestar[3]{A\in\mathcal A}{\rA}
    \proofstepwidestar[1,3]{A\in\mathbb R}{%
      \FormulaRefAuto{A\subseteq B\dsep x\in A\vdash x\in B}{1,3}}
    \proofstepwidestar[3]{A\subseteq\bigcup\mathcal A}{%
      \FormulaRefAuto{FamilyMemberSubsetUnion}[thm]{3}}
    \proofstepwidestar[2,3]{A\subseteq S}{\rIE{2,5}}
    \proofstepwidestar[1,2,3]{A\RealLe S}{%
      \FormulaRefAuto{RealCutOrder}[def]{4,6}}
    \proofstepwidestar[1,2]{\forall A\in\mathcal A\,(A\RealLe S)}{%
      \rUI{\rRI{3,7}}}
  \closeproofpartwideindR

  \proofpartwideindR[Vereinigung liegt unter jeder oberen Schranke]{%
    \begin{aligned}[t]
      &\mathcal A\subseteq\mathbb R\dsep V\in\mathbb R\dsep
        \forall A\in\mathcal A\,(A\RealLe V)\\[-2pt]
      &\quad\vdash\bigcup\mathcal A\RealLe V
    \end{aligned}
  }{%
    \mathcal A\subseteq\mathbb R\dsep V\in\mathbb R\dsep
    \forall A\in\mathcal A\,(A\RealLe V)
    \vdash\bigcup\mathcal A\RealLe V
  }[RealCutUnionLeastUpperBound]
    \proofstepwidestar[1]{\mathcal A\subseteq\mathbb R}{\rA}
    \proofstepwidestar[2]{V\in\mathbb R}{\rA}
    \proofstepwidestar[3]{\forall A\in\mathcal A\,(A\RealLe V)}{\rA}
    \proofstepwidestar[4]{x\in\bigcup\mathcal A}{\rA}
    \proofstepwidestar[4]{\exists A\in\mathcal A\,(x\in A)}{%
      \FormulaRefAuto{x\in\bigcup A\eqvdash\exists B\in A\,(x\in B)}{4}}
    \proofstepwidestar[4]{A\in\mathcal A\land x\in A}{\rA}
    \proofstepwidestar[1,4]{A\in\mathbb R}{%
      \FormulaRefAuto{A\subseteq B\dsep x\in A\vdash x\in B}{1,6}}
    \proofstepwidestar[3,4]{A\RealLe V}{\rUE{3,6}}
    \proofstepwidestar[1,2,3,4]{A\subseteq V}{%
      \FormulaRefAuto{RealCutOrderCharacterization}[thm]{7,2,8}}
    \proofstepwidestar[1,2,3,4]{x\in V}{%
      \FormulaRefAuto{A\subseteq B\dsep x\in A\vdash x\in B}{9,6}}
    \proofstepwidestar[1,2,3,4]{x\in V}{\rEE{5,6,10}}
    \proofstepwidestar[1,2,3]{\bigcup\mathcal A\subseteq V}{%
      \FormulaRefAuto{SubsetIntroductionRule}[thm]{[4]\vdots 11}}
    \proofstepwidestar[1,2,3]{\bigcup\mathcal A\RealLe V}{%
      \FormulaRefAuto{RealCutOrder}[def]{2,12}}
  \closeproofpartwideindR

  \proofpartwide{Existenz der kleinsten oberen Schranke}
    \proofstepwidestarbreakkeep[1]{\varnothing\neq\mathcal A\subseteq\mathbb R}{\rA}
    \proofstepwidestarbreakkeep[2]{%
      \exists U\in\mathbb R\,\forall A\in\mathcal A\,(A\RealLe U)}{\rA}
    \proofstepwidestarbreakkeep[2]{%
      U\in\mathbb R\land\forall A\in\mathcal A\,(A\RealLe U)}{\rA}
    \proofstepwidestarbreakkeep[1,2]{\RealCut{\bigcup\mathcal A}}{%
      \FormulaRefAuto{RealBoundedCutUnionIsCut}[thm]{1,3}}
    \proofstepwidestarbreakkeep[1,2]{\bigcup\mathcal A\in\mathbb R}{%
      \FormulaRefAuto{RealCutMembership}[thm]{4}}
    \proofstepwidestarbreakkeep[1]{%
      \forall A\in\mathcal A\,(A\RealLe\bigcup\mathcal A)}{%
      \FormulaRefAuto{RealCutUnionUpperBound}[thm]{1}}
    \proofstepwidestarbreakkeep[1]{%
      \forall V\in\mathbb R\,
      (\forall A\in\mathcal A\,(A\RealLe V)
       \rightarrow\bigcup\mathcal A\RealLe V)}{%
      \FormulaRefAuto{RealCutUnionLeastUpperBound}[thm]{1}}
    \proofstepwidestarbreakkeep[1,2]{%
      \exists S\in\mathbb R\,
      (\forall A\in\mathcal A\,(A\RealLe S)\land
       \forall V\in\mathbb R\,
       (\forall A\in\mathcal A\,(A\RealLe V)\rightarrow S\RealLe V))}{%
      \rEI{5,\rAI{6,7}}}
    \proofstepwidestarbreakkeep[1,2]{%
      \exists S\in\mathbb R\,
      (\forall A\in\mathcal A\,(A\RealLe S)\land
       \forall V\in\mathbb R\,
       (\forall A\in\mathcal A\,(A\RealLe V)\rightarrow S\RealLe V))}{%
      \rEE{2,3,8}}
  \closeproofpartwide
\end{tabproofsplitwide}

Erst die vorstehende Existenz- und Kleinste-Obergrenzen-Aussage erfüllt die
Definitionsprämisse des allgemeinen Supremumsbegriffs aus Band~13. Daher ist
der folgende Supremumsterm wohldefiniert.

\FormulaThmDeltaK[Supremum einer beschränkten Schnittfamilie]{%
  \begin{aligned}[t]
    &\varnothing\neq\mathcal A\subseteq\mathbb R\dsep
      U\in\mathbb R\dsep
      \forall A\in\mathcal A\,(A\RealLe U)\\[-2pt]
    &\hspace{20mm}\vdash
      \sup_{\mathbb R,\RealLe}(\mathcal A)=\bigcup\mathcal A.
  \end{aligned}%
}{RealCutSupremumUnion}{%
  \DeltaRow{Mengen}{\mathcal A\dsep A\dsep U}%
  \DeltaRow{Relationsoperatoren}{\RealLe}%
}
\begin{tabproofwide}
  \proofstepwidestarbreakkeep[1]{\varnothing\neq\mathcal A\subseteq\mathbb R}{\rA}
  \proofstepwidestarbreakkeep[2]{U\in\mathbb R}{\rA}
  \proofstepwidestarbreakkeep[3]{\forall A\in\mathcal A\,(A\RealLe U)}{\rA}
  \proofstepwidestar[1,2,3]{\RealCut { \bigcup \mathcal A }}{%
      \FormulaRefAuto{RealBoundedCutUnionIsCut}[thm]{1,2,3}}
  \proofstepwidestarbreakkeep[1,2,3]{\bigcup\mathcal A\in\mathbb R}{%
    \FormulaRefAuto{RealCutMembership}[thm]{%
      4}}
  \proofstepwidestarbreakkeep[1]{%
    \forall A\in\mathcal A\,(A\RealLe\bigcup\mathcal A)}{%
    \FormulaRefAuto{RealCutUnionUpperBound}[thm]{1}}
  \proofstepwidestarbreakkeep[1]{%
    \forall V\in\mathbb R\,
    (\forall A\in\mathcal A\,(A\RealLe V)
      \rightarrow\bigcup\mathcal A\RealLe V)}{%
    \FormulaRefAuto{RealCutUnionLeastUpperBound}[thm]{1}}
  \proofstepwidestarbreakkeep[1,2,3]{%
    \bigcup\mathcal A\in\UB_{\mathbb R,\RealLe}(\mathcal A)}{%
    \FormulaRefAuto{UpperBoundSetIntro}[thm]{5,6}}
  \proofstepwidestarbreakkeep[1,2,3]{%
    \forall V\in\UB_{\mathbb R,\RealLe}(\mathcal A)\,
      (\bigcup\mathcal A\RealLe V)}{%
    \FormulaRefAuto{UpperBoundSetElim}[thm]{7}}
  \proofstepwidestarbreakkeep[1,2,3]{%
    \bigcup\mathcal A=\sup_{\mathbb R,\RealLe}(\mathcal A)}{%
    \FormulaRefAuto{SupremumCandidateEqualsSupremum}[thm]{1,8,9}}
  \proofstepwidestarbreakkeep[1,2,3]{%
    \sup_{\mathbb R,\RealLe}(\mathcal A)=\bigcup\mathcal A}{%
    \FormulaRefAuto{a=b\vdash b=a}{10}}
\end{tabproofwide}

\FormulaThmDeltaK[Infimumsvollständigkeit der reellen Ordnung]{%
  \begin{aligned}[t]
    &\varnothing\neq\mathcal A\subseteq\mathbb R\dsep
      \exists L\in\mathbb R\,
        \forall A\in\mathcal A\,(L\RealLe A)\\[-2pt]
    &\quad\vdash \exists I\in\mathbb R\,
      \left(
        \forall A\in\mathcal A\,(I\RealLe A)\land
        \forall V\in\mathbb R\,
          \bigl(\forall A\in\mathcal A\,(V\RealLe A)
          \rightarrow V\RealLe I\bigr)
      \right).
  \end{aligned}%
}{RealGreatestLowerBoundProperty}{%
  \DeltaRow{Mengen}{\mathcal A\dsep A\dsep L\dsep V}%
  \DeltaRow{Relationsoperatoren}{\RealLe}%
}
\begin{tabproofsplitwide}
\proofpartwide{Nichtleerheit der Schrankenmenge}
\proofstepwidestarbreakkeep[1]{\varnothing\neq\mathcal A\subseteq\mathbb R}{\rA}
    \proofstepwidestarbreakkeep[2]{L\in\mathbb R}{\rA}
    \proofstepwidestarbreakkeep[3]{\forall A\in\mathcal A\,(L\RealLe A)}{\rA}
    \proofstepwidestarbreakkeep[1,2,3]{%
      L\in\LB_{\mathbb R,\RealLe}(\mathcal A)}{%
      \FormulaRefAuto{LowerBoundSetIntro}[thm]{1,2,3}}
    \proofstepwidestarbreakkeep[1,2,3]{%
      \LB_{\mathbb R,\RealLe}(\mathcal A)\neq\varnothing}{%
      \FormulaRefAuto{x\in A\vdash A\neq\varnothing}{4}}

\closeproofpartwide

\proofpartwide{Trägerinklusion}
  \setcounter{proofstepnr}{5}% Fortlaufende Schritte dieses Beweises.
\proofstepwidestarbreakkeep[1]{%
      \LB_{\mathbb R,\RealLe}(\mathcal A)\subseteq\mathbb R}{%
      \FormulaRefAuto{LowerBoundSetSubsetCarrier}[thm]{1}}

\closeproofpartwide

\proofpartwide{Obere Beschränktheit}
  \setcounter{proofstepnr}{6}% Fortlaufende Schritte dieses Beweises.
\proofstepwidestarbreakkeep[1]{\exists A_0\in\mathcal A}{%
      \FormulaRefAuto{A\neq\varnothing\eqvdash\exists x\,(x\in A)}{1}}
    \proofstepwidestarbreakkeep[1]{A_0\in\mathcal A}{\rA}
    \proofstepwidestarbreakkeep[1]{A_0\in\mathbb R}{%
      \FormulaRefAuto{A\subseteq B\dsep x\in A\vdash x\in B}{1,8}}
    \proofstepwidestarbreakkeep[10]{V\in\LB_{\mathbb R,\RealLe}(\mathcal A)}{\rA}
    \proofstepwidestarbreakkeep[1,10]{V\RealLe A_0}{%
      \FormulaRefAuto{LowerBoundSetElim}[thm]{1,10,8}}
    \proofstepwidestarbreakkeep[1]{%
      \forall V\in\LB_{\mathbb R,\RealLe}(\mathcal A)\,(V\RealLe A_0)}{%
      \rUI{\rRI{10,11}}}
    \proofstepwidestarbreakkeep[1]{%
      \exists A_0\in\mathbb R\,
      \forall V\in\LB_{\mathbb R,\RealLe}(\mathcal A)\,(V\RealLe A_0)}{%
      \rEI{9,12}}
    \proofstepwidestarbreakkeep[1]{%
      \exists A_0\in\mathbb R\,
      \forall V\in\LB_{\mathbb R,\RealLe}(\mathcal A)\,(V\RealLe A_0)}{%
      \rEE{7,8,13}}

\closeproofpartwide

\proofpartwideindR[Schrankenmenge ist nichtleer und beschränkt]{%
    \begin{aligned}[t]
      &\varnothing\neq\mathcal A\subseteq\mathbb R\dsep
        L\in\mathbb R\dsep\forall A\in\mathcal A\,(L\RealLe A)\\[-2pt]
      &\quad\vdash
        \varnothing\neq\LB_{\mathbb R,\RealLe}(\mathcal A)
        \subseteq\mathbb R\land{}\\[-2pt]
      &\qquad\exists A_0\in\mathbb R\,
        \forall V\in\LB_{\mathbb R,\RealLe}(\mathcal A)\,
          (V\RealLe A_0)
    \end{aligned}
  }{%
    \varnothing\neq\mathcal A\subseteq\mathbb R\dsep
    L\in\mathbb R\dsep\forall A\in\mathcal A\,(L\RealLe A)
    \vdash
    \varnothing\neq\LB_{\mathbb R,\RealLe}(\mathcal A)\subseteq\mathbb R
    \land
    \exists A_0\in\mathbb R\,
    \forall V\in\LB_{\mathbb R,\RealLe}(\mathcal A)\,(V\RealLe A_0)
  }[RealLowerBoundFamilyBounded]
  \setcounter{proofstepnr}{14}% Fortlaufende Schritte dieses Beweises.
\proofstepwidestarbreakkeep[1,2,3]{%
      \varnothing\neq\LB_{\mathbb R,\RealLe}(\mathcal A)\subseteq\mathbb R
      \land
      \exists A_0\in\mathbb R\,
      \forall V\in\LB_{\mathbb R,\RealLe}(\mathcal A)\,(V\RealLe A_0)}{%
      \rAI{\rAI{5,6},14}}

\closeproofpartwideindR


  \proofpartwide{Supremum der unteren Schranken}
    \proofstepwidestarbreakkeep[1]{\varnothing\neq\mathcal A\subseteq\mathbb R}{\rA}
    \proofstepwidestarbreakkeep[2]{%
      \exists L\in\mathbb R\,\forall A\in\mathcal A\,(L\RealLe A)}{\rA}
    \proofstepwidestarbreakkeep[2]{L\in\mathbb R\land
      \forall A\in\mathcal A\,(L\RealLe A)}{\rA}
    \proofstepwidestarbreakkeep[1,2]{%
      \begin{aligned}[t]
        &\varnothing\neq\LB_{\mathbb R,\RealLe}(\mathcal A)\subseteq\mathbb R
          \land{}\\[-2pt]
        &\exists A_0\in\mathbb R\,
          \forall V\in\LB_{\mathbb R,\RealLe}(\mathcal A)\,(V\RealLe A_0)
      \end{aligned}}{%
      \FormulaRefAuto{RealLowerBoundFamilyBounded}[thm]{1,3}}
    \proofstepwidestarbreakkeep[1,2]{%
      \begin{aligned}[t]
        &\exists I\in\mathbb R\,\bigl(\\[-2pt]
        &\quad \forall V\in\LB_{\mathbb R,\RealLe}(\mathcal A)\,(V\RealLe I)
          \land{}\\[-2pt]
        &\quad \forall W\in\mathbb R\,\bigl(\\[-2pt]
        &\qquad \forall V\in\LB_{\mathbb R,\RealLe}(\mathcal A)\,(V\RealLe W)
          \rightarrow I\RealLe W\bigr)\bigr)
      \end{aligned}}{%
      \FormulaRefAuto{RealLeastUpperBoundProperty}[thm]{4}}
    \proofstepwidestarbreakkeep[1,2]{%
      \begin{aligned}[t]
        &I\in\mathbb R\land{}\\[-2pt]
        &\forall V\in\LB_{\mathbb R,\RealLe}(\mathcal A)\,(V\RealLe I)
          \land{}\\[-2pt]
        &\forall W\in\mathbb R\,\bigl(\\[-2pt]
        &\quad \forall V\in\LB_{\mathbb R,\RealLe}(\mathcal A)\,(V\RealLe W)
          \rightarrow I\RealLe W\bigr)
      \end{aligned}}{\rA}
    \proofstepwidestarbreakkeep[7]{A\in\mathcal A}{\rA}
    \proofstepwidestarbreakkeep[1,7]{A\in\mathbb R}{%
      \FormulaRefAuto{A\subseteq B\dsep x\in A\vdash x\in B}{1,7}}
    \proofstepwidestarbreakkeep[1,7]{%
      \forall V\in\LB_{\mathbb R,\RealLe}(\mathcal A)\,(V\RealLe A)}{%
      \FormulaRefAuto{LowerBoundSetElim}[thm]{1,7}}
    \proofstepwidestarbreakkeep[1,2,7]{I\RealLe A}{\rUE{6,8,9}}
    \proofstepwidestarbreakkeep[1,2]{\forall A\in\mathcal A\,(I\RealLe A)}{%
      \rUI{\rRI{7,10}}}
    \proofstepwidestarbreakkeep[12]{V\in\mathbb R}{\rA}
    \proofstepwidestarbreakkeep[13]{\forall A\in\mathcal A\,(V\RealLe A)}{\rA}
    \proofstepwidestarbreakkeep[1,12,13]{%
      V\in\LB_{\mathbb R,\RealLe}(\mathcal A)}{%
      \FormulaRefAuto{LowerBoundSetIntro}[thm]{1,12,13}}
    \proofstepwidestarbreakkeep[1,2,12,13]{V\RealLe I}{\rUE{6,14}}
    \proofstepwidestarbreakkeep[1,2]{%
      \forall V\in\mathbb R\,
      (\forall A\in\mathcal A\,(V\RealLe A)\rightarrow V\RealLe I)}{%
      \rUI{\rRI{12,\rRI{13,15}}}}
    \proofstepwidestarbreakkeep[1,2]{%
      \begin{aligned}[t]
        &\exists I\in\mathbb R\,\bigl(\\[-2pt]
        &\quad \forall A\in\mathcal A\,(I\RealLe A)\land{}\\[-2pt]
        &\quad \forall V\in\mathbb R\,\bigl(
          \forall A\in\mathcal A\,(V\RealLe A)\\[-2pt]
        &\hspace{38mm}{}\rightarrow V\RealLe I\bigr)\bigr)
      \end{aligned}}{%
      \rEI{6,\rAI{11,16}}}
    \proofstepwidestarbreakkeep[1,2]{%
      \begin{aligned}[t]
        &\exists I\in\mathbb R\,\bigl(\\[-2pt]
        &\quad \forall A\in\mathcal A\,(I\RealLe A)\land{}\\[-2pt]
        &\quad \forall V\in\mathbb R\,\bigl(
          \forall A\in\mathcal A\,(V\RealLe A)\\[-2pt]
        &\hspace{38mm}{}\rightarrow V\RealLe I\bigr)\bigr)
      \end{aligned}}{%
      \rEE{5,6,17}}
    \proofstepwidestarbreakkeep[1,2]{%
      \begin{aligned}[t]
        &\exists I\in\mathbb R\,\bigl(\\[-2pt]
        &\quad \forall A\in\mathcal A\,(I\RealLe A)\land{}\\[-2pt]
        &\quad \forall V\in\mathbb R\,\bigl(
          \forall A\in\mathcal A\,(V\RealLe A)\\[-2pt]
        &\hspace{38mm}{}\rightarrow V\RealLe I\bigr)\bigr)
      \end{aligned}}{%
      \rEE{2,3,18}}
  \closeproofpartwide
\end{tabproofsplitwide}

Auch hier wird der allgemeine Infimumsterm erst nach diesem Existenzsatz
eingeführt.

\FormulaThmDeltaK[Charakterisierung des reellen Infimums]{%
  \begin{aligned}[t]
    &\varnothing\neq\mathcal A\subseteq\mathbb R\dsep
      \exists L\in\mathbb R\,
        \forall A\in\mathcal A\,(L\RealLe A)\\[-2pt]
    &\quad\vdash
      \inf_{\mathbb R,\RealLe}(\mathcal A)\in\mathbb R\land{}\\[-2pt]
    &\qquad
      \forall A\in\mathcal A\,
        \bigl(\inf_{\mathbb R,\RealLe}(\mathcal A)\RealLe A\bigr)\land{}\\[-2pt]
    &\qquad
      \forall V\in\mathbb R\,
        \bigl(\forall A\in\mathcal A\,(V\RealLe A)
        \rightarrow
        V\RealLe\inf_{\mathbb R,\RealLe}(\mathcal A)\bigr).
  \end{aligned}%
}{RealInfimumCharacterization}{%
  \DeltaRow{Mengen}{\mathcal A\dsep A\dsep L\dsep V}%
  \DeltaRow{Relationsoperatoren}{\RealLe}%
}
\begin{tabproofsplitwide}
  \proofpartwideindR[Existenz eines größten Elements der Schrankenmenge]{%
    \begin{aligned}[t]
      &\varnothing\neq\mathcal A\subseteq\mathbb R\dsep
        \exists L\in\mathbb R\,\forall A\in\mathcal A\,(L\RealLe A)\\[-2pt]
      &\quad\vdash
        \exists I\in\LB_{\mathbb R,\RealLe}(\mathcal A)\,
        \forall V\in\LB_{\mathbb R,\RealLe}(\mathcal A)\,(V\RealLe I)
    \end{aligned}
  }{%
    \varnothing\neq\mathcal A\subseteq\mathbb R\dsep
    \exists L\in\mathbb R\,\forall A\in\mathcal A\,(L\RealLe A)
    \vdash
    \exists I\in\LB_{\mathbb R,\RealLe}(\mathcal A)\,
    \forall V\in\LB_{\mathbb R,\RealLe}(\mathcal A)\,(V\RealLe I)
  }[RealLowerBoundSetHasGreatest]
    \proofstepwidestarbreakkeep[1]{\varnothing\neq\mathcal A\subseteq\mathbb R}{\rA}
    \proofstepwidestarbreakkeep[2]{%
      \exists L\in\mathbb R\,\forall A\in\mathcal A\,(L\RealLe A)}{\rA}
    \proofstepwidestarbreakkeep[1,2]{%
      \exists I\in\mathbb R\,
      (\forall A\in\mathcal A\,(I\RealLe A)\land
       \forall V\in\mathbb R\,
       (\forall A\in\mathcal A\,(V\RealLe A)\rightarrow V\RealLe I))}{%
      \FormulaRefAuto{RealGreatestLowerBoundProperty}[thm]{1,2}}
    \proofstepwidestarbreakkeep[1,2]{%
      I\in\mathbb R\land
      \forall A\in\mathcal A\,(I\RealLe A)\land
      \forall V\in\mathbb R\,
      (\forall A\in\mathcal A\,(V\RealLe A)\rightarrow V\RealLe I)}{\rA}
    \proofstepwidestarbreakkeep[1,2]{%
      I\in\LB_{\mathbb R,\RealLe}(\mathcal A)}{%
      \FormulaRefAuto{LowerBoundSetIntro}[thm]{1,4}}
    \proofstepwidestarbreakkeep[6]{V\in\LB_{\mathbb R,\RealLe}(\mathcal A)}{\rA}
    \proofstepwidestarbreakkeep[1,6]{V\in\mathbb R\land
      \forall A\in\mathcal A\,(V\RealLe A)}{%
      \FormulaRefAuto{LowerBoundSetCharacterization}[thm]{1,6}}
    \proofstepwidestarbreakkeep[1,2,6]{V\RealLe I}{\rUE{4,7}}
    \proofstepwidestarbreakkeep[1,2]{%
      \forall V\in\LB_{\mathbb R,\RealLe}(\mathcal A)\,(V\RealLe I)}{%
      \rUI{\rRI{6,8}}}
    \proofstepwidestarbreakkeep[1,2]{%
      \exists I\in\LB_{\mathbb R,\RealLe}(\mathcal A)\,
      \forall V\in\LB_{\mathbb R,\RealLe}(\mathcal A)\,(V\RealLe I)}{%
      \rEI{5,9}}
    \proofstepwidestarbreakkeep[1,2]{%
      \exists I\in\LB_{\mathbb R,\RealLe}(\mathcal A)\,
      \forall V\in\LB_{\mathbb R,\RealLe}(\mathcal A)\,(V\RealLe I)}{%
      \rEE{3,4,10}}
  \closeproofpartwideindR

  \proofpartwide{Eigenschaften des definierten Infimums}
    \proofstepwidestarbreakkeep[1]{\varnothing\neq\mathcal A\subseteq\mathbb R}{\rA}
    \proofstepwidestarbreakkeep[2]{%
      \exists L\in\mathbb R\,\forall A\in\mathcal A\,(L\RealLe A)}{\rA}
    \proofstepwidestarbreakkeep[1,2]{%
      \exists I\in\LB_{\mathbb R,\RealLe}(\mathcal A)\,
      \forall V\in\LB_{\mathbb R,\RealLe}(\mathcal A)\,(V\RealLe I)}{%
      \FormulaRefAuto{RealLowerBoundSetHasGreatest}[thm]{1,2}}
    \proofstepwidestarbreakkeep[1,2]{%
      \inf_{\mathbb R,\RealLe}(\mathcal A)\in\mathbb R}{%
      \FormulaRefAuto{InfimumInCarrier}[thm]{1,3}}
    \proofstepwidestarbreakkeep[1,2]{%
      \forall A\in\mathcal A\,
      (\inf_{\mathbb R,\RealLe}(\mathcal A)\RealLe A)}{%
      \FormulaRefAuto{InfimumBelowMember}[thm]{1,3}}
    \proofstepwidestarbreakkeep[1,2]{%
      \forall V\in\mathbb R\,
      (\forall A\in\mathcal A\,(V\RealLe A)
        \rightarrow V\RealLe\inf_{\mathbb R,\RealLe}(\mathcal A))}{%
      \FormulaRefAuto{InfimumGreatestLowerBound}[thm]{1,3}}
    \proofstepwidestarbreakkeep[1,2]{%
      \begin{aligned}[t]
        &\inf_{\mathbb R,\RealLe}(\mathcal A)\in\mathbb R\land{}\\[-2pt]
        &\forall A\in\mathcal A\,
          (\inf_{\mathbb R,\RealLe}(\mathcal A)\RealLe A)\land{}\\[-2pt]
        &\forall V\in\mathbb R\,\bigl(\forall A\in\mathcal A\,(V\RealLe A)\\[-2pt]
        &\hspace{30mm}{}\rightarrow
          V\RealLe\inf_{\mathbb R,\RealLe}(\mathcal A)\bigr)
      \end{aligned}}{%
      \rAI{4,\rAI{5,6}}}
  \closeproofpartwide
\end{tabproofsplitwide}

\chapter{Arithmetik der Schnitte}

\section{Addition und additive Inverse}

\FormulaDefDeltaK[Addition reeller Schnitte]{%
  A,B\in\mathbb R\vdash
  A\RealAdd B\coloneqq
  \bigl\{q\in\mathbb Q\bigm|
    \exists a\in A\,\exists b\in B\,(q<a+b)\bigr\}%
}{RealCutAddition}{%
  \DeltaRow{Mengen}{A\dsep B\dsep a\dsep b\dsep q}%
  \DeltaRow{Funktionsoperatoren}{\RealAdd}%
  \DeltaRow{\textbf{Neues Symbol}}{}%
  \DeltaRow{Reelle Addition}{\RealAdd}%
}

\FormulaDefDeltaK[Additive Null der reellen Zahlen]{%
  \RealZero\coloneqq\RealHat{\RatZero}%
}{RealCutZero}{%
  \DeltaRow{Mengen}{\RealZero}%
  \DeltaRow{\textbf{Neues Symbol}}{}%
  \DeltaRow{Reelle Null}{\RealZero}%
}

\FormulaDefDeltaK[Additives Inverses eines reellen Schnittes]{%
  A\in\mathbb R\vdash
  \RealNeg{A}\coloneqq
  \bigl\{q\in\mathbb Q\bigm|
    \exists r\in\mathbb Q\,(r\notin A\land q<-r)\bigr\}%
}{RealCutNegation}{%
  \DeltaRow{Mengen}{A\dsep q\dsep r}%
  \DeltaRow{Funktionsoperatoren}{\RealNeg{\,\cdot\,}}%
  \DeltaRow{\textbf{Neues Symbol}}{}%
  \DeltaRow{Reelle Negation}{\RealNeg{\,\cdot\,}}%
}

\FormulaThmDeltaK[Abgeschlossenheit der additiven Schnittoperationen]{%
  A,B\in\mathbb R\vdash
  A\RealAdd B\in\mathbb R\land
  \RealNeg{A}\in\mathbb R\land
  \RealZero\in\mathbb R%
}{RealCutAdditiveClosure}{%
  \DeltaRow{Mengen}{A\dsep B}%
  \DeltaRow{Funktionsoperatoren}{\RealAdd\dsep\RealNeg{\,\cdot\,}}%
}
\begin{tabproofsplitwide}
  \proofpartwideindR[Summe zweier Schnitte]{%
    A,B\in\mathbb R\vdash\RealCut{A\RealAdd B}%
  }{%
    A,B\in\mathbb R\vdash\RealCut{A\RealAdd B}%
  }[RealCutAdditionIsCut]
    \proofstepwidestar[1]{A,B\in\mathbb R}{\rA}
    \proofstepwidestar[1]{\RealCut A\land\RealCut B}{%
      \FormulaRefAuto{RealCutMembership}[thm]{1}}
    \proofstepwidestarbreakkeep[1]{\exists a_0\in A\land\exists b_0\in B}{%
      \FormulaRefAuto{DedekindCutCharacterization}[thm]{2}}
    \proofstepwidestar[1]{a_0\in A\land b_0\in B}{\rA}
    \proofstepwidestarbreakkeep[1]{%
      a_0+b_0-\RatOne\in\mathbb Q\land
      a_0+b_0-\RatOne<a_0+b_0}{%
      \FormulaRefAuto{RationalOrderedFieldAxioms}[def]{2,4}}
    \proofstepwidestarbreakkeep[1]{a_0+b_0-\RatOne\in A\RealAdd B}{%
      \FormulaRefAuto{RealCutAddition}[def]{4,5}}
    \proofstepwidestar[1]{A\RealAdd B\neq\varnothing}{%
      \FormulaRefAuto{x\in A\vdash A\neq\varnothing}{6}}
    \proofstepwidestar[1]{A\RealAdd B\neq\varnothing}{%
      \rEE{3,4,7}}
    \proofstepwidestarbreakkeep[1]{%
      \exists\alpha\in\mathbb Q\setminus A\land
      \exists\beta\in\mathbb Q\setminus B}{%
      \FormulaRefAuto{DedekindCutCharacterization}[thm]{2}}
    \proofstepwidestarbreakkeep[10]{%
      \alpha\in\mathbb Q\setminus A\land
      \beta\in\mathbb Q\setminus B}{\rA}
    \proofstepwidestarbreakkeep[1,10]{%
      \forall a\in A\,(a\leq\alpha)\land
      \forall b\in B\,(b\leq\beta)}{%
      \FormulaRefAuto{DedekindCutCharacterization}[thm]{2,10}}
    \proofstepwidestar[1]{A\RealAdd B=\{q\in\mathbb Q\mid\exists a\in A\,\exists b\in B\,(q<a+b)\}}{%
      \FormulaRefAuto{RealCutAddition}[def]{1}}
    \proofstepwidestarbreakkeep[1,10]{\alpha+\beta\notin A\RealAdd B}{%
      \FormulaRefAuto{RationalOrderTranslation}[thm]{%
        11,12}}
    \proofstepwidestar[1,10]{\alpha+\beta\in\mathbb Q}{%
      \FormulaRefAuto{RationalAddClosure}[thm]{10}}
    \proofstepwidestarbreakkeep[1]{A\RealAdd B\neq\mathbb Q}{%
      \rEE{9,10,\FormulaRefAuto{x\in B, x\not\in A\vdash A\neq B}{%
        14,13}}}
    \proofstepwidestar[16]{q\in A\RealAdd B}{\rA}
    \proofstepwidestar[17]{s\in\mathbb Q\land s<q}{\rA}
    \proofstepwidestarbreakkeep[16]{%
      \exists a\in A\,\exists b\in B\,(q<a+b)}{%
      \FormulaRefAuto{RealCutAddition}[def]{16}}
    \proofstepwidestarbreakkeep[16]{a\in A\land b\in B\land q<a+b}{\rA}
    \proofstepwidestar[16,17]{s<a+b}{%
      \FormulaRefAuto{RationalOrderTotal}[thm]{17,19}}
    \proofstepwidestarbreakkeep[16,17]{s\in A\RealAdd B}{%
      \FormulaRefAuto{RealCutAddition}[def]{17,19,20}}
    \proofstepwidestar[16,17]{s\in A\RealAdd B}{%
      \rEE{18,19,21}}
    \proofstepwidestarbreakkeep[1]{%
      \forall q\in A\RealAdd B\,\forall s\in\mathbb Q\,
      (s<q\rightarrow s\in A\RealAdd B)}{%
      \rUI{\rRI{16,\rRI{17,22}}}}
    \proofstepwidestar[24]{q\in A\RealAdd B}{\rA}
    \proofstepwidestarbreakkeep[24]{%
      \exists a\in A\,\exists b\in B\,(q<a+b)}{%
      \FormulaRefAuto{RealCutAddition}[def]{24}}
    \proofstepwidestarbreakkeep[24]{a\in A\land b\in B\land q<a+b}{\rA}
    \proofstepwidestarbreakkeep[24]{%
      \exists r\in\mathbb Q\,(q<r\land r<a+b)}{%
      \FormulaRefAuto{RationalOrderDense}[thm]{26}}
    \proofstepwidestarbreakkeep[24]{r\in\mathbb Q\land q<r\land r<a+b}{\rA}
    \proofstepwidestarbreakkeep[24]{r\in A\RealAdd B}{%
      \FormulaRefAuto{RealCutAddition}[def]{26,28}}
    \proofstepwidestarbreakkeep[24]{\exists r\in A\RealAdd B\,(q<r)}{%
      \rEI{28,29}}
    \proofstepwidestar[24]{\exists r\in A\RealAdd B\,(q<r)}{%
      \rEE{27,28,30}}
    \proofstepwidestar[24]{\exists r\in A\RealAdd B\,(q<r)}{%
      \rEE{25,26,31}}
    \proofstepwidestarbreakkeep[1]{%
      \forall q\in A\RealAdd B\,\exists r\in A\RealAdd B\,(q<r)}{%
      \rUI{\rRI{24,32}}}
    \proofstepwidestarbreakkeep[1]{A\RealAdd B\subseteq\mathbb Q}{%
      \FormulaRefAuto{RealCutAddition}[def]{1}}
    \proofstepwidestarbreakkeep[1]{\RealCut{A\RealAdd B}}{%
      \FormulaRefAuto{DedekindCutCharacterization}[thm]{34,8,15,23,33}}
  \closeproofpartwideindR

  \proofpartwideindR[Negation eines Schnittes]{%
    A\in\mathbb R\vdash\RealCut{\RealNeg A}%
  }{%
    A\in\mathbb R\vdash\RealCut{\RealNeg A}%
  }[RealCutNegationIsCut]
    \proofstepwidestar[1]{A\in\mathbb R}{\rA}
    \proofstepwidestar[1]{\RealCut A}{%
      \FormulaRefAuto{RealCutMembership}[thm]{1}}
    \proofstepwidestar[1]{\exists r_0\in\mathbb Q\setminus A}{%
      \FormulaRefAuto{DedekindCutCharacterization}[thm]{2}}
    \proofstepwidestar[1]{r_0\in\mathbb Q\land r_0\notin A}{\rA}
    \proofstepwidestar[1]{-r_0-\RatOne<-r_0}{%
      \FormulaRefAuto{RationalOrderedFieldAxioms}[def]{4}}
    \proofstepwidestar[1]{-r_0-\RatOne\in\RealNeg A}{%
      \FormulaRefAuto{RealCutNegation}[def]{4,5}}
    \proofstepwidestar[1]{\RealNeg A\neq\varnothing}{%
      \FormulaRefAuto{x\in A\vdash A\neq\varnothing}{6}}
    \proofstepwidestar[1]{\RealNeg A\neq\varnothing}{\rEE{3,4,7}}
    \proofstepwidestarbreakkeep[1]{\exists a_0\in A}{%
      \FormulaRefAuto{DedekindCutCharacterization}[thm]{2}}
    \proofstepwidestar[10]{a_0\in A}{\rA}
    \proofstepwidestar[1,10]{-a_0\notin\RealNeg A}{%
      \FormulaRefAuto{RealCutNegation}[def]{2,10}}
    \proofstepwidestar[1,10]{-a_0\in\mathbb Q}{%
      \FormulaRefAuto{RationalNegationLaws}[thm]{10}}
    \proofstepwidestarbreakkeep[1,10]{\RealNeg A\neq\mathbb Q}{%
      \FormulaRefAuto{x\in B, x\not\in A\vdash A\neq B}{%
        12,11}}
    \proofstepwidestarbreakkeep[1]{\RealNeg A\neq\mathbb Q}{\rEE{9,10,13}}
    \proofstepwidestar[15]{q\in\RealNeg A}{\rA}
    \proofstepwidestar[16]{s\in\mathbb Q\land s<q}{\rA}
    \proofstepwidestarbreakkeep[15]{%
      \exists r\in\mathbb Q\,(r\notin A\land q<-r)}{%
      \FormulaRefAuto{RealCutNegation}[def]{15}}
    \proofstepwidestar[15]{r\in\mathbb Q\land r\notin A\land q<-r}{\rA}
    \proofstepwidestar[15,16]{s<-r}{%
      \FormulaRefAuto{RationalOrderTotal}[thm]{16,18}}
    \proofstepwidestar[15,16]{s\in\RealNeg A}{%
      \FormulaRefAuto{RealCutNegation}[def]{16,18,19}}
    \proofstepwidestar[15,16]{s\in\RealNeg A}{\rEE{17,18,20}}
    \proofstepwidestarbreakkeep[1]{%
      \forall q\in\RealNeg A\,\forall s\in\mathbb Q\,
      (s<q\rightarrow s\in\RealNeg A)}{%
      \rUI{\rRI{15,\rRI{16,21}}}}
    \proofstepwidestar[23]{q\in\RealNeg A}{\rA}
    \proofstepwidestarbreakkeep[23]{%
      \exists r\in\mathbb Q\,(r\notin A\land q<-r)}{%
      \FormulaRefAuto{RealCutNegation}[def]{23}}
    \proofstepwidestar[23]{r\in\mathbb Q\land r\notin A\land q<-r}{\rA}
    \proofstepwidestarbreakkeep[23]{%
      \exists s\in\mathbb Q\,(q<s\land s<-r)}{%
      \FormulaRefAuto{RationalOrderDense}[thm]{25}}
    \proofstepwidestar[23]{s\in\mathbb Q\land q<s\land s<-r}{\rA}
    \proofstepwidestar[23]{s\in\RealNeg A}{%
      \FormulaRefAuto{RealCutNegation}[def]{25,27}}
    \proofstepwidestar[23]{\exists s\in\RealNeg A\,(q<s)}{\rEI{27,28}}
    \proofstepwidestar[23]{\exists s\in\RealNeg A\,(q<s)}{\rEE{26,27,29}}
    \proofstepwidestar[23]{\exists s\in\RealNeg A\,(q<s)}{\rEE{24,25,30}}
    \proofstepwidestarbreakkeep[1]{%
      \forall q\in\RealNeg A\,\exists s\in\RealNeg A\,(q<s)}{%
      \rUI{\rRI{23,31}}}
    \proofstepwidestar[1]{\RealNeg A\subseteq\mathbb Q}{%
      \FormulaRefAuto{RealCutNegation}[def]{1}}
    \proofstepwidestar[1]{\RealCut{\RealNeg A}}{%
      \FormulaRefAuto{DedekindCutCharacterization}[thm]{33,8,14,22,32}}
  \closeproofpartwideindR

  \proofpartwide{Nullschnitt und Zusammenfassung}
    \proofstepwidestar[1]{A,B\in\mathbb R}{\rA}
    \proofstepwidestar[1]{\RealCut{A\RealAdd B}}{%
      \FormulaRefAuto{RealCutAdditionIsCut}[thm]{1}}
    \proofstepwidestar[1]{A\RealAdd B\in\mathbb R}{%
      \FormulaRefAuto{RealCutMembership}[thm]{%
        2}}
    \proofstepwidestar[1]{\RealCut{\RealNeg A}}{%
      \FormulaRefAuto{RealCutNegationIsCut}[thm]{1}}
    \proofstepwidestar[1]{\RealNeg A\in\mathbb R}{%
      \FormulaRefAuto{RealCutMembership}[thm]{%
        4}}
    \proofstepwidestar[]{\RatZero\in\mathbb Q}{%
      \FormulaRefAuto{RationalZeroOneCarrier}[thm]}
    \proofstepwidestar[]{\RealCut{\RealHat{\RatZero}}}{%
      \FormulaRefAuto{RealPrincipalCutIsCut}[thm]{6}}
    \proofstepwidestar[]{\RealHat{\RatZero}\in\mathbb R}{%
      \FormulaRefAuto{RealCutMembership}[thm]{7}}
    \proofstepwidestar[]{\RealZero\in\mathbb R}{%
      \FormulaRefAuto{RealCutZero}[def]{%
        8}}
    \proofstepwidestar[1]{%
      A\RealAdd B\in\mathbb R\land
      \RealNeg A\in\mathbb R\land\RealZero\in\mathbb R}{%
      \rAI{3,\rAI{5,9}}}
  \closeproofpartwide
\end{tabproofsplitwide}

\FormulaThmDeltaK[Rationale Randapproximation eines Schnittes]{%
  \begin{aligned}[t]
  &A\in\mathbb R\dsep h\in\mathbb Q\dsep\RatZero<h\vdash{}\\[-2pt]
  &\quad\exists a\in A\,\exists r\in\mathbb Q\,
    \bigl(r\notin A\land a<r\land r-a<h\bigr).
  \end{aligned}%
}{RealCutAdditiveBoundaryApproximation}{%
  \DeltaRow{Mengen}{A\dsep a\dsep r\dsep h}%
}
\begin{tabproofsplitwide}
  \proofpartwideindR[Ein Gitterpunkt liegt außerhalb des Schnittes]{%
    \begin{aligned}[t]
      &A\in\mathbb R\dsep h\in\mathbb Q\dsep\RatZero<h\dsep
        p\in A\dsep u\in\mathbb Q\setminus A\\[-2pt]
      &\quad\vdash\exists m\in\mathbb N\,
        \left(u<p+\RatEmbed(\IntEmbed(m))\frac h2\right)
    \end{aligned}
  }{%
    A\in\mathbb R\dsep h\in\mathbb Q\dsep\RatZero<h\dsep
    p\in A\dsep u\in\mathbb Q\setminus A
    \vdash\exists m\in\mathbb N\,
    \left(u<p+\RatEmbed(\IntEmbed(m))\frac h2\right)
  }[RealCutAdditiveBoundaryGridEscape]
    \proofstepwidestarbreakkeep[1]{A\in\mathbb R}{\rA}
    \proofstepwidestarbreakkeep[2]{h\in\mathbb Q}{\rA}
    \proofstepwidestarbreakkeep[3]{\RatZero<h}{\rA}
    \proofstepwidestarbreakkeep[4]{p\in A}{\rA}
    \proofstepwidestarbreakkeep[5]{u\in\mathbb Q\setminus A}{\rA}
    \proofstepwidestar[1,4]{\RealCut A}{%
      \FormulaRefAuto{RealCutMembership}[thm]{1}}
    \proofstepwidestarbreakkeep[1,4]{p\in\mathbb Q}{%
      \FormulaRefAuto{DedekindCutCharacterization}[thm]{%
        6,4}}
    \proofstepwidestarbreakkeep[2,3]{%
      \frac h2\in\mathbb Q\land\RatZero<\frac h2}{%
      \FormulaRefAuto{RationalOrderedFieldAxioms}[def]{2,3}}
    \proofstepwidestarbreakkeep[1,2,3,4,5]{%
      (u-p)\left(\frac h2\right)^{-1}\in\mathbb Q}{%
      \FormulaRefAuto{RationalOrderedFieldAxioms}[def]{2,5,7,8}}
    \proofstepwidestarbreakkeep[1,2,3,4,5]{%
      \exists m\in\mathbb N\,
      (u-p)\left(\frac h2\right)^{-1}
      <\RatEmbed(\IntEmbed(m))}{%
      \FormulaRefAuto{RationalArchimedean}[thm]{9}}
    \proofstepwidestarbreakkeep[1,2,3,4,5]{m\in\mathbb N\land
      (u-p)(h/2)^{-1}<\RatEmbed(\IntEmbed(m))}{\rA}
    \proofstepwidestarbreakkeep[1,2,3,4,5]{%
      u<p+\RatEmbed(\IntEmbed(m))\frac h2}{%
      \FormulaRefAuto{RationalPositiveMulStrictOrder}[thm]{8,11}}
    \proofstepwidestarbreakkeep[1,2,3,4,5]{%
      \exists m\in\mathbb N\,
      \left(u<p+\RatEmbed(\IntEmbed(m))\frac h2\right)}{%
      \rEI{11,12}}
    \proofstepwidestarbreakkeep[1,2,3,4,5]{%
      \exists m\in\mathbb N\,
      \left(u<p+\RatEmbed(\IntEmbed(m))\frac h2\right)}{%
      \rEE{10,11,13}}
  \closeproofpartwideindR

  \proofpartwideindR[Erster äußerer Gitterpunkt]{%
    \begin{aligned}[t]
      &A\in\mathbb R\dsep h\in\mathbb Q\dsep\RatZero<h\dsep
        p\in A\dsep u\in\mathbb Q\setminus A\\[-2pt]
      &\quad\vdash\exists m_0\in\mathbb N\setminus\{0\}\,
        \bigl(
          p+\RatEmbed(\IntEmbed(m_0))\frac h2\notin A\land{}\\[-2pt]
      &\hspace{36mm}
          p+\RatEmbed(\IntEmbed(m_0-1))\frac h2\in A
        \bigr)
    \end{aligned}
  }{%
    A\in\mathbb R\dsep h\in\mathbb Q\dsep\RatZero<h\dsep
    p\in A\dsep u\in\mathbb Q\setminus A
    \vdash\exists m_0\in\mathbb N\setminus\{0\}\,
    \left(
      p+\RatEmbed(\IntEmbed(m_0))\frac h2\notin A\land
      p+\RatEmbed(\IntEmbed(m_0-1))\frac h2\in A
    \right)
  }[RealCutAdditiveBoundaryFirstOutside]
    \proofstepwidestarbreakkeep[1]{A\in\mathbb R}{\rA}
    \proofstepwidestarbreakkeep[2]{h\in\mathbb Q}{\rA}
    \proofstepwidestarbreakkeep[3]{\RatZero<h}{\rA}
    \proofstepwidestarbreakkeep[4]{p\in A}{\rA}
    \proofstepwidestarbreakkeep[5]{u\in\mathbb Q\setminus A}{\rA}
    \proofstepwidestarbreakkeep[1,2,3,4,5]{%
      \exists m\in\mathbb N\,
      \left(u<p+\RatEmbed(\IntEmbed(m))\frac h2\right)}{%
      \FormulaRefAuto{RealCutAdditiveBoundaryGridEscape}[thm]{1,2,3,4,5}}
    \proofstepwidestarbreakkeep[1,2,3,4,5]{%
      M\coloneqq\left\{n\in\mathbb N\mid
        p+\RatEmbed(\IntEmbed(n))\frac h2\notin A\right\}}{\rII}
    \proofstepwidestarbreakkeep[1,2,3,4,5]{M\neq\varnothing}{%
      \FormulaRefAuto{DedekindCutCharacterization}[thm]{5,6,7}}
    \proofstepwidestarbreakkeep[1,2,3,4,5]{%
      \exists m_0\in M\,\forall n\in M\,(m_0\leq n)}{%
      \FormulaRefAuto{NaturalWellOrderingPrinciple}[thm]{7,8}}
    \proofstepwidestarbreakkeep[1,2,3,4,5]{%
      m_0\in M\land\forall n\in M\,(m_0\leq n)}{\rA}
    \proofstepwidestarbreakkeep[1,2,3,4,5]{m_0\neq0}{%
      \FormulaRefAuto{PeanoAddRightZero}[thm]{4,7,10}}
    \proofstepwidestarbreakkeep[1,2,3,4,5]{(m_0-1)\in\mathbb N}{%
      \FormulaRefAuto{PeanoMinusOneNatural}[thm]{10,11}}
    \proofstepwidestarbreakkeep[1,2,3,4,5]{(m_0-1)<m_0}{%
      \FormulaRefAuto{PeanoMinusOneLt}[thm]{10,11}}
    \proofstepwidestarbreakkeep[1,2,3,4,5]{%
      p+\RatEmbed(\IntEmbed(m_0-1))\frac h2\in A}{%
      \FormulaRefAuto{NaturalWellOrderingPrinciple}[thm]{7,10,13}}
    \proofstepwidestarbreakkeep[1,2,3,4,5]{%
      p+\RatEmbed(\IntEmbed(m_0))\frac h2\notin A}{%
      \rAEa{10}}
    \proofstepwidestarbreakkeep[1,2,3,4,5]{m_0\in\mathbb N\setminus\{0\}}{%
      \FormulaRefAuto{x\in A\setminus B\eqvdash x\in A\land x\notin B}{10,11}}
    \proofstepwidestarbreakkeep[1,2,3,4,5]{\begin{aligned}[t]
      &\exists m_0\in\mathbb N\setminus\{0\}\,\bigl(\\[-2pt]
      &\quad p+\RatEmbed(\IntEmbed(m_0))\frac h2\notin A\land{}\\[-2pt]
      &\quad p+\RatEmbed(\IntEmbed(m_0-1))\frac h2\in A\bigr)
      \end{aligned}}{%
      \rEI{16,\rAI{15,14}}}
    \proofstepwidestarbreakkeep[1,2,3,4,5]{\begin{aligned}[t]
      &\exists m_0\in\mathbb N\setminus\{0\}\,\bigl(\\[-2pt]
      &\quad p+\RatEmbed(\IntEmbed(m_0))\frac h2\notin A\land{}\\[-2pt]
      &\quad p+\RatEmbed(\IntEmbed(m_0-1))\frac h2\in A\bigr)
      \end{aligned}}{%
      \rEE{9,10,17}}
  \closeproofpartwideindR

  \proofpartwide{Schluss}
    \proofstepwidestarbreakkeep[1]{A\in\mathbb R}{\rA}
    \proofstepwidestarbreakkeep[2]{h\in\mathbb Q}{\rA}
    \proofstepwidestarbreakkeep[3]{\RatZero<h}{\rA}
    \proofstepwidestar[1]{\RealCut A}{%
      \FormulaRefAuto{RealCutMembership}[thm]{1}}
    \proofstepwidestarbreakkeep[1]{\exists p\in A}{%
      \FormulaRefAuto{DedekindCutCharacterization}[thm]{%
        4}}
    \proofstepwidestar[1]{\RealCut A}{%
      \FormulaRefAuto{RealCutMembership}[thm]{1}}
    \proofstepwidestarbreakkeep[1]{\exists u\in\mathbb Q\setminus A}{%
      \FormulaRefAuto{DedekindCutCharacterization}[thm]{%
        6}}
    \proofstepwidestarbreakkeep[1]{p\in A}{\rA}
    \proofstepwidestarbreakkeep[1]{u\in\mathbb Q\setminus A}{\rA}
    \proofstepwidestarbreakkeep[1,2,3]{\begin{aligned}[t]
      &\exists m_0\in\mathbb N\setminus\{0\}\,\bigl(\\[-2pt]
      &\quad p+\RatEmbed(\IntEmbed(m_0))\frac h2\notin A\land{}\\[-2pt]
      &\quad p+\RatEmbed(\IntEmbed(m_0-1))\frac h2\in A\bigr)
      \end{aligned}}{%
      \FormulaRefAuto{RealCutAdditiveBoundaryFirstOutside}[thm]{1,2,3,8,9}}
    \proofstepwidestarbreakkeep[1,2,3]{%
      m_0\in\mathbb N\setminus\{0\}\land
      p+\RatEmbed(\IntEmbed(m_0))\frac h2\notin A\land
      p+\RatEmbed(\IntEmbed(m_0-1))\frac h2\in A}{\rA}
    \proofstepwidestar[1,2,3]{( { m _ 0 } - 1 ) + 1 = { m _ 0 }}{%
      \FormulaRefAuto{PeanoMinusOneAddOneOnDomain}[thm]{\rAEa{11}}}
    \proofstepwidestar[1,2,3]{{ m _ 0 } - 1 \in \mathbb N}{%
      \FormulaRefAuto{PeanoMinusOneNaturalOnDomain}[thm]{\rAEa{11}}}
    \proofstepwidestar[1,2,3]{\IntEmbed ( { ( m _ 0 - 1 ) } ) + \IntOne = \IntEmbed ( { ( m _ 0 - 1 ) } + 1 )}{%
      \FormulaRefAuto{IntegerSuccessorOnEmbedding}[thm]{%
            13}}
    \proofstepwidestarbreakkeep[1,2,3]{%
      \IntEmbed(m_0)=\IntEmbed(m_0-1)+\IntOne}{%
      \FormulaRefAuto{a=b\vdash b=a}[thm]{%
        \rIE{12,
          14}}}
    \proofstepwidestarbreakkeep[1,2,3]{%
      \RatEmbed(\IntEmbed(m_0))=
      \RatEmbed(\IntEmbed(m_0-1))+\RatOne}{%
      \FormulaRefAuto{IntegerRationalEmbeddingAdditive}[thm]{15}}
    \proofstepwidestarbreakkeep[1,2,3]{%
      a\coloneqq p+\RatEmbed(\IntEmbed(m_0-1))\frac h2\in A}{%
      \rIE{11}}
    \proofstepwidestarbreakkeep[1,2,3]{%
      r\coloneqq p+\RatEmbed(\IntEmbed(m_0))\frac h2\notin A}{%
      \rIE{11}}
    \proofstepwidestarbreakkeep[1,2,3]{r-a=h/2}{%
      \FormulaRefAuto{RationalOrderedFieldAxioms}[def]{16,17,18}}
    \proofstepwidestarbreakkeep[1,2,3]{\RatZero<r-a<h}{%
      \FormulaRefAuto{RationalOrderedFieldAxioms}[def]{2,3,19}}
    \proofstepwidestarbreakkeep[1,2,3]{a<r}{%
      \FormulaRefAuto{RationalOrderTranslation}[thm]{20}}
    \proofstepwidestarbreakkeep[1,2,3]{%
      \exists a\in A\,\exists r\in\mathbb Q\,
      (r\notin A\land a<r\land r-a<h)}{%
      \rEI{17,\rEI{18,\rAI{18,\rAI{21,20}}}}}
    \proofstepwidestarbreakkeep[1,2,3]{%
      \exists a\in A\,\exists r\in\mathbb Q\,
      (r\notin A\land a<r\land r-a<h)}{%
      \rEE{10,11,22}}
    \proofstepwidestarbreakkeep[1,2,3]{%
      \exists a\in A\,\exists r\in\mathbb Q\,
      (r\notin A\land a<r\land r-a<h)}{%
      \rEE{7,9,23}}
    \proofstepwidestarbreakkeep[1,2,3]{%
      \exists a\in A\,\exists r\in\mathbb Q\,
      (r\notin A\land a<r\land r-a<h)}{%
      \rEE{5,8,24}}
  \closeproofpartwide
\end{tabproofsplitwide}

\FormulaThmDeltaKR[Gesetze der reellen Addition]{%
  \begin{aligned}[t]
    A,B,C\in\mathbb R\vdash{}&
    \text{(i)}\quad (A\RealAdd B)\RealAdd C=A\RealAdd(B\RealAdd C),\\[-2pt]
    &\text{(ii)}\quad A\RealAdd B=B\RealAdd A,\\[-2pt]
    &\text{(iii)}\quad A\RealAdd\RealZero=A,\\[-2pt]
    &\text{(iv)}\quad A\RealAdd\RealNeg{A}=\RealZero.
  \end{aligned}%
}{%
  \begin{aligned}[t]
  A,B,C\in\mathbb R\vdash{}&
    (A\RealAdd B)\RealAdd C=A\RealAdd(B\RealAdd C)\dsep{}\\[-2pt]
  &A\RealAdd B=B\RealAdd A\dsep
    A\RealAdd\RealZero=A\dsep
    A\RealAdd\RealNeg{A}=\RealZero.
  \end{aligned}%
}{RealCutAdditiveLaws}{%
  \DeltaRow{Mengen}{A\dsep B\dsep C}%
  \DeltaRow{Funktionsoperatoren}{\RealAdd\dsep\RealNeg{\,\cdot\,}}%
}
\begin{tabproofsplitwide}
  \proofpartwideindR[Assoziativität]{%
    A,B,C\in\mathbb R\vdash
    (A\RealAdd B)\RealAdd C=A\RealAdd(B\RealAdd C)
  }{%
    A,B,C\in\mathbb R\vdash
    (A\RealAdd B)\RealAdd C=A\RealAdd(B\RealAdd C)
  }[RealCutAdditionAssociative]
    \proofstepwidestarbreakkeep[1]{A,B,C\in\mathbb R}{\rA}
    \proofstepwidestarbreakkeep[1]{%
      A\RealAdd B\in\mathbb R\land B\RealAdd C\in\mathbb R}{%
      \FormulaRefAuto{RealCutAdditiveClosure}[thm]{1}}
    \proofstepwidestarbreakkeep[3]{q\in(A\RealAdd B)\RealAdd C}{\rA}
    \proofstepwidestarbreakkeep[1,3]{%
      \exists x\in A\RealAdd B\,\exists c\in C\,(q<x+c)}{%
      \FormulaRefAuto{RealCutAddition}[def]{2,3}}
    \proofstepwidestarbreakkeep[1,3]{%
      \exists a\in A\,\exists b\in B\,\exists c\in C\,(q<(a+b)+c)}{%
      \FormulaRefAuto{RealCutAddition}[def]{4}}
    \proofstepwidestarbreakkeep[1,3]{%
      \exists a\in A\,\exists y\in B\RealAdd C\,(q<a+y)}{%
      \FormulaRefAuto{RationalAddAssociative}[thm]{5}}
    \proofstepwidestarbreakkeep[1,3]{q\in A\RealAdd(B\RealAdd C)}{%
      \FormulaRefAuto{RealCutAddition}[def]{2,6}}
    \proofstepwidestarbreakkeep[1]{%
      (A\RealAdd B)\RealAdd C\subseteq A\RealAdd(B\RealAdd C)}{%
      \FormulaRefAuto{SubsetIntroductionRule}[thm]{[3]\vdots 7}}
    \proofstepwidestarbreakkeep[9]{q\in A\RealAdd(B\RealAdd C)}{\rA}
    \proofstepwidestarbreakkeep[1,9]{%
      \exists a\in A\,\exists b\in B\,\exists c\in C\,(q<a+(b+c))}{%
      \FormulaRefAuto{RealCutAddition}[def]{2,9}}
    \proofstepwidestarbreakkeep[1,9]{%
      \exists x\in A\RealAdd B\,\exists c\in C\,(q<x+c)}{%
      \FormulaRefAuto{RationalAddAssociative}[thm]{10}}
    \proofstepwidestarbreakkeep[1,9]{q\in(A\RealAdd B)\RealAdd C}{%
      \FormulaRefAuto{RealCutAddition}[def]{2,11}}
    \proofstepwidestarbreakkeep[1]{%
      A\RealAdd(B\RealAdd C)\subseteq(A\RealAdd B)\RealAdd C}{%
      \FormulaRefAuto{SubsetIntroductionRule}[thm]{[9]\vdots 12}}
    \proofstepwidestarbreakkeep[1]{%
      (A\RealAdd B)\RealAdd C=A\RealAdd(B\RealAdd C)}{%
      \FormulaRefAuto{A\subseteq B\dsep B\subseteq A\vdash A=B}{8,13}}
  \closeproofpartwideindR

  \proofpartwideindR[Kommutativität]{%
    A,B\in\mathbb R\vdash A\RealAdd B=B\RealAdd A
  }{%
    A,B\in\mathbb R\vdash A\RealAdd B=B\RealAdd A
  }[RealCutAdditionCommutative]
    \proofstepwidestar[1]{A,B\in\mathbb R}{\rA}
    \proofstepwidestar[2]{q\in A\RealAdd B}{\rA}
    \proofstepwidestar[1,2]{%
      \exists a\in A\,\exists b\in B\,(q<a+b)}{%
      \FormulaRefAuto{RealCutAddition}[def]{1,2}}
    \proofstepwidestar[1,2]{%
      \exists b\in B\,\exists a\in A\,(q<b+a)}{%
      \FormulaRefAuto{RationalAddCommutative}[thm]{3}}
    \proofstepwidestar[1,2]{q\in B\RealAdd A}{%
      \FormulaRefAuto{RealCutAddition}[def]{1,4}}
    \proofstepwidestar[1]{A\RealAdd B\subseteq B\RealAdd A}{%
      \FormulaRefAuto{SubsetIntroductionRule}[thm]{[2]\vdots 5}}
    \proofstepwidestar[1]{B\RealAdd A\subseteq A\RealAdd B}{%
      \rIE{6}}
    \proofstepwidestar[1]{A\RealAdd B=B\RealAdd A}{%
      \FormulaRefAuto{A\subseteq B\dsep B\subseteq A\vdash A=B}{6,7}}
  \closeproofpartwideindR

  \proofpartwideindR[Nullgesetz]{%
    A\in\mathbb R\vdash A\RealAdd\RealZero=A
  }{%
    A\in\mathbb R\vdash A\RealAdd\RealZero=A
  }[RealCutAdditionZero]
    \proofstepwidestar[1]{A\in\mathbb R}{\rA}
    \proofstepwidestar[2]{q\in A\RealAdd\RealZero}{\rA}
    \proofstepwidestar[1,2]{\RealZero = \RealHat \RatZero}{%
      \FormulaRefAuto{RealCutZero}[def]}
    \proofstepwidestar[1,2]{%
      \exists a\in A\,\exists r<\RatZero\,(q<a+r)}{%
      \FormulaRefAuto{RealCutAddition}[def]{%
        3,2}}
    \proofstepwidestar[1,2]{q<a}{%
      \FormulaRefAuto{RationalOrderTranslation}[thm]{4}}
    \proofstepwidestar[1,2]{\RealCut A}{%
      \FormulaRefAuto{RealCutMembership}[thm]{1}}
    \proofstepwidestar[1,2]{q\in A}{%
      \FormulaRefAuto{DedekindCutCharacterization}[thm]{%
        6,4,5}}
    \proofstepwidestar[1]{A\RealAdd\RealZero\subseteq A}{%
      \FormulaRefAuto{SubsetIntroductionRule}[thm]{[2]\vdots 7}}
    \proofstepwidestar[9]{q\in A}{\rA}
    \proofstepwidestar[1,9]{\RealCut A}{%
      \FormulaRefAuto{RealCutMembership}[thm]{1}}
    \proofstepwidestar[1,9]{\exists a\in A\,(q<a)}{%
      \FormulaRefAuto{DedekindCutCharacterization}[thm]{%
        10,9}}
    \proofstepwidestar[1,9]{a\in A\land q<a}{\rA}
    \proofstepwidestar[1,9]{q-a<\RatZero}{%
      \FormulaRefAuto{RationalOrderTranslation}[thm]{12}}
    \proofstepwidestar[1,9]{%
      \exists r\in\mathbb Q\,(q-a<r\land r<\RatZero)}{%
      \FormulaRefAuto{RationalOrderDense}[thm]{13}}
    \proofstepwidestar[1,9]{r\in\mathbb Q\land q-a<r\land r<\RatZero}{\rA}
    \proofstepwidestar[1,9]{r\in\RealHat{\RatZero}}{%
      \FormulaRefAuto{RealPrincipalCut}[def]{15}}
    \proofstepwidestar[1,9]{r\in\RealZero}{%
      \FormulaRefAuto{RealCutZero}[def]{%
        16}}
    \proofstepwidestar[1,9]{q<a+r}{%
      \FormulaRefAuto{RationalOrderTranslation}[thm]{15}}
    \proofstepwidestar[1,9]{q\in A\RealAdd\RealZero}{%
      \FormulaRefAuto{RealCutAddition}[def]{12,15,17,18}}
    \proofstepwidestar[1,9]{q\in A\RealAdd\RealZero}{\rEE{14,15,19}}
    \proofstepwidestar[1,9]{q\in A\RealAdd\RealZero}{\rEE{11,12,20}}
    \proofstepwidestar[1]{A\subseteq A\RealAdd\RealZero}{%
      \FormulaRefAuto{SubsetIntroductionRule}[thm]{[9]\vdots 21}}
    \proofstepwidestar[1]{A\RealAdd\RealZero=A}{%
      \FormulaRefAuto{A\subseteq B\dsep B\subseteq A\vdash A=B}{8,22}}
  \closeproofpartwideindR

  \proofpartwideindR[Additives Inverses]{%
    A\in\mathbb R\vdash A\RealAdd\RealNeg A=\RealZero
  }{%
    A\in\mathbb R\vdash A\RealAdd\RealNeg A=\RealZero
  }[RealCutAdditionInverse]
    \proofstepwidestarbreakkeep[1]{A\in\mathbb R}{\rA}
    \proofstepwidestarbreakkeep[2]{q\in A\RealAdd\RealNeg A}{\rA}
    \proofstepwidestarbreakkeep[1,2]{%
      \exists a\in A\,\exists b\in\RealNeg A\,(q<a+b)}{%
      \FormulaRefAuto{RealCutAddition}[def]{1,2}}
    \proofstepwidestarbreakkeep[1,2]{%
      \exists a\in A\,\exists r\notin A\,(q<a-r)}{%
      \FormulaRefAuto{RealCutNegation}[def]{3}}
    \proofstepwidestar[1,2]{\RealCut A}{%
      \FormulaRefAuto{RealCutMembership}[thm]{1}}
    \proofstepwidestarbreakkeep[1,2]{q<\RatZero}{%
      \FormulaRefAuto{DedekindCutCharacterization}[thm]{%
        5,4}}
    \proofstepwidestar[1,2]{q\in\RealHat{\RatZero}}{%
      \FormulaRefAuto{RealPrincipalCut}[def]{6}}
    \proofstepwidestarbreakkeep[1,2]{q\in\RealZero}{%
      \FormulaRefAuto{RealCutZero}[def]{%
        7}}
    \proofstepwidestarbreakkeep[1]{A\RealAdd\RealNeg A\subseteq\RealZero}{%
      \FormulaRefAuto{SubsetIntroductionRule}[thm]{[2]\vdots 8}}
    \proofstepwidestarbreakkeep[10]{q\in\RealZero}{\rA}
    \proofstepwidestar[10]{q<\RatZero}{%
      \FormulaRefAuto{RealPrincipalCut}[def]{10}}
    \proofstepwidestarbreakkeep[10]{q<\RatZero}{%
      \FormulaRefAuto{RealCutZero}[def]{%
        11}}
    \proofstepwidestarbreakkeep[10]{%
      \exists h\in\mathbb Q\,(\RatZero<h\land h<-q)}{%
      \FormulaRefAuto{RationalOrderDense}[thm]{12}}
    \proofstepwidestarbreakkeep[10]{h\in\mathbb Q\land\RatZero<h\land h<-q}{\rA}
    \proofstepwidestarbreakkeep[1,10]{%
      \exists a\in A\,\exists r\in\mathbb Q\,
      (r\notin A\land a<r\land r-a<h)}{%
      \FormulaRefAuto{RealCutAdditiveBoundaryApproximation}[thm]{1,14}}
    \proofstepwidestarbreakkeep[1,10]{%
      a\in A\land r\in\mathbb Q\land r\notin A\land a<r\land r-a<h}{\rA}
    \proofstepwidestarbreakkeep[1,10]{q<a-r}{%
      \FormulaRefAuto{RationalOrderTranslation}[thm]{12,14,16}}
    \proofstepwidestarbreakkeep[1,10]{%
      \exists b\in\mathbb Q\,(q-a<b\land b<-r)}{%
      \FormulaRefAuto{RationalOrderDense}[thm]{17}}
    \proofstepwidestarbreakkeep[1,10]{b\in\mathbb Q\land q-a<b\land b<-r}{\rA}
    \proofstepwidestarbreakkeep[1,10]{b\in\RealNeg A}{%
      \FormulaRefAuto{RealCutNegation}[def]{16,19}}
    \proofstepwidestarbreakkeep[1,10]{q<a+b}{%
      \FormulaRefAuto{RationalOrderTranslation}[thm]{19}}
    \proofstepwidestarbreakkeep[1,10]{q\in A\RealAdd\RealNeg A}{%
      \FormulaRefAuto{RealCutAddition}[def]{16,20,21}}
    \proofstepwidestarbreakkeep[1,10]{q\in A\RealAdd\RealNeg A}{\rEE{18,19,22}}
    \proofstepwidestarbreakkeep[1,10]{q\in A\RealAdd\RealNeg A}{\rEE{15,16,23}}
    \proofstepwidestarbreakkeep[1,10]{q\in A\RealAdd\RealNeg A}{\rEE{13,14,24}}
    \proofstepwidestarbreakkeep[1]{\RealZero\subseteq A\RealAdd\RealNeg A}{%
      \FormulaRefAuto{SubsetIntroductionRule}[thm]{[10]\vdots 25}}
    \proofstepwidestarbreakkeep[1]{A\RealAdd\RealNeg A=\RealZero}{%
      \FormulaRefAuto{A\subseteq B\dsep B\subseteq A\vdash A=B}{9,26}}
  \closeproofpartwideindR

  \proofpartwide{Zusammenfassung}
    \proofstepwidestarbreakkeep[1]{A,B,C\in\mathbb R}{\rA}
    \proofstepwidestarbreakkeep[1]{(A\RealAdd B)\RealAdd C=A\RealAdd(B\RealAdd C)}{%
      \FormulaRefAuto{RealCutAdditionAssociative}[thm]{1}}
    \proofstepwidestarbreakkeep[1]{A\RealAdd B=B\RealAdd A}{%
      \FormulaRefAuto{RealCutAdditionCommutative}[thm]{1}}
    \proofstepwidestarbreakkeep[1]{A\RealAdd\RealZero=A}{%
      \FormulaRefAuto{RealCutAdditionZero}[thm]{1}}
    \proofstepwidestarbreakkeep[1]{A\RealAdd\RealNeg A=\RealZero}{%
      \FormulaRefAuto{RealCutAdditionInverse}[thm]{1}}
    \proofstepwidestarbreakkeep[1]{%
      (A\RealAdd B)\RealAdd C=A\RealAdd(B\RealAdd C)\land
      A\RealAdd B=B\RealAdd A\land
      A\RealAdd\RealZero=A\land
      A\RealAdd\RealNeg A=\RealZero}{%
      \rAI{2,\rAI{3,\rAI{4,5}}}}
  \closeproofpartwide
\end{tabproofsplitwide}

\FormulaThmDeltaKR[Negation kehrt die Schnittordnung um]{%
  \begin{aligned}[t]
    A,B\in\mathbb R\vdash{}&
    \text{(i)}\quad A\RealLe B\leftrightarrow\RealNeg B\RealLe\RealNeg A,\\[-2pt]
    &\text{(ii)}\quad A\RealLt B\leftrightarrow\RealNeg B\RealLt\RealNeg A,\\[-2pt]
    &\text{(iii)}\quad \RealNeg{\RealNeg A}=A,\\[-2pt]
    &\text{(iv)}\quad \RealNeg{\RealZero}=\RealZero.
  \end{aligned}%
}{%
  A,B\in\mathbb R\vdash
  \begin{aligned}[t]
    &\bigl(A\RealLe B\leftrightarrow
      \RealNeg B\RealLe\RealNeg A\bigr)\dsep
      \bigl(A\RealLt B\leftrightarrow
      \RealNeg B\RealLt\RealNeg A\bigr)\dsep{}\\[-2pt]
    &\RealNeg{\RealNeg A}=A\dsep
      \RealNeg{\RealZero}=\RealZero.
  \end{aligned}%
}{RealCutNegationOrder}{%
  \DeltaRow{Mengen}{A\dsep B}%
  \DeltaRow{Relationsoperatoren}{\RealLe\dsep\RealLt}%
  \DeltaRow{Funktionsoperatoren}{\RealNeg{\,\cdot\,}}%
}
\begin{tabproofsplitwide}
\proofpartwideindR[Eindeutigkeit des additiven Inversen]{%
    X,Y\in\mathbb R\dsep X\RealAdd Y=\RealZero
    \vdash Y=\RealNeg X
  }{%
    X,Y\in\mathbb R\dsep X\RealAdd Y=\RealZero
    \vdash Y=\RealNeg X
  }[RealCutAdditiveInverseUnique]
    \proofstepwidestar[1]{X,Y\in\mathbb R}{\rA}
    \proofstepwidestar[2]{X\RealAdd Y=\RealZero}{\rA}
    \proofstepwide[1]{Y}{=}{Y\RealAdd\RealZero}{%
      \FormulaRefAuto{RealCutAdditiveLaws}[thm]{1}}
    \proofstepwide[1]{Y\RealAdd\RealZero}{=}{%
      Y\RealAdd(X\RealAdd\RealNeg X)}{%
      \FormulaRefAuto{RealCutAdditiveLaws}[thm]{1}}
    \proofstepwide[1]{Y\RealAdd(X\RealAdd\RealNeg X)}{=}{%
      (Y\RealAdd X)\RealAdd\RealNeg X}{%
      \FormulaRefAuto{RealCutAdditionAssociative}[thm]{1}}
    \proofstepwide[1]{(Y\RealAdd X)\RealAdd\RealNeg X}{=}{%
      (X\RealAdd Y)\RealAdd\RealNeg X}{%
      \FormulaRefAuto{RealCutAdditionCommutative}[thm]{1}}
    \proofstepwide[1,2]{(X\RealAdd Y)\RealAdd\RealNeg X}{=}{%
      \RealZero\RealAdd\RealNeg X}{\rIE{2}}
    \proofstepwidestar[1]{(\RealNeg X)\RealAdd\RealZero=\RealNeg X}{%
      \FormulaRefAuto{RealCutAdditionZero}[thm]{1}}
    \proofstepwide[1]{\RealZero\RealAdd\RealNeg X}{=}{\RealNeg X}{%
      \FormulaRefAuto{RealCutAdditionCommutative}[thm]{%
        8}}
    \proofstepwide[1,2]{Y}{=}{\RealNeg X}{\rChain{3,4,5,6,7,9}}
  \closeproofpartwideindR

  \proofpartwide{Involutivität}
\proofstepwidestar[1]{A\in\mathbb R}{\rA}
    \proofstepwidestar[1]{\RealNeg A\in\mathbb R}{%
      \FormulaRefAuto{RealCutAdditiveClosure}[thm]{1}}
    \proofstepwide[1]{\RealNeg A\RealAdd A}{=}{A\RealAdd\RealNeg A}{%
      \FormulaRefAuto{RealCutAdditionCommutative}[thm]{1}}
    \proofstepwide[1]{A\RealAdd\RealNeg A}{=}{\RealZero}{%
      \FormulaRefAuto{RealCutAdditionInverse}[thm]{1}}
    \proofstepwidestar[1]{A=\RealNeg{\RealNeg A}}{%
      \FormulaRefAuto{RealCutAdditiveInverseUnique}[thm]{2,3,4}}
    \proofstepwidestar[1]{\RealNeg{\RealNeg A}=A}{%
      \FormulaRefAuto{a=b\vdash b=a}{5}}

\closeproofpartwide

\proofpartwide{Negation des Nullschnitts}
  \setcounter{proofstepnr}{6}% Fortlaufende Schritte dieses Beweises.
\proofstepwidestar[]{\RealZero\in\mathbb R}{%
      \FormulaRefAuto{RealCutAdditiveClosure}[thm]}
    \proofstepwide[]{\RealZero\RealAdd\RealZero}{=}{\RealZero}{%
      \FormulaRefAuto{RealCutAdditionZero}[thm]{7}}
    \proofstepwidestar[]{\RealZero=\RealNeg{\RealZero}}{%
      \FormulaRefAuto{RealCutAdditiveInverseUnique}[thm]{7,8}}
    \proofstepwidestar[]{\RealNeg{\RealZero}=\RealZero}{%
      \FormulaRefAuto{a=b\vdash b=a}{9}}

\closeproofpartwide

\proofpartwideindR[Involutivität und Null]{%
    A\in\mathbb R\vdash
    \RealNeg{\RealNeg A}=A\land\RealNeg{\RealZero}=\RealZero
  }{%
    A\in\mathbb R\vdash
    \RealNeg{\RealNeg A}=A\land\RealNeg{\RealZero}=\RealZero
  }[RealCutNegationInvolutionZero]
  \setcounter{proofstepnr}{10}% Fortlaufende Schritte dieses Beweises.
\proofstepwidestar[1]{%
      \RealNeg{\RealNeg A}=A\land\RealNeg{\RealZero}=\RealZero}{%
      \rAI{6,10}}

\closeproofpartwideindR


  \proofpartwideindR[Antitonie]{%
    A,B\in\mathbb R\dsep A\subseteq B
    \vdash\RealNeg B\subseteq\RealNeg A
  }{%
    A,B\in\mathbb R\dsep A\subseteq B
    \vdash\RealNeg B\subseteq\RealNeg A
  }[RealCutNegationAntitone]
    \proofstepwidestarbreakkeep[1]{A,B\in\mathbb R}{\rA}
    \proofstepwidestarbreakkeep[2]{A\subseteq B}{\rA}
    \proofstepwidestarbreakkeep[3]{q\in\RealNeg B}{\rA}
    \proofstepwidestarbreakkeep[1,3]{%
      \exists r\in\mathbb Q\,(r\notin B\land q<-r)}{%
      \FormulaRefAuto{RealCutNegation}[def]{1,3}}
    \proofstepwidestarbreakkeep[1,3]{r\in\mathbb Q\land r\notin B\land q<-r}{\rA}
    \proofstepwidestar[2,3]{\forall x\,(x\in A\rightarrow x\in B)}{%
      \FormulaRefAuto{A\subseteq B\coloneq
          \forall x(x\in A\rightarrow x\in B)}{2}}
    \proofstepwidestarbreakkeep[2,3]{r\notin A}{%
      \FormulaRefAuto{P\rightarrow Q\dsep\neg Q\vdash\neg P}[thm]{%
        6,5}}
    \proofstepwidestarbreakkeep[1,2,3]{q\in\RealNeg A}{%
      \FormulaRefAuto{RealCutNegation}[def]{1,5,7}}
    \proofstepwidestarbreakkeep[1,2,3]{q\in\RealNeg A}{\rEE{4,5,8}}
    \proofstepwidestarbreakkeep[1,2]{\RealNeg B\subseteq\RealNeg A}{%
      \FormulaRefAuto{SubsetIntroductionRule}[thm]{[3]\vdots 9}}
  \closeproofpartwideindR

  \proofpartwide{Ordnungsäquivalenzen}
    \proofstepwidestarbreakkeep[1]{A,B\in\mathbb R}{\rA}
    \proofstepwidestarbreakkeep[2]{A\subseteq B}{\rA}
    \proofstepwidestarbreakkeep[1,2]{\RealNeg B\subseteq\RealNeg A}{%
      \FormulaRefAuto{RealCutNegationAntitone}[thm]{1,2}}
    \proofstepwidestarbreakkeep[4]{\RealNeg B\subseteq\RealNeg A}{\rA}
    \proofstepwidestarbreakkeep[1,4]{\RealNeg{\RealNeg A}\subseteq\RealNeg{\RealNeg B}}{%
      \FormulaRefAuto{RealCutNegationAntitone}[thm]{1,4}}
    \proofstepwidestarbreakkeep[1,4]{A\subseteq B}{%
      \FormulaRefAuto{RealCutNegationInvolutionZero}[thm]{1,5}}
    \proofstepwide[1]{A\subseteq B}{\leftrightarrow}{%
      \RealNeg B\subseteq\RealNeg A}{\rLRI{3,6}}
    \proofstepwide[1]{A\RealLe B}{\leftrightarrow}{%
      \RealNeg B\RealLe\RealNeg A}{%
      \FormulaRefAuto{RealCutOrderCharacterization}[thm]{1,7}}
    \proofstepwidestarbreakkeep[9]{A\subsetneq B}{\rA}
    \proofstepwidestarbreakkeep[1,9]{A\subseteq B\land A\neq B}{%
      \FormulaRefAuto{A\subset B\coloneq A\subseteq B\land A\neq B}{9}}
    \proofstepwidestarbreakkeep[1,9]{\RealNeg B\subseteq\RealNeg A}{%
      \FormulaRefAuto{RealCutNegationAntitone}[thm]{1,10}}
    \proofstepwidestarbreakkeep[12]{\RealNeg B=\RealNeg A}{\rA}
    \proofstepwidestarbreakkeep[1,12]{A=B}{%
      \FormulaRefAuto{RealCutNegationInvolutionZero}[thm]{1,12}}
    \proofstepwidestarbreakkeep[1,9,12]{\bot}{\rBI{10,13}}
    \proofstepwidestarbreakkeep[1,9]{\RealNeg B\neq\RealNeg A}{\rCI{12,14}}
    \proofstepwidestarbreakkeep[1,9]{\RealNeg B\subsetneq\RealNeg A}{%
      \FormulaRefAuto{A\subset B\coloneq A\subseteq B\land A\neq B}{11,15}}
    \proofstepwidestarbreakkeep[17]{\RealNeg B\subsetneq\RealNeg A}{\rA}
    \proofstepwidestarbreakkeep[1,17]{\RealNeg B\subseteq\RealNeg A}{%
      \FormulaRefAuto{A\subset B\vdash A\subseteq B}{17}}
    \proofstepwidestarbreakkeep[1,17]{A\subseteq B}{\rLREb{7,18}}
    \proofstepwidestarbreakkeep[20]{A=B}{\rA}
    \proofstepwidestarbreakkeep[1,20]{\RealNeg B=\RealNeg A}{\rIE{20}}
    \proofstepwidestarbreakkeep[1,17,20]{\bot}{%
      \FormulaRefAuto{A\subset B\vdash A\neq B}{17,21}}
    \proofstepwidestarbreakkeep[1,17]{A\neq B}{\rCI{20,22}}
    \proofstepwidestarbreakkeep[1,17]{A\subsetneq B}{%
      \FormulaRefAuto{A\subset B\coloneq A\subseteq B\land A\neq B}{19,23}}
    \proofstepwide[1]{A\subsetneq B}{\leftrightarrow}{%
      \RealNeg B\subsetneq\RealNeg A}{\rLRI{16,24}}
    \proofstepwide[1]{A\RealLt B}{\leftrightarrow}{%
      \RealNeg B\RealLt\RealNeg A}{%
      \FormulaRefAuto{RealCutOrderCharacterization}[thm]{1,25}}
    \proofstepwidestarbreakkeep[1]{\RealNeg{\RealNeg A}=A}{%
      \FormulaRefAuto{RealCutNegationInvolutionZero}[thm]{1}}
    \proofstepwidestarbreakkeep[]{\RealNeg{\RealZero}=\RealZero}{%
      \FormulaRefAuto{RealCutNegationInvolutionZero}[thm]}
    \proofstepwidestarbreakkeep[1]{%
      \begin{aligned}[t]
        &(A\RealLe B\leftrightarrow\RealNeg B\RealLe\RealNeg A)\land{}\\[-2pt]
        &(A\RealLt B\leftrightarrow\RealNeg B\RealLt\RealNeg A)\land{}\\[-2pt]
        &\RealNeg{\RealNeg A}=A\land{}\\[-2pt]
        &\RealNeg{\RealZero}=\RealZero
      \end{aligned}}{%
      \rAI{8,\rAI{26,\rAI{27,28}}}}
  \closeproofpartwide
\end{tabproofsplitwide}

\section{Multiplikation und multiplikative Inverse}

Zunächst wird das Produkt nichtnegativer Schnitte definiert. Die rationale
negative Halbachse \(\RealZero\) wird ausdrücklich beigefügt. Dadurch ist
das Produkt mit \(\RealZero\) wieder \(\RealZero\); der rationale Nullpunkt
gehört genau dann zum positiven Produkt, wenn er zu beiden Faktoren gehört.

\FormulaDefDeltaK[Positives Produkt reeller Schnitte]{%
  \begin{aligned}[t]
  &A,B\in\mathbb R\dsep
    \RealZero\RealLe A\dsep\RealZero\RealLe B\\[-2pt]
  &\quad\vdash
  A\RealPosMul B\coloneqq
  \RealZero\cup
  \bigl\{q\in\mathbb Q\bigm|
    \exists a\in A\,\exists b\in B\,
      (\RatZero<a\land\RatZero<b\land q<ab)\bigr\}.
  \end{aligned}%
}{RealCutPositiveProduct}{%
  \DeltaRow{Mengen}{A\dsep B\dsep a\dsep b\dsep q}%
  \DeltaRow{Relationsoperatoren}{\RealLe}%
  \DeltaRow{Funktionsoperatoren}{\RealPosMul}%
  \DeltaRow{\textbf{Neues Symbol}}{}%
  \DeltaRow{Positives Schnittprodukt}{\RealPosMul}%
}

\FormulaDefDeltaK[Multiplikation beliebiger reeller Schnitte]{%
  A,B\in\mathbb R\vdash
  A\RealMul B\coloneqq
  \begin{cases}
    A\RealPosMul B,
      &\RealZero\RealLe A\land\RealZero\RealLe B,\\
    \RealNeg{\bigl((\RealNeg A)\RealPosMul B\bigr)},
      &A\RealLt\RealZero\land\RealZero\RealLe B,\\
    \RealNeg{\bigl(A\RealPosMul(\RealNeg B)\bigr)},
      &\RealZero\RealLe A\land B\RealLt\RealZero,\\
    (\RealNeg A)\RealPosMul(\RealNeg B),
      &A\RealLt\RealZero\land B\RealLt\RealZero.
  \end{cases}%
}{RealCutMultiplication}{%
  \DeltaRow{Mengen}{A\dsep B}%
  \DeltaRow{Relationsoperatoren}{\RealLe\dsep\RealLt}%
  \DeltaRow{Funktionsoperatoren}{\RealMul\dsep\RealPosMul
    \dsep\RealNeg{\,\cdot\,}}%
  \DeltaRow{\textbf{Neues Symbol}}{}%
  \DeltaRow{Reelle Multiplikation}{\RealMul}%
}

\FormulaDefDeltaK[Multiplikative Eins der reellen Zahlen]{%
  \RealOne\coloneqq\RealHat{\RatOne}%
}{RealCutOne}{%
  \DeltaRow{Mengen}{\RealOne}%
  \DeltaRow{\textbf{Neues Symbol}}{}%
  \DeltaRow{Reelle Eins}{\RealOne}%
}

\FormulaDefDeltaK[Positives Reziprokes eines Schnittes]{%
  \begin{aligned}[t]
  &A\in\mathbb R\dsep\RealZero\RealLt A\\[-2pt]
  &\quad\vdash
  \RealPosInv A\coloneqq
  \RealZero\cup
  \bigl\{q\in\mathbb Q\bigm|
    \exists r\in\mathbb Q\,
      (\RatZero<r\land r\notin A\land q<r^{-1})\bigr\}.
  \end{aligned}%
}{RealCutPositiveReciprocal}{%
  \DeltaRow{Mengen}{A\dsep q\dsep r}%
  \DeltaRow{Relationsoperatoren}{\RealLt}%
  \DeltaRow{Funktionsoperatoren}{\RealPosInv{\,\cdot\,}}%
  \DeltaRow{\textbf{Neues Symbol}}{}%
  \DeltaRow{Positives Reziprokes}{\RealPosInv{\,\cdot\,}}%
}

\FormulaDefDeltaK[Reziprokes einer von Null verschiedenen reellen Zahl]{%
  A\in\mathbb R\dsep A\neq\RealZero\vdash
  \RealInv A\coloneqq
  \begin{cases}
    \RealPosInv A,&\RealZero\RealLt A,\\
    \RealNeg{\bigl(\RealPosInv{\RealNeg A}\bigr)},&A\RealLt\RealZero.
  \end{cases}%
}{RealCutReciprocal}{%
  \DeltaRow{Mengen}{A}%
  \DeltaRow{Relationsoperatoren}{\RealLt}%
  \DeltaRow{Funktionsoperatoren}{\RealInv{\,\cdot\,}
    \dsep\RealPosInv{\,\cdot\,}\dsep\RealNeg{\,\cdot\,}}%
  \DeltaRow{\textbf{Neues Symbol}}{}%
  \DeltaRow{Reelles Reziprokes}{\RealInv{\,\cdot\,}}%
}

\FormulaThmDeltaKR[Abgeschlossenheit der multiplikativen Schnittoperationen]{%
  \begin{aligned}[t]
    &\text{(i)}\quad A,B\in\mathbb R\dsep
      \RealZero\RealLe A\dsep\RealZero\RealLe B
      \vdash A\RealPosMul B\in\mathbb R,\\[-2pt]
    &\text{(ii)}\quad A,B\in\mathbb R
      \vdash A\RealMul B\in\mathbb R\land\RealOne\in\mathbb R,\\[-2pt]
    &\text{(iii)}\quad A\in\mathbb R\dsep\RealZero\RealLt A
      \vdash\RealPosInv A\in\mathbb R,\\[-2pt]
    &\text{(iv)}\quad A\in\mathbb R\dsep A\neq\RealZero
      \vdash\RealInv A\in\mathbb R.
  \end{aligned}%
}{%
  \begin{aligned}[t]
    &A,B\in\mathbb R\dsep
      \RealZero\RealLe A\dsep\RealZero\RealLe B
      \vdash A\RealPosMul B\in\mathbb R,\\[-2pt]
    &A,B\in\mathbb R
      \vdash A\RealMul B\in\mathbb R\land\RealOne\in\mathbb R,\\[-2pt]
    &A\in\mathbb R\dsep\RealZero\RealLt A
      \vdash\RealPosInv A\in\mathbb R,\\[-2pt]
    &A\in\mathbb R\dsep A\neq\RealZero
      \vdash\RealInv A\in\mathbb R.
  \end{aligned}%
}{RealCutMultiplicativeClosure}{%
  \DeltaRow{Mengen}{A\dsep B}%
  \DeltaRow{Funktionsoperatoren}{\RealPosMul\dsep\RealMul
    \dsep\RealPosInv{\,\cdot\,}\dsep\RealInv{\,\cdot\,}}%
}
\begin{tabproofsplitwide}
  \proofpartwideindR[Schnittpunkt unter rationalem Au\ss enpunkt]{%
    A\in\mathbb R\dsep a\in A\dsep r\in\mathbb Q\dsep r\notin A
    \vdash a<r
  }{%
    A\in\mathbb R\dsep a\in A\dsep r\in\mathbb Q\dsep r\notin A
    \vdash a<r
  }[RealCutMemberBelowExterior]
    \proofstepwidestar[1]{A\in\mathbb R}{\rA}
    \proofstepwidestar[2]{a\in A}{\rA}
    \proofstepwidestar[3]{r\in\mathbb Q}{\rA}
    \proofstepwidestar[4]{r\notin A}{\rA}
    \proofstepwidestar[1]{\RealCut A}{%
      \FormulaRefAuto{RealCutMembership}[thm]{1}}
    \proofstepwidestar[1,2]{a\in\mathbb Q}{%
      \FormulaRefAuto{DedekindCutCharacterization}[thm]{5,2}}
    \proofstepwidestar[7]{r<a}{\rA}
    \proofstepwidestar[1,2,3,7]{r\in A}{%
      \FormulaRefAuto{DedekindCutCharacterization}[thm]{5,2,3,7}}
    \proofstepwidestar[1,2,3,4,7]{\bot}{\rBI{4,8}}
    \proofstepwidestar[1,2,3,4]{\neg(r<a)}{\rCI{7,9}}
    \proofstepwidestar[1,2,3,4]{a\leq r}{%
      \FormulaRefAuto{RationalOrderTotal}[thm]{6,3,10}}
    \proofstepwidestar[2,4]{a\neq r}{%
      \FormulaRefAuto{x\in A\dsep y\notin A \vdash x\neq y}[thm]{2,4}}
    \proofstepwidestar[1,2,3,4]{a<r}{%
      \FormulaRefAuto{RationalStrictOrderDef}[def]{6,3,11,12}}
  \closeproofpartwideindR

  \proofpartwideindR[Positiver Punkt eines positiven Schnittes]{%
    A\in\mathbb R\dsep\RealZero\RealLt A
    \vdash\exists a\in A\,(\RatZero<a)
  }{%
    A\in\mathbb R\dsep\RealZero\RealLt A
    \vdash\exists a\in A\,(\RatZero<a)
  }[RealCutPositiveMember]
    \proofstepwidestar[1]{A\in\mathbb R}{\rA}
    \proofstepwidestar[2]{\RealZero\RealLt A}{\rA}
    \proofstepwidestar[1]{\RealCut A}{%
      \FormulaRefAuto{RealCutMembership}[thm]{1}}
    \proofstepwidestarbreakkeep[1,2]{\RealZero\subsetneq A}{%
      \FormulaRefAuto{RealCutOrderCharacterization}[thm]{1,2}}
    \proofstepwidestarbreakkeep[1,2]{%
      \RealZero\subseteq A\land\RealZero\neq A}{%
      \FormulaRefAuto{A\subset B\coloneq A\subseteq B\land A\neq B}{4}}
    \proofstepwidestarbreakkeep[1,2]{\exists x\in A\setminus\RealZero}{%
      \FormulaRefAuto{A\subseteq B\dsep A\neq B\vdash
        \exists x\in B\setminus A}[thm]{\rAEa{5},\rAEb{5}}}
    \proofstepwidestar[1,2]{x\in A\setminus\RealZero}{\rA}
    \proofstepwidestarbreakkeep[1,2]{x\in A\land x\notin\RealZero}{%
      \FormulaRefAuto{x\in A\setminus B\eqvdash
        x\in A\land x\notin B}{7}}
    \proofstepwidestarbreakkeep[1,2]{x\in\mathbb Q}{%
      \FormulaRefAuto{DedekindCutCharacterization}[thm]{3,\rAEa{8}}}
    \proofstepwidestar[1,2]{\neg(x<\RatZero)}{%
      \FormulaRefAuto{RealPrincipalCut}[def]{9,\rAEb{8}}}
    \proofstepwidestarbreakkeep[1,2]{\neg(x<\RatZero)}{%
      \FormulaRefAuto{RealCutZero}[def]{%
        10}}
    \proofstepwidestar[1,2]{\RatZero \in \mathbb Q \land \RatOne \in \mathbb Q}{%
      \FormulaRefAuto{RationalZeroOneCarrier}[thm]}
    \proofstepwidestarbreakkeep[1,2]{\RatZero\leq x}{%
      \FormulaRefAuto{RationalOrderTotal}[thm]{%
        12,9,11}}
    \proofstepwidestarbreakkeep[1,2]{\exists a\in A\,(x<a)}{%
      \FormulaRefAuto{DedekindCutCharacterization}[thm]{3,\rAEa{8}}}
    \proofstepwidestar[1,2]{a\in A\land x<a}{\rA}
    \proofstepwidestarbreakkeep[1,2]{a\in\mathbb Q}{%
      \FormulaRefAuto{DedekindCutCharacterization}[thm]{3,\rAEa{15}}}
    \proofstepwidestarbreakkeep[1,2]{\RatZero<a}{%
      \FormulaRefAuto{RationalOrderTotal}[thm]{13,\rAEb{15}}}
    \proofstepwidestarbreakkeep[1,2]{\exists a\in A\,(\RatZero<a)}{%
      \rEI{\rAI{\rAEa{15},17}}}
    \proofstepwidestar[1,2]{\exists a\in A\,(\RatZero<a)}{%
      \rEE{14,15,18}}
    \proofstepwidestar[1,2]{\exists a\in A\,(\RatZero<a)}{%
      \rEE{6,7,19}}
  \closeproofpartwideindR

  \proofpartwideindR[Positive Au\ss engrenzen begrenzen das Schnittprodukt]{%
    \begin{aligned}[t]
      &A,B\in\mathbb R\dsep
        \RealZero\RealLe A\dsep\RealZero\RealLe B\dsep
        \alpha,\beta\in\mathbb Q\dsep{}\\[-2pt]
      &\RatZero<\alpha\dsep\alpha\notin A\dsep
        \RatZero<\beta\dsep\beta\notin B\\[-2pt]
      &\quad\vdash\alpha\beta\in\mathbb Q\land
        \alpha\beta\notin A\RealPosMul B
    \end{aligned}
  }{%
    \begin{aligned}[t]
      &A,B\in\mathbb R\dsep
        \RealZero\RealLe A\dsep\RealZero\RealLe B\dsep
        \alpha,\beta\in\mathbb Q\dsep{}\\[-2pt]
      &\RatZero<\alpha\dsep\alpha\notin A\dsep
        \RatZero<\beta\dsep\beta\notin B\\[-2pt]
      &\quad\vdash\alpha\beta\in\mathbb Q\land
        \alpha\beta\notin A\RealPosMul B
    \end{aligned}
  }[RealCutPositiveProductOutsideBound]
    \proofstepwidestar[1]{A,B\in\mathbb R}{\rA}
    \proofstepwidestar[2]{\RealZero\RealLe A}{\rA}
    \proofstepwidestar[3]{\RealZero\RealLe B}{\rA}
    \proofstepwidestar[4]{\alpha,\beta\in\mathbb Q}{\rA}
    \proofstepwidestar[5]{\RatZero<\alpha}{\rA}
    \proofstepwidestar[6]{\alpha\notin A}{\rA}
    \proofstepwidestar[7]{\RatZero<\beta}{\rA}
    \proofstepwidestar[8]{\beta\notin B}{\rA}
    \proofstepwidestar[4]{\alpha\beta\in\mathbb Q}{%
      \FormulaRefAuto{RationalMulClosure}[thm]{4}}
    \proofstepwidestarbreakkeep[4,5,7]{\RatZero<\alpha\beta}{%
      \FormulaRefAuto{RationalPositiveMulStrictOrder}[thm]{4,5,7},%
      \FormulaRefAuto{RationalMulZero}[thm]{\rAEa{4}}}
    \proofstepwidestar[11]{\alpha\beta\in A\RealPosMul B}{\rA}
    \proofstepwidestarbreakkeep[1,2,3,11]{%
      \alpha\beta\in\RealZero\lor
      \exists a\in A\,\exists b\in B\,
      (\RatZero<a\land\RatZero<b\land\alpha\beta<ab)}{%
      \FormulaRefAuto{RealCutPositiveProduct}[def]{1,2,3,11}}
    \proofstepwidestar[13]{\alpha\beta\in\RealZero}{\rA}
    \proofstepwidestar[9,13]{\alpha\beta<\RatZero}{%
      \FormulaRefAuto{RealPrincipalCut}[def]{9,13}}
    \proofstepwidestarbreakkeep[9,13]{\alpha\beta<\RatZero}{%
      \FormulaRefAuto{RealCutZero}[def]{%
        14}}
    \proofstepwidestarbreakkeep[4,5,7,13]{\bot}{%
      \FormulaRefAuto{RationalOrderTotal}[thm]{9,10,15}}
    \proofstepwidestarbreakkeep[17]{%
      \exists a\in A\,\exists b\in B\,
      (\RatZero<a\land\RatZero<b\land\alpha\beta<ab)}{\rA}
    \proofstepwidestarbreakkeep[17]{%
      a\in A\land b\in B\land\RatZero<a\land\RatZero<b\land
      \alpha\beta<ab}{\rA}
    \proofstepwidestar[1,17]{\RealCut A}{%
      \FormulaRefAuto{RealCutMembership}[thm]{\rAEa{1}}}
    \proofstepwidestarbreakkeep[1,17]{a\in\mathbb Q}{%
      \FormulaRefAuto{DedekindCutCharacterization}[thm]{%
        19,18}}
    \proofstepwidestar[1,17]{\RealCut B}{%
      \FormulaRefAuto{RealCutMembership}[thm]{\rAEb{1}}}
    \proofstepwidestarbreakkeep[1,17]{b\in\mathbb Q}{%
      \FormulaRefAuto{DedekindCutCharacterization}[thm]{%
        21,18}}
    \proofstepwidestarbreakkeep[1,4,6,17]{a<\alpha}{%
      \FormulaRefAuto{RealCutMemberBelowExterior}[thm]{%
        \rAEa{1},18,\rAEa{4},6}}
    \proofstepwidestarbreakkeep[1,4,8,17]{b<\beta}{%
      \FormulaRefAuto{RealCutMemberBelowExterior}[thm]{%
        \rAEb{1},18,\rAEb{4},8}}
    \proofstepwidestarbreakkeep[1,4,7,17]{ab<\alpha b}{%
      \FormulaRefAuto{RationalPositiveMulStrictOrder}[thm]{%
        20,22,23,\rAEb{18}},%
      \FormulaRefAuto{RationalMulCommutative}[thm]{20,22,\rAEa{4}}}
    \proofstepwidestarbreakkeep[4,5,17]{\alpha b<\alpha\beta}{%
      \FormulaRefAuto{RationalPositiveMulStrictOrder}[thm]{%
        4,22,24,5}}
    \proofstepwidestarbreakkeep[1,4,5,7,17]{ab<\alpha\beta}{%
      \FormulaRefAuto{RationalOrderTotal}[thm]{25,26}}
    \proofstepwidestarbreakkeep[1,4,5,7,17]{\bot}{%
      \FormulaRefAuto{RationalOrderTotal}[thm]{18,27}}
    \proofstepwidestarbreakkeep[1,4\text{--}8,17]{\bot}{%
      \rEE{17,18,28}}
    \proofstepwidestarbreakkeep[1\text{--}8,11]{\bot}{%
      \rOE{12,13,16,17,29}}
    \proofstepwidestarbreakkeep[1\text{--}8]{%
      \alpha\beta\notin A\RealPosMul B}{\rCI{11,30}}
    \proofstepwidestarbreakkeep[1\text{--}8]{%
      \alpha\beta\in\mathbb Q\land
      \alpha\beta\notin A\RealPosMul B}{\rAI{9,31}}
  \closeproofpartwideindR

  \proofpartwideindR[Positives Produkt ist ein Schnitt]{%
    \begin{aligned}[t]
      &A,B\in\mathbb R\dsep\RealZero\RealLe A\dsep\RealZero\RealLe B\\[-2pt]
      &\quad\vdash\RealCut{A\RealPosMul B}
    \end{aligned}
  }{%
    A,B\in\mathbb R\dsep\RealZero\RealLe A\dsep\RealZero\RealLe B
    \vdash\RealCut{A\RealPosMul B}
  }[RealCutPositiveProductIsCut]
    \proofstepwidestar[1]{A,B\in\mathbb R}{\rA}
    \proofstepwidestar[2]{\RealZero\RealLe A}{\rA}
    \proofstepwidestar[3]{\RealZero\RealLe B}{\rA}
    \multicolumn{6}{@{}l@{}}{\itshape Nichtleerheit}\\*
    \proofstepwidestarbreakkeep[]{\RealZero\in\mathbb R}{%
      \FormulaRefAuto{RealCutAdditiveClosure}[thm]}
    \proofstepwidestarbreakkeep[]{\RealCut{\RealZero}}{%
      \FormulaRefAuto{RealCutMembership}[thm]{4}}
    \proofstepwidestarbreakkeep[]{%
      \RealZero\neq\varnothing\land\RealZero\neq\mathbb Q}{%
      \FormulaRefAuto{DedekindCutCharacterization}[thm]{5}}
    \proofstepwidestarbreakkeep[1,2,3]{\RealZero\subseteq A\RealPosMul B}{%
      \FormulaRefAuto{RealCutPositiveProduct}[def]{1,2,3}}
    \proofstepwidestarbreakkeep[]{\exists z\in\RealZero}{%
      \FormulaRefAuto{DedekindCutCharacterization}[thm]{5}}
    \proofstepwidestarbreakkeep[]{z\in\RealZero}{\rA}
    \proofstepwidestarbreakkeep[1,2,3]{z\in A\RealPosMul B}{%
      \FormulaRefAuto{A\subseteq B\dsep x\in A\vdash x\in B}[thm]{7,9}}
    \proofstepwidestarbreakkeep[1,2,3]{A\RealPosMul B\neq\varnothing}{%
      \FormulaRefAuto{x\in A\vdash A\neq\varnothing}[thm]{10}}
    \proofstepwidestarbreakkeep[1,2,3]{A\RealPosMul B\neq\varnothing}{%
      \rEE{8,9,11}}
    \multicolumn{6}{@{}l@{}}{\itshape Echtheit}\\*
    \proofstepwidestarbreakkeep[1,2,3]{%
      A=\RealZero\lor B=\RealZero\lor
      (\RealZero\RealLt A\land\RealZero\RealLt B)}{%
      \FormulaRefAuto{RealCutTotalOrder}[thm]{1,2,3}}
    \proofstepwidestarbreakkeep[14]{A=\RealZero}{\rA}
    \proofstepwidestarbreakkeep[1,2,3,14]{A\RealPosMul B=\RealZero}{%
      \FormulaRefAuto{RealCutPositiveProduct}[def]{1,2,3,14},%
      \FormulaRefAuto{RealCutZero}[def],%
      \FormulaRefAuto{RealPrincipalCut}[def]}
    \proofstepwidestarbreakkeep[1,2,3,14]{A\RealPosMul B\neq\mathbb Q}{%
      \rIE{15,\rAEb{6}}}
    \proofstepwidestarbreakkeep[17]{B=\RealZero}{\rA}
    \proofstepwidestarbreakkeep[1,2,3,17]{A\RealPosMul B=\RealZero}{%
      \FormulaRefAuto{RealCutPositiveProduct}[def]{1,2,3,17},%
      \FormulaRefAuto{RealCutZero}[def],%
      \FormulaRefAuto{RealPrincipalCut}[def]}
    \proofstepwidestarbreakkeep[1,2,3,17]{A\RealPosMul B\neq\mathbb Q}{%
      \rIE{18,\rAEb{6}}}
    \proofstepwidestarbreakkeep[20]{%
      \RealZero\RealLt A\land\RealZero\RealLt B}{\rA}
    \proofstepwidestarbreakkeep[1,2,3,20]{%
      \exists a\in A\,(\RatZero<a)}{%
      \FormulaRefAuto{RealCutPositiveMember}[thm]{\rAEa{1},\rAEa{20}}}
    \proofstepwidestarbreakkeep[22]{%
      a\in A\land\RatZero<a}{\rA}
    \proofstepwidestar[1]{\RealCut A}{%
      \FormulaRefAuto{RealCutMembership}[thm]{\rAEa{1}}}
    \proofstepwidestarbreakkeep[1]{A\subseteq\mathbb Q\land A\neq\mathbb Q}{%
      \FormulaRefAuto{DedekindCutCharacterization}[thm]{%
        23}}
    \proofstepwidestarbreakkeep[1,2,3,20]{%
      \exists\alpha\in\mathbb Q\setminus A}{%
      \FormulaRefAuto{A\subseteq B\dsep A\neq B\vdash
        \exists x\in B\setminus A}[thm]{\rAEa{24},\rAEb{24}}}
    \proofstepwidestarbreakkeep[26]{%
      \alpha\in\mathbb Q\setminus A}{\rA}
    \proofstepwidestarbreakkeep[26]{%
      \alpha\in\mathbb Q\land\alpha\notin A}{%
      \FormulaRefAuto{x\in A\setminus B\eqvdash
        x\in A\land x\notin B}{26}}
    \proofstepwidestarbreakkeep[1,22,26]{a<\alpha}{%
      \FormulaRefAuto{RealCutMemberBelowExterior}[thm]{%
        \rAEa{1},\rAEa{22},\rAEa{27},\rAEb{27}}}
    \proofstepwidestarbreakkeep[1,22,26]{\RatZero<\alpha}{%
      \FormulaRefAuto{RationalOrderTotal}[thm]{22,28}}
    \proofstepwidestarbreakkeep[1,2,3,20]{%
      \exists b\in B\,(\RatZero<b)}{%
      \FormulaRefAuto{RealCutPositiveMember}[thm]{\rAEb{1},\rAEb{20}}}
    \proofstepwidestarbreakkeep[31]{%
      b\in B\land\RatZero<b}{\rA}
    \proofstepwidestar[1]{\RealCut B}{%
      \FormulaRefAuto{RealCutMembership}[thm]{\rAEb{1}}}
    \proofstepwidestarbreakkeep[1]{B\subseteq\mathbb Q\land B\neq\mathbb Q}{%
      \FormulaRefAuto{DedekindCutCharacterization}[thm]{%
        32}}
    \proofstepwidestarbreakkeep[1,2,3,20]{%
      \exists\beta\in\mathbb Q\setminus B}{%
      \FormulaRefAuto{A\subseteq B\dsep A\neq B\vdash
        \exists x\in B\setminus A}[thm]{\rAEa{33},\rAEb{33}}}
    \proofstepwidestarbreakkeep[35]{%
      \beta\in\mathbb Q\setminus B}{\rA}
    \proofstepwidestarbreakkeep[35]{%
      \beta\in\mathbb Q\land\beta\notin B}{%
      \FormulaRefAuto{x\in A\setminus B\eqvdash
        x\in A\land x\notin B}{35}}
    \proofstepwidestarbreakkeep[1,31,35]{b<\beta}{%
      \FormulaRefAuto{RealCutMemberBelowExterior}[thm]{%
        \rAEb{1},\rAEa{31},\rAEa{36},\rAEb{36}}}
    \proofstepwidestarbreakkeep[1,31,35]{\RatZero<\beta}{%
      \FormulaRefAuto{RationalOrderTotal}[thm]{31,37}}
    \proofstepwidestarbreakkeep[1,2,3,22,26,31,35]{%
      \alpha\beta\in\mathbb Q\land
      \alpha\beta\notin A\RealPosMul B}{%
      \FormulaRefAuto{RealCutPositiveProductOutsideBound}[thm]{%
        1,2,3,27,29,36,38}}
    \proofstepwidestarbreakkeep[1,2,3,22,26,31,35]{%
      A\RealPosMul B\neq\mathbb Q}{%
      \FormulaRefAuto{x\in B, x\not\in A\vdash A\neq B}{%
        \rAEa{39},\rAEb{39}}}
    \proofstepwidestarbreakkeep[1,2,3,20,22,26,31]{%
      A\RealPosMul B\neq\mathbb Q}{\rEE{34,35,40}}
    \proofstepwidestarbreakkeep[1,2,3,20,22,26]{%
      A\RealPosMul B\neq\mathbb Q}{\rEE{30,31,41}}
    \proofstepwidestarbreakkeep[1,2,3,20,22]{%
      A\RealPosMul B\neq\mathbb Q}{\rEE{25,26,42}}
    \proofstepwidestarbreakkeep[1,2,3,20]{%
      A\RealPosMul B\neq\mathbb Q}{\rEE{21,22,43}}
    \proofstepwidestarbreakkeep[1,2,3]{A\RealPosMul B\neq\mathbb Q}{%
      \rOE{13,14,16,17,19,20,44}}
    \multicolumn{6}{@{}l@{}}{\itshape Abw\"artsabgeschlossenheit}\\*
    \proofstepwidestarbreakkeep[46]{q\in A\RealPosMul B}{\rA}
    \proofstepwidestarbreakkeep[47]{s\in\mathbb Q\land s<q}{\rA}
    \proofstepwidestarbreakkeep[1,2,3,46]{%
      q\in\RealZero\lor
      \exists a\in A\,\exists b\in B\,
      (\RatZero<a\land\RatZero<b\land q<ab)}{%
      \FormulaRefAuto{RealCutPositiveProduct}[def]{1,2,3,46}}
    \proofstepwidestarbreakkeep[49]{q\in\RealZero}{\rA}
    \proofstepwidestar[47,49]{\RealZero = \RealHat \RatZero}{%
      \FormulaRefAuto{RealCutZero}[def]{49}}
    \proofstepwidestar[47,49]{\RatZero \in \mathbb Q \land \RatOne \in \mathbb Q}{%
      \FormulaRefAuto{RationalZeroOneCarrier}[thm]}
    \proofstepwidestar[47,49]{s\in\RealHat{\RatZero}}{%
      \FormulaRefAuto{RealPrincipalCutDownwardClosed}[thm]{%
          51,%
          50,47}}
    \proofstepwidestarbreakkeep[47,49]{s\in\RealZero}{%
      \FormulaRefAuto{RealCutZero}[def]{%
        52}}
    \proofstepwidestarbreakkeep[1,2,3,47,49]{s\in A\RealPosMul B}{%
      \FormulaRefAuto{A\subseteq B\dsep x\in A\vdash x\in B}[thm]{7,53}}
    \proofstepwidestarbreakkeep[55]{%
      \exists a\in A\,\exists b\in B\,
      (\RatZero<a\land\RatZero<b\land q<ab)}{\rA}
    \proofstepwidestarbreakkeep[55]{%
      a\in A\land b\in B\land\RatZero<a\land\RatZero<b\land q<ab}{\rA}
    \proofstepwidestarbreakkeep[47,55]{s<ab}{%
      \FormulaRefAuto{RationalOrderTotal}[thm]{47,56}}
    \proofstepwidestarbreakkeep[1,2,3,47,55]{s\in A\RealPosMul B}{%
      \FormulaRefAuto{RealCutPositiveProduct}[def]{1,2,3,47,56,57}}
    \proofstepwidestarbreakkeep[1,2,3,47,55]{s\in A\RealPosMul B}{%
      \rEE{55,56,58}}
    \proofstepwidestarbreakkeep[1,2,3,46,47]{s\in A\RealPosMul B}{%
      \rOE{48,49,54,55,59}}
    \proofstepwidestarbreakkeep[1,2,3]{%
      \forall q\in A\RealPosMul B\,\forall s\in\mathbb Q\,
      (s<q\rightarrow s\in A\RealPosMul B)}{%
      \rUI{\rRI{46,\rRI{47,60}}}}
    \multicolumn{6}{@{}l@{}}{\itshape Kein gr\"o\ss tes Element}\\*
    \proofstepwidestarbreakkeep[62]{q\in A\RealPosMul B}{\rA}
    \proofstepwidestarbreakkeep[1,2,3,62]{%
      q\in\RealZero\lor
      \exists a\in A\,\exists b\in B\,
      (\RatZero<a\land\RatZero<b\land q<ab)}{%
      \FormulaRefAuto{RealCutPositiveProduct}[def]{1,2,3,62}}
    \proofstepwidestarbreakkeep[64]{q\in\RealZero}{\rA}
    \proofstepwidestar[64]{\RealZero = \RealHat \RatZero}{%
      \FormulaRefAuto{RealCutZero}[def]{64}}
    \proofstepwidestar[64]{\RatZero \in \mathbb Q \land \RatOne \in \mathbb Q}{%
      \FormulaRefAuto{RationalZeroOneCarrier}[thm]}
    \proofstepwidestar[64]{\exists r\in\RealHat{\RatZero}\,(q<r)}{%
      \FormulaRefAuto{RealPrincipalCutNoMaximum}[thm]{%
          66,%
          65}}
    \proofstepwidestarbreakkeep[64]{%
      \exists r\in\RealZero\,(q<r)}{%
      \FormulaRefAuto{RealCutZero}[def]{%
        67}}
    \proofstepwidestarbreakkeep[64]{r\in\RealZero\land q<r}{\rA}
    \proofstepwidestarbreakkeep[1,2,3,64]{r\in A\RealPosMul B}{%
      \FormulaRefAuto{A\subseteq B\dsep x\in A\vdash x\in B}[thm]{7,69}}
    \proofstepwidestarbreakkeep[1,2,3,64]{%
      \exists r\in A\RealPosMul B\,(q<r)}{%
      \rEI{\rAI{70,\rAEb{69}}}}
    \proofstepwidestarbreakkeep[1,2,3,64]{%
      \exists r\in A\RealPosMul B\,(q<r)}{\rEE{68,69,71}}
    \proofstepwidestarbreakkeep[73]{%
      \exists a\in A\,\exists b\in B\,
      (\RatZero<a\land\RatZero<b\land q<ab)}{\rA}
    \proofstepwidestarbreakkeep[73]{%
      a\in A\land b\in B\land\RatZero<a\land\RatZero<b\land q<ab}{\rA}
    \proofstepwidestarbreakkeep[1,2,3,62]{q\in\mathbb Q}{%
      \FormulaRefAuto{A\subseteq B\dsep x\in A\vdash x\in B}[thm]{7,62}}
    \proofstepwidestarbreakkeep[1,73]{a,b\in\mathbb Q}{%
      \rAIRepeat{2}{%
        \FormulaRefAuto{DedekindCutCharacterization}[thm]\circ
        \FormulaRefAuto{RealCutMembership}[thm]}{%
        \rAEa{1},74\mid\rAEb{1},74}}
    \proofstepwidestarbreakkeep[1,73]{ab\in\mathbb Q}{%
      \FormulaRefAuto{RationalMulClosure}[thm]{76}}
    \proofstepwidestarbreakkeep[1,2,3,62,73]{%
      \exists r\in\mathbb Q\,(q<r\land r<ab)}{%
      \FormulaRefAuto{RationalOrderDense}[thm]{75,77,74}}
    \proofstepwidestarbreakkeep[73]{%
      r\in\mathbb Q\land q<r\land r<ab}{\rA}
    \proofstepwidestarbreakkeep[1,2,3,73]{r\in A\RealPosMul B}{%
      \FormulaRefAuto{RealCutPositiveProduct}[def]{1,2,3,74,79}}
    \proofstepwidestarbreakkeep[1,2,3,73]{%
      \exists r\in A\RealPosMul B\,(q<r)}{%
      \rEI{\rAI{80,\rAEa{\rAEb{79}}}}}
    \proofstepwidestarbreakkeep[1,2,3,73]{%
      \exists r\in A\RealPosMul B\,(q<r)}{\rEE{78,79,81}}
    \proofstepwidestarbreakkeep[1,2,3,73]{%
      \exists r\in A\RealPosMul B\,(q<r)}{\rEE{73,74,82}}
    \proofstepwidestarbreakkeep[1,2,3,62]{%
      \exists r\in A\RealPosMul B\,(q<r)}{%
      \rOE{63,64,72,73,83}}
    \proofstepwidestarbreakkeep[1,2,3]{%
      \forall q\in A\RealPosMul B\,
      \exists r\in A\RealPosMul B\,(q<r)}{%
      \rUI{\rRI{62,84}}}
    \multicolumn{6}{@{}l@{}}{\itshape Zusammensetzung der Schnitteigenschaften}\\*
    \proofstepwidestarbreakkeep[1,2,3]{A\RealPosMul B\subseteq\mathbb Q}{%
      \FormulaRefAuto{RealCutPositiveProduct}[def]{1,2,3}}
    \proofstepwidestarbreakkeep[1,2,3]{\RealCut{A\RealPosMul B}}{%
      \FormulaRefAuto{DedekindCutCharacterization}[thm]{86,12,45,61,85}}
  \closeproofpartwideindR

  \proofpartwideindR[Positives Reziprokes ist ein Schnitt]{%
    A\in\mathbb R\dsep\RealZero\RealLt A
    \vdash\RealCut{\RealPosInv A}
  }{%
    A\in\mathbb R\dsep\RealZero\RealLt A
    \vdash\RealCut{\RealPosInv A}
  }[RealCutPositiveReciprocalIsCut]
    \proofstepwidestar[1]{A\in\mathbb R}{\rA}
    \proofstepwidestar[2]{\RealZero\RealLt A}{\rA}
    \multicolumn{6}{@{}l@{}}{\itshape Nichtleerheit}\\*
    \proofstepwidestarbreakkeep[]{\RealZero\in\mathbb R}{%
      \FormulaRefAuto{RealCutAdditiveClosure}[thm]}
    \proofstepwidestarbreakkeep[]{\RealCut{\RealZero}}{%
      \FormulaRefAuto{RealCutMembership}[thm]{3}}
    \proofstepwidestarbreakkeep[]{%
      \RealZero\neq\varnothing\land\RealZero\neq\mathbb Q}{%
      \FormulaRefAuto{DedekindCutCharacterization}[thm]{4}}
    \proofstepwidestarbreakkeep[1,2]{\RealZero\subseteq\RealPosInv A}{%
      \FormulaRefAuto{RealCutPositiveReciprocal}[def]{1,2}}
    \proofstepwidestarbreakkeep[]{\exists z\in\RealZero}{%
      \FormulaRefAuto{DedekindCutCharacterization}[thm]{4}}
    \proofstepwidestarbreakkeep[]{z\in\RealZero}{\rA}
    \proofstepwidestarbreakkeep[1,2]{z\in\RealPosInv A}{%
      \FormulaRefAuto{A\subseteq B\dsep x\in A\vdash x\in B}[thm]{6,8}}
    \proofstepwidestarbreakkeep[1,2]{\RealPosInv A\neq\varnothing}{%
      \FormulaRefAuto{x\in A\vdash A\neq\varnothing}[thm]{9}}
    \proofstepwidestarbreakkeep[1,2]{\RealPosInv A\neq\varnothing}{%
      \rEE{7,8,10}}
    \multicolumn{6}{@{}l@{}}{\itshape Echtheit}\\*
    \proofstepwidestarbreakkeep[1,2]{\exists a\in A\,(\RatZero<a)}{%
      \FormulaRefAuto{RealCutPositiveMember}[thm]{1,2}}
    \proofstepwidestarbreakkeep[13]{a\in A\land\RatZero<a}{\rA}
    \proofstepwidestar[1]{\RealCut A}{%
      \FormulaRefAuto{RealCutMembership}[thm]{1}}
    \proofstepwidestarbreakkeep[1]{A\subseteq\mathbb Q}{%
      \FormulaRefAuto{DedekindCutCharacterization}[thm]{%
        14}}
    \proofstepwidestarbreakkeep[1,2,13]{a\in\mathbb Q}{%
      \FormulaRefAuto{A\subseteq B\dsep x\in A\vdash x\in B}[thm]{15,13}}
    \proofstepwidestar[1,2,13]{\RatZero\leq a\land\RatZero\neq a}{%
      \FormulaRefAuto{RationalStrictOrderDef}[def]{\rAEb{13}}}
    \proofstepwidestarbreakkeep[1,2,13]{a\neq\RatZero}{%
      \FormulaRefAuto{a\neq b\vdash b\neq a}[thm]{%
        17}}
    \proofstepwidestar[1,2,13]{a ^ { - 1 } \in \mathbb Q}{%
      \FormulaRefAuto{RationalReciprocalClosureAxiom}[ax]{16,18}}
    \proofstepwidestar[1,2,13]{\RatZero < a ^ { - 1 }}{%
      \FormulaRefAuto{RationalPositiveReciprocalAxiom}[ax]{16,\rAEb{13}}}
    \proofstepwidestarbreakkeep[1,2,13]{%
      a^{-1}\in\mathbb Q\land\RatZero<a^{-1}}{%
      \rAI{%
        19,%
        20}}
    \proofstepwidestarbreakkeep[22]{a^{-1}\in\RealPosInv A}{\rA}
    \proofstepwidestarbreakkeep[1,2,22]{%
      a^{-1}\in\RealZero\lor
      \exists r\in\mathbb Q\,(\RatZero<r\land r\notin A\land
        a^{-1}<r^{-1})}{%
      \FormulaRefAuto{RealCutPositiveReciprocal}[def]{1,2,22}}
    \proofstepwidestarbreakkeep[24]{a^{-1}\in\RealZero}{\rA}
    \proofstepwidestar[24]{a^{-1}<\RatZero}{%
      \FormulaRefAuto{RealPrincipalCut}[def]{24}}
    \proofstepwidestarbreakkeep[24]{a^{-1}<\RatZero}{%
      \FormulaRefAuto{RealCutZero}[def]{%
        25}}
    \proofstepwidestarbreakkeep[1,2,13,24]{\bot}{%
      \FormulaRefAuto{RationalOrderTotal}[thm]{21,26}}
    \proofstepwidestarbreakkeep[28]{%
      \exists r\in\mathbb Q\,(\RatZero<r\land r\notin A\land
        a^{-1}<r^{-1})}{\rA}
    \proofstepwidestarbreakkeep[29]{%
      r\in\mathbb Q\land\RatZero<r\land r\notin A\land
      a^{-1}<r^{-1}}{\rA}
    \proofstepwidestarbreakkeep[1,2,13,29]{r<a}{%
      \FormulaRefAuto{RationalPositiveReciprocalOrderReversal}[thm]{%
        16,\rAEb{13},29}}
    \proofstepwidestar[1]{\RealCut A}{%
      \FormulaRefAuto{RealCutMembership}[thm]{1}}
    \proofstepwidestarbreakkeep[1,2,13,29]{r\in A}{%
      \FormulaRefAuto{DedekindCutCharacterization}[thm]{%
        31,13,29,30}}
    \proofstepwidestarbreakkeep[1,2,13,29]{\bot}{%
      \rBI{\rAEa{\rAEb{\rAEb{29}}},32}}
    \proofstepwidestarbreakkeep[1,2,13,28]{\bot}{\rEE{28,29,33}}
    \proofstepwidestarbreakkeep[1,2,13,22]{\bot}{%
      \rOE{23,24,27,28,34}}
    \proofstepwidestarbreakkeep[1,2,13]{a^{-1}\notin\RealPosInv A}{%
      \rCI{22,35}}
    \proofstepwidestarbreakkeep[1,2,13]{\RealPosInv A\neq\mathbb Q}{%
      \FormulaRefAuto{x\in B, x\not\in A\vdash A\neq B}{%
        \rAEa{21},36}}
    \proofstepwidestarbreakkeep[1,2]{\RealPosInv A\neq\mathbb Q}{%
      \rEE{12,13,37}}
    \multicolumn{6}{@{}l@{}}{\itshape Abw\"artsabgeschlossenheit}\\*
    \proofstepwidestarbreakkeep[39]{q\in\RealPosInv A}{\rA}
    \proofstepwidestarbreakkeep[40]{s\in\mathbb Q\land s<q}{\rA}
    \proofstepwidestarbreakkeep[1,2,39]{%
      q\in\RealZero\lor
      \exists r\in\mathbb Q\,
      (\RatZero<r\land r\notin A\land q<r^{-1})}{%
      \FormulaRefAuto{RealCutPositiveReciprocal}[def]{1,2,39}}
    \proofstepwidestarbreakkeep[42]{q\in\RealZero}{\rA}
    \proofstepwidestar[40,42]{\RealZero = \RealHat \RatZero}{%
      \FormulaRefAuto{RealCutZero}[def]{42}}
    \proofstepwidestar[40,42]{\RatZero \in \mathbb Q \land \RatOne \in \mathbb Q}{%
      \FormulaRefAuto{RationalZeroOneCarrier}[thm]}
    \proofstepwidestar[40,42]{s\in\RealHat{\RatZero}}{%
      \FormulaRefAuto{RealPrincipalCutDownwardClosed}[thm]{%
          44,%
          43,40}}
    \proofstepwidestarbreakkeep[40,42]{s\in\RealZero}{%
      \FormulaRefAuto{RealCutZero}[def]{%
        45}}
    \proofstepwidestarbreakkeep[1,2,40,42]{s\in\RealPosInv A}{%
      \FormulaRefAuto{A\subseteq B\dsep x\in A\vdash x\in B}[thm]{6,46}}
    \proofstepwidestarbreakkeep[48]{%
      \exists r\in\mathbb Q\,
      (\RatZero<r\land r\notin A\land q<r^{-1})}{\rA}
    \proofstepwidestarbreakkeep[48]{%
      r\in\mathbb Q\land\RatZero<r\land r\notin A\land q<r^{-1}}{\rA}
    \proofstepwidestarbreakkeep[40,48]{s<r^{-1}}{%
      \FormulaRefAuto{RationalOrderTotal}[thm]{40,49}}
    \proofstepwidestarbreakkeep[1,2,40,48]{s\in\RealPosInv A}{%
      \FormulaRefAuto{RealCutPositiveReciprocal}[def]{1,2,40,49,50}}
    \proofstepwidestarbreakkeep[1,2,40,48]{s\in\RealPosInv A}{%
      \rEE{48,49,51}}
    \proofstepwidestarbreakkeep[1,2,39,40]{s\in\RealPosInv A}{%
      \rOE{41,42,47,48,52}}
    \proofstepwidestarbreakkeep[1,2]{%
      \forall q\in\RealPosInv A\,\forall s\in\mathbb Q\,
      (s<q\rightarrow s\in\RealPosInv A)}{%
      \rUI{\rRI{39,\rRI{40,53}}}}
    \multicolumn{6}{@{}l@{}}{\itshape Kein gr\"o\ss tes Element}\\*
    \proofstepwidestarbreakkeep[55]{q\in\RealPosInv A}{\rA}
    \proofstepwidestarbreakkeep[1,2,55]{%
      q\in\RealZero\lor
      \exists r\in\mathbb Q\,
      (\RatZero<r\land r\notin A\land q<r^{-1})}{%
      \FormulaRefAuto{RealCutPositiveReciprocal}[def]{1,2,55}}
    \proofstepwidestarbreakkeep[57]{q\in\RealZero}{\rA}
    \proofstepwidestar[57]{\RealZero = \RealHat \RatZero}{%
      \FormulaRefAuto{RealCutZero}[def]{57}}
    \proofstepwidestar[57]{\RatZero \in \mathbb Q \land \RatOne \in \mathbb Q}{%
      \FormulaRefAuto{RationalZeroOneCarrier}[thm]}
    \proofstepwidestar[57]{\exists s\in\RealHat{\RatZero}\,(q<s)}{%
      \FormulaRefAuto{RealPrincipalCutNoMaximum}[thm]{%
          59,%
          58}}
    \proofstepwidestarbreakkeep[57]{%
      \exists s\in\RealZero\,(q<s)}{%
      \FormulaRefAuto{RealCutZero}[def]{%
        60}}
    \proofstepwidestarbreakkeep[57]{s\in\RealZero\land q<s}{\rA}
    \proofstepwidestarbreakkeep[1,2,57]{s\in\RealPosInv A}{%
      \FormulaRefAuto{A\subseteq B\dsep x\in A\vdash x\in B}[thm]{6,62}}
    \proofstepwidestarbreakkeep[1,2,57]{%
      \exists s\in\RealPosInv A\,(q<s)}{\rEI{\rAI{63,\rAEb{62}}}}
    \proofstepwidestarbreakkeep[1,2,57]{%
      \exists s\in\RealPosInv A\,(q<s)}{\rEE{61,62,64}}
    \proofstepwidestarbreakkeep[66]{%
      \exists r\in\mathbb Q\,
      (\RatZero<r\land r\notin A\land q<r^{-1})}{\rA}
    \proofstepwidestarbreakkeep[66]{%
      r\in\mathbb Q\land\RatZero<r\land r\notin A\land q<r^{-1}}{\rA}
    \proofstepwidestarbreakkeep[1,2,55]{q\in\mathbb Q}{%
      \FormulaRefAuto{RealCutPositiveReciprocal}[def]{1,2,55}}
    \proofstepwidestar[66]{\RatZero\leq r\land\RatZero\neq r}{%
      \FormulaRefAuto{RationalStrictOrderDef}[def]{67}}
    \proofstepwidestarbreakkeep[66]{r\neq\RatZero}{%
      \FormulaRefAuto{a\neq b\vdash b\neq a}[thm]{%
        69}}
    \proofstepwidestarbreakkeep[66]{r^{-1}\in\mathbb Q}{%
      \FormulaRefAuto{RationalReciprocalClosure}[thm]{67,70}}
    \proofstepwidestarbreakkeep[1,2,55,66]{%
      \exists s\in\mathbb Q\,(q<s\land s<r^{-1})}{%
      \FormulaRefAuto{RationalOrderDense}[thm]{68,71,67}}
    \proofstepwidestarbreakkeep[66]{%
      s\in\mathbb Q\land q<s\land s<r^{-1}}{\rA}
    \proofstepwidestarbreakkeep[1,2,66]{s\in\RealPosInv A}{%
      \FormulaRefAuto{RealCutPositiveReciprocal}[def]{1,2,67,73}}
    \proofstepwidestarbreakkeep[1,2,66]{%
      \exists s\in\RealPosInv A\,(q<s)}{%
      \rEI{\rAI{74,\rAEa{\rAEb{73}}}}}
    \proofstepwidestarbreakkeep[1,2,66]{%
      \exists s\in\RealPosInv A\,(q<s)}{\rEE{72,73,75}}
    \proofstepwidestarbreakkeep[1,2,66]{%
      \exists s\in\RealPosInv A\,(q<s)}{\rEE{66,67,76}}
    \proofstepwidestarbreakkeep[1,2,55]{%
      \exists s\in\RealPosInv A\,(q<s)}{\rOE{56,57,65,66,77}}
    \proofstepwidestarbreakkeep[1,2]{%
      \forall q\in\RealPosInv A\,\exists s\in\RealPosInv A\,(q<s)}{%
      \rUI{\rRI{55,78}}}
    \multicolumn{6}{@{}l@{}}{\itshape Zusammensetzung der Schnitteigenschaften}\\*
    \proofstepwidestarbreakkeep[1,2]{\RealPosInv A\subseteq\mathbb Q}{%
      \FormulaRefAuto{RealCutPositiveReciprocal}[def]{1,2}}
    \proofstepwidestarbreakkeep[1,2]{\RealCut{\RealPosInv A}}{%
      \FormulaRefAuto{DedekindCutCharacterization}[thm]{80,11,38,54,79}}
  \closeproofpartwideindR

  \proofpartwide{Fallzerlegung und Zusammenfassung}
    \proofstepwidestarbreakkeep[1]{A,B\in\mathbb R}{\rA}
    \proofstepwidestarbreakkeep[2]{\RealZero\RealLe A\land\RealZero\RealLe B}{\rA}
    \proofstepwidestar[1,2]{\RealCut{A\RealPosMul B}}{%
      \FormulaRefAuto{RealCutPositiveProductIsCut}[thm]{1,2}}
    \proofstepwidestarbreakkeep[1,2]{A\RealPosMul B\in\mathbb R}{%
      \FormulaRefAuto{RealCutMembership}[thm]{%
        3}}
    \proofstepwidestarbreakkeep[1]{%
      \begin{aligned}[t]
      &(A\RealLt\RealZero\rightarrow
        (\RealZero\RealLt\RealNeg A\land\RealNeg A\in\mathbb R))\land{}\\[-2pt]
      &(B\RealLt\RealZero\rightarrow
        (\RealZero\RealLt\RealNeg B\land\RealNeg B\in\mathbb R))
      \end{aligned}}{%
      \FormulaRefAuto{RealCutNegationOrder}[thm]{1},
      \FormulaRefAuto{RealCutAdditiveClosure}[thm]{1}}
    \proofstepwidestar[1]{\TotOrd { \mathbb R , \RealLe }}{%
      \FormulaRefAuto{RealCutTotalOrder}[thm]{1}}
    \proofstepwidestar[1]{\RealNeg A,\RealNeg B\in\mathbb R}{%
      \FormulaRefAuto{RealCutAdditiveClosure}[thm]{1}}
    \proofstepwidestar[1]{\begin{aligned}[t]&A,B\in\mathbb R\dsep\RealZero\RealLe A\dsep\RealZero\RealLe B\\[-2pt]&\quad\vdash\RealCut{A\RealPosMul B}\end{aligned}}{%
      \FormulaRefAuto{RealCutPositiveProductIsCut}[thm]}
    \proofstepwidestarbreakkeep[1]{%
      A\RealMul B\in\mathbb R}{%
      \FormulaRefAuto{RealCutMultiplication}[def]{%
        6,4,5,
        8,
        7}}
    \proofstepwidestarbreakkeep[]{\RatOne\in\mathbb Q}{%
      \FormulaRefAuto{RationalZeroOneCarrier}[thm]}
    \proofstepwidestar[]{\RealCut { \RealHat \RatOne }}{%
      \FormulaRefAuto{RealPrincipalCutIsCut}[thm]{10}}
    \proofstepwidestar[]{\RealHat{\RatOne}\in\mathbb R}{%
      \FormulaRefAuto{RealCutMembership}[thm]{%
          11}}
    \proofstepwidestarbreakkeep[]{\RealOne\in\mathbb R}{%
      \FormulaRefAuto{RealCutOne}[def]{%
        12}}
    \proofstepwidestarbreakkeep[14]{A\in\mathbb R\land\RealZero\RealLt A}{\rA}
    \proofstepwidestar[14]{\RealCut{\RealPosInv A}}{%
      \FormulaRefAuto{RealCutPositiveReciprocalIsCut}[thm]{14}}
    \proofstepwidestarbreakkeep[14]{\RealPosInv A\in\mathbb R}{%
      \FormulaRefAuto{RealCutMembership}[thm]{%
        15}}
    \proofstepwidestarbreakkeep[17]{A\in\mathbb R\land A\neq\RealZero}{\rA}
    \proofstepwidestarbreakkeep[17]{%
      \RealZero\RealLt A\lor A\RealLt\RealZero}{%
      \FormulaRefAuto{RealCutTotalOrder}[thm]{17}}
    \proofstepwidestar[17]{A\RealLt\RealZero\leftrightarrow\RealZero\RealLt\RealNeg A}{%
      \FormulaRefAuto{RealCutNegationOrder}[thm]{17}}
    \proofstepwidestar[17]{\RealNeg A\in\mathbb R}{%
      \FormulaRefAuto{RealCutAdditiveClosure}[thm]{17}}
    \proofstepwidestarbreakkeep[17]{\RealInv A\in\mathbb R}{%
      \FormulaRefAuto{RealCutReciprocal}[def]{%
        17,18,16,19,
        20}}
    \proofstepwidestarbreakkeep[1,2]{%
      A\RealPosMul B\in\mathbb R\land
      A\RealMul B\in\mathbb R\land\RealOne\in\mathbb R}{%
      \rAI{4,\rAI{9,13}}}
    \proofstepwidestarbreakkeep[14]{\RealPosInv A\in\mathbb R}{\rIE{16}}
    \proofstepwidestarbreakkeep[17]{\RealInv A\in\mathbb R}{\rIE{21}}
  \closeproofpartwide
\end{tabproofsplitwide}

\FormulaThmDeltaK[Multiplikative Randapproximation eines positiven Schnittes]{%
  \begin{aligned}[t]
  &A\in\mathbb R\dsep\RealZero\RealLt A\dsep
    q\in\mathbb Q\dsep\RatZero<q<\RatOne\vdash{}\\[-2pt]
  &\quad\exists a\in A\,\exists r\in\mathbb Q\,
    \bigl(\RatZero<a<r\land r\notin A\land q<a r^{-1}\bigr).
  \end{aligned}%
}{RealCutMultiplicativeBoundaryApproximation}{%
  \DeltaRow{Mengen}{A\dsep a\dsep r\dsep q}%
  \DeltaRow{Relationsoperatoren}{\RealLt}%
}
\begin{tabproofsplitwide}
  \proofpartwideindR[Wahl einer positiven Maschenweite]{%
    \begin{aligned}[t]
      &p,q\in\mathbb Q\dsep\RatZero<p\dsep
        \RatZero<q<\RatOne\\[-2pt]
      &\quad\vdash\exists h\in\mathbb Q\,
        (\RatZero<h\land qh<(\RatOne-q)p)
    \end{aligned}
  }{%
    p,q\in\mathbb Q\dsep\RatZero<p\dsep
    \RatZero<q<\RatOne
    \vdash\exists h\in\mathbb Q\,
    (\RatZero<h\land qh<(\RatOne-q)p)
  }[RealCutMultiplicativeMeshChoice]
    \proofstepwidestar[1]{p,q\in\mathbb Q}{\rA}
    \proofstepwidestar[2]{\RatZero<p}{\rA}
    \proofstepwidestar[3]{\RatZero<q<\RatOne}{\rA}
    \proofstepwidestar[1,2,3]{%
      \RatZero<(\RatOne-q)pq^{-1}}{%
      \FormulaRefAuto{RationalPositiveReciprocal}[thm]{1,2,3}}
    \proofstepwidestar[1,2,3]{%
      \exists h\in\mathbb Q\,
      (\RatZero<h\land h<(\RatOne-q)pq^{-1})}{%
      \FormulaRefAuto{RationalOrderDense}[thm]{4}}
    \proofstepwidestar[1,2,3]{%
      h\in\mathbb Q\land\RatZero<h\land
      h<(\RatOne-q)pq^{-1}}{\rA}
    \proofstepwidestar[1,2,3]{qh<(\RatOne-q)p}{%
      \FormulaRefAuto{RationalPositiveMulStrictOrder}[thm]{3,6}}
    \proofstepwidestar[1,2,3]{%
      \exists h\in\mathbb Q\,
      (\RatZero<h\land qh<(\RatOne-q)p)}{%
      \rEI{6,\rAI{6,7}}}
    \proofstepwidestar[1,2,3]{%
      \exists h\in\mathbb Q\,
      (\RatZero<h\land qh<(\RatOne-q)p)}{%
      \rEE{5,6,8}}
  \closeproofpartwideindR

  \proofpartwide{Gitterrand und Quotientenvergleich}
    \proofstepwidestarbreakkeep[1]{A\in\mathbb R}{\rA}
    \proofstepwidestarbreakkeep[2]{\RealZero\RealLt A}{\rA}
    \proofstepwidestarbreakkeep[3]{q\in\mathbb Q}{\rA}
    \proofstepwidestarbreakkeep[4]{\RatZero<q<\RatOne}{\rA}
    \proofstepwidestarbreakkeep[1,2]{\exists p\in A\,(\RatZero<p)}{%
      \FormulaRefAuto{RealCutOrderCharacterization}[thm]{1,2}}
    \proofstepwidestarbreakkeep[1,2]{p\in A\land\RatZero<p}{\rA}
    \proofstepwidestar[1]{\RealCut A}{%
      \FormulaRefAuto{RealCutMembership}[thm]{1}}
    \proofstepwidestarbreakkeep[1]{\exists u\in\mathbb Q\setminus A}{%
      \FormulaRefAuto{DedekindCutCharacterization}[thm]{%
        7}}
    \proofstepwidestarbreakkeep[1]{u\in\mathbb Q\setminus A}{\rA}
    \proofstepwidestarbreakkeep[1,2,3,4]{%
      \exists h\in\mathbb Q\,
      (\RatZero<h\land qh<(\RatOne-q)p)}{%
      \FormulaRefAuto{RealCutMultiplicativeMeshChoice}[thm]{3,4,6}}
    \proofstepwidestarbreakkeep[1,2,3,4]{%
      h\in\mathbb Q\land\RatZero<h\land qh<(\RatOne-q)p}{\rA}
    \proofstepwidestarbreakkeep[1,2,3,4]{%
      \begin{aligned}[t]
        &\exists m_0\in\mathbb N\setminus\{0\}\,\bigl(\\[-2pt]
        &\quad p+\RatEmbed(\IntEmbed(m_0))\frac h2\notin A\land{}\\[-2pt]
        &\quad p+\RatEmbed(\IntEmbed(m_0-1))\frac h2\in A\bigr)
      \end{aligned}}{%
      \FormulaRefAuto{RealCutAdditiveBoundaryFirstOutside}[thm]{1,11,6,9}}
    \proofstepwidestarbreakkeep[1,2,3,4]{%
      a\in A\land r\notin A\land p\leq a<r\land r-a<h}{%
      \FormulaRefAuto{RationalOrderedFieldAxioms}[def]{11,12}}
    \proofstepwidestarbreakkeep[1,2,3,4]{qr=qa+q(r-a)}{%
      \FormulaRefAuto{RationalDistributive}[thm]{3,13}}
    \proofstepwidestarbreakkeep[1,2,3,4]{%
      qa+q(r-a)<qa+(\RatOne-q)p}{%
      \FormulaRefAuto{RationalPositiveMulStrictOrder}[thm]{11,13}}
    \proofstepwidestarbreakkeep[1,2,3,4]{%
      qa+(\RatOne-q)p\leq qa+(\RatOne-q)a}{%
      \FormulaRefAuto{RationalOrderTranslation}[thm]{4,13}}
    \proofstepwidestarbreakkeep[1,2,3,4]{qa+(\RatOne-q)a=a}{%
      \FormulaRefAuto{RationalDistributive}[thm]{3,13}}
    \proofstepwidestarbreakkeep[1,2,3,4]{qr<a}{\rChain{14,15,16,17}}
    \proofstepwidestarbreakkeep[1,2,3,4]{q<ar^{-1}}{%
      \FormulaRefAuto{RationalPositiveMulStrictOrder}[thm]{13,18}}
    \proofstepwidestarbreakkeep[1,2,3,4]{%
      \exists a\in A\,\exists r\in\mathbb Q\,
      (\RatZero<a<r\land r\notin A\land q<ar^{-1})}{%
      \rEI{13,\rEI{13,\rAI{13,\rAI{13,19}}}}}
    \proofstepwidestarbreakkeep[1,2,3,4]{%
      \exists a\in A\,\exists r\in\mathbb Q\,
      (\RatZero<a<r\land r\notin A\land q<ar^{-1})}{%
      \rEE{12,13,20}}
    \proofstepwidestarbreakkeep[1,2,3,4]{%
      \exists a\in A\,\exists r\in\mathbb Q\,
      (\RatZero<a<r\land r\notin A\land q<ar^{-1})}{%
      \rEE{10,11,21}}
    \proofstepwidestarbreakkeep[1,2,3,4]{%
      \exists a\in A\,\exists r\in\mathbb Q\,
      (\RatZero<a<r\land r\notin A\land q<ar^{-1})}{%
      \rEE{8,9,22}}
    \proofstepwidestarbreakkeep[1,2,3,4]{%
      \exists a\in A\,\exists r\in\mathbb Q\,
      (\RatZero<a<r\land r\notin A\land q<ar^{-1})}{%
      \rEE{5,6,23}}
  \closeproofpartwide
\end{tabproofsplitwide}

\FormulaThmDeltaK[Zerlegung eines Punktes des Nullschnitts]{%
  q\in\RealZero\vdash\exists a\in\RealZero\,\exists b\in\RealZero\,(q<a+b)
}{RealCutZeroAdditionWitnesses}{%
  \DeltaRow{Mengen}{q\dsep a\dsep b}
}
\begin{tabproofwide}
    \proofstepwidestar[1]{q\in\RealZero}{%
      \rA}
    \proofstepwidestar[]{\RatZero\in\mathbb Q}{%
      \FormulaRefAuto{RationalZeroOneCarrier}[thm]}
    \proofstepwidestar[]{\RealCut{\RealHat{\RatZero}}}{%
      \FormulaRefAuto{RealPrincipalCutIsCut}[thm]{2}}
    \proofstepwidestar[]{\RealHat{\RatZero}\in\mathbb R}{%
      \FormulaRefAuto{RealCutMembership}[thm]{3}}
    \proofstepwidestar[]{\RealZero\in\mathbb R}{%
      \FormulaRefAuto{RealCutZero}[def]{4}}
    \proofstepwidestar[]{\RealZero\RealAdd\RealZero=\RealZero}{%
      \FormulaRefAuto{RealCutAdditionZero}[thm]{5}}
    \proofstepwidestar[1]{q\in\RealZero\RealAdd\RealZero}{%
      \rIE{\rIE{6,\rII},1}}
    \proofstepwidestar[1]{\exists a\in\RealZero\,\exists b\in\RealZero\,(q<a+b)}{%
      \FormulaRefAuto{RealCutAddition}[def]{5,5,7}}
\end{tabproofwide}

\FormulaThmDeltaKR[Gesetze der reellen Multiplikation]{%
  \begin{aligned}[t]
    A,B,C\in\mathbb R\vdash{}&
    \text{(i)}\quad (A\RealMul B)\RealMul C=A\RealMul(B\RealMul C),\\[-2pt]
    &\text{(ii)}\quad A\RealMul B=B\RealMul A,\\[-2pt]
    &\text{(iii)}\quad A\RealMul\RealOne=A,\\[-2pt]
    &\text{(iv)}\quad A\RealMul(B\RealAdd C)=(A\RealMul B)\RealAdd(A\RealMul C),\\[-2pt]
    &\text{(v)}\quad A\neq\RealZero\rightarrow A\RealMul\RealInv A=\RealOne.
  \end{aligned}%
}{%
  \begin{aligned}[t]
  A,B,C\in\mathbb R\vdash{}&
    (A\RealMul B)\RealMul C=A\RealMul(B\RealMul C)\dsep{}\\[-2pt]
  &A\RealMul B=B\RealMul A\dsep
    A\RealMul\RealOne=A\dsep{}\\[-2pt]
  &A\RealMul(B\RealAdd C)
      =(A\RealMul B)\RealAdd(A\RealMul C),\\[-2pt]
  &A\neq\RealZero\rightarrow A\RealMul\RealInv A=\RealOne.
  \end{aligned}%
}{RealCutMultiplicativeLaws}{%
  \DeltaRow{Mengen}{A\dsep B\dsep C}%
  \DeltaRow{Funktionsoperatoren}{\RealAdd\dsep\RealMul
    \dsep\RealInv{\,\cdot\,}}%
}
\begin{tabproofsplitwide}
\proofpartwide{Nichtnegativität der Summe}
\proofstepwidestar[1]{A,B\in\mathbb R}{\rA}
    \proofstepwidestar[2]{%
      \RealZero\RealLe A\land\RealZero\RealLe B}{\rA}
    \proofstepwidestarbreakkeep[1,2]{%
      \RealZero\subseteq A\land\RealZero\subseteq B}{%
      \FormulaRefAuto{RealCutOrderCharacterization}[thm]{1,2}}
    \proofstepwidestar[4]{q\in\RealZero}{\rA}
    \proofstepwidestarbreakkeep[4]{%
      \exists a\in\RealZero\,\exists b\in\RealZero\,(q<a+b)}{\FormulaRefAuto{RealCutZeroAdditionWitnesses}[thm]{4}}
    \proofstepwidestarbreakkeep[1,2,4]{%
      \exists a\in A\,\exists b\in B\,(q<a+b)}{%
      \FormulaRefAuto{A\subseteq B\dsep x\in A\vdash x\in B}[thm]{3,5}}
    \proofstepwidestar[1,2,4]{q\in A\RealAdd B}{%
      \FormulaRefAuto{RealCutAddition}[def]{1,6}}
    \proofstepwidestarbreakkeep[1,2]{%
      \RealZero\subseteq A\RealAdd B}{%
      \FormulaRefAuto{SubsetIntroductionRule}[thm]{[4]\vdots 7}}
    \proofstepwidestarbreakkeep[1]{A\RealAdd B\in\mathbb R}{%
      \FormulaRefAuto{RealCutAdditiveClosure}[thm]{1}}
    \proofstepwidestarbreakkeep[1,2]{\RealZero\RealLe A\RealAdd B}{%
      \FormulaRefAuto{RealCutOrderCharacterization}[thm]{8,9}}

\closeproofpartwide

\proofpartwide{Nullpunktkriterium}
  \setcounter{proofstepnr}{10}% Fortlaufende Schritte dieses Beweises.
\proofstepwidestar[11]{\RatZero\in A\RealAdd B}{\rA}
    \proofstepwidestarbreakkeep[1,11]{%
      \exists a\in A\,\exists b\in B\,(\RatZero<a+b)}{%
      \FormulaRefAuto{RealCutAddition}[def]{1,11}}
    \proofstepwidestar[1]{\RealCut A}{%
      \FormulaRefAuto{RealCutMembership}[thm]{\rAEa{1}}}
    \proofstepwidestarbreakkeep[1]{%
      a\in A\dsep\RatZero\notin A\vdash a<\RatZero}{\text{\parbox[t]{\linewidth}{\raggedright \ensuremath{a\in\mathbb Q.} Aus \ensuremath{\RatZero\leq\allowbreak  a} folgte bei \ensuremath{a=\allowbreak \RatZero} sofort, bei \ensuremath{\RatZero<\allowbreak a} durch Abwärtsabschluss jeweils \ensuremath{\RatZero\in A.} \FormulaRefAuto{DedekindCutCharacterization}[thm]{
        13}, \FormulaRefAuto{RationalOrderTotal}[thm]}}}
    \proofstepwidestar[1]{\RealCut B}{%
      \FormulaRefAuto{RealCutMembership}[thm]{\rAEb{1}}}
    \proofstepwidestarbreakkeep[1]{%
      b\in B\dsep\RatZero\notin B\vdash b<\RatZero}{\text{\parbox[t]{\linewidth}{\raggedright \ensuremath{b\in\mathbb Q.} Aus \ensuremath{\RatZero\leq\allowbreak  b} folgte wie in \ensuremath{14} durch Gleichheit oder Abwärtsabschluss \ensuremath{\RatZero\in B.} \FormulaRefAuto{DedekindCutCharacterization}[thm]{
        15}, \FormulaRefAuto{RationalOrderTotal}[thm]}}}
    \proofstepwidestarbreakkeep[1,11]{%
      \RatZero\notin A\land\RatZero\notin B\vdash\bot}{\text{\parbox[t]{\linewidth}{\raggedright Für die Zeugen aus \ensuremath{12} liefern \ensuremath{14} und \ensuremath{16} die Ungleichungen \ensuremath{a<\allowbreak 0,\allowbreak  b<\allowbreak 0,\allowbreak } also \ensuremath{a+b<\allowbreak 0}, im Widerspruch zu \ensuremath{0<\allowbreak a+b.} \FormulaRefAuto{RationalOrderedFieldAxioms}[def]{12,14,16}}}}
    \proofstepwidestarbreakkeep[1,11]{%
      \RatZero\in A\lor\RatZero\in B}{%
      \FormulaRefAuto{\neg(P\land Q) \eqvdash \neg P\lor\neg Q}[thm]{17},
      \FormulaRefAuto{rDN}[thm]}
    \proofstepwidestar[19]{\RatZero\in A}{\rA}
    \proofstepwidestar[1,19]{\RealCut A}{%
      \FormulaRefAuto{RealCutMembership}[thm]{\rAEa{1}}}
    \proofstepwidestarbreakkeep[1,19]{%
      \exists a\in A\,(\RatZero<a)}{%
      \FormulaRefAuto{DedekindCutCharacterization}[thm]{%
        20,19}}
    \proofstepwidestar[19]{a\in A\land\RatZero<a}{\rA}
    \proofstepwidestarbreakkeep[19]{%
      \exists b\in\RealZero\,(-a<b\land\RatZero<a+b)}{\text{\parbox[t]{\linewidth}{\raggedright \ensuremath{-a<\allowbreak \RatZero.} Wähle rational \ensuremath{b} mit \ensuremath{-a<\allowbreak b<\allowbreak \RatZero;} Addition von \ensuremath{a} ergibt \ensuremath{\RatZero<\allowbreak a+b.} \FormulaRefAuto{RationalOrderDense}[thm]{\rAEb{22}}, \FormulaRefAuto{RationalOrderTranslation}[thm] , \FormulaRefAuto{RealCutZero}[def] , \FormulaRefAuto{RealPrincipalCut}[def]}}}
    \proofstepwidestar[19]{%
      b\in\RealZero\land-a<b\land\RatZero<a+b}{\rA}
    \proofstepwidestarbreakkeep[1,2,19]{b\in B}{%
      \FormulaRefAuto{A\subseteq B\dsep x\in A\vdash x\in B}[thm]{%
        \rAEb{3},\rAEa{24}}}
    \proofstepwidestarbreakkeep[1,2,19]{\RatZero\in A\RealAdd B}{%
      \FormulaRefAuto{RealCutAddition}[def]{\rAEa{22},25,\rAEb{\rAEb{24}}}}
    \proofstepwidestarbreakkeep[1,2]{%
      \RatZero\in A\rightarrow\RatZero\in A\RealAdd B}{%
      \rRI{19,\rEE{21,22,\rEE{23,24,26}}}}
    \proofstepwidestarbreakkeep[1,2]{%
      \RatZero\in B\rightarrow\RatZero\in A\RealAdd B}{%
      \FormulaRefAuto{RealCutAdditionCommutative}[thm]{1,27}}
    \proofstepwidestarbreakkeep[1,2]{%
      (\RatZero\in A\lor\RatZero\in B)
      \rightarrow\RatZero\in A\RealAdd B}{%
      \rOE{27,28}}
    \proofstepwidestarbreakkeep[1,2]{%
      \RatZero\in A\RealAdd B\leftrightarrow
      (\RatZero\in A\lor\RatZero\in B)}{%
      \rBI{\rRI{11,18},29}}

\closeproofpartwide

\proofpartwideindR[Nullpunkt einer nichtnegativen Summe]{%
    \begin{aligned}[t]
    &A,B\in\mathbb R\dsep\RealZero\RealLe A\dsep\RealZero\RealLe B\\[-2pt]
    &\quad\vdash\RealZero\RealLe A\RealAdd B\land
      \bigl(\RatZero\in A\RealAdd B\leftrightarrow
        (\RatZero\in A\lor\RatZero\in B)\bigr)
    \end{aligned}
  }{%
    A,B\in\mathbb R\dsep\RealZero\RealLe A\dsep\RealZero\RealLe B
    \vdash\RealZero\RealLe A\RealAdd B\land
    (\RatZero\in A\RealAdd B\leftrightarrow
      (\RatZero\in A\lor\RatZero\in B))
  }[RealCutNonnegativeAdditionZeroMembership]
  \setcounter{proofstepnr}{30}% Fortlaufende Schritte dieses Beweises.
\proofstepwidestarbreakkeep[1,2]{%
      \RealZero\RealLe A\RealAdd B\land
      (\RatZero\in A\RealAdd B\leftrightarrow
       (\RatZero\in A\lor\RatZero\in B))}{\rAI{10,30}}

\closeproofpartwideindR


  \proofpartwideindR[Nullpunkt eines positiven Produkts]{%
    \begin{aligned}[t]
    &A,B\in\mathbb R\dsep\RealZero\RealLe A\dsep\RealZero\RealLe B\\[-2pt]
    &\quad\vdash
      \bigl(\RatZero\in A\RealPosMul B\leftrightarrow
        (\RatZero\in A\land\RatZero\in B)\bigr)
    \end{aligned}
  }{%
    A,B\in\mathbb R\dsep\RealZero\RealLe A\dsep\RealZero\RealLe B
    \vdash
    (\RatZero\in A\RealPosMul B\leftrightarrow
      (\RatZero\in A\land\RatZero\in B))
  }[RealCutPositiveProductZeroMembership]
    \proofstepwidestar[1]{A,B\in\mathbb R}{%
      \rA}
    \proofstepwidestar[2]{\RealZero\RealLe A\land\RealZero\RealLe B}{%
      \rA}
    \proofstepwidestar[1]{\RealCut A}{%
      \FormulaRefAuto{RealCutMembership}[thm]{\rAEa{1}}}
    \proofstepwidestar[1]{\RealCut B}{%
      \FormulaRefAuto{RealCutMembership}[thm]{\rAEb{1}}}
    \proofstepwidestar[]{\RatZero\in\mathbb Q}{%
      \FormulaRefAuto{RationalZeroCarrierAxiom}[ax]}
    \proofstepwidestar[6]{\RatZero<\RatZero}{%
      \rA}
    \proofstepwidestar[6]{\RatZero\neq\RatZero}{%
      \rAEb{\FormulaRefAuto{RationalStrictOrderDef}[def]{6}}}
    \proofstepwidestar[6]{\bot}{%
      \rBI{\rII,7}}
    \proofstepwidestar[]{\neg(\RatZero<\RatZero)}{%
      \rCI{6,8}}
    \proofstepwidestar[]{\RatZero\notin\RealHat{\RatZero}}{%
      \FormulaRefAuto{RealPrincipalCut}[def]{5,9}}
    \proofstepwidestar[]{\RatZero\notin\RealZero}{%
      \FormulaRefAuto{RealCutZero}[def]{10}}
    \proofstepwidestar[12]{\RatZero\in A\RealPosMul B}{%
      \rA}
    \proofstepwidestar[1,2,12]{\exists a\in A\,\exists b\in B\,(\RatZero<a\land\RatZero<b\land\RatZero<ab)}{%
      \FormulaRefAuto{RealCutPositiveProduct}[def]{1,2,11,12}}
    \proofstepwidestar[14]{a\in A\land\exists b\in B\,(\RatZero<a\land\RatZero<b\land\RatZero<ab)}{%
      \rA}
    \proofstepwidestar[15]{b\in B\land\RatZero<a\land\RatZero<b\land\RatZero<ab}{%
      \rA}
    \proofstepwidestar[1]{\forall p\in A\,\forall q\in\mathbb Q\,(q<p\rightarrow q\in A)}{%
      \FormulaRefAuto{DedekindCutDownwardAxiom}[ax]{3}}
    \proofstepwidestar[1]{\forall p\in B\,\forall q\in\mathbb Q\,(q<p\rightarrow q\in B)}{%
      \FormulaRefAuto{DedekindCutDownwardAxiom}[ax]{4}}
    \proofstepwidestar[1,14,15]{\RatZero\in A}{%
      \rRE{\rRE{\rUE{\rRE{\rUE{16},\rAEa{14}}},5},\rAEa{\rAEb{15}}}}
    \proofstepwidestar[1,15]{\RatZero\in B}{%
      \rRE{\rRE{\rUE{\rRE{\rUE{17},\rAEa{15}}},5},\rAEa{\rAEb{\rAEb{15}}}}}
    \proofstepwidestar[1,14,15]{\RatZero\in A\land\RatZero\in B}{%
      \rAI{18,19}}
    \proofstepwidestar[1,14]{\RatZero\in A\land\RatZero\in B}{%
      \rEE{\rAEb{14},15,20}}
    \proofstepwidestar[1,2,12]{\RatZero\in A\land\RatZero\in B}{%
      \rEE{13,14,21}}
    \proofstepwidestar[1,2]{\RatZero\in A\RealPosMul B\rightarrow(\RatZero\in A\land\RatZero\in B)}{%
      \rRI{12,22}}
    \proofstepwidestar[24]{\RatZero\in A\land\RatZero\in B}{%
      \rA}
    \proofstepwidestar[1]{\forall p\in A\,\exists a\in A\,(p<a)}{%
      \FormulaRefAuto{DedekindCutNoMaximumAxiom}[ax]{3}}
    \proofstepwidestar[1]{\forall p\in B\,\exists b\in B\,(p<b)}{%
      \FormulaRefAuto{DedekindCutNoMaximumAxiom}[ax]{4}}
    \proofstepwidestar[1,24]{\exists a\in A\,(\RatZero<a)}{%
      \rRE{\rUE{25},\rAEa{24}}}
    \proofstepwidestar[1,24]{\exists b\in B\,(\RatZero<b)}{%
      \rRE{\rUE{26},\rAEb{24}}}
    \proofstepwidestar[29]{a\in A\land\RatZero<a}{%
      \rA}
    \proofstepwidestar[30]{b\in B\land\RatZero<b}{%
      \rA}
    \proofstepwidestar[1,29]{a\in\mathbb Q}{%
      \FormulaRefAuto{DedekindCutCharacterization}[thm]{3,\rAEa{29}}}
    \proofstepwidestar[1,30]{b\in\mathbb Q}{%
      \FormulaRefAuto{DedekindCutCharacterization}[thm]{4,\rAEa{30}}}
    \proofstepwidestar[1,29,30]{a\RatZero<ab}{%
      \FormulaRefAuto{RationalPositiveMulStrictOrderAxiom}[ax]{31,5,32,\rAEb{29},\rAEb{30}}}
    \proofstepwidestar[1,29]{a\RatZero=\RatZero}{%
      \FormulaRefAuto{RationalMulZero}[thm]{31}}
    \proofstepwidestar[1,29,30]{\RatZero<ab}{%
      \rIE{34,33}}
    \proofstepwidestar[1,29,30]{\exists a\in A\,\exists b\in B\,(\RatZero<a\land\RatZero<b\land\RatZero<ab)}{%
      \rEI{\rAI{\rAEa{29},\rEI{\rAI{\rAEa{30},\rAI{\rAEb{29},\rAI{\rAEb{30},35}}}}}}}
    \proofstepwidestar[1,2,29,30]{\RatZero\in A\RealPosMul B}{%
      \FormulaRefAuto{RealCutPositiveProduct}[def]{1,2,5,36}}
    \proofstepwidestar[1,2,24,29]{\RatZero\in A\RealPosMul B}{%
      \rEE{28,30,37}}
    \proofstepwidestar[1,2,24]{\RatZero\in A\RealPosMul B}{%
      \rEE{27,29,38}}
    \proofstepwidestar[1,2]{(\RatZero\in A\land\RatZero\in B)\rightarrow\RatZero\in A\RealPosMul B}{%
      \rRI{24,39}}
    \proofstepwidestar[1,2]{\RatZero\in A\RealPosMul B\leftrightarrow(\RatZero\in A\land\RatZero\in B)}{%
      \rLRI{23,40}}
  \closeproofpartwideindR

  \proofpartwide{Assoziativität des positiven Produkts}
\proofstepwidestar[1]{A,B,C\in\mathbb R}{\rA}
    \proofstepwidestar[2]{%
      \RealZero\RealLe A\land\RealZero\RealLe B\land\RealZero\RealLe C}{\rA}
    \proofstepwidestarbreakkeep[1,2]{%
      \begin{aligned}[t]
      &A\RealPosMul B\in\mathbb R\land
        B\RealPosMul C\in\mathbb R\land{}\\[-2pt]
      &(A\RealPosMul B)\RealPosMul C\in\mathbb R\land
        A\RealPosMul(B\RealPosMul C)\in\mathbb R,\\[-2pt]
      &\RealZero\RealLe A\RealPosMul B\land
        \RealZero\RealLe B\RealPosMul C,\\[-2pt]
      &\RealZero\subseteq(A\RealPosMul B)\RealPosMul C\land
        \RealZero\subseteq A\RealPosMul(B\RealPosMul C)
      \end{aligned}}{%
      \FormulaRefAuto{RealCutMultiplicativeClosure}[thm]{1,2},
      \FormulaRefAuto{RealCutPositiveProduct}[def]{1,2},
      \FormulaRefAuto{RealCutOrderCharacterization}[thm]{1}}
    \proofstepwidestarbreakkeep[1,2]{%
      q<\RatZero\vdash
      \bigl(q\in(A\RealPosMul B)\RealPosMul C\leftrightarrow
        q\in A\RealPosMul(B\RealPosMul C)\bigr)}{%
      \FormulaRefAuto{RealCutPositiveProduct}[def]{3},
      \FormulaRefAuto{RealCutZero}[def]}
    \proofstepwidestarbreakkeep[1,2]{%
      \begin{aligned}[t]
      &\RatZero\in(A\RealPosMul B)\RealPosMul C\\[-2pt]
      &\quad\leftrightarrow
        ((\RatZero\in A\land\RatZero\in B)\land\RatZero\in C)\\[-2pt]
      &\quad\leftrightarrow
        (\RatZero\in A\land(\RatZero\in B\land\RatZero\in C))\\[-2pt]
      &\quad\leftrightarrow
        \RatZero\in A\RealPosMul(B\RealPosMul C)
      \end{aligned}}{%
      \FormulaRefAuto{RealCutPositiveProductZeroMembership}[thm]{1,2,3}}
    \proofstepwidestarbreakkeep[1,2]{%
      \begin{aligned}[t]
      &q\in(A\RealPosMul B)\RealPosMul C\dsep\RatZero<q\vdash{}\\[-2pt]
      &\quad q\in A\RealPosMul(B\RealPosMul C)
      \end{aligned}}{\text{\parbox[t]{\linewidth}{\raggedright Wähle zunächst \ensuremath{u\in A\RealPosMul\allowbreak  B,\allowbreak \ c\in C} mit \ensuremath{q<\allowbreak uc,\allowbreak } danach \ensuremath{a\in A,\allowbreak \ b\in B} mit \ensuremath{u<\allowbreak ab.} Aus \ensuremath{q<\allowbreak uc<\allowbreak abc} liefert die rationale Dichte \ensuremath{d<\allowbreak bc} mit \ensuremath{q<\allowbreak ad;} damit ist \ensuremath{d\in B\RealPosMul\allowbreak  C.} \FormulaRefAuto{RealCutPositiveProduct}[def]{1,2}, \FormulaRefAuto{RationalMulAssociative}[thm] , \FormulaRefAuto{RationalPositiveMulStrictOrder}[thm] , \FormulaRefAuto{RationalOrderDense}[thm]}}}
    \proofstepwidestarbreakkeep[1,2]{%
      \begin{aligned}[t]
      &q\in A\RealPosMul(B\RealPosMul C)\dsep\RatZero<q\vdash{}\\[-2pt]
      &\quad q\in(A\RealPosMul B)\RealPosMul C
      \end{aligned}}{\text{\parbox[t]{\linewidth}{\raggedright Wähle zunächst \ensuremath{a\in A,\allowbreak \ d\in B\RealPosMul\allowbreak  C} mit \ensuremath{q<\allowbreak ad,\allowbreak } danach \ensuremath{b\in B,\allowbreak \ c\in C} mit \ensuremath{d<\allowbreak bc.} Aus \ensuremath{q<\allowbreak ad<\allowbreak abc} liefert die rationale Dichte \ensuremath{u<\allowbreak ab} mit \ensuremath{q<\allowbreak uc;} damit ist \ensuremath{u\in A\RealPosMul\allowbreak  B.} \FormulaRefAuto{RealCutPositiveProduct}[def]{1,2}, \FormulaRefAuto{RationalMulAssociative}[thm] , \FormulaRefAuto{RationalPositiveMulStrictOrder}[thm] , \FormulaRefAuto{RationalOrderDense}[thm]}}}
    \proofstepwidestarbreakkeep[1,2]{%
      \forall q\in\mathbb Q\,
      \left(q\in(A\RealPosMul B)\RealPosMul C\leftrightarrow
      q\in A\RealPosMul(B\RealPosMul C)\right)}{%
      \FormulaRefAuto{RationalOrderTotal}[thm]{4,5,6,7},
      \FormulaRefAuto{RealCutPositiveProduct}[def]{1,2}}
    \proofstepwidestarbreakkeep[1,2]{%
      (A\RealPosMul B)\RealPosMul C=A\RealPosMul(B\RealPosMul C)}{%
      \FormulaRefAuto{\forall x(x\in A\leftrightarrow x\in B)\vdash A=B}{8}}

\closeproofpartwide

\proofpartwide{Kommutativität des positiven Produkts}
  \setcounter{proofstepnr}{9}% Fortlaufende Schritte dieses Beweises.
\proofstepwidestar[1,2]{\begin{aligned}[t]&A\RealPosMul B=\RealZero\cup{}\\[-2pt]&\{q\in\mathbb Q\mid\exists a\in A\,\exists b\in B\,(\RatZero<a\land\RatZero<b\land q<ab)\}\end{aligned}}{%
      \FormulaRefAuto{RealCutPositiveProduct}[def]{1,2}}
\proofstepwidestarbreakkeep[1,2]{%
      \forall q\in\mathbb Q\,
      (q\in A\RealPosMul B\leftrightarrow q\in B\RealPosMul A)}{%
      \FormulaRefAuto{RationalMulCommutative}[thm]{%
        10}}
    \proofstepwidestarbreakkeep[1,2]{A\RealPosMul B=B\RealPosMul A}{%
      \FormulaRefAuto{\forall x(x\in A\leftrightarrow x\in B)\vdash A=B}{11}}

\closeproofpartwide

\proofpartwide{Einheit des positiven Produkts}
  \setcounter{proofstepnr}{12}% Fortlaufende Schritte dieses Beweises.
\proofstepwidestarbreakkeep[1,2]{%
      \RealOne\in\mathbb R\land\RealZero\RealLe\RealOne\land
      \RatZero\in\RealOne}{%
      \FormulaRefAuto{RealCutOne}[def],
      \FormulaRefAuto{RationalZeroOneClasses}[thm],
      \FormulaRefAuto{RationalOnePositiveAxiom}[ax],
      \FormulaRefAuto{RationalToRealEmbeddingOrder}[thm]}
    \proofstepwidestarbreakkeep[1,2]{%
      q<\RatZero\vdash
      (q\in A\RealPosMul\RealOne\leftrightarrow q\in A)}{%
      \FormulaRefAuto{RealCutPositiveProduct}[def]{1,2},
      \FormulaRefAuto{RealCutOrderCharacterization}[thm]{1,2},
      \FormulaRefAuto{RealCutZero}[def]}
    \proofstepwidestarbreakkeep[1,2]{%
      \RatZero\in A\RealPosMul\RealOne\leftrightarrow\RatZero\in A}{%
      \FormulaRefAuto{RealCutPositiveProductZeroMembership}[thm]{1,2,13}}
    \proofstepwidestar[1,2]{\RealCut A}{%
      \FormulaRefAuto{RealCutMembership}[thm]{\rAEa{1}}}
    \proofstepwidestarbreakkeep[1,2]{%
      q\in A\RealPosMul\RealOne\dsep\RatZero<q\vdash q\in A}{\text{\parbox[t]{\linewidth}{\raggedright Die Produktdefinition liefert \ensuremath{a\in A} und \ensuremath{0<\allowbreak b<\allowbreak 1} mit \ensuremath{q<\allowbreak ab<\allowbreak a;} Abwärtsabschluss liefert \ensuremath{q\in A.} \FormulaRefAuto{RealCutPositiveProduct}[def]{1,2}, \FormulaRefAuto{RealCutOne}[def] , \FormulaRefAuto{RationalMulOne}[thm] , \FormulaRefAuto{DedekindCutCharacterization}[thm]{
           16}}}}
    \proofstepwidestar[1,2]{\RealCut A}{%
      \FormulaRefAuto{RealCutMembership}[thm]{\rAEa{1}}}
    \proofstepwidestarbreakkeep[1,2]{%
      q\in A\dsep\RatZero<q\vdash q\in A\RealPosMul\RealOne}{\text{\parbox[t]{\linewidth}{\raggedright Wähle \ensuremath{a\in A} mit \ensuremath{q<\allowbreak a} und dann rational \ensuremath{b} mit \ensuremath{qa^{-1}<\allowbreak b<\allowbreak \RatOne;} daraus folgt \ensuremath{q<\allowbreak ab.} \FormulaRefAuto{DedekindCutCharacterization}[thm]{
           18}, \FormulaRefAuto{RationalOrderDense}[thm] , \FormulaRefAuto{RealCutOne}[def] , \FormulaRefAuto{RationalPositiveMulStrictOrder}[thm]}}}
    \proofstepwidestarbreakkeep[1,2]{%
      \forall q\in\mathbb Q\,
      (q\in A\RealPosMul\RealOne\leftrightarrow q\in A)}{%
      \FormulaRefAuto{RationalOrderTotal}[thm]{14,15,17,19},
      \FormulaRefAuto{RealCutPositiveProduct}[def]{1,2}}
    \proofstepwidestarbreakkeep[1,2]{A\RealPosMul\RealOne=A}{%
      \FormulaRefAuto{\forall x(x\in A\leftrightarrow x\in B)\vdash A=B}{20}}

\closeproofpartwide

\proofpartwideindR[Positive Produktgesetze]{%
    \begin{aligned}[t]
      &A,B,C\in\mathbb R\dsep
        \RealZero\RealLe A\dsep\RealZero\RealLe B\dsep\RealZero\RealLe C\\[-2pt]
      &\quad\vdash
        (A\RealPosMul B)\RealPosMul C=A\RealPosMul(B\RealPosMul C)\land{}\\[-2pt]
      &\qquad A\RealPosMul B=B\RealPosMul A\land
        A\RealPosMul\RealOne=A
    \end{aligned}
  }{%
    A,B,C\in\mathbb R\dsep
    \RealZero\RealLe A\dsep\RealZero\RealLe B\dsep\RealZero\RealLe C
    \vdash
    (A\RealPosMul B)\RealPosMul C=A\RealPosMul(B\RealPosMul C)\land
    A\RealPosMul B=B\RealPosMul A\land
    A\RealPosMul\RealOne=A
  }[RealCutPositiveProductLaws]
  \setcounter{proofstepnr}{21}% Fortlaufende Schritte dieses Beweises.
\proofstepwidestarbreakkeep[1,2]{%
      \begin{aligned}[t]
      &(A\RealPosMul B)\RealPosMul C=A\RealPosMul(B\RealPosMul C)\land{}\\[-2pt]
      &A\RealPosMul B=B\RealPosMul A\land{}\\[-2pt]
      &A\RealPosMul\RealOne=A
      \end{aligned}}{\rAI{9,\rAI{12,21}}}

\closeproofpartwideindR


  \proofpartwideindR[Positive Distributivität]{%
    \begin{aligned}[t]
      &A,B,C\in\mathbb R\dsep
        \RealZero\RealLe A\dsep\RealZero\RealLe B\dsep\RealZero\RealLe C\\[-2pt]
      &\quad\vdash
        A\RealPosMul(B\RealAdd C)=
        (A\RealPosMul B)\RealAdd(A\RealPosMul C)
    \end{aligned}
  }{%
    A,B,C\in\mathbb R\dsep
    \RealZero\RealLe A\dsep\RealZero\RealLe B\dsep\RealZero\RealLe C
    \vdash
    A\RealPosMul(B\RealAdd C)=
    (A\RealPosMul B)\RealAdd(A\RealPosMul C)
  }[RealCutPositiveDistributive]
    \proofstepwidestar[1]{A,B,C\in\mathbb R}{\rA}
    \proofstepwidestar[2]{%
      \RealZero\RealLe A\land\RealZero\RealLe B\land\RealZero\RealLe C}{\rA}
    \proofstepwidestarbreakkeep[1,2]{%
      A=\RealZero\lor B=\RealZero\lor C=\RealZero\vdash
      A\RealPosMul(B\RealAdd C)=
      (A\RealPosMul B)\RealAdd(A\RealPosMul C)}{\text{\parbox[t]{\linewidth}{\raggedright Ist \ensuremath{A=\allowbreak 0,\allowbreak } sind beide Produkte rechts und das Produkt links \ensuremath{\RealZero.} Ist \ensuremath{B=\allowbreak 0} oder \ensuremath{C=\allowbreak 0,\allowbreak } verschwinden genau die
        entsprechenden Faktoren; die übrigen Summanden werden durch das additive Nullgesetz
        identifiziert. \FormulaRefAuto{RealCutPositiveProduct}[def]{1,2}, \FormulaRefAuto{RealCutAdditiveLaws}[thm]{1}}}}
    \proofstepwidestarbreakkeep[1,2]{\begin{aligned}[t]
      &B\RealAdd C\in\mathbb R\land A\RealPosMul B\in\mathbb R\land{}\\[-2pt]
      &A\RealPosMul C\in\mathbb R\land{}\\[-2pt]
      &A\RealPosMul(B\RealAdd C)\in\mathbb R\land{}\\[-2pt]
      &(A\RealPosMul B)\RealAdd(A\RealPosMul C)\in\mathbb R,\\[-2pt]
      &\RealZero\RealLe B\RealAdd C\land{}\\[-2pt]
      &\RealZero\RealLe A\RealPosMul B\land\RealZero\RealLe A\RealPosMul C\land{}\\[-2pt]
      &\RealZero\RealLe A\RealPosMul(B\RealAdd C)\land{}\\[-2pt]
      &\RealZero\RealLe(A\RealPosMul B)\RealAdd(A\RealPosMul C)
      \end{aligned}}{%
      \FormulaRefAuto{RealCutNonnegativeAdditionZeroMembership}[thm]{1,2},
      \FormulaRefAuto{RealCutMultiplicativeClosure}[thm]{1,2},
      \FormulaRefAuto{RealCutPositiveProduct}[def]{1,2},
      \FormulaRefAuto{RealCutOrderCharacterization}[thm]{1}}
    \proofstepwidestarbreakkeep[1,2]{%
      q<\RatZero\vdash
      \bigl(q\in A\RealPosMul(B\RealAdd C)\leftrightarrow
        q\in(A\RealPosMul B)\RealAdd(A\RealPosMul C)\bigr)}{%
      \FormulaRefAuto{RealCutOrderCharacterization}[thm]{4},
      \FormulaRefAuto{RealCutZero}[def]}
    \proofstepwidestarbreakkeep[1,2]{%
      \begin{aligned}[t]
      &\RatZero\in A\RealPosMul(B\RealAdd C)\\[-2pt]
      &\quad\leftrightarrow
        (\RatZero\in A\land
          (\RatZero\in B\lor\RatZero\in C))\\[-2pt]
      &\quad\leftrightarrow
        ((\RatZero\in A\land\RatZero\in B)\lor
          (\RatZero\in A\land\RatZero\in C))\\[-2pt]
      &\quad\leftrightarrow
        \RatZero\in(A\RealPosMul B)\RealAdd(A\RealPosMul C)
      \end{aligned}}{%
      \FormulaRefAuto{RealCutPositiveProductZeroMembership}[thm]{1,2,4},
      \FormulaRefAuto{RealCutNonnegativeAdditionZeroMembership}[thm]{1,2,4}}
    \proofstepwidestarbreakkeep[1,2]{%
      \begin{aligned}[t]
      &\RealZero\RealLt A\dsep\RealZero\RealLt B\dsep
        \RealZero\RealLt C\dsep{}\\[-2pt]
      &q\in A\RealPosMul(B\RealAdd C)\dsep\RatZero<q\vdash{}\\[-2pt]
      &\quad q\in(A\RealPosMul B)\RealAdd(A\RealPosMul C)
      \end{aligned}}{\text{\parbox[t]{\linewidth}{\raggedright Aus \ensuremath{q<\allowbreak a d} mit \ensuremath{a\in A,\allowbreak \ d\in B\RealAdd\allowbreak  C} wähle \ensuremath{b\in B,\allowbreak c\in C} so, dass \ensuremath{d<\allowbreak b+c.} Dann ist \ensuremath{q<\allowbreak a(b+c)=\allowbreak ab+ac.} Rationale Dichte und Offenheit liefern \ensuremath{u>\allowbreak ab,\allowbreak \ v>\allowbreak ac} mit \ensuremath{u\in A\RealPosMul\allowbreak  B,\allowbreak \ v\in A\RealPosMul\allowbreak  C} und weiterhin \ensuremath{q<\allowbreak u+v.} \FormulaRefAuto{RealCutPositiveProduct}[def]{1,2}, \FormulaRefAuto{RealCutAddition}[def]{1}, \FormulaRefAuto{RationalDistributive}[thm] , \FormulaRefAuto{RationalOrderDense}[thm]}}}
    \proofstepwidestarbreakkeep[1,2]{%
      \begin{aligned}[t]
      &\RealZero\RealLt A\dsep\RealZero\RealLt B\dsep
        \RealZero\RealLt C\dsep{}\\[-2pt]
      &q\in(A\RealPosMul B)\RealAdd(A\RealPosMul C)
        \dsep\RatZero<q\vdash{}\\[-2pt]
      &\quad q\in A\RealPosMul(B\RealAdd C)
      \end{aligned}}{\text{\parbox[t]{\linewidth}{\raggedright Schreibe \ensuremath{q<\allowbreak u+v} mit \ensuremath{u<\allowbreak ab} und \ensuremath{v<\allowbreak a'c.} Wähle wegen der Offenheit von \ensuremath{A} ein \ensuremath{a''\in A} mit \ensuremath{a,\allowbreak a'<\allowbreak a'';} wähle ferner \ensuremath{b'>\allowbreak b} in \ensuremath{B,\allowbreak
          \ c'>\allowbreak c} in \ensuremath{C} und \ensuremath{d\in B\RealAdd\allowbreak  C} mit \ensuremath{b'+c'<\allowbreak d.} Dann folgt \ensuremath{q<\allowbreak ab+a'c<\allowbreak a''(b'+c')<\allowbreak a''d.} \FormulaRefAuto{RealCutPositiveProduct}[def]{1,2}, \FormulaRefAuto{RealCutAddition}[def]{1}, \FormulaRefAuto{DedekindCutCharacterization}[thm] , \FormulaRefAuto{RationalDistributive}[thm] , \FormulaRefAuto{RationalOrderDense}[thm] , \FormulaRefAuto{RationalPositiveMulStrictOrder}[thm]}}}
    \proofstepwidestarbreakkeep[1,2]{\begin{aligned}[t]
      &\RealZero\RealLt A\dsep\RealZero\RealLt B\dsep\RealZero\RealLt C\vdash{}\\[-2pt]
      &\forall q\in\mathbb Q\,\bigl(q\in A\RealPosMul(B\RealAdd C)\\[-2pt]
      &\quad\leftrightarrow q\in(A\RealPosMul B)\RealAdd(A\RealPosMul C)\bigr)
      \end{aligned}}{%
      \FormulaRefAuto{RationalOrderTotal}[thm]{5,6,7,8},
      \FormulaRefAuto{RealCutPositiveProduct}[def]{1,2},
      \FormulaRefAuto{RealCutAddition}[def]{1}}
    \proofstepwidestarbreakkeep[1,2]{%
      \begin{aligned}[t]
      &\RealZero\RealLt A\dsep\RealZero\RealLt B\dsep
        \RealZero\RealLt C\vdash{}\\[-2pt]
      &A\RealPosMul(B\RealAdd C)=
        (A\RealPosMul B)\RealAdd(A\RealPosMul C)
      \end{aligned}}{%
      \FormulaRefAuto{\forall x(x\in A\leftrightarrow x\in B)\vdash A=B}{9}}
    \proofstepwidestarbreakkeep[1,2]{%
      \begin{aligned}[t]
      &(A=\RealZero\lor\RealZero\RealLt A)\land
        (B=\RealZero\lor\RealZero\RealLt B)\land{}\\[-2pt]
      &(C=\RealZero\lor\RealZero\RealLt C)
      \end{aligned}}{%
      \FormulaRefAuto{RealCutTotalOrder}[thm]{1,2}}
    \proofstepwidestarbreakkeep[1,2]{%
      A\RealPosMul(B\RealAdd C)=
      (A\RealPosMul B)\RealAdd(A\RealPosMul C)}{\text{\parbox[t]{\linewidth}{\raggedright \ensuremath{11} zerlegt jeden der drei nichtnegativen Faktoren in Null- und
        strikt positiven Fall. Sobald ein Nullfall eintritt, gilt \ensuremath{3;} der einzige
        verbleibende Fall ist \ensuremath{10.}}}}
  \closeproofpartwideindR

  \proofpartwide{Ein negatives Vorzeichen}
\proofstepwidestar[1]{A,B\in\mathbb R}{\rA}
    \proofstepwidestar[1]{\RealNeg A,\RealNeg B\in\mathbb R}{%
      \FormulaRefAuto{RealCutAdditiveClosure}[thm]{1}}
    \proofstepwidestar[1]{%
      A=\RealZero\lor\RealZero\RealLt A\lor A\RealLt\RealZero}{%
      \FormulaRefAuto{RealCutTotalOrder}[thm]{1}}
    \proofstepwidestar[1]{%
      B=\RealZero\lor\RealZero\RealLt B\lor B\RealLt\RealZero}{%
      \FormulaRefAuto{RealCutTotalOrder}[thm]{1}}
    \proofstepwidestarbreakkeep[1]{%
      \begin{aligned}[t]
      A=\RealZero\lor B=\RealZero\quad\vdash\quad
      &A\RealMul B=\RealZero,\\[-2pt]
      &(\RealNeg A)\RealMul B=\RealNeg{(A\RealMul B)},\\[-2pt]
      &A\RealMul(\RealNeg B)=\RealNeg{(A\RealMul B)}
      \end{aligned}}{%
      \FormulaRefAuto{RealCutMultiplication}[def]{1},
      \FormulaRefAuto{RealCutNegationOrder}[thm],
      \FormulaRefAuto{RealCutPositiveProduct}[def]}
    \proofstepwidestarbreakkeep[1]{%
      \begin{aligned}[t]
      &\RealZero\RealLt A\dsep\RealZero\RealLt B\vdash{}\\[-2pt]
      &(\RealNeg A)\RealMul B
        =\RealNeg{(A\RealPosMul B)}
        =\RealNeg{(A\RealMul B)},\\[-2pt]
      &A\RealMul(\RealNeg B)
        =\RealNeg{(A\RealPosMul B)}
        =\RealNeg{(A\RealMul B)}
      \end{aligned}}{%
      \FormulaRefAuto{RealCutMultiplication}[def]{1},
      \FormulaRefAuto{RealCutNegationOrder}[thm]{1}}
    \proofstepwidestarbreakkeep[1]{%
      \begin{aligned}[t]
      &\RealZero\RealLt A\dsep B\RealLt\RealZero\vdash{}\\[-2pt]
      &(\RealNeg A)\RealMul B
        =(A\RealPosMul\RealNeg B)
        =\RealNeg{(A\RealMul B)},\\[-2pt]
      &A\RealMul(\RealNeg B)
        =A\RealPosMul(\RealNeg B)
        =\RealNeg{(A\RealMul B)}
      \end{aligned}}{%
      \FormulaRefAuto{RealCutMultiplication}[def]{1},
      \FormulaRefAuto{RealCutNegationOrder}[thm]{1}}
    \proofstepwidestarbreakkeep[1]{%
      \begin{aligned}[t]
      &A\RealLt\RealZero\dsep\RealZero\RealLt B\vdash{}\\[-2pt]
      &(\RealNeg A)\RealMul B
        =(\RealNeg A)\RealPosMul B
        =\RealNeg{(A\RealMul B)},\\[-2pt]
      &A\RealMul(\RealNeg B)
        =(\RealNeg A)\RealPosMul B
        =\RealNeg{(A\RealMul B)}
      \end{aligned}}{%
      \FormulaRefAuto{RealCutMultiplication}[def]{1},
      \FormulaRefAuto{RealCutNegationOrder}[thm]{1},
      \FormulaRefAuto{RealCutPositiveProductLaws}[thm]{1}}
    \proofstepwidestarbreakkeep[1]{%
      \begin{aligned}[t]
      &A\RealLt\RealZero\dsep B\RealLt\RealZero\vdash{}\\[-2pt]
      &(\RealNeg A)\RealMul B
        =\RealNeg{((\RealNeg A)\RealPosMul(\RealNeg B))}
        =\RealNeg{(A\RealMul B)},\\[-2pt]
      &A\RealMul(\RealNeg B)
        =\RealNeg{((\RealNeg A)\RealPosMul(\RealNeg B))}
        =\RealNeg{(A\RealMul B)}
      \end{aligned}}{%
      \FormulaRefAuto{RealCutMultiplication}[def]{1},
      \FormulaRefAuto{RealCutNegationOrder}[thm]{1}}
    \proofstepwidestarbreakkeep[1]{%
      (\RealNeg A)\RealMul B=\RealNeg{(A\RealMul B)}\land
      A\RealMul(\RealNeg B)=\RealNeg{(A\RealMul B)}}{%
      \rOE{3,4,5,6,7,8,9}}

\closeproofpartwide

\proofpartwide{Zwei negative Vorzeichen}
  \setcounter{proofstepnr}{10}% Fortlaufende Schritte dieses Beweises.
\proofstepwidestar[1]{A\RealMul B\in\mathbb R}{%
      \FormulaRefAuto{RealCutMultiplicativeClosure}[thm]{1}}
\proofstepwidestarbreakkeep[1]{%
      \RealNeg{\RealNeg{(A\RealMul B)}}=A\RealMul B}{%
      \FormulaRefAuto{RealCutNegationOrder}[thm]{%
        11}}
    \proofstepwidestarbreakkeep[1]{%
      (\RealNeg A)\RealMul(\RealNeg B)=A\RealMul B}{%
      \rIE{10,12}}

\closeproofpartwide

\proofpartwide{Multiplikation mit dem Nullschnitt}
  \setcounter{proofstepnr}{13}% Fortlaufende Schritte dieses Beweises.
\proofstepwidestarbreakkeep[1]{A\RealMul\RealZero=\RealZero}{%
      \FormulaRefAuto{RealCutMultiplication}[def]{1},
      \FormulaRefAuto{RealCutPositiveProduct}[def],
      \FormulaRefAuto{RealCutNegationOrder}[thm]}

\closeproofpartwide

\proofpartwideindR[Vorzeichenidentitäten]{%
    A,B\in\mathbb R\vdash
    (\RealNeg A)\RealMul B=\RealNeg{(A\RealMul B)}\land
    (\RealNeg A)\RealMul(\RealNeg B)=A\RealMul B\land
    A\RealMul\RealZero=\RealZero
  }{%
    A,B\in\mathbb R\vdash
    (\RealNeg A)\RealMul B=\RealNeg{(A\RealMul B)}\land
    (\RealNeg A)\RealMul(\RealNeg B)=A\RealMul B\land
    A\RealMul\RealZero=\RealZero
  }[RealCutMultiplicationSignRules]
  \setcounter{proofstepnr}{14}% Fortlaufende Schritte dieses Beweises.
\proofstepwidestarbreakkeep[1]{%
      (\RealNeg A)\RealMul B=\RealNeg{(A\RealMul B)}\land
      (\RealNeg A)\RealMul(\RealNeg B)=A\RealMul B\land
      A\RealMul\RealZero=\RealZero}{%
      \rAI{\rAEa{10},\rAI{13,14}}}

\closeproofpartwideindR


  \proofpartwideindR[Positives Produktinverses]{%
    A\in\mathbb R\dsep\RealZero\RealLt A
    \vdash A\RealPosMul\RealPosInv A=\RealOne
  }{%
    A\in\mathbb R\dsep\RealZero\RealLt A
    \vdash A\RealPosMul\RealPosInv A=\RealOne
  }[RealCutPositiveProductInverse]
    \proofstepwidestarbreakkeep[1]{A\in\mathbb R}{\rA}
    \proofstepwidestarbreakkeep[2]{\RealZero\RealLt A}{\rA}
    \proofstepwidestarbreakkeep[3]{q\in A\RealPosMul\RealPosInv A}{\rA}
    \proofstepwidestar[1,2,3]{p,q,r\in\mathbb Q\dsep q<r\dsep\RatZero<p\vdash p\cdot q<p\cdot r}{%
      \FormulaRefAuto{RationalPositiveMulStrictOrder}[thm]}
    \proofstepwidestar[1,2,3]{\begin{aligned}[t]&\RealPosInv A=\RealZero\cup{}\\[-2pt]&\{q\in\mathbb Q\mid\exists r\in\mathbb Q\,(\RatZero<r\land r\notin A\land q<r^{-1})\}\end{aligned}}{%
      \FormulaRefAuto{RealCutPositiveReciprocal}[def]{1,2,3}}
    \proofstepwidestarbreakkeep[1,2,3]{q\in\RealOne}{%
      \FormulaRefAuto{RealCutPositiveProduct}[def]{%
        5,
        4}}
    \proofstepwidestarbreakkeep[1,2]{A\RealPosMul\RealPosInv A\subseteq\RealOne}{%
      \FormulaRefAuto{SubsetIntroductionRule}[thm]{[3]\vdots 6}}
    \proofstepwidestarbreakkeep[8]{q\in\RealOne}{\rA}
    \proofstepwidestar[8]{q<\RatOne}{%
      \FormulaRefAuto{RealPrincipalCut}[def]{8}}
    \proofstepwidestarbreakkeep[8]{q<\RatOne}{%
      \FormulaRefAuto{RealCutOne}[def]{%
        9}}
    \proofstepwidestarbreakkeep[1,2,8]{%
      q<\RatZero\lor q=\RatZero\lor(\RatZero<q<\RatOne)}{%
      \FormulaRefAuto{RationalOrderTotal}[thm]{10}}
    \proofstepwidestarbreakkeep[12]{q<\RatZero}{\rA}
    \proofstepwidestarbreakkeep[1,2,12]{q\in A\RealPosMul\RealPosInv A}{%
      \FormulaRefAuto{RealCutPositiveProduct}[def]{1,2},
      \FormulaRefAuto{RealCutZero}[def]{12}}
    \proofstepwidestarbreakkeep[14]{q=\RatZero}{\rA}
    \proofstepwidestar[1,2,14]{\RealCut A}{%
      \FormulaRefAuto{RealCutMembership}[thm]{1}}
    \proofstepwidestarbreakkeep[1,2,14]{%
      \exists a\in A\,\exists r\in\mathbb Q\setminus A\,
      (\RatZero<a<r)}{\text{\parbox[t]{\linewidth}{\raggedright \FormulaRefAuto{RealCutPositiveMember}[thm]{1,2}, \FormulaRefAuto{DedekindCutCharacterization}[thm]{
          15}, \FormulaRefAuto{RealCutMemberBelowExterior}[thm]{1}}}}
    \proofstepwidestarbreakkeep[16]{%
      a\in A\land r\in\mathbb Q\setminus A\land\RatZero<a<r}{\rA}
    \proofstepwidestarbreakkeep[16]{%
      \exists b\in\mathbb Q\,(\RatZero<b<r^{-1})}{%
      \FormulaRefAuto{RationalPositiveReciprocal}[thm]{17},
      \FormulaRefAuto{RationalOrderDense}[thm]{17}}
    \proofstepwidestarbreakkeep[16]{%
      b\in\mathbb Q\land\RatZero<b<r^{-1}}{\rA}
    \proofstepwidestarbreakkeep[1,2,16]{%
      b\in\RealPosInv A\land\RatZero<ab}{%
      \FormulaRefAuto{RealCutPositiveReciprocal}[def]{17,19},
      \FormulaRefAuto{RationalOrderedFieldAxioms}[def]{17,19}}
    \proofstepwidestarbreakkeep[1,2,14,16]{%
      q\in A\RealPosMul\RealPosInv A}{%
      \FormulaRefAuto{RealCutPositiveProduct}[def]{14,17,20}}
    \proofstepwidestarbreakkeep[22]{\RatZero<q<\RatOne}{\rA}
    \proofstepwidestarbreakkeep[1,2,22]{%
      \exists a\in A\,\exists r\notin A\,(\RatZero<a<r\land q<ar^{-1})}{%
      \FormulaRefAuto{RealCutMultiplicativeBoundaryApproximation}[thm]{1,2,22}}
    \proofstepwidestarbreakkeep[1,2,22]{%
      \exists b\in\mathbb Q\,(qa^{-1}<b\land b<r^{-1})}{%
      \FormulaRefAuto{RationalOrderDense}[thm]{23}}
    \proofstepwidestarbreakkeep[1,2,22]{b\in\RealPosInv A\land q<ab}{%
      \FormulaRefAuto{RealCutPositiveReciprocal}[def]{23,24}}
    \proofstepwidestarbreakkeep[1,2,22]{q\in A\RealPosMul\RealPosInv A}{%
      \FormulaRefAuto{RealCutPositiveProduct}[def]{23,25}}
    \proofstepwidestarbreakkeep[1,2,8]{q\in A\RealPosMul\RealPosInv A}{%
      \rOE{11,12,13,14,21,22,26}}
    \proofstepwidestarbreakkeep[1,2]{\RealOne\subseteq A\RealPosMul\RealPosInv A}{%
      \FormulaRefAuto{SubsetIntroductionRule}[thm]{[8]\vdots 27}}
    \proofstepwidestarbreakkeep[1,2]{A\RealPosMul\RealPosInv A=\RealOne}{%
      \FormulaRefAuto{A\subseteq B\dsep B\subseteq A\vdash A=B}{7,28}}
  \closeproofpartwideindR

  \proofpartwide{Assoziativität}
\proofstepwidestar[1]{A,B,C\in\mathbb R}{\rA}
    \proofstepwidestarbreakkeep[1]{%
      \begin{aligned}[t]
      &(A=\RealZero\lor\RealZero\RealLt A\lor A\RealLt\RealZero)\land{}\\[-2pt]
      &(B=\RealZero\lor\RealZero\RealLt B\lor B\RealLt\RealZero)\land{}\\[-2pt]
      &(C=\RealZero\lor\RealZero\RealLt C\lor C\RealLt\RealZero)
      \end{aligned}}{%
      \FormulaRefAuto{RealCutTotalOrder}[thm]{1}}
    \proofstepwidestarbreakkeep[1]{%
      A=\RealZero\lor B=\RealZero\lor C=\RealZero\vdash
      (A\RealMul B)\RealMul C=A\RealMul(B\RealMul C)}{%
      \FormulaRefAuto{RealCutMultiplicationSignRules}[thm]{1}}
    \proofstepwidestarbreakkeep[1]{%
      \begin{aligned}[t]
      &\RealZero\RealLt A\dsep\RealZero\RealLt B\dsep
        \RealZero\RealLt C\vdash{}\\[-2pt]
      &(A\RealMul B)\RealMul C
       =(A\RealPosMul B)\RealPosMul C\\[-2pt]
      &\quad=A\RealPosMul(B\RealPosMul C)
       =A\RealMul(B\RealMul C)
      \end{aligned}}{%
      \FormulaRefAuto{RealCutMultiplication}[def]{1},
      \FormulaRefAuto{RealCutPositiveProductLaws}[thm]{1}}
    \proofstepwidestarbreakkeep[1]{%
      \begin{aligned}[t]
      &A\RealLt\RealZero\dsep\RealZero\RealLt B\dsep
        \RealZero\RealLt C\vdash{}\\[-2pt]
      &(A\RealMul B)\RealMul C
       =\RealNeg{(((\RealNeg A)\RealPosMul B)\RealPosMul C)}\\[-2pt]
      &\quad=\RealNeg{((\RealNeg A)\RealPosMul(B\RealPosMul C))}
       =A\RealMul(B\RealMul C)
      \end{aligned}}{%
      \FormulaRefAuto{RealCutMultiplication}[def]{1},
      \FormulaRefAuto{RealCutPositiveProductLaws}[thm]{1},
      \FormulaRefAuto{RealCutMultiplicationSignRules}[thm]{1}}
    \proofstepwidestarbreakkeep[1]{%
      \begin{aligned}[t]
      &B\RealLt\RealZero\dsep\RealZero\RealLt A\dsep
        \RealZero\RealLt C\vdash{}\\[-2pt]
      &(A\RealMul B)\RealMul C
       =\RealNeg{((A\RealPosMul\RealNeg B)\RealPosMul C)}\\[-2pt]
      &\quad=\RealNeg{(A\RealPosMul((\RealNeg B)\RealPosMul C))}
       =A\RealMul(B\RealMul C)
      \end{aligned}}{%
      \FormulaRefAuto{RealCutMultiplication}[def]{1},
      \FormulaRefAuto{RealCutPositiveProductLaws}[thm]{1},
      \FormulaRefAuto{RealCutMultiplicationSignRules}[thm]{1}}
    \proofstepwidestarbreakkeep[1]{%
      \begin{aligned}[t]
      &C\RealLt\RealZero\dsep\RealZero\RealLt A\dsep
        \RealZero\RealLt B\vdash{}\\[-2pt]
      &(A\RealMul B)\RealMul C
       =\RealNeg{((A\RealPosMul B)\RealPosMul\RealNeg C)}\\[-2pt]
      &\quad=\RealNeg{(A\RealPosMul(B\RealPosMul\RealNeg C))}
       =A\RealMul(B\RealMul C)
      \end{aligned}}{%
      \FormulaRefAuto{RealCutMultiplication}[def]{1},
      \FormulaRefAuto{RealCutPositiveProductLaws}[thm]{1},
      \FormulaRefAuto{RealCutMultiplicationSignRules}[thm]{1}}
    \proofstepwidestarbreakkeep[1]{%
      \begin{aligned}[t]
      &A\RealLt\RealZero\dsep B\RealLt\RealZero\dsep
        \RealZero\RealLt C\vdash{}\\[-2pt]
      &(A\RealMul B)\RealMul C
       =((\RealNeg A)\RealPosMul(\RealNeg B))\RealPosMul C\\[-2pt]
      &\quad=(\RealNeg A)\RealPosMul((\RealNeg B)\RealPosMul C)
       =A\RealMul(B\RealMul C)
      \end{aligned}}{%
      \FormulaRefAuto{RealCutMultiplication}[def]{1},
      \FormulaRefAuto{RealCutPositiveProductLaws}[thm]{1}}
    \proofstepwidestarbreakkeep[1]{%
      \begin{aligned}[t]
      &A\RealLt\RealZero\dsep C\RealLt\RealZero\dsep
        \RealZero\RealLt B\vdash{}\\[-2pt]
      &(A\RealMul B)\RealMul C
       =((\RealNeg A)\RealPosMul B)\RealPosMul(\RealNeg C)\\[-2pt]
      &\quad=(\RealNeg A)\RealPosMul(B\RealPosMul\RealNeg C)
       =A\RealMul(B\RealMul C)
      \end{aligned}}{%
      \FormulaRefAuto{RealCutMultiplication}[def]{1},
      \FormulaRefAuto{RealCutPositiveProductLaws}[thm]{1}}
    \proofstepwidestarbreakkeep[1]{%
      \begin{aligned}[t]
      &B\RealLt\RealZero\dsep C\RealLt\RealZero\dsep
        \RealZero\RealLt A\vdash{}\\[-2pt]
      &(A\RealMul B)\RealMul C
       =(A\RealPosMul\RealNeg B)\RealPosMul(\RealNeg C)\\[-2pt]
      &\quad=A\RealPosMul((\RealNeg B)\RealPosMul(\RealNeg C))
       =A\RealMul(B\RealMul C)
      \end{aligned}}{%
      \FormulaRefAuto{RealCutMultiplication}[def]{1},
      \FormulaRefAuto{RealCutPositiveProductLaws}[thm]{1}}
    \proofstepwidestarbreakkeep[1]{\begin{aligned}[t]
      &A\RealLt\RealZero\dsep B\RealLt\RealZero\dsep C\RealLt\RealZero\vdash{}\\[-2pt]
      &(A\RealMul B)\RealMul C\\[-2pt]
      &\quad=\RealNeg{(((\RealNeg A)\RealPosMul(\RealNeg B))\RealPosMul(\RealNeg C))}\\[-2pt]
      &\quad=\RealNeg{((\RealNeg A)\RealPosMul((\RealNeg B)\RealPosMul(\RealNeg C)))}\\[-2pt]
      &\quad=A\RealMul(B\RealMul C)
      \end{aligned}}{%
      \FormulaRefAuto{RealCutMultiplication}[def]{1},
      \FormulaRefAuto{RealCutPositiveProductLaws}[thm]{1},
      \FormulaRefAuto{RealCutMultiplicationSignRules}[thm]{1}}
    \proofstepwidestarbreakkeep[1]{%
      (A\RealMul B)\RealMul C=A\RealMul(B\RealMul C)}{%
      \rOE{2,3,4,5,6,7,8,9,10,11}}

\closeproofpartwide

\proofpartwide{Kommutativität}
  \setcounter{proofstepnr}{12}% Fortlaufende Schritte dieses Beweises.
\proofstepwidestarbreakkeep[1]{%
      \begin{aligned}[t]
      &\RealZero\RealLe A\land\RealZero\RealLe B,
        \quad A\RealLt\RealZero\land\RealZero\RealLe B,\\[-2pt]
      &\RealZero\RealLe A\land B\RealLt\RealZero,
        \quad A\RealLt\RealZero\land B\RealLt\RealZero
      \end{aligned}}{%
      \FormulaRefAuto{RealCutTotalOrder}[thm]{1}}
    \proofstepwidestarbreakkeep[1]{A\RealMul B=B\RealMul A}{%
      \FormulaRefAuto{RealCutMultiplication}[def]{1,13},
      \FormulaRefAuto{RealCutPositiveProductLaws}[thm]{1},
      \FormulaRefAuto{RealCutMultiplicationSignRules}[thm]{1}}

\closeproofpartwide

\proofpartwide{Einheit}
  \setcounter{proofstepnr}{14}% Fortlaufende Schritte dieses Beweises.
\proofstepwidestarbreakkeep[1]{%
      \begin{aligned}[t]
      \RealZero\RealLe A&\vdash
        A\RealMul\RealOne=A\RealPosMul\RealOne=A,\\[-2pt]
      A\RealLt\RealZero&\vdash
        A\RealMul\RealOne
        =\RealNeg{((\RealNeg A)\RealPosMul\RealOne)}=A
      \end{aligned}}{%
      \FormulaRefAuto{RealCutMultiplication}[def]{1},
      \FormulaRefAuto{RealCutPositiveProductLaws}[thm]{1},
      \FormulaRefAuto{RealCutNegationOrder}[thm]{1}}
    \proofstepwidestarbreakkeep[1]{A\RealMul\RealOne=A}{%
      \FormulaRefAuto{RealCutTotalOrder}[thm]{1,15}}

\closeproofpartwide

\proofpartwide{Multiplikatives Inverses}
  \setcounter{proofstepnr}{16}% Fortlaufende Schritte dieses Beweises.
\proofstepwidestar[17]{A\neq\RealZero}{\rA}
    \proofstepwidestarbreakkeep[1,17]{%
      \RealZero\RealLt A\lor A\RealLt\RealZero}{%
      \FormulaRefAuto{RealCutTotalOrder}[thm]{1,17}}
    \proofstepwidestar[1,17]{\RealZero\RealLt A}{\rA}
    \proofstepwidestarbreakkeep[1,17]{%
      A\RealMul\RealInv A=A\RealPosMul\RealPosInv A}{%
      \FormulaRefAuto{RealCutMultiplication}[def]{1,19},
      \FormulaRefAuto{RealCutReciprocal}[def]{1,17,19}}
    \proofstepwidestarbreakkeep[1,17]{%
      A\RealPosMul\RealPosInv A=\RealOne}{%
      \FormulaRefAuto{RealCutPositiveProductInverse}[thm]{1,19}}
    \proofstepwidestar[22]{A\RealLt\RealZero}{\rA}
    \proofstepwidestarbreakkeep[1,17]{%
      A\RealMul\RealInv A=
      (\RealNeg A)\RealPosMul\RealPosInv{\RealNeg A}}{%
      \FormulaRefAuto{RealCutMultiplication}[def]{1,22},
      \FormulaRefAuto{RealCutReciprocal}[def]{1,17,22},
      \FormulaRefAuto{RealCutNegationOrder}[thm]{1}}
    \proofstepwidestar[1,17]{\RealZero\RealLt\RealNeg A}{%
      \FormulaRefAuto{RealCutNegationOrder}[thm]{1,22}}
    \proofstepwidestarbreakkeep[1,17]{%
      (\RealNeg A)\RealPosMul\RealPosInv{\RealNeg A}=\RealOne}{%
      \FormulaRefAuto{RealCutPositiveProductInverse}[thm]{%
        1,24}}
    \proofstepwidestarbreakkeep[1,17]{A\RealMul\RealInv A=\RealOne}{%
      \rOE{18,19,20,21,22,23,25}}
    \proofstepwidestarbreakkeep[1]{%
      A\neq\RealZero\rightarrow A\RealMul\RealInv A=\RealOne}{%
      \rRI{17,26}}

\closeproofpartwide

\proofpartwideindR[Allgemeine Assoziativität, Einheit und Inverses]{%
    \begin{aligned}[t]
      A,B,C\in\mathbb R\vdash{}&
      (A\RealMul B)\RealMul C=A\RealMul(B\RealMul C)\land{}\\[-2pt]
      &A\RealMul B=B\RealMul A\land A\RealMul\RealOne=A\land{}\\[-2pt]
      &(A\neq\RealZero\rightarrow A\RealMul\RealInv A=\RealOne)
    \end{aligned}
  }{%
    A,B,C\in\mathbb R\vdash
    (A\RealMul B)\RealMul C=A\RealMul(B\RealMul C)\land
    A\RealMul B=B\RealMul A\land A\RealMul\RealOne=A\land
    (A\neq\RealZero\rightarrow A\RealMul\RealInv A=\RealOne)
  }[RealCutMultiplicationCoreLaws]
  \setcounter{proofstepnr}{27}% Fortlaufende Schritte dieses Beweises.
\proofstepwidestarbreakkeep[1]{%
      \begin{aligned}[t]
        &(A\RealMul B)\RealMul C=A\RealMul(B\RealMul C)\land{}\\[-2pt]
        &A\RealMul B=B\RealMul A\land{}\\[-2pt]
        &A\RealMul\RealOne=A\land{}\\[-2pt]
        &(A\neq\RealZero\rightarrow A\RealMul\RealInv A=\RealOne)
      \end{aligned}}{%
      \rAI{12,\rAI{14,\rAI{16,27}}}}

\closeproofpartwideindR


  \proofpartwideindR[Nichtnegative Differenz]{%
    X,Y\in\mathbb R\dsep X\RealLe Y
    \vdash\RealZero\RealLe Y\RealAdd\RealNeg X
  }{%
    X,Y\in\mathbb R\dsep X\RealLe Y
    \vdash\RealZero\RealLe Y\RealAdd\RealNeg X
  }[RealCutDifferenceNonnegative]
    \proofstepwidestar[1]{X,Y\in\mathbb R}{\rA}
    \proofstepwidestar[2]{X\RealLe Y}{\rA}
    \proofstepwidestarbreakkeep[1,2]{X\subseteq Y}{%
      \FormulaRefAuto{RealCutOrderCharacterization}[thm]{1,2}}
    \proofstepwidestarbreakkeep[1]{X\RealAdd\RealNeg X=\RealZero}{%
      \FormulaRefAuto{RealCutAdditiveLaws}[thm]{1}}
    \proofstepwidestar[5]{q\in\RealZero}{\rA}
    \proofstepwidestarbreakkeep[1,5]{q\in X\RealAdd\RealNeg X}{%
      \rIE{4,5}}
    \proofstepwidestarbreakkeep[1,5]{%
      \exists x\in X\,\exists z\in\RealNeg X\,(q<x+z)}{%
      \FormulaRefAuto{RealCutAddition}[def]{1,6}}
    \proofstepwidestarbreakkeep[1,2,5]{%
      \exists y\in Y\,\exists z\in\RealNeg X\,(q<y+z)}{%
      \FormulaRefAuto{A\subseteq B\dsep x\in A\vdash x\in B}[thm]{3,7}}
    \proofstepwidestarbreakkeep[1,2,5]{q\in Y\RealAdd\RealNeg X}{%
      \FormulaRefAuto{RealCutAddition}[def]{1,8}}
    \proofstepwidestarbreakkeep[1,2]{%
      \RealZero\subseteq Y\RealAdd\RealNeg X}{%
      \FormulaRefAuto{SubsetIntroductionRule}[thm]{[5]\vdots 9}}
    \proofstepwidestarbreakkeep[1]{Y\RealAdd\RealNeg X\in\mathbb R}{%
      \FormulaRefAuto{RealCutAdditiveClosure}[thm]{1}}
    \proofstepwidestarbreakkeep[1,2]{%
      \RealZero\RealLe Y\RealAdd\RealNeg X}{%
      \FormulaRefAuto{RealCutOrderCharacterization}[thm]{10,11}}
  \closeproofpartwideindR

  \proofpartwideindR[Allgemeine Distributivität]{%
    A,B,C\in\mathbb R\vdash
    A\RealMul(B\RealAdd C)=
    (A\RealMul B)\RealAdd(A\RealMul C)
  }{%
    A,B,C\in\mathbb R\vdash
    A\RealMul(B\RealAdd C)=
    (A\RealMul B)\RealAdd(A\RealMul C)
  }[RealCutMultiplicationDistributive]
    \proofstepwidestar[1]{A,B,C\in\mathbb R}{\rA}
    \proofstepwidestarbreakkeep[1]{%
      \begin{aligned}[t]
      &(\RealZero\RealLe A\lor A\RealLt\RealZero)\land{}\\[-2pt]
      &(\RealZero\RealLe B\lor B\RealLt\RealZero)\land{}\\[-2pt]
      &(\RealZero\RealLe C\lor C\RealLt\RealZero)
      \end{aligned}}{%
      \FormulaRefAuto{RealCutTotalOrder}[thm]{1}}
    \proofstepwidestarbreakkeep[1]{%
      \RealZero\RealLe A\dsep\RealZero\RealLe B\dsep\RealZero\RealLe C
      \vdash A\RealMul(B\RealAdd C)=
      (A\RealMul B)\RealAdd(A\RealMul C)}{%
      \FormulaRefAuto{RealCutPositiveDistributive}[thm]{1},
      \FormulaRefAuto{RealCutMultiplication}[def]{1}}
    \proofstepwidestarbreakkeep[1]{%
      \begin{aligned}[t]
      &\RealZero\RealLe A\dsep B\RealLt\RealZero\dsep C\RealLt\RealZero
        \vdash{}\\[-2pt]
      &A\RealMul(B\RealAdd C)
       =\RealNeg{\bigl(A\RealPosMul
          ((\RealNeg B)\RealAdd(\RealNeg C))\bigr)}\\[-2pt]
      &\quad=\RealNeg{\bigl((A\RealPosMul\RealNeg B)
          \RealAdd(A\RealPosMul\RealNeg C)\bigr)}\\[-2pt]
      &\quad=(A\RealMul B)\RealAdd(A\RealMul C)
      \end{aligned}}{%
      \FormulaRefAuto{RealCutPositiveDistributive}[thm]{1},
      \FormulaRefAuto{RealCutMultiplicationSignRules}[thm]{1},
      \FormulaRefAuto{RealCutAdditiveLaws}[thm]{1},
      \FormulaRefAuto{RealCutNegationOrder}[thm]{1}}
    \proofstepwidestarbreakkeep[1]{%
      \begin{aligned}[t]
      &\RealZero\RealLe A\dsep\RealZero\RealLe B\dsep C\RealLt\RealZero
        \vdash{}\\[-2pt]
      &D\coloneqq\RealNeg C\dsep\RealZero\RealLt D\dsep
        (D\RealLe B\lor B\RealLe D)
      \end{aligned}}{%
      \FormulaRefAuto{RealCutNegationOrder}[thm]{1},
      \FormulaRefAuto{RealCutTotalOrder}[thm]{1}}
    \proofstepwidestarbreakkeep[1]{%
      \begin{aligned}[t]
      &D\RealLe B\dsep E\coloneqq B\RealAdd C\vdash{}\\[-2pt]
      &\RealZero\RealLe E\dsep B=E\RealAdd D\dsep{}\\[-2pt]
      &A\RealMul B=(A\RealMul E)\RealAdd(A\RealMul D)
      \end{aligned}}{%
      \FormulaRefAuto{RealCutDifferenceNonnegative}[thm]{1},
      \FormulaRefAuto{RealCutAdditiveLaws}[thm]{1},
      \FormulaRefAuto{RealCutPositiveDistributive}[thm]{1}}
    \proofstepwidestarbreakkeep[1]{%
      \begin{aligned}[t]
      &D\RealLe B\vdash{}\\[-2pt]
      &A\RealMul(B\RealAdd C)=A\RealMul E\\[-2pt]
      &\quad=(A\RealMul B)\RealAdd\RealNeg{(A\RealMul D)}
       =(A\RealMul B)\RealAdd(A\RealMul C)
      \end{aligned}}{%
      \FormulaRefAuto{RealCutAdditiveInverseUnique}[thm]{6},
      \FormulaRefAuto{RealCutMultiplicationSignRules}[thm]{1}}
    \proofstepwidestarbreakkeep[1]{%
      \begin{aligned}[t]
      &B\RealLe D\dsep E\coloneqq\RealNeg{(B\RealAdd C)}\vdash{}\\[-2pt]
      &\RealZero\RealLe E\dsep D=B\RealAdd E\dsep{}\\[-2pt]
      &A\RealMul D=(A\RealMul B)\RealAdd(A\RealMul E)
      \end{aligned}}{%
      \FormulaRefAuto{RealCutDifferenceNonnegative}[thm]{1},
      \FormulaRefAuto{RealCutAdditiveLaws}[thm]{1},
      \FormulaRefAuto{RealCutNegationOrder}[thm]{1},
      \FormulaRefAuto{RealCutPositiveDistributive}[thm]{1}}
    \proofstepwidestarbreakkeep[1]{%
      \begin{aligned}[t]
      &B\RealLe D\vdash{}\\[-2pt]
      &A\RealMul(B\RealAdd C)=\RealNeg{(A\RealMul E)}\\[-2pt]
      &\quad=(A\RealMul B)\RealAdd\RealNeg{(A\RealMul D)}
       =(A\RealMul B)\RealAdd(A\RealMul C)
      \end{aligned}}{%
      \FormulaRefAuto{RealCutAdditiveInverseUnique}[thm]{8},
      \FormulaRefAuto{RealCutMultiplicationSignRules}[thm]{1}}
    \proofstepwidestarbreakkeep[1]{%
      \RealZero\RealLe A\dsep\RealZero\RealLe B\dsep C\RealLt\RealZero
      \vdash A\RealMul(B\RealAdd C)=
      (A\RealMul B)\RealAdd(A\RealMul C)}{%
      \rOE{5,6,7,8,9}}
    \proofstepwidestarbreakkeep[1]{%
      \RealZero\RealLe A\dsep B\RealLt\RealZero\dsep\RealZero\RealLe C
      \vdash A\RealMul(B\RealAdd C)=
      (A\RealMul B)\RealAdd(A\RealMul C)}{%
      \FormulaRefAuto{RealCutAdditionCommutative}[thm]{1,10},
      \FormulaRefAuto{RealCutMultiplicationCoreLaws}[thm]{1}}
    \proofstepwidestarbreakkeep[1]{%
      \RealZero\RealLe A\vdash
      A\RealMul(B\RealAdd C)=
      (A\RealMul B)\RealAdd(A\RealMul C)}{%
      \rOE{2,3,4,10,11}}
    \proofstepwidestarbreakkeep[1]{%
      \begin{aligned}[t]
      &A\RealLt\RealZero\vdash{}\\[-2pt]
      &A\RealMul(B\RealAdd C)
       =\RealNeg{\bigl((\RealNeg A)\RealMul(B\RealAdd C)\bigr)}\\[-2pt]
      &\quad=\RealNeg{\bigl(((\RealNeg A)\RealMul B)
          \RealAdd((\RealNeg A)\RealMul C)\bigr)}\\[-2pt]
      &\quad=(A\RealMul B)\RealAdd(A\RealMul C)
      \end{aligned}}{%
      \FormulaRefAuto{RealCutMultiplicationSignRules}[thm]{1},
      \FormulaRefAuto{RealCutNegationOrder}[thm]{1},12,
      \FormulaRefAuto{RealCutAdditiveLaws}[thm]{1}}
    \proofstepwidestarbreakkeep[1]{%
      A\RealMul(B\RealAdd C)=
      (A\RealMul B)\RealAdd(A\RealMul C)}{%
      \rOE{2,12,13}}
  \closeproofpartwideindR

  \proofpartwide{Zusammenfassung}
    \proofstepwidestarbreakkeep[1]{A,B,C\in\mathbb R}{\rA}
    \proofstepwidestarbreakkeep[1]{%
      \begin{aligned}[t]
        &(A\RealMul B)\RealMul C=A\RealMul(B\RealMul C)\land{}\\[-2pt]
        &A\RealMul B=B\RealMul A\land{}\\[-2pt]
        &A\RealMul\RealOne=A\land{}\\[-2pt]
        &(A\neq\RealZero\rightarrow A\RealMul\RealInv A=\RealOne)
      \end{aligned}}{%
      \FormulaRefAuto{RealCutMultiplicationCoreLaws}[thm]{1}}
    \proofstepwidestarbreakkeep[1]{%
      A\RealMul(B\RealAdd C)=
      (A\RealMul B)\RealAdd(A\RealMul C)}{%
      \FormulaRefAuto{RealCutMultiplicationDistributive}[thm]{1}}
    \proofstepwidestarbreakkeep[1]{%
      \begin{aligned}[t]
        &(A\RealMul B)\RealMul C=A\RealMul(B\RealMul C)\land{}\\[-2pt]
        &A\RealMul B=B\RealMul A\land{}\\[-2pt]
        &A\RealMul\RealOne=A\land{}\\[-2pt]
        &A\RealMul(B\RealAdd C)=
          (A\RealMul B)\RealAdd(A\RealMul C)\land{}\\[-2pt]
        &(A\neq\RealZero\rightarrow A\RealMul\RealInv A=\RealOne)
      \end{aligned}}{%
      \rAI{2,3}}
  \closeproofpartwide
\end{tabproofsplitwide}

\FormulaThmDeltaK[Verträglichkeit von Ordnung und Arithmetik]{%
  \begin{aligned}[t]
  A,B,C\in\mathbb R\vdash{}&
    \bigl(A\RealLe B\rightarrow
      A\RealAdd C\RealLe B\RealAdd C\bigr)\land{}\\[-2pt]
  &\bigl(\RealZero\RealLe A\land\RealZero\RealLe B
      \rightarrow\RealZero\RealLe A\RealMul B\bigr).
  \end{aligned}%
}{RealCutOrderArithmeticCompatibility}{%
  \DeltaRow{Mengen}{A\dsep B\dsep C}%
  \DeltaRow{Relationsoperatoren}{\RealLe}%
  \DeltaRow{Funktionsoperatoren}{\RealAdd\dsep\RealMul}%
}
\begin{tabproofsplitwide}
  \proofpartwideindR[Translationsmonotonie]{%
    A,B,C\in\mathbb R\dsep A\RealLe B
    \vdash A\RealAdd C\RealLe B\RealAdd C
  }{%
    A,B,C\in\mathbb R\dsep A\RealLe B
    \vdash A\RealAdd C\RealLe B\RealAdd C
  }[RealCutAdditionMonotone]
    \proofstepwidestar[1]{A,B,C\in\mathbb R}{\rA}
    \proofstepwidestar[2]{A\RealLe B}{\rA}
    \proofstepwidestar[1,2]{A\subseteq B}{%
      \FormulaRefAuto{RealCutOrderCharacterization}[thm]{1,2}}
    \proofstepwidestar[4]{q\in A\RealAdd C}{\rA}
    \proofstepwidestar[1,4]{%
      \exists a\in A\,\exists c\in C\,(q<a+c)}{%
      \FormulaRefAuto{RealCutAddition}[def]{1,4}}
    \proofstepwidestar[1,2,4]{%
      \exists a\in B\,\exists c\in C\,(q<a+c)}{%
      \FormulaRefAuto{A\subseteq B\dsep x\in A\vdash x\in B}{3,5}}
    \proofstepwidestar[1,2,4]{q\in B\RealAdd C}{%
      \FormulaRefAuto{RealCutAddition}[def]{1,6}}
    \proofstepwidestar[1,2]{A\RealAdd C\subseteq B\RealAdd C}{%
      \FormulaRefAuto{SubsetIntroductionRule}[thm]{[4]\vdots 7}}
    \proofstepwidestar[1,2]{A\RealAdd C\in\mathbb R\land B\RealAdd C\in\mathbb R}{%
      \FormulaRefAuto{RealCutAdditiveClosure}[thm]{1}}
    \proofstepwidestar[1,2]{A\RealAdd C\RealLe B\RealAdd C}{%
      \FormulaRefAuto{RealCutOrder}[def]{%
        9,8}}
  \closeproofpartwideindR

  \proofpartwideindR[Nichtnegative Produkte]{%
    A,B\in\mathbb R\dsep\RealZero\RealLe A\dsep\RealZero\RealLe B
    \vdash\RealZero\RealLe A\RealMul B
  }{%
    A,B\in\mathbb R\dsep\RealZero\RealLe A\dsep\RealZero\RealLe B
    \vdash\RealZero\RealLe A\RealMul B
  }[RealCutNonnegativeProduct]
    \proofstepwidestar[1]{A,B\in\mathbb R}{\rA}
    \proofstepwidestar[2]{\RealZero\RealLe A}{\rA}
    \proofstepwidestar[3]{\RealZero\RealLe B}{\rA}
    \proofstepwidestar[1,2,3]{A\RealMul B=A\RealPosMul B}{%
      \FormulaRefAuto{RealCutMultiplication}[def]{1,2,3}}
    \proofstepwidestar[1,2,3]{\RealZero\subseteq A\RealPosMul B}{%
      \FormulaRefAuto{RealCutPositiveProduct}[def]{1,2,3}}
    \proofstepwidestar[1,2,3]{\RealZero\subseteq A\RealMul B}{\rIE{4,5}}
    \proofstepwidestar[1,2,3]{A\RealMul B\in\mathbb R}{%
      \FormulaRefAuto{RealCutMultiplicativeClosure}[thm]{1}}
    \proofstepwidestar[1,2,3]{\RealZero\RealLe A\RealMul B}{%
      \FormulaRefAuto{RealCutOrder}[def]{%
        7,6}}
  \closeproofpartwideindR

  \proofpartwide{Zusammenfassung}
    \proofstepwidestarbreakkeep[1]{A,B,C\in\mathbb R}{\rA}
    \proofstepwidestarbreakkeep[1]{%
      A\RealLe B\rightarrow A\RealAdd C\RealLe B\RealAdd C}{%
      \FormulaRefAuto{RealCutAdditionMonotone}[thm]{1}}
    \proofstepwidestarbreakkeep[1]{%
      (\RealZero\RealLe A\land\RealZero\RealLe B)
      \rightarrow\RealZero\RealLe A\RealMul B}{%
      \FormulaRefAuto{RealCutNonnegativeProduct}[thm]{1}}
    \proofstepwidestarbreakkeep[1]{%
      (A\RealLe B\rightarrow A\RealAdd C\RealLe B\RealAdd C)\land
      ((\RealZero\RealLe A\land\RealZero\RealLe B)
       \rightarrow\RealZero\RealLe A\RealMul B)}{\rAI{2,3}}
  \closeproofpartwide
\end{tabproofsplitwide}

\section{Erhaltung der rationalen Arithmetik}

\FormulaThmDeltaKR[Die rationale Einbettung erhält die Körperoperationen]{%
  \begin{aligned}[t]
    p,q\in\mathbb Q\vdash{}&
    \text{(i)}\quad \RealEmb(p+q)=\RealEmb(p)\RealAdd\RealEmb(q),\\[-2pt]
    &\text{(ii)}\quad \RealEmb(pq)=\RealEmb(p)\RealMul\RealEmb(q),\\[-2pt]
    &\text{(iii)}\quad \RealEmb(-p)=\RealNeg{\RealEmb(p)},\\[-2pt]
    &\text{(iv)}\quad p\neq\RatZero\rightarrow\RealEmb(p^{-1})=\RealInv{\RealEmb(p)}.
  \end{aligned}%
}{%
  \begin{aligned}[t]
  p,q\in\mathbb Q\vdash{}&
    \RealEmb(p+q)=\RealEmb(p)\RealAdd\RealEmb(q)\dsep{}\\[-2pt]
  &\RealEmb(pq)=\RealEmb(p)\RealMul\RealEmb(q)\dsep{}\\[-2pt]
  &\RealEmb(-p)=\RealNeg{\RealEmb(p)}\dsep{}\\[-2pt]
  &p\neq\RatZero\rightarrow
    \RealEmb(p^{-1})=\RealInv{\RealEmb(p)}.
  \end{aligned}%
}{RationalToRealEmbeddingArithmetic}{%
  \DeltaRow{Mengen}{p\dsep q}%
  \DeltaRow{Funktionen}{\RealEmb}%
  \DeltaRow{Funktionsoperatoren}{\RealAdd\dsep\RealMul
    \dsep\RealNeg{\,\cdot\,}\dsep\RealInv{\,\cdot\,}}%
}
\begin{tabproofsplitwide}
\proofpartwideindR[Erhaltung der Addition]{%
    p,q\in\mathbb Q\vdash
    \RealEmb(p+q)=\RealEmb(p)\RealAdd\RealEmb(q)
  }{%
    p,q\in\mathbb Q\vdash
    \RealEmb(p+q)=\RealEmb(p)\RealAdd\RealEmb(q)
  }[RationalToRealEmbeddingAdditive]
    \proofstepwidestar[1]{p,q\in\mathbb Q}{\rA}
    \proofstepwidestar[2]{t\in\RealEmb(p)\RealAdd\RealEmb(q)}{\rA}
    \proofstepwidestar[1,2]{\exists a\in\RealEmb(p)\,\exists b\in\RealEmb(q)\,(t<a+b)}{%
      \FormulaRefAuto{RealCutAddition}[def]{1,2}}
    \proofstepwidestar[1,2]{%
      \exists a<p\,\exists b<q\,(t<a+b)}{%
      \FormulaRefAuto{RationalToRealEmbedding}[def]{%
        3}}
    \proofstepwidestar[1,2]{t<p+q}{%
      \FormulaRefAuto{RationalOrderTranslation}[thm]{4}}
    \proofstepwidestar[1,2]{t\in\RealHat{p+q}}{%
      \FormulaRefAuto{RealPrincipalCut}[def]{1,5}}
    \proofstepwidestar[1,2]{t\in\RealEmb(p+q)}{%
      \FormulaRefAuto{RationalToRealEmbedding}[def]{%
        6}}
    \proofstepwidestar[1]{%
      \RealEmb(p)\RealAdd\RealEmb(q)\subseteq\RealEmb(p+q)}{%
      \FormulaRefAuto{SubsetIntroductionRule}[thm]{[2]\vdots 7}}
    \proofstepwidestar[9]{t\in\RealEmb(p+q)}{\rA}
    \proofstepwidestar[1,9]{t<p+q}{%
      \FormulaRefAuto{RealPrincipalCut}[def]{1,9}}
    \proofstepwidestar[1,9]{t<p+q}{%
      \FormulaRefAuto{RationalToRealEmbedding}[def]{%
        10}}
    \proofstepwidestar[1,9]{%
      \exists a<p\,\exists b<q\,(t<a+b)}{%
      \FormulaRefAuto{RationalOrderDense}[thm]{11}}
    \proofstepwidestar[1,9]{t\in\RealEmb(p)\RealAdd\RealEmb(q)}{%
      \FormulaRefAuto{RealCutAddition}[def]{1,12}}
    \proofstepwidestar[1,9]{t\in\RealEmb(p)\RealAdd\RealEmb(q)}{%
      \FormulaRefAuto{RationalToRealEmbedding}[def]{%
        13}}
    \proofstepwidestar[1]{%
      \RealEmb(p+q)\subseteq\RealEmb(p)\RealAdd\RealEmb(q)}{%
      \FormulaRefAuto{SubsetIntroductionRule}[thm]{[9]\vdots 14}}
    \proofstepwidestar[1]{%
      \RealEmb(p+q)=\RealEmb(p)\RealAdd\RealEmb(q)}{%
      \FormulaRefAuto{A\subseteq B\dsep B\subseteq A\vdash A=B}{15,8}}
  \closeproofpartwideindR

  \proofpartwide{Erhaltung der Negation}
\proofstepwidestar[1]{p,q\in\mathbb Q}{\rA}
    \proofstepwidestar[1]{\RealNeg{\RealEmb(p)}=\{t\in\mathbb Q\mid\exists r\in\mathbb Q\,(r\notin\RealEmb(p)\land t<-r)\}}{%
      \FormulaRefAuto{RealCutNegation}[def]{1}}
    \proofstepwidestar[1]{\RealEmb(-p)=\RealHat{-p}}{%
      \FormulaRefAuto{RationalToRealEmbedding}[def]{1}}
    \proofstepwidestarbreakkeep[1]{%
      \forall t\in\mathbb Q\,
      (t\in\RealEmb(-p)\leftrightarrow t\in\RealNeg{\RealEmb(p)})}{%
      \FormulaRefAuto{RationalNegationLaws}[thm]{%
        3,
        2}}
    \proofstepwidestar[1]{\RealEmb(-p)=\RealNeg{\RealEmb(p)}}{%
      \FormulaRefAuto{\forall x(x\in A\leftrightarrow x\in B)\vdash A=B}{4}}

\closeproofpartwide

\proofpartwide{Erhaltung der Multiplikation}
  \setcounter{proofstepnr}{5}% Fortlaufende Schritte dieses Beweises.
\proofstepwidestarbreakkeep[1]{%
      \begin{aligned}[t]
      &(p=\RatZero\lor\RatZero<p\lor p<\RatZero)\land{}\\[-2pt]
      &(q=\RatZero\lor\RatZero<q\lor q<\RatZero)
      \end{aligned}}{%
      \FormulaRefAuto{RationalOrderTotal}[thm]{1}}
    \proofstepwidestar[1]{\RealEmb(p),\RealEmb(q)\in\mathbb R}{%
      \FormulaRefAuto{RationalToRealEmbeddingFunction}[thm]{1}}
    \proofstepwidestarbreakkeep[1]{%
      p=\RatZero\lor q=\RatZero\vdash
      \RealEmb(pq)=\RealEmb(p)\RealMul\RealEmb(q)}{%
      \FormulaRefAuto{RationalMulZero}[thm]{1},
      \FormulaRefAuto{RealCutZero}[def],
      \FormulaRefAuto{RealCutMultiplicationSignRules}[thm]{%
        7}}
    \proofstepwidestarbreakkeep[1]{%
      \forall t\in\mathbb Q\,
      \bigl(\RatZero<p\land\RatZero<q\rightarrow
       (t\in\RealEmb(pq)\leftrightarrow
        t\in\RealEmb(p)\RealPosMul\RealEmb(q))\bigr)}{\text{\parbox[t]{\linewidth}{\raggedright \FormulaRefAuto{RationalToRealEmbedding}[def]{1}, \FormulaRefAuto{RealCutPositiveProduct}[def]{1}, \FormulaRefAuto{RationalMulAssociative}[thm]{1}, \FormulaRefAuto{RationalOrderDense}[thm]}}}
    \proofstepwidestarbreakkeep[1]{%
      \RatZero<p\dsep\RatZero<q\vdash
      \RealEmb(pq)=\RealEmb(p)\RealMul\RealEmb(q)}{%
      \FormulaRefAuto{\forall x(x\in A\leftrightarrow x\in B)\vdash A=B}{9},
      \FormulaRefAuto{RationalToRealEmbeddingOrder}[thm]{1},
      \FormulaRefAuto{RealCutMultiplication}[def]{1}}
    \proofstepwidestarbreakkeep[1]{%
      \begin{aligned}[t]
      &\RatZero<p\dsep q<\RatZero\vdash{}\\[-2pt]
      &\RealEmb(pq)=\RealNeg{\RealEmb(p(-q))}
       =\RealNeg{(\RealEmb(p)\RealMul\RealEmb(-q))}\\[-2pt]
      &\quad=\RealEmb(p)\RealMul\RealEmb(q)
      \end{aligned}}{%
      \FormulaRefAuto{RationalNegationLaws}[thm]{1},5,10,
      \FormulaRefAuto{RationalToRealEmbeddingOrder}[thm]{1},
      \FormulaRefAuto{RealCutMultiplicationSignRules}[thm]{1}}
    \proofstepwidestarbreakkeep[1]{%
      \begin{aligned}[t]
      &p<\RatZero\dsep\RatZero<q\vdash{}\\[-2pt]
      &\RealEmb(pq)=\RealNeg{\RealEmb((-p)q)}
       =\RealNeg{(\RealEmb(-p)\RealMul\RealEmb(q))}\\[-2pt]
      &\quad=\RealEmb(p)\RealMul\RealEmb(q)
      \end{aligned}}{%
      \FormulaRefAuto{RationalNegationLaws}[thm]{1},5,10,
      \FormulaRefAuto{RationalToRealEmbeddingOrder}[thm]{1},
      \FormulaRefAuto{RealCutMultiplicationSignRules}[thm]{1}}
    \proofstepwidestarbreakkeep[1]{%
      \begin{aligned}[t]
      &p<\RatZero\dsep q<\RatZero\vdash{}\\[-2pt]
      &\RealEmb(pq)=\RealEmb((-p)(-q))\\[-2pt]
      &\quad=\RealEmb(-p)\RealMul\RealEmb(-q)
       =\RealEmb(p)\RealMul\RealEmb(q)
      \end{aligned}}{%
      \FormulaRefAuto{RationalNegationLaws}[thm]{1},5,10,
      \FormulaRefAuto{RationalToRealEmbeddingOrder}[thm]{1},
      \FormulaRefAuto{RealCutMultiplicationSignRules}[thm]{1}}
    \proofstepwidestarbreakkeep[1]{\RealEmb(pq)=\RealEmb(p)\RealMul\RealEmb(q)}{%
      \rOE{6,8,10,11,12,13}}

\closeproofpartwide

\proofpartwideindR[Erhaltung von Negation und Produkt]{%
    \begin{aligned}[t]
      p,q\in\mathbb Q\vdash{}&
      \RealEmb(-p)=\RealNeg{\RealEmb(p)}\land{}\\[-2pt]
      &\RealEmb(pq)=\RealEmb(p)\RealMul\RealEmb(q)
    \end{aligned}
  }{%
    p,q\in\mathbb Q\vdash
    \RealEmb(-p)=\RealNeg{\RealEmb(p)}\land
    \RealEmb(pq)=\RealEmb(p)\RealMul\RealEmb(q)
  }[RationalToRealEmbeddingNegationProduct]
  \setcounter{proofstepnr}{14}% Fortlaufende Schritte dieses Beweises.
\proofstepwidestarbreakkeep[1]{%
      \RealEmb(-p)=\RealNeg{\RealEmb(p)}\land
      \RealEmb(pq)=\RealEmb(p)\RealMul\RealEmb(q)}{\rAI{5,14}}

\closeproofpartwideindR


  \proofpartwideindR[Erhaltung des Kehrwertes]{%
    p\in\mathbb Q\dsep p\neq\RatZero\vdash
    \RealEmb(p^{-1})=\RealInv{\RealEmb(p)}
  }{%
    p\in\mathbb Q\dsep p\neq\RatZero\vdash
    \RealEmb(p^{-1})=\RealInv{\RealEmb(p)}
  }[RationalToRealEmbeddingReciprocal]
    \proofstepwidestar[1]{p\in\mathbb Q}{\rA}
    \proofstepwidestar[2]{p\neq\RatZero}{\rA}
    \proofstepwidestar[1,2]{\RealZero = \RealHat \RatZero}{%
      \FormulaRefAuto{RealCutZero}[def]}
    \proofstepwidestar[1,2]{\RealEmb(p)\neq\RealZero}{%
      \FormulaRefAuto{RationalToRealEmbeddingInjective}[thm]{%
        1,2,3}}
    \proofstepwidestar[1,2]{%
      \RealEmb(p)\RealMul\RealInv{\RealEmb(p)}=\RealOne}{%
      \FormulaRefAuto{RealCutMultiplicativeLaws}[thm]{4}}
    \proofstepwidestar[1,2]{p \cdot p ^ { - 1 } = \RatOne}{%
      \FormulaRefAuto{RationalMulInverse}[thm]{1,2}}
    \proofstepwidestar[1,2]{%
      \RealEmb(p)\RealMul\RealEmb(p^{-1})=\RealOne}{%
      \FormulaRefAuto{RationalToRealEmbeddingNegationProduct}[thm]{%
        6}}
    \proofstepwidestar[1,2]{%
      \RealEmb(p^{-1})=\RealInv{\RealEmb(p)}}{%
      \FormulaRefAuto{RealCutMultiplicationCoreLaws}[thm]{4,5,7}}
  \closeproofpartwideindR

  \proofpartwide{Zusammenfassung}
    \proofstepwidestarbreakkeep[1]{p,q\in\mathbb Q}{\rA}
    \proofstepwidestarbreakkeep[1]{\RealEmb(p+q)=\RealEmb(p)\RealAdd\RealEmb(q)}{%
      \FormulaRefAuto{RationalToRealEmbeddingAdditive}[thm]{1}}
    \proofstepwidestarbreakkeep[1]{%
      \begin{aligned}[t]
        &\RealEmb(-p)=\RealNeg{\RealEmb(p)}\land{}\\[-2pt]
        &\RealEmb(pq)=\RealEmb(p)\RealMul\RealEmb(q)
      \end{aligned}}{%
      \FormulaRefAuto{RationalToRealEmbeddingNegationProduct}[thm]{1}}
    \proofstepwidestarbreakkeep[4]{p\neq\RatZero}{\rA}
    \proofstepwidestarbreakkeep[1,4]{\RealEmb(p^{-1})=\RealInv{\RealEmb(p)}}{%
      \FormulaRefAuto{RationalToRealEmbeddingReciprocal}[thm]{1,4}}
    \proofstepwidestarbreakkeep[1]{%
      p\neq\RatZero\rightarrow
      \RealEmb(p^{-1})=\RealInv{\RealEmb(p)}}{\rRI{4,5}}
    \proofstepwidestarbreakkeep[1]{%
      \begin{aligned}[t]
        &\RealEmb(p+q)=\RealEmb(p)\RealAdd\RealEmb(q)\land{}\\[-2pt]
        &\RealEmb(pq)=\RealEmb(p)\RealMul\RealEmb(q)\land{}\\[-2pt]
        &\RealEmb(-p)=\RealNeg{\RealEmb(p)}\land{}\\[-2pt]
        &(p\neq\RatZero\rightarrow
          \RealEmb(p^{-1})=\RealInv{\RealEmb(p)})
      \end{aligned}}{%
      \rAI{2,\rAI{3,6}}}
  \closeproofpartwide
\end{tabproofsplitwide}

Insbesondere dürfen die rationalen Zahlen von nun an über \(\RealEmb\) mit
ihren Hauptschnitten identifiziert werden. Diese Identifikation ändert weder
Gleichungen noch Ungleichungen.

\chapter{Betrag, Abstand und Intervalle}

\section{Betrag und Abstand}

\FormulaDefDeltaK[Absolutbetrag einer reellen Zahl]{%
  A\in\mathbb R\vdash
  \RealAbs A\coloneqq
  \begin{cases}
    A,&\RealZero\RealLe A,\\
    \RealNeg A,&A\RealLt\RealZero.
  \end{cases}%
}{RealAbsoluteValue}{%
  \DeltaRow{Mengen}{A}%
  \DeltaRow{Relationsoperatoren}{\RealLe\dsep\RealLt}%
  \DeltaRow{Funktionsoperatoren}{\RealNeg{\,\cdot\,}}%
  \DeltaRow{\textbf{Neues Symbol}}{}%
  \DeltaRow{Reeller Betrag}{\RealAbs{\,\cdot\,}}%
}

\FormulaDefDeltaK[Abstand zweier reeller Zahlen]{%
  A,B\in\mathbb R\vdash
  \RealDist{A}{B}\coloneqq\RealAbs{A\RealAdd\RealNeg{B}}%
}{RealDistance}{%
  \DeltaRow{Mengen}{A\dsep B}%
  \DeltaRow{Funktionsoperatoren}{\RealAdd\dsep\RealNeg{\,\cdot\,}
    \dsep\RealAbs{\,\cdot\,}\dsep\RealDist{\,\cdot\,}{\,\cdot\,}}%
  \DeltaRow{\textbf{Neues Symbol}}{}%
  \DeltaRow{Reeller Abstand}{\RealDist{\,\cdot\,}{\,\cdot\,}}%
}

\FormulaThmDeltaKR[Grundgesetze des reellen Betrags]{%
  \begin{aligned}[t]
    A,B\in\mathbb R\vdash{}&
    \text{(i)}\quad \RealZero\RealLe\RealAbs A,\\[-2pt]
    &\text{(ii)}\quad \RealAbs A=\RealZero\leftrightarrow A=\RealZero,\\[-2pt]
    &\text{(iii)}\quad \RealAbs{\RealNeg A}=\RealAbs A,\\[-2pt]
    &\text{(iv)}\quad \RealAbs{A\RealMul B}=\RealAbs A\RealMul\RealAbs B,\\[-2pt]
    &\text{(v)}\quad \RealAbs{A\RealAdd B}\RealLe\RealAbs A\RealAdd\RealAbs B.
  \end{aligned}%
}{%
  \begin{aligned}[t]
  A,B\in\mathbb R\vdash{}&
    \RealZero\RealLe\RealAbs A\dsep
    \bigl(\RealAbs A=\RealZero\leftrightarrow A=\RealZero\bigr)\dsep{}\\[-2pt]
  &\RealAbs{\RealNeg A}=\RealAbs A\dsep
    \RealAbs{A\RealMul B}=\RealAbs A\RealMul\RealAbs B\dsep{}\\[-2pt]
  &\RealAbs{A\RealAdd B}\RealLe\RealAbs A\RealAdd\RealAbs B.
  \end{aligned}%
}{RealAbsoluteValueLaws}{%
  \DeltaRow{Mengen}{A\dsep B}%
  \DeltaRow{Relationsoperatoren}{\RealLe}%
  \DeltaRow{Funktionsoperatoren}{\RealAdd\dsep\RealMul
    \dsep\RealNeg{\,\cdot\,}\dsep\RealAbs{\,\cdot\,}}%
}
\begin{tabproofsplitwide}
\proofpartwide{Nichtnegativität}
\proofstepwidestar[1]{A\in\mathbb R}{\rA}
    \proofstepwidestar[1]{\RealZero\RealLe A\lor A\RealLt\RealZero}{%
      \FormulaRefAuto{RealCutTotalOrder}[thm]{1}}
    \proofstepwidestar[3]{\RealZero\RealLe A}{\rA}
    \proofstepwidestar[1,3]{\RealAbs A=A}{%
      \FormulaRefAuto{RealAbsoluteValue}[def]{1,3}}
    \proofstepwidestar[5]{A\RealLt\RealZero}{\rA}
    \proofstepwidestar[1,5]{\RealZero\RealLt\RealNeg A}{%
      \FormulaRefAuto{RealCutNegationOrder}[thm]{1,5}}
    \proofstepwidestar[1,5]{\RealAbs A=\RealNeg A}{%
      \FormulaRefAuto{RealAbsoluteValue}[def]{1,5}}
    \proofstepwidestar[1]{\RealZero\RealLe\RealAbs A}{%
      \rOE{2,3,4,5,6,7}}

\closeproofpartwide

\proofpartwide{Nullkriterium}
  \setcounter{proofstepnr}{8}% Fortlaufende Schritte dieses Beweises.
\proofstepwidestar[9]{\RealAbs A=\RealZero}{\rA}
    \proofstepwidestar[1,9]{\RealNeg{\RealNeg A}=A\land\RealNeg\RealZero=\RealZero}{%
      \FormulaRefAuto{RealCutNegationOrder}[thm]{1}}
    \proofstepwidestar[1,9]{A=\RealZero}{%
      \FormulaRefAuto{RealAbsoluteValue}[def]{%
        1,2,9,10}}
    \proofstepwidestar[12]{A=\RealZero}{\rA}
    \proofstepwidestar[1,12]{\RealNeg\RealZero=\RealZero}{%
      \FormulaRefAuto{RealCutNegationOrder}[thm]}
    \proofstepwidestar[1,12]{\RealAbs A=\RealZero}{%
      \FormulaRefAuto{RealAbsoluteValue}[def]{%
        1,12,13}}
    \proofstepwidestar[1]{\RealAbs A=\RealZero\leftrightarrow A=\RealZero}{%
      \rLRI{\rRI{9,11},\rRI{12,14}}}

\closeproofpartwide

\proofpartwideindR[Positivität und Nullkriterium]{%
    A\in\mathbb R\vdash
    \RealZero\RealLe\RealAbs A\land
    (\RealAbs A=\RealZero\leftrightarrow A=\RealZero)
  }{%
    A\in\mathbb R\vdash
    \RealZero\RealLe\RealAbs A\land
    (\RealAbs A=\RealZero\leftrightarrow A=\RealZero)
  }[RealAbsoluteValueNonnegativeZero]
  \setcounter{proofstepnr}{15}% Fortlaufende Schritte dieses Beweises.
\proofstepwidestar[1]{%
      \RealZero\RealLe\RealAbs A\land
      (\RealAbs A=\RealZero\leftrightarrow A=\RealZero)}{\rAI{8,15}}

\closeproofpartwideindR


  \proofpartwide{Verträglichkeit mit Negation}
\proofstepwidestarbreakkeep[1]{A,B\in\mathbb R}{\rA}
    \proofstepwidestarbreakkeep[1]{%
      \RealZero\RealLe A\lor A\RealLt\RealZero}{%
      \FormulaRefAuto{RealCutTotalOrder}[thm]{1}}
    \proofstepwidestar[1]{\RealNeg{\RealNeg A}=A\land\RealNeg\RealZero=\RealZero}{%
      \FormulaRefAuto{RealCutNegationOrder}[thm]{1}}
    \proofstepwidestarbreakkeep[1]{\RealAbs{\RealNeg A}=\RealAbs A}{%
      \FormulaRefAuto{RealAbsoluteValue}[def]{%
        1,2,3}}

\closeproofpartwide

\proofpartwide{Verträglichkeit mit Multiplikation}
  \setcounter{proofstepnr}{4}% Fortlaufende Schritte dieses Beweises.
\proofstepwidestarbreakkeep[1]{\text{\parbox[t]{\linewidth}{\raggedright die vier Vorzeichenkombinationen von \ensuremath{A,\allowbreak B} sind vollständig}}}{%
      \FormulaRefAuto{RealCutTotalOrder}[thm]{1}}
    \proofstepwidestar[1]{\begin{aligned}[t]&(\RealNeg A)\RealMul B=\RealNeg{(A\RealMul B)}\land{}\\[-2pt]&(\RealNeg A)\RealMul(\RealNeg B)=A\RealMul B\land A\RealMul\RealZero=\RealZero\end{aligned}}{%
      \FormulaRefAuto{RealCutMultiplicationSignRules}[thm]{1}}
    \proofstepwidestarbreakkeep[1]{%
      \RealAbs{A\RealMul B}=\RealAbs A\RealMul\RealAbs B}{%
      \FormulaRefAuto{RealAbsoluteValue}[def]{%
        1,5,6}}

\closeproofpartwide

\proofpartwideindR[Negation und Produkt]{%
    A,B\in\mathbb R\vdash
    \RealAbs{\RealNeg A}=\RealAbs A\land
    \RealAbs{A\RealMul B}=\RealAbs A\RealMul\RealAbs B
  }{%
    A,B\in\mathbb R\vdash
    \RealAbs{\RealNeg A}=\RealAbs A\land
    \RealAbs{A\RealMul B}=\RealAbs A\RealMul\RealAbs B
  }[RealAbsoluteValueNegationProduct]
  \setcounter{proofstepnr}{7}% Fortlaufende Schritte dieses Beweises.
\proofstepwidestarbreakkeep[1]{%
      \RealAbs{\RealNeg A}=\RealAbs A\land
      \RealAbs{A\RealMul B}=\RealAbs A\RealMul\RealAbs B}{\rAI{4,7}}

\closeproofpartwideindR


  \proofpartwideindR[Dreiecksungleichung]{%
    A,B\in\mathbb R\vdash
    \RealAbs{A\RealAdd B}\RealLe\RealAbs A\RealAdd\RealAbs B
  }{%
    A,B\in\mathbb R\vdash
    \RealAbs{A\RealAdd B}\RealLe\RealAbs A\RealAdd\RealAbs B
  }[RealAbsoluteValueTriangle]
    \proofstepwidestar[1]{A,B\in\mathbb R}{\rA}
    \proofstepwidestar[1]{\TotOrd { \mathbb R , \RealLe }}{%
      \FormulaRefAuto{RealCutTotalOrder}[thm]{1}}
    \proofstepwidestar[1]{A\RealLe\RealAbs A\land B\RealLe\RealAbs B}{%
      \FormulaRefAuto{RealAbsoluteValue}[def]{%
        1,2}}
    \proofstepwidestar[1]{%
      A\RealAdd B\RealLe\RealAbs A\RealAdd\RealAbs B}{%
      \FormulaRefAuto{RealCutAdditionMonotone}[thm]{1,3}}
    \proofstepwidestar[1]{\RealNeg{(A\RealAdd B)}=(\RealNeg A)\RealAdd(\RealNeg B)}{%
      \FormulaRefAuto{RealCutAdditiveLaws}[thm]{1}}
    \proofstepwidestar[1]{%
      \RealNeg{(A\RealAdd B)}\RealLe
      \RealAbs A\RealAdd\RealAbs B}{%
      \FormulaRefAuto{RealCutNegationOrder}[thm]{%
        1,3,5}}
    \proofstepwidestar[1]{\TotOrd { \mathbb R , \RealLe }}{%
      \FormulaRefAuto{RealCutTotalOrder}[thm]{1}}
    \proofstepwidestar[1]{%
      \RealAbs{A\RealAdd B}\RealLe\RealAbs A\RealAdd\RealAbs B}{%
      \FormulaRefAuto{RealAbsoluteValue}[def]{%
        1,4,6,7}}
  \closeproofpartwideindR

  \proofpartwide{Zusammenfassung}
    \proofstepwidestarbreakkeep[1]{A,B\in\mathbb R}{\rA}
    \proofstepwidestarbreakkeep[1]{%
      \RealZero\RealLe\RealAbs A\land
      (\RealAbs A=\RealZero\leftrightarrow A=\RealZero)}{%
      \FormulaRefAuto{RealAbsoluteValueNonnegativeZero}[thm]{1}}
    \proofstepwidestarbreakkeep[1]{%
      \RealAbs{\RealNeg A}=\RealAbs A\land
      \RealAbs{A\RealMul B}=\RealAbs A\RealMul\RealAbs B}{%
      \FormulaRefAuto{RealAbsoluteValueNegationProduct}[thm]{1}}
    \proofstepwidestarbreakkeep[1]{%
      \RealAbs{A\RealAdd B}\RealLe\RealAbs A\RealAdd\RealAbs B}{%
      \FormulaRefAuto{RealAbsoluteValueTriangle}[thm]{1}}
    \proofstepwidestarbreakkeep[1]{%
      \begin{aligned}[t]
        &\RealZero\RealLe\RealAbs A\land{}\\[-2pt]
        &(\RealAbs A=\RealZero\leftrightarrow A=\RealZero)\land{}\\[-2pt]
        &\RealAbs{\RealNeg A}=\RealAbs A\land{}\\[-2pt]
        &\RealAbs{A\RealMul B}=\RealAbs A\RealMul\RealAbs B\land{}\\[-2pt]
        &\RealAbs{A\RealAdd B}\RealLe\RealAbs A\RealAdd\RealAbs B
      \end{aligned}}{%
      \rAI{2,\rAI{3,4}}}
  \closeproofpartwide
\end{tabproofsplitwide}

\FormulaThmDeltaKR[Grundgesetze des reellen Abstands]{%
  \begin{aligned}[t]
    A,B,C\in\mathbb R\vdash{}&
    \text{(i)}\quad \RealZero\RealLe\RealDist{A}{B},\\[-2pt]
    &\text{(ii)}\quad \RealDist{A}{B}=\RealZero\leftrightarrow A=B,\\[-2pt]
    &\text{(iii)}\quad \RealDist{A}{B}=\RealDist{B}{A},\\[-2pt]
    &\text{(iv)}\quad \RealDist{A}{C}\RealLe\RealDist{A}{B}\RealAdd\RealDist{B}{C}.
  \end{aligned}%
}{%
  \begin{aligned}[t]
  A,B,C\in\mathbb R\vdash{}&
    \RealZero\RealLe\RealDist{A}{B}\dsep
    \bigl(\RealDist{A}{B}=\RealZero\leftrightarrow A=B\bigr)\dsep{}\\[-2pt]
  &\RealDist{A}{B}=\RealDist{B}{A}\dsep
    \RealDist{A}{C}\RealLe
      \RealDist{A}{B}\RealAdd\RealDist{B}{C}.
  \end{aligned}%
}{RealDistanceLaws}{%
  \DeltaRow{Mengen}{A\dsep B\dsep C}%
  \DeltaRow{Relationsoperatoren}{\RealLe}%
  \DeltaRow{Funktionsoperatoren}{\RealAdd\dsep\RealDist{\,\cdot\,}{\,\cdot\,}}%
}
\begin{tabproofsplitwide}
\proofpartwide{Nichtnegativität}
\proofstepwidestarbreakkeep[1]{A,B\in\mathbb R}{\rA}
    \proofstepwidestarbreakkeep[1]{\RealDist AB=\RealAbs{A\RealAdd\RealNeg B}}{%
      \FormulaRefAuto{RealDistance}[def]{1}}
    \proofstepwidestarbreakkeep[1]{\RealZero\RealLe\RealDist AB}{%
      \FormulaRefAuto{RealAbsoluteValueNonnegativeZero}[thm]{1,2}}

\closeproofpartwide

\proofpartwide{Definitheit}
  \setcounter{proofstepnr}{3}% Fortlaufende Schritte dieses Beweises.
\proofstepwidestarbreakkeep[1]{%
      \RealDist AB=\RealZero\leftrightarrow
      A\RealAdd\RealNeg B=\RealZero}{%
      \FormulaRefAuto{RealAbsoluteValueNonnegativeZero}[thm]{1,2}}
    \proofstepwidestarbreakkeep[1]{%
      A\RealAdd\RealNeg B=\RealZero\leftrightarrow A=B}{%
      \FormulaRefAuto{RealCutAdditiveInverseUnique}[thm]{1}}
    \proofstepwidestarbreakkeep[1]{\RealDist AB=\RealZero\leftrightarrow A=B}{%
      \rChain{4,5}}

\closeproofpartwide

\proofpartwide{Symmetrie}
  \setcounter{proofstepnr}{6}% Fortlaufende Schritte dieses Beweises.
\proofstepwidestar[1]{\RealNeg{\RealNeg B}=B}{%
      \FormulaRefAuto{RealCutNegationOrder}[thm]{1}}
\proofstepwide[1]{B\RealAdd\RealNeg A}{=}{%
      \RealNeg{(A\RealAdd\RealNeg B)}}{%
      \FormulaRefAuto{RealCutAdditiveLaws}[thm]{%
        1,7}}
    \proofstepwidestar[1]{\RealDist A B=\RealAbs{A\RealAdd\RealNeg B}\land\RealDist B A=\RealAbs{B\RealAdd\RealNeg A}}{%
      \FormulaRefAuto{RealDistance}[def]{1}}
    \proofstepwidestarbreakkeep[1]{\RealDist AB=\RealDist BA}{%
      \FormulaRefAuto{RealAbsoluteValueNegationProduct}[thm]{%
        9,8}}

\closeproofpartwide

\proofpartwideindR[Positivität, Definitheit und Symmetrie]{%
    \begin{aligned}[t]
      A,B\in\mathbb R\vdash{}&
      \RealZero\RealLe\RealDist AB\land{}\\[-2pt]
      &(\RealDist AB=\RealZero\leftrightarrow A=B)\land
      \RealDist AB=\RealDist BA
    \end{aligned}
  }{%
    A,B\in\mathbb R\vdash
    \RealZero\RealLe\RealDist AB\land
    (\RealDist AB=\RealZero\leftrightarrow A=B)\land
    \RealDist AB=\RealDist BA
  }[RealDistanceBasicLaws]
  \setcounter{proofstepnr}{10}% Fortlaufende Schritte dieses Beweises.
\proofstepwidestarbreakkeep[1]{%
      \begin{aligned}[t]
        &\RealZero\RealLe\RealDist AB\land{}\\[-2pt]
        &(\RealDist AB=\RealZero\leftrightarrow A=B)\land{}\\[-2pt]
        &\RealDist AB=\RealDist BA
      \end{aligned}}{\rAI{3,\rAI{6,10}}}

\closeproofpartwideindR


  \proofpartwideindR[Dreiecksungleichung des Abstands]{%
    A,B,C\in\mathbb R\vdash
    \RealDist AC\RealLe\RealDist AB\RealAdd\RealDist BC
  }{%
    A,B,C\in\mathbb R\vdash
    \RealDist AC\RealLe\RealDist AB\RealAdd\RealDist BC
  }[RealDistanceTriangle]
    \proofstepwidestar[1]{A,B,C\in\mathbb R}{\rA}
    \proofstepwide[1]{A\RealAdd\RealNeg C}{=}{%
      (A\RealAdd\RealNeg B)\RealAdd(B\RealAdd\RealNeg C)}{%
      \FormulaRefAuto{RealCutAdditiveLaws}[thm]{1}}
    \proofstepwidestar[1]{%
      \RealAbs{A\RealAdd\RealNeg C}\RealLe
      \RealAbs{A\RealAdd\RealNeg B}\RealAdd
      \RealAbs{B\RealAdd\RealNeg C}}{%
      \FormulaRefAuto{RealAbsoluteValueTriangle}[thm]{1,2}}
    \proofstepwidestar[1]{%
      \RealDist AC\RealLe\RealDist AB\RealAdd\RealDist BC}{%
      \FormulaRefAuto{RealDistance}[def]{3}}
  \closeproofpartwideindR

  \proofpartwide{Zusammenfassung}
    \proofstepwidestarbreakkeep[1]{A,B,C\in\mathbb R}{\rA}
    \proofstepwidestarbreakkeep[1]{%
      \begin{aligned}[t]
        &\RealZero\RealLe\RealDist AB\land{}\\[-2pt]
        &(\RealDist AB=\RealZero\leftrightarrow A=B)\land{}\\[-2pt]
        &\RealDist AB=\RealDist BA
      \end{aligned}}{%
      \FormulaRefAuto{RealDistanceBasicLaws}[thm]{1}}
    \proofstepwidestarbreakkeep[1]{%
      \RealDist AC\RealLe\RealDist AB\RealAdd\RealDist BC}{%
      \FormulaRefAuto{RealDistanceTriangle}[thm]{1}}
    \proofstepwidestarbreakkeep[1]{%
      \begin{aligned}[t]
        &\RealZero\RealLe\RealDist AB\land{}\\[-2pt]
        &(\RealDist AB=\RealZero\leftrightarrow A=B)\land{}\\[-2pt]
        &\RealDist AB=\RealDist BA\land{}\\[-2pt]
        &\RealDist AC\RealLe\RealDist AB\RealAdd\RealDist BC
      \end{aligned}}{\rAI{2,3}}
  \closeproofpartwide
\end{tabproofsplitwide}

\section{Intervalle und Umgebungen}

\FormulaDefDeltaK[Beschränkte reelle Intervalle]{%
  \begin{aligned}[t]
  A,B\in\mathbb R\vdash\quad
  [A,B]_{\mathbb R}&\coloneqq
    \{x\in\mathbb R\mid A\RealLe x\land x\RealLe B\},\\
  (A,B)_{\mathbb R}&\coloneqq
    \{x\in\mathbb R\mid A\RealLt x\land x\RealLt B\},\\
  [A,B)_{\mathbb R}&\coloneqq
    \{x\in\mathbb R\mid A\RealLe x\land x\RealLt B\},\\
  (A,B]_{\mathbb R}&\coloneqq
    \{x\in\mathbb R\mid A\RealLt x\land x\RealLe B\}.
  \end{aligned}%
}{RealBoundedIntervals}{%
  \DeltaRow{Mengen}{A\dsep B\dsep x}%
  \DeltaRow{Relationsoperatoren}{\RealLe\dsep\RealLt}%
  \DeltaRow{\textbf{Neue Symbole}}{}%
  \DeltaRow{Beschränkte Intervalle}
    {[A,B]_{\mathbb R}\dsep(A,B)_{\mathbb R}
      \dsep[A,B)_{\mathbb R}\dsep(A,B]_{\mathbb R}}%
}

\FormulaDefDeltaK[Unbeschränkte reelle Intervalle]{%
  \begin{aligned}[t]
  A,B\in\mathbb R\vdash\quad
  (-\infty,B)_{\mathbb R}&\coloneqq
    \{x\in\mathbb R\mid x\RealLt B\},\\
  (-\infty,B]_{\mathbb R}&\coloneqq
    \{x\in\mathbb R\mid x\RealLe B\},\\
  (A,\infty)_{\mathbb R}&\coloneqq
    \{x\in\mathbb R\mid A\RealLt x\},\\
  [A,\infty)_{\mathbb R}&\coloneqq
    \{x\in\mathbb R\mid A\RealLe x\}.
  \end{aligned}%
}{RealUnboundedIntervals}{%
  \DeltaRow{Mengen}{A\dsep B\dsep x}%
  \DeltaRow{Relationsoperatoren}{\RealLe\dsep\RealLt}%
  \DeltaRow{\textbf{Neue Symbole}}{}%
  \DeltaRow{Unbeschränkte Intervalle}
    {(-\infty,B)_{\mathbb R}\dsep(-\infty,B]_{\mathbb R}}%
  \DeltaRow{Unbeschränkte Intervalle}
    {(A,\infty)_{\mathbb R}\dsep[A,\infty)_{\mathbb R}}%
}

\FormulaDefDeltaK[Offene Epsilon-Umgebung]{%
  \begin{aligned}[t]
  &x,\varepsilon\in\mathbb R\dsep\RealZero\RealLt\varepsilon\\[-2pt]
  &\quad\vdash
  B_{\varepsilon}(x)\coloneqq
    \{y\in\mathbb R\mid\RealDist{y}{x}\RealLt\varepsilon\}.
  \end{aligned}%
}{RealEpsilonNeighborhood}{%
  \DeltaRow{Mengen}{x\dsep y\dsep\varepsilon}%
  \DeltaRow{Relationsoperatoren}{\RealLt}%
  \DeltaRow{Funktionsoperatoren}{\RealDist{\,\cdot\,}{\,\cdot\,}}%
  \DeltaRow{\textbf{Neues Symbol}}{}%
  \DeltaRow{Offene Umgebung}{B_{\varepsilon}(x)}%
}

\FormulaThmDeltaK[Intervallform der Epsilon-Umgebung]{%
  x,\varepsilon\in\mathbb R\dsep\RealZero\RealLt\varepsilon
  \vdash
  B_{\varepsilon}(x)=
  \bigl(x\RealAdd\RealNeg\varepsilon,
        x\RealAdd\varepsilon\bigr)_{\mathbb R}%
}{RealEpsilonNeighborhoodInterval}{%
  \DeltaRow{Mengen}{x\dsep\varepsilon}%
  \DeltaRow{Relationsoperatoren}{\RealLt}%
  \DeltaRow{Funktionsoperatoren}{\RealAdd\dsep\RealNeg{\,\cdot\,}
    \dsep\RealDist{\,\cdot\,}{\,\cdot\,}}%
}
\begin{tabproofwide}
  \proofstepwidestarbreakkeep[1]{x,\varepsilon\in\mathbb R}{\rA}
  \proofstepwidestarbreakkeep[2]{\RealZero\RealLt\varepsilon}{\rA}
  \proofstepwidestarbreakkeep[1,2]{%
    y\in B_\varepsilon(x)\leftrightarrow\RealDist yx\RealLt\varepsilon}{%
    \FormulaRefAuto{RealEpsilonNeighborhood}[def]{1,2}}
  \proofstepwidestarbreakkeep[1,2]{%
    \RealDist yx\RealLt\varepsilon\leftrightarrow
    \RealAbs{y\RealAdd\RealNeg x}\RealLt\varepsilon}{%
    \FormulaRefAuto{RealDistance}[def]{1}}
  \proofstepwidestar[1,2]{\TotOrd { \mathbb R , \RealLe }}{%
      \FormulaRefAuto{RealCutTotalOrder}[thm]{1}}
  \proofstepwidestarbreakkeep[1,2]{%
    \begin{aligned}[t]
      &\RealAbs{y\RealAdd\RealNeg x}\RealLt\varepsilon\leftrightarrow{}\\[-2pt]
      &\quad \RealNeg\varepsilon\RealLt y\RealAdd\RealNeg x
        \land y\RealAdd\RealNeg x\RealLt\varepsilon
    \end{aligned}}{%
    \FormulaRefAuto{RealAbsoluteValue}[def]{%
      1,2,5}}
  \proofstepwidestar[1,2]{\RealNeg{\RealNeg x}=x\land\RealNeg{\RealNeg\varepsilon}=\varepsilon}{%
      \FormulaRefAuto{RealCutNegationOrder}[thm]{1}}
  \proofstepwidestarbreakkeep[1,2]{%
    \begin{aligned}[t]
      &\RealNeg\varepsilon\RealLt y\RealAdd\RealNeg x
        \land y\RealAdd\RealNeg x\RealLt\varepsilon\\[-2pt]
      &\quad\leftrightarrow
        x\RealAdd\RealNeg\varepsilon\RealLt y
        \land y\RealLt x\RealAdd\varepsilon
    \end{aligned}}{%
    \FormulaRefAuto{RealCutAdditionMonotone}[thm]{%
      1,7}}
  \proofstepwidestarbreakkeep[1,2]{%
    \begin{aligned}[t]
      &x\RealAdd\RealNeg\varepsilon\RealLt y
        \land y\RealLt x\RealAdd\varepsilon\\[-2pt]
      &\quad\leftrightarrow
        y\in(x\RealAdd\RealNeg\varepsilon,
          x\RealAdd\varepsilon)_{\mathbb R}
    \end{aligned}}{%
    \FormulaRefAuto{RealBoundedIntervals}[def]{1}}
  \proofstepwidestarbreakkeep[1,2]{%
    \begin{aligned}[t]
      &y\in B_\varepsilon(x)\leftrightarrow{}\\[-2pt]
      &\quad y\in(x\RealAdd\RealNeg\varepsilon,
        x\RealAdd\varepsilon)_{\mathbb R}
    \end{aligned}}{%
    \rChain{3,4,6,8,9}}
  \proofstepwidestarbreakkeep[1,2]{%
    \begin{aligned}[t]
      B_\varepsilon(x)
        &=\bigl(x\RealAdd\RealNeg\varepsilon,\\[-2pt]
        &\hspace{12mm}x\RealAdd\varepsilon\bigr)_{\mathbb R}
    \end{aligned}}{%
    \FormulaRefAuto{\forall x(x\in A\leftrightarrow x\in B)\vdash A=B}{%
      \rUI{10}}}
\end{tabproofwide}

\chapter{Archimedische Grundbegriffe}

Die kanonische reelle Einbettung einer natürlichen Zahl ist die Komposition
\[
  n\longmapsto
  \RealEmb\bigl(\RatEmbed(\IntEmbed(n))\bigr),
\]
diejenige einer ganzen Zahl ist
\(z\mapsto\RealEmb(\RatEmbed(z))\). In den folgenden Formeln werden diese
Kompositionen ausgeschrieben; damit bleiben alle Zwischentypen aus den
Bänden~17 und~18 sichtbar.

\FormulaDefDeltaK[Archimedische Einbettung]{%
  \begin{aligned}[t]
  &\eta\colon\mathbb N\to K\dsep\TotOrd{K,\preccurlyeq}\\[-2pt]
  &\quad\vdash
  \operatorname{Arch}(K,\preccurlyeq,\eta)
  \coloneqq
  \forall x\in K\,\exists n\in\mathbb N\,
    \bigl(x\preccurlyeq\eta(n)\land x\neq\eta(n)\bigr).
  \end{aligned}%
}{ArchimedeanOrderedEmbedding}{%
  \DeltaRow{Mengen}{K\dsep x\dsep n}%
  \DeltaRow{Funktionen}{\eta}%
  \DeltaRow{Relationsoperatoren}{\preccurlyeq}%
  \DeltaRow{\textbf{Neues Symbol}}{}%
  \DeltaRow{Archimedisches Prädikat}
    {\operatorname{Arch}(K,\preccurlyeq,\eta)}%
}

\FormulaThmDeltaK[Archimedische Eigenschaft der reellen Zahlen]{%
  \forall x\in\mathbb R\,\exists n\in\mathbb N\,
    \bigl(x\RealLt
      \RealEmb(\RatEmbed(\IntEmbed(n)))\bigr)%
}{RealArchimedeanProperty}{%
  \DeltaRow{Mengen}{x\dsep n}%
  \DeltaRow{Funktionen}{\IntEmbed\dsep\RatEmbed\dsep\RealEmb}%
  \DeltaRow{Relationsoperatoren}{\RealLt}%
}
\begin{tabproofwide}
  \proofstepwidestar[1]{x\in\mathbb R}{\rA}
  \proofstepwidestar[1]{\RealCut x}{%
    \FormulaRefAuto{RealCutMembership}[thm]{1}}
  \proofstepwidestar[1]{\exists q\in\mathbb Q\setminus x}{%
    \FormulaRefAuto{DedekindCutCharacterization}[thm]{2}}
  \proofstepwidestar[4]{q\in\mathbb Q\land q\notin x}{\rA}
  \proofstepwidestar[1,4]{%
    \exists n\in\mathbb N\,
    q<\RatEmbed(\IntEmbed(n))}{%
    \FormulaRefAuto{RationalArchimedean}[thm]{4}}
  \proofstepwidestar[6]{%
    n\in\mathbb N\land q<\RatEmbed(\IntEmbed(n))}{\rA}
  \proofstepwidestar[1,4,6]{%
    x\subseteq\RealHat{\RatEmbed(\IntEmbed(n))}}{%
    \FormulaRefAuto{DedekindCutCharacterization}[thm]{2,4,6}}
  \proofstepwidestar[4,6]{%
    q\in\RealHat{\RatEmbed(\IntEmbed(n))}\setminus x}{%
    \FormulaRefAuto{RealPrincipalCut}[def]{4,6}}
  \proofstepwidestar[1,4,6]{%
    x\subsetneq\RealHat{\RatEmbed(\IntEmbed(n))}}{%
    \FormulaRefAuto{A\subset B\coloneq A\subseteq B\land A\neq B}{%
      \rAI{7,\FormulaRefAuto{x\in B, x\not\in A\vdash A\neq B}{%
        \FormulaRefAuto{x\in A\setminus B\vdash x\in A}[thm]{8},%
        \FormulaRefAuto{x\in A\setminus B\vdash x\notin B}[thm]{8}}}}}
  \proofstepwidestar[6]{\RealEmb(\RatEmbed(\IntEmbed(n)))=\RealHat{\RatEmbed(\IntEmbed(n))}}{%
      \FormulaRefAuto{RationalToRealEmbedding}[def]{6}}
  \proofstepwidestar[1,4,6]{%
    x\RealLt\RealEmb(\RatEmbed(\IntEmbed(n)))}{%
    \FormulaRefAuto{RealCutStrictOrder}[def]{%
      10,9}}
  \proofstepwidestar[1,4,6]{%
    \exists n\in\mathbb N\,
    x\RealLt\RealEmb(\RatEmbed(\IntEmbed(n)))}{\rEI{6,11}}
  \proofstepwidestar[1,4]{%
    \exists n\in\mathbb N\,
    x\RealLt\RealEmb(\RatEmbed(\IntEmbed(n)))}{\rEE{5,6,12}}
  \proofstepwidestar[1]{%
    \exists n\in\mathbb N\,
    x\RealLt\RealEmb(\RatEmbed(\IntEmbed(n)))}{\rEE{3,4,13}}
  \proofstepwidestar[]{%
    \forall x\in\mathbb R\,\exists n\in\mathbb N\,
    x\RealLt\RealEmb(\RatEmbed(\IntEmbed(n)))}{\rUI{\rRI{1,14}}}
\end{tabproofwide}

\FormulaThmDeltaK[Ganzzahliger Abschnitt einer reellen Zahl]{%
  \forall x\in\mathbb R\,\exists z\in\mathbb Z\,
  \bigl(\RealEmb(\RatEmbed(z))\RealLe x\land
        x\RealLt\RealEmb(\RatEmbed(z+\IntOne))\bigr)%
}{RealIntegerPartProperty}{%
  \DeltaRow{Mengen}{x\dsep z}%
  \DeltaRow{Funktionen}{\RatEmbed\dsep\RealEmb}%
  \DeltaRow{Relationsoperatoren}{\RealLe\dsep\RealLt}%
}
\begin{tabproofsplitwide}
  \proofpartwideindR[Ganzzahlige obere und untere Vergleichswerte]{%
    x\in\mathbb R\vdash
    \exists w,u\in\mathbb Z\,
    (\RealEmb(\RatEmbed(w))\RealLe x\land
     x\RealLt\RealEmb(\RatEmbed(u)))
  }{%
    x\in\mathbb R\vdash
    \exists w,u\in\mathbb Z\,
    (\RealEmb(\RatEmbed(w))\RealLe x\land
     x\RealLt\RealEmb(\RatEmbed(u)))
  }[RealIntegerPartBounds]
    \proofstepwidestarbreakkeep[1]{x\in\mathbb R}{\rA}
    \proofstepwidestarbreakkeep[1]{\RealNeg x\in\mathbb R}{%
      \FormulaRefAuto{RealCutAdditiveClosure}[thm]{1}}
    \proofstepwidestarbreakkeep[1]{%
      \exists m,n\in\mathbb N\,
      (x\RealLt\RealEmb(\RatEmbed(\IntEmbed(n)))\land
       \RealNeg x\RealLt\RealEmb(\RatEmbed(\IntEmbed(m))))}{%
      \FormulaRefAuto{RealArchimedeanProperty}[thm]{1,2}}
    \proofstepwidestarbreakkeep[1]{m,n\in\mathbb N\land
      x\RealLt\RealEmb(\RatEmbed(\IntEmbed(n)))\land
      \RealNeg x\RealLt\RealEmb(\RatEmbed(\IntEmbed(m)))}{\rA}
    \proofstepwidestar[1]{\RealEmb(-\RatEmbed(\IntEmbed(m)))=\RealNeg{\RealEmb(\RatEmbed(\IntEmbed(m)))}}{%
      \FormulaRefAuto{RationalToRealEmbeddingArithmetic}[thm]{4}}
    \proofstepwidestarbreakkeep[1]{%
      \RealEmb(\RatEmbed(-\IntEmbed(m)))\RealLe x}{%
      \FormulaRefAuto{RealCutNegationOrder}[thm]{%
        5}}
    \proofstepwidestarbreakkeep[1]{%
      \exists w,u\in\mathbb Z\,
      (\RealEmb(\RatEmbed(w))\RealLe x\land
       x\RealLt\RealEmb(\RatEmbed(u)))}{%
      \rEI{-\IntEmbed(m),\rEI{\IntEmbed(n),\rAI{6,4}}}}
    \proofstepwidestarbreakkeep[1]{%
      \exists w,u\in\mathbb Z\,
      (\RealEmb(\RatEmbed(w))\RealLe x\land
       x\RealLt\RealEmb(\RatEmbed(u)))}{\rEE{3,4,7}}
  \closeproofpartwideindR

  \proofpartwide{Kleinster Überschreitungsindex}
    \proofstepwidestarbreakkeep[1]{x\in\mathbb R}{\rA}
    \proofstepwidestarbreakkeep[1]{%
      \exists w,u\in\mathbb Z\,
      (\RealEmb(\RatEmbed(w))\RealLe x\land
       x\RealLt\RealEmb(\RatEmbed(u)))}{%
      \FormulaRefAuto{RealIntegerPartBounds}[thm]{1}}
    \proofstepwidestarbreakkeep[1]{w,u\in\mathbb Z\land
      \RealEmb(\RatEmbed(w))\RealLe x\land
      x\RealLt\RealEmb(\RatEmbed(u))}{\rA}
    \proofstepwidestarbreakkeep[1]{%
      M\coloneqq\{n\in\mathbb N\mid
      x\RealLt\RealEmb(\RatEmbed(w+\IntEmbed(n)))\}}{\rII}
    \proofstepwidestarbreakkeep[1]{M\neq\varnothing}{%
      \FormulaRefAuto{IntegerSignedNormalForm}[thm]{3,4}}
    \proofstepwidestarbreakkeep[1]{%
      \exists m_0\in M\,\forall n\in M\,(m_0\leq n)}{%
      \FormulaRefAuto{NaturalWellOrderingPrinciple}[thm]{4,5}}
    \proofstepwidestarbreakkeep[1]{m_0\in M\land\forall n\in M\,(m_0\leq n)}{\rA}
    \proofstepwidestarbreakkeep[1]{m_0\neq0}{%
      \FormulaRefAuto{RealCutTotalOrder}[thm]{3,4,7}}
    \proofstepwidestarbreakkeep[1]{k\coloneqq(m_0-1)\in\mathbb N}{%
      \FormulaRefAuto{PeanoMinusOneNatural}[thm]{7,8}}
    \proofstepwidestarbreakkeep[1]{m_0=(k+1)}{%
      \FormulaRefAuto{PeanoMinusOneRecoverValue}[thm]{7,8,9}}
    \proofstepwidestarbreakkeep[1]{%
      \RealEmb(\RatEmbed(w+\IntEmbed(k)))\RealLe x}{%
      \FormulaRefAuto{NaturalWellOrderingPrinciple}[thm]{4,7,9,10}}
    \proofstepwidestarbreakkeep[1]{%
      x\RealLt\RealEmb(\RatEmbed(w+\IntEmbed(k)+\IntOne))}{%
      \FormulaRefAuto{IntegerSuccessorOnEmbedding}[thm]{7,10}}
    \proofstepwidestarbreakkeep[1]{z\coloneqq w+\IntEmbed(k)\in\mathbb Z}{%
      \FormulaRefAuto{IntAddClosure}[thm]{3,9}}
    \proofstepwidestarbreakkeep[1]{%
      \exists z\in\mathbb Z\,
      (\RealEmb(\RatEmbed(z))\RealLe x\land
       x\RealLt\RealEmb(\RatEmbed(z+\IntOne)))}{%
      \rEI{13,\rAI{11,12}}}
    \proofstepwidestarbreakkeep[1]{%
      \exists z\in\mathbb Z\,
      (\RealEmb(\RatEmbed(z))\RealLe x\land
       x\RealLt\RealEmb(\RatEmbed(z+\IntOne)))}{\rEE{6,7,14}}
    \proofstepwidestarbreakkeep[1]{%
      \exists z\in\mathbb Z\,
      (\RealEmb(\RatEmbed(z))\RealLe x\land
       x\RealLt\RealEmb(\RatEmbed(z+\IntOne)))}{\rEE{2,3,15}}
    \proofstepwidestarbreakkeep[]{%
      \forall x\in\mathbb R\,\exists z\in\mathbb Z\,
      (\RealEmb(\RatEmbed(z))\RealLe x\land
       x\RealLt\RealEmb(\RatEmbed(z+\IntOne)))}{\rUI{\rRI{1,16}}}
  \closeproofpartwide
\end{tabproofsplitwide}

\FormulaThmDeltaK[Dichtheit der rationalen Zahlen in den reellen Zahlen]{%
  x,y\in\mathbb R\dsep x\RealLt y
  \vdash
  \exists q\in\mathbb Q\,
    \bigl(x\RealLt\RealEmb(q)\land\RealEmb(q)\RealLt y\bigr)%
}{RationalDenseInReal}{%
  \DeltaRow{Mengen}{x\dsep y\dsep q}%
  \DeltaRow{Funktionen}{\RealEmb}%
  \DeltaRow{Relationsoperatoren}{\RealLt}%
}
\begin{tabproofwide}
  \proofstepwidestar[1]{x,y\in\mathbb R}{\rA}
  \proofstepwidestar[2]{x\RealLt y}{\rA}
  \proofstepwidestar[1,2]{x\subsetneq y}{%
    \FormulaRefAuto{RealCutStrictOrder}[def]{1,2}}
  \proofstepwidestar[1,2]{\exists b\in y\setminus x}{%
    \FormulaRefAuto{A\subset B\vdash\exists x\in B\setminus A}{3}}
  \proofstepwidestar[5]{b\in y\land b\notin x}{\rA}
  \proofstepwidestar[1]{\RealCut y}{%
      \FormulaRefAuto{RealCutMembership}[thm]{1}}
  \proofstepwidestar[1,5]{\exists q\in y\,(b<q)}{%
    \FormulaRefAuto{DedekindCutCharacterization}[thm]{%
      6,5}}
  \proofstepwidestar[8]{q\in y\land b<q}{\rA}
  \proofstepwidestar[1]{\RealCut x}{%
      \FormulaRefAuto{RealCutMembership}[thm]{1}}
  \proofstepwidestar[1,2,5,8]{x\subseteq\RealHat q}{%
    \FormulaRefAuto{DedekindCutCharacterization}[thm]{%
      9,5,8}}
  \proofstepwidestar[5,8]{b\in\RealHat q\setminus x}{%
    \FormulaRefAuto{RealPrincipalCut}[def]{5,8}}
  \proofstepwidestar[1,2,5,8]{x\subsetneq\RealHat q}{%
    \FormulaRefAuto{A\subset B\coloneq A\subseteq B\land A\neq B}{%
      \rAI{10,\FormulaRefAuto{x\in B, x\not\in A\vdash A\neq B}{%
        \FormulaRefAuto{x\in A\setminus B\vdash x\in A}[thm]{11},%
        \FormulaRefAuto{x\in A\setminus B\vdash x\notin B}[thm]{11}}}}}
  \proofstepwidestar[1]{\RealCut y}{%
      \FormulaRefAuto{RealCutMembership}[thm]{1}}
  \proofstepwidestar[1,8]{\RealHat q\subseteq y}{%
    \FormulaRefAuto{DedekindCutCharacterization}[thm]{%
      13,8}}
  \proofstepwidestar[1,8]{q\in y\setminus\RealHat q}{%
    \FormulaRefAuto{RealPrincipalCut}[def]{8}}
  \proofstepwidestar[1,8]{\RealHat q\subsetneq y}{%
    \FormulaRefAuto{A\subset B\coloneq A\subseteq B\land A\neq B}{%
      \rAI{14,\FormulaRefAuto{x\in B, x\not\in A\vdash A\neq B}{%
        \FormulaRefAuto{x\in A\setminus B\vdash x\in A}[thm]{15},%
        \FormulaRefAuto{x\in A\setminus B\vdash x\notin B}[thm]{15}}}}}
  \proofstepwidestar[1,2,5,8]{x\RealLt\RealHat q\land\RealHat q\RealLt y}{%
      \FormulaRefAuto{RealCutStrictOrder}[def]{12,16}}
  \proofstepwidestar[1,2,5,8]{%
    x\RealLt\RealEmb(q)\land\RealEmb(q)\RealLt y}{%
    \FormulaRefAuto{RationalToRealEmbedding}[def]{%
      17}}
  \proofstepwidestar[1,2,5,8]{%
    \exists q\in\mathbb Q\,
    (x\RealLt\RealEmb(q)\land\RealEmb(q)\RealLt y)}{\rEI{8,18}}
  \proofstepwidestar[1,2,5]{%
    \exists q\in\mathbb Q\,
    (x\RealLt\RealEmb(q)\land\RealEmb(q)\RealLt y)}{\rEE{7,8,19}}
  \proofstepwidestar[1,2]{%
    \exists q\in\mathbb Q\,
    (x\RealLt\RealEmb(q)\land\RealEmb(q)\RealLt y)}{\rEE{4,5,20}}
\end{tabproofwide}

\FormulaThmDeltaK[Beliebig kleine positive rationale Hauptschnitte]{%
  \begin{aligned}[t]
  &\varepsilon\in\mathbb R\dsep\RealZero\RealLt\varepsilon
    \vdash \exists n\in\mathbb N\,\bigl(0<n\land{}\\[-2pt]
  &\qquad \RealZero\RealLt
      \RealEmb\bigl((\RatEmbed(\IntEmbed(n)))^{-1}\bigr)
    \land{}\\[-2pt]
  &\qquad \RealEmb\bigl((\RatEmbed(\IntEmbed(n)))^{-1}\bigr)
        \RealLt\varepsilon\bigr).
  \end{aligned}%
}{RealSmallPositiveRational}{%
  \DeltaRow{Mengen}{\varepsilon\dsep n\dsep q\dsep m\dsep M_m}%
  \DeltaRow{Funktionen}{\IntEmbed\dsep\RatEmbed\dsep\RealEmb}%
  \DeltaRow{Relationsoperatoren}{\RealLt}%
}
\begin{tabproofwide}
  \proofstepwidestarbreakkeep[1]{\varepsilon\in\mathbb R}{\rA}
  \proofstepwidestarbreakkeep[2]{\RealZero\RealLt\varepsilon}{\rA}
  \proofstepwidestarbreakkeep[]{\RealZero\in\mathbb R}{%
    \FormulaRefAuto{RealCutAdditiveClosure}[thm]}
  \proofstepwidestarbreakkeep[1,2]{%
    \exists q\in\mathbb Q\,
    (\RealZero\RealLt\RealEmb(q)\land
     \RealEmb(q)\RealLt\varepsilon)}{%
    \FormulaRefAuto{RationalDenseInReal}[thm]{3,1,2}}
  \proofstepwidestarbreakkeep[1,2]{q\in\mathbb Q\land
    \RealZero\RealLt\RealEmb(q)\land
    \RealEmb(q)\RealLt\varepsilon}{\rA}
  \proofstepwidestar[1,2]{\RealZero=\RealHat{\RatZero}}{%
      \FormulaRefAuto{RealCutZero}[def]}
  \proofstepwidestar[1,2]{\RatZero\in\mathbb Q}{%
      \FormulaRefAuto{RationalZeroCarrierAxiom}[ax]}
  \proofstepwidestarbreakkeep[1,2]{%
    q\in\mathbb Q\land\RatZero<q\land q\neq\RatZero\land
    q^{-1}\in\mathbb Q\land\RatZero<q^{-1}}{\text{\parbox[t]{\linewidth}{\raggedright \FormulaRefAuto{RationalToRealEmbeddingOrder}[thm]{
        7,\rAEa{5},
        6,\rAEa{\rAEb{5}}}, \FormulaRefAuto{RationalStrictOrderDef}[def] , \FormulaRefAuto{RationalReciprocalClosureAxiom}[ax] , \FormulaRefAuto{RationalPositiveReciprocalAxiom}[ax]}}}
  \proofstepwidestarbreakkeep[1,2]{%
    \exists m\in\mathbb N\,
    q^{-1}<\RatEmbed(\IntEmbed(m))}{%
    \FormulaRefAuto{RationalArchimedean}[thm]{8}}
  \proofstepwidestarbreakkeep[1,2]{m\in\mathbb N\land
    q^{-1}<\RatEmbed(\IntEmbed(m))}{\rA}
  \proofstepwidestarbreakkeep[1,2]{%
    \begin{aligned}[t]
      &M_m\coloneqq\RatEmbed(\IntEmbed(m))\in\mathbb Q,\\[-2pt]
      &0<m,\qquad \RatZero<M_m,
       \qquad q^{-1}<M_m
    \end{aligned}}{\text{\parbox[t]{\linewidth}{\raggedright \FormulaRefAuto{NaturalIntegerEmbeddingFunction}[thm] , \FormulaRefAuto{IntegerRationalEmbeddingFunction}[thm] , \FormulaRefAuto{IntegerRationalEmbeddingOrder}[thm] , \FormulaRefAuto{IntegerOrderClassChar}[thm] , \FormulaRefAuto{NaturalIntegerEmbeddingValue}[thm] , \FormulaRefAuto{IntegerZeroDef}[def] , \FormulaRefAuto{RationalOrderTotal}[thm]{8,10}}}}
  \proofstepwidestarbreakkeep[1,2]{%
    M_m^{-1}\in\mathbb Q\land
    \RatZero<M_m^{-1}\land M_m^{-1}<q}{\text{\parbox[t]{\linewidth}{\raggedright \FormulaRefAuto{RationalStrictOrderDef}[def] , \FormulaRefAuto{RationalReciprocalClosureAxiom}[ax] , \FormulaRefAuto{RationalPositiveReciprocalAxiom}[ax] , \FormulaRefAuto{RationalPositiveReciprocalOrderReversal}[thm]{11,8}, \FormulaRefAuto{RationalOrderedFieldAxioms}[def]{\rAEa{8},8}}}}
  \proofstepwidestar[1,2]{\RealZero=\RealHat{\RatZero}}{%
      \FormulaRefAuto{RealCutZero}[def]}
  \proofstepwidestar[1,2]{\RatZero\in\mathbb Q}{%
      \FormulaRefAuto{RationalZeroCarrierAxiom}[ax]}
  \proofstepwidestarbreakkeep[1,2]{%
    \RealZero\RealLt\RealEmb(M_m^{-1})
    \land\RealEmb(M_m^{-1})\RealLt\RealEmb(q)}{\text{\parbox[t]{\linewidth}{\raggedright \FormulaRefAuto{RationalToRealEmbeddingOrder}[thm]{
        14,12,
        13}, \FormulaRefAuto{RationalToRealEmbeddingOrder}[thm]{12,8}}}}
  \proofstepwidestarbreakkeep[1,2]{%
    \RealEmb(M_m^{-1})
      \RealLt\varepsilon}{%
    \FormulaRefAuto{RealCutTotalOrder}[thm]{15,\rAEb{\rAEb{5}}}}
  \proofstepwidestarbreakkeep[1,2]{%
    \exists n\in\mathbb N\,(0<n\land
    \RealZero\RealLt\RealEmb((\RatEmbed(\IntEmbed(n)))^{-1})\land
    \RealEmb((\RatEmbed(\IntEmbed(n)))^{-1})\RealLt\varepsilon)}{%
    \rEI{m,11,15,16}}
  \proofstepwidestarbreakkeep[1,2]{%
    \exists n\in\mathbb N\,(0<n\land
    \RealZero\RealLt\RealEmb((\RatEmbed(\IntEmbed(n)))^{-1})\land
    \RealEmb((\RatEmbed(\IntEmbed(n)))^{-1})\RealLt\varepsilon)}{%
    \rEE{9,10,17}}
  \proofstepwidestarbreakkeep[1,2]{%
    \exists n\in\mathbb N\,(0<n\land
    \RealZero\RealLt\RealEmb((\RatEmbed(\IntEmbed(n)))^{-1})\land
    \RealEmb((\RatEmbed(\IntEmbed(n)))^{-1})\RealLt\varepsilon)}{%
    \rEE{4,5,18}}
\end{tabproofwide}

Die letzte Aussage ist die Form der archimedischen Eigenschaft, die in
späteren Epsilon-Argumenten am häufigsten verwendet wird. Sie setzt nur die
positive reelle Zahl \(\varepsilon\) voraus und liefert eine rationale
Vergleichsgröße innerhalb derselben Zielmenge \(\mathbb R\).

\chapter{Axiomatische Zusammenfassung}

\section{Vollständig angeordnete Körper}

Die folgende Definition fasst genau die Eigenschaften zusammen, welche die
Schnittkonstruktion bereitstellt. Die Vollständigkeitszeile quantifiziert über
Teilmengen des Trägers und ist daher in dieser Form zweitstufig zu lesen.

\FormulaDefDeltaK[Vollständig angeordneter Körper]{%
  \CompleteOrderedField{\begin{gathered}
    K,\boxplus,\boxtimes,\boxminus,\\[-2pt]
    \operatorname{inv},0_K,1_K,\preccurlyeq
  \end{gathered}}%
}{CompleteOrderedFieldAxioms}{%
  \DeltaRow{Mengen}{K\dsep x\dsep y\dsep z}%
  \DeltaRow{Mengen}{S\dsep s\dsep u}%
  \DeltaRow{Relationsoperatoren}{\preccurlyeq}%
  \DeltaRow{Funktionsoperatoren}{\boxplus\dsep\boxtimes}%
  \DeltaRow{Funktionsoperatoren}{\boxminus\dsep\operatorname{inv}}%
  \multicolumn{3}{@{}l@{}}{\textbf{Typen}}\\[-2pt]%
  \DeltaRow{Operationen}
    {\boxplus,\boxtimes\colon K\times K\to K}%
  \DeltaRow{Negation}{\boxminus\colon K\to K}%
  \DeltaRow{Kehrwert}
    {\operatorname{inv}\colon K\setminus\{0_K\}\to K}%
  \DeltaRow{Konstanten}{0_K,1_K\in K\dsep0_K\neq1_K}%
  \multicolumn{3}{@{}l@{}}{\textbf{Additive abelsche Gruppe}}\\[-2pt]%
  \DeltaRow{Ass. Addition}
    {x,y,z\in K\vdash
      (x\boxplus y)\boxplus z=x\boxplus(y\boxplus z)}%
  \DeltaRow{Komm. Addition}
    {x,y\in K\vdash x\boxplus y=y\boxplus x}%
  \DeltaRow{Null und Inverses}
    {x\in K\vdash x\boxplus0_K=x\dsep
      x\boxplus(\boxminus x)=0_K}%
  \multicolumn{3}{@{}l@{}}{%
    \textbf{Kommutatives multiplikatives Monoid auf \(K\);
      Inversen auf \(K^\times\)}}\\[-2pt]%
  \DeltaRow{Ass. Produkt}
    {x,y,z\in K\vdash
      (x\boxtimes y)\boxtimes z=x\boxtimes(y\boxtimes z)}%
  \DeltaRow{Komm. Produkt}
    {x,y\in K\vdash x\boxtimes y=y\boxtimes x}%
  \DeltaRow{Eins}
    {x\in K\vdash x\boxtimes1_K=x}%
  \DeltaRow{Produktinverses}
    {x\in K\setminus\{0_K\}\vdash
      x\boxtimes\operatorname{inv}(x)=1_K}%
  \multicolumn{3}{@{}l@{}}{\textbf{Feldkopplung}}\\[-2pt]%
  \DeltaRow{Distributivität}
    {x,y,z\in K\vdash x\boxtimes(y\boxplus z)=
      (x\boxtimes y)\boxplus(x\boxtimes z)}%
  \multicolumn{3}{@{}l@{}}{\textbf{Totale Ordnung}}\\[-2pt]%
  \DeltaRow{Totalität}{\TotOrd{K,\preccurlyeq}}%
  \multicolumn{3}{@{}l@{}}{%
    \textbf{Verträglichkeit von Ordnung und Körper}}\\[-2pt]%
  \DeltaRow{Positive Eins}{0_K\preccurlyeq1_K}%
  \DeltaRow{Translation}
    {x,y,z\in K\dsep x\preccurlyeq y\vdash
      x\boxplus z\preccurlyeq y\boxplus z}%
  \DeltaRow{Positive Produkte}
    {x,y\in K\dsep0_K\preccurlyeq x\dsep0_K\preccurlyeq y
      \vdash0_K\preccurlyeq x\boxtimes y}%
  \multicolumn{3}{@{}l@{}}{\textbf{Dedekind-Vollständigkeit}}\\[-2pt]%
  \DeltaRow{Vollständigkeit}
    {\begin{aligned}[t]
      &\varnothing\neq S\subseteq K\dsep
        \UB_{K,\preccurlyeq}(S)\neq\varnothing\vdash{}\\[-2pt]
      &\exists s\in\UB_{K,\preccurlyeq}(S)\,
        \forall u\in\UB_{K,\preccurlyeq}(S)\,(s\preccurlyeq u)
     \end{aligned}}%
  \multicolumn{3}{@{}l@{}}{\textbf{Struktur}}\\[-2pt]%
  \DeltaRow{Prädikat}{\mathsf{COF}}%
}

\FormulaThmDeltaK[Das Schnittmodell ist ein vollständig angeordneter Körper]{%
  \CompleteOrderedField{\mathbb R,\RealAdd,\RealMul,
    \RealNeg{\,\cdot\,},\RealInv{\,\cdot\,},
    \RealZero,\RealOne,\RealLe}%
}{RealCutsCompleteOrderedField}{%
  \DeltaRow{Mengen}{\mathbb R}%
  \DeltaRow{Relationsoperatoren}{\RealLe}%
  \DeltaRow{Funktionsoperatoren}{\RealAdd\dsep\RealMul
    \dsep\RealNeg{\,\cdot\,}\dsep\RealInv{\,\cdot\,}}%
}
\begin{tabproofsplitwide}
  \proofpartwide{Träger und Körpergesetze}
    \proofstepwidestarbreakkeep[]{\RealZero,\RealOne\in\mathbb R}{\FormulaRefAuto{RealCutAdditiveClosure}[thm], %
      \FormulaRefAuto{RealCutMultiplicativeClosure}[thm]}
    \proofstepwidestar[]{\RealOne = \RealHat \RatOne}{%
      \FormulaRefAuto{RealCutOne}[def]}
    \proofstepwidestar[]{\RealZero = \RealHat \RatZero}{%
      \FormulaRefAuto{RealCutZero}[def]}
    \proofstepwidestar[]{\RatZero<\RatOne}{%
      \FormulaRefAuto{RationalOnePositiveAxiom}[ax]}
    \proofstepwidestar[]{\RatZero\leq\RatOne\land\RatZero\neq\RatOne}{\FormulaRefAuto{RationalStrictOrderDef}[def]{4}}
    \proofstepwidestarbreakkeep[]{\RealZero\neq\RealOne}{%
      \FormulaRefAuto{RationalToRealEmbeddingInjective}[thm]{%
        3,2,
        5}}
    \proofstepwidestarbreakkeep[]{\RealAdd\colon\mathbb R\times\mathbb R\to\mathbb R}{%
      \FormulaRefAuto{RealCutAdditiveClosure}[thm]}
    \proofstepwidestarbreakkeep[]{\RealMul\colon\mathbb R\times\mathbb R\to\mathbb R}{%
      \FormulaRefAuto{RealCutMultiplicativeClosure}[thm]}
    \proofstepwidestarbreakkeep[]{\RealNeg{\,\cdot\,}\colon\mathbb R\to\mathbb R}{%
      \FormulaRefAuto{RealCutAdditiveClosure}[thm]}
    \proofstepwidestarbreakkeep[]{%
      \RealInv{\,\cdot\,}\colon\mathbb R\setminus\{\RealZero\}\to\mathbb R}{%
      \FormulaRefAuto{RealCutMultiplicativeClosure}[thm]}
    \proofstepwidestarbreakkeep[]{\text{alle additiven Körperzeilen gelten}}{%
      \FormulaRefAuto{RealCutAdditiveLaws}[thm]}
    \proofstepwidestarbreakkeep[]{\text{alle multiplikativen Körperzeilen gelten}}{%
      \FormulaRefAuto{RealCutMultiplicativeLaws}[thm]}
  \closeproofpartwide

  \proofpartwide{Ordnung und Vollständigkeit}
    \setcounter{proofstepnr}{12}
    \proofstepwidestarbreakkeep[]{\TotOrd{\mathbb R,\RealLe}}{%
      \FormulaRefAuto{RealCutTotalOrder}[thm]}
    \proofstepwidestar[]{\RatZero<\RatOne}{%
      \FormulaRefAuto{RationalOnePositiveAxiom}[ax]}
    \proofstepwidestar[]{\RatZero\leq\RatOne\land\RatZero\neq\RatOne}{%
      \FormulaRefAuto{RationalStrictOrderDef}[def]{%
          14}}
    \proofstepwidestarbreakkeep[]{\RealZero\RealLe\RealOne}{\text{\parbox[t]{\linewidth}{\raggedright \FormulaRefAuto{RationalToRealEmbeddingOrder}[thm] , \FormulaRefAuto{RationalZeroCarrierAxiom}[ax] , \FormulaRefAuto{RationalOneCarrierAxiom}[ax] , \ensuremath{\rAEa{15},\allowbreak } \FormulaRefAuto{RealCutZero}[def] , \FormulaRefAuto{RealCutOne}[def]}}}
    \proofstepwidestarbreakkeep[]{\text{\parbox[t]{\linewidth}{\raggedright Translation und positive Produkte sind ordnungstreu}}}{%
      \FormulaRefAuto{RealCutOrderArithmeticCompatibility}[thm]}
    \proofstepwidestarbreakkeep[18]{\varnothing\neq S\subseteq\mathbb R}{\rA}
    \proofstepwidestarbreakkeep[19]{\UB_{\mathbb R,\RealLe}(S)\neq\varnothing}{\rA}
    \proofstepwidestarbreakkeep[18,19]{%
      \exists U\in\mathbb R\,\forall s\in S\,(s\RealLe U)}{%
      \FormulaRefAuto{UpperBoundSetCharacterization}[thm]{18,19}}
    \proofstepwidestarbreakkeep[18,19]{%
      \exists s\in\UB_{\mathbb R,\RealLe}(S)\,
      \forall u\in\UB_{\mathbb R,\RealLe}(S)\,(s\RealLe u)}{%
      \FormulaRefAuto{RealLeastUpperBoundProperty}[thm]{18,20}}
    \proofstepwidestarbreakkeep[]{%
      \CompleteOrderedField{\mathbb R,\RealAdd,\RealMul,
        \RealNeg{\,\cdot\,},\RealInv{\,\cdot\,},
        \RealZero,\RealOne,\RealLe}}{\FormulaRefAuto{CompleteOrderedFieldAxioms}[def]{1,2,3,4,5,6,7,8,9,10,11,12,13,14,15,16,17,\rUI{\rRI{18,\rRI{19,21}}}}}
  \closeproofpartwide
\end{tabproofsplitwide}

\FormulaThmDeltaK[Kanonische natürliche Abbildung eines angeordneten Körpers]{%
  \begin{aligned}[t]
  &\CompleteOrderedField{K,\boxplus,\boxtimes,
    \boxminus,\operatorname{inv},0_K,1_K,\preccurlyeq}\\[-2pt]
  &\quad\vdash
    \exists!\eta\colon\mathbb N\to K\,
    \left(
      \eta(0)=0_K\land
      \forall n\in\mathbb N\,
        \bigl(\eta(n+1)=\eta(n)\boxplus1_K\bigr)
    \right).
  \end{aligned}%
}{CompleteOrderedFieldNaturalEmbedding}{%
  \DeltaRow{Mengen}{K\dsep n}%
  \DeltaRow{Funktionen}{\eta}%
  \DeltaRow{Funktionsoperatoren}{\boxplus}%
}
\begin{tabproofwide}
  \proofstepwidestarbreakkeep[1]{%
    \CompleteOrderedField{K,\boxplus,\boxtimes,
      \boxminus,\operatorname{inv},0_K,1_K,\preccurlyeq}}{\rA}
  \proofstepwidestarbreakkeep[1]{0_K\in K\land1_K\in K}{%
    \FormulaRefAuto{CompleteOrderedFieldAxioms}[def]{1}}
  \proofstepwidestarbreakkeep[1]{\boxplus\colon K\times K\to K}{%
    \FormulaRefAuto{CompleteOrderedFieldAxioms}[def]{1}}
  \proofstepwidestarbreakkeep[1]{%
    \forall x\in K\,(x\boxplus1_K\in K)}{%
    \FormulaRefAuto{CompleteOrderedFieldAxioms}[def]{1,3,\rAEb{2}}}
  \proofstepwidestarbreakkeep[1]{%
    \exists!s_K\,\bigl(s_K\colon K\to K\land
      \forall x\in K\,(s_K(x)=x\boxplus1_K)\bigr)}{%
    \FormulaRefAuto{TypedFunctionDefinitionPrinciple}[thm]{4}}
  \proofstepwidestarbreakkeep[1]{%
    s_K\colon K\to K\land
      \forall x\in K\,(s_K(x)=x\boxplus1_K)}{\rA}
  \proofstepwidestarbreakkeep[1]{%
    \exists!\eta\colon\mathbb N\to K\,
    (\eta(0)=0_K\land
     \forall n\in\mathbb N\,(\eta(n+1)=s_K(\eta(n))))}{%
    \FormulaRefAuto{DedekindRecursionTheoremAddOne}[thm]{\rAEa{2},\rAEa{6}}}
  \proofstepwidestarbreakkeep[1]{%
    \exists!\eta\colon\mathbb N\to K\,
    (\eta(0)=0_K\land
     \forall n\in\mathbb N\,
     (\eta(n+1)=\eta(n)\boxplus1_K))}{%
    \rIE{\rAEb{6},7}}
  \proofstepwidestarbreakkeep[1]{%
    \exists!\eta\colon\mathbb N\to K\,
    (\eta(0)=0_K\land
     \forall n\in\mathbb N\,
     (\eta(n+1)=\eta(n)\boxplus1_K))}{\UEE{5,6,8}}
\end{tabproofwide}

Wir schreiben \(\eta_K\) für die durch die beiden Rekursionsgleichungen
eindeutig bestimmte Funktion.

\FormulaThmDeltaK[Archimedizität vollständig angeordneter Körper]{%
  \begin{aligned}[t]
  &\CompleteOrderedField{\begin{gathered}
      K,\boxplus,\boxtimes,\boxminus,\\[-2pt]
      \operatorname{inv},0_K,1_K,\preccurlyeq
    \end{gathered}}\dsep{}\\[-2pt]
  &\eta_K\colon\mathbb N\to K\dsep
    \eta_K(0)=0_K\dsep{}\\[-2pt]
  &\forall n\in\mathbb N\,
      \bigl(\eta_K(n+1)=\eta_K(n)\boxplus1_K\bigr)\\[-2pt]
  &\quad\vdash
    \operatorname{Arch}(K,\preccurlyeq,\eta_K).
  \end{aligned}%
}{CompleteOrderedFieldArchimedean}{%
  \DeltaRow{Mengen}{K\dsep n\dsep s\dsep u\dsep x\dsep y\dsep z}%
  \DeltaRow{Funktionen}{\eta_K}%
  \DeltaRow{Relationsoperatoren}{\preccurlyeq}%
}
\begin{tabproofsplitwide}
  \proofpartwideindR[Strikte Translation]{%
    \begin{aligned}[t]
      &\CompleteOrderedField{K,\boxplus,\boxtimes,\boxminus,
        \operatorname{inv},0_K,1_K,\preccurlyeq}\dsep
        x,y,z\in K\dsep{}\\[-2pt]
      &x\preccurlyeq y\dsep x\neq y
        \vdash
        x\boxplus z\preccurlyeq y\boxplus z\land
        x\boxplus z\neq y\boxplus z
    \end{aligned}
  }{%
    \CompleteOrderedField{K,\boxplus,\boxtimes,\boxminus,
      \operatorname{inv},0_K,1_K,\preccurlyeq}\dsep
    x,y,z\in K\dsep x\preccurlyeq y\dsep x\neq y
    \vdash
    x\boxplus z\preccurlyeq y\boxplus z\land
    x\boxplus z\neq y\boxplus z
  }[CompleteOrderedFieldStrictTranslation]
    \proofstepwidestar[1]{\CompleteOrderedField{K,\boxplus,\boxtimes,
      \boxminus,\operatorname{inv},0_K,1_K,\preccurlyeq}}{\rA}
    \proofstepwidestar[2]{x,y,z\in K}{\rA}
    \proofstepwidestar[3]{x\preccurlyeq y}{\rA}
    \proofstepwidestar[4]{x\neq y}{\rA}
    \proofstepwidestar[1,2,3]{%
      x\boxplus z\preccurlyeq y\boxplus z}{%
      \FormulaRefAuto{CompleteOrderedFieldAxioms}[def]{1,2,3}}
    \proofstepwidestar[1,2]{\boxminus z\in K}{%
      \FormulaRefAuto{CompleteOrderedFieldAxioms}[def]{1,2}}
    \proofstepwidestar[7]{x\boxplus z=y\boxplus z}{\rA}
    \proofstepwidestar[1,2,7]{%
      (x\boxplus z)\boxplus(\boxminus z)
      =(y\boxplus z)\boxplus(\boxminus z)}{\rIE{7,\rII}}
    \proofstepwidestar[1,2,7]{\begin{aligned}[t]&(x\boxplus z)\boxplus(\boxminus z)=x\land{}\\[-2pt]&(y\boxplus z)\boxplus(\boxminus z)=y\end{aligned}}{%
      \FormulaRefAuto{CompleteOrderedFieldAxioms}[def]{1,2,6}}
    \proofstepwidestar[1,2,7]{x=y}{%
      \rIE{8,9}}
    \proofstepwidestar[1,2,4,7]{\bot}{\rBI{4,10}}
    \proofstepwidestar[1,2,4]{x\boxplus z\neq y\boxplus z}{%
      \rCI{7,11}}
    \proofstepwidestar[1,2,3,4]{%
      x\boxplus z\preccurlyeq y\boxplus z\land
      x\boxplus z\neq y\boxplus z}{\rAI{5,12}}
  \closeproofpartwideindR

  \proofpartwideindR[Nichtoberschranken liefern strikte Bildzeugen]{%
    \begin{aligned}[t]
      &\TotOrd{K,\preccurlyeq}\dsep
        \eta_K\colon\mathbb N\to K\dsep u\in K\dsep
        u\notin\UB_{K,\preccurlyeq}(\eta_K[\mathbb N])\\[-2pt]
      &\quad\vdash
        \exists n\in\mathbb N\,
        (u\preccurlyeq\eta_K(n)\land u\neq\eta_K(n))
    \end{aligned}
  }{%
    \TotOrd{K,\preccurlyeq}\dsep
    \eta_K\colon\mathbb N\to K\dsep u\in K\dsep
    u\notin\UB_{K,\preccurlyeq}(\eta_K[\mathbb N])
    \vdash
    \exists n\in\mathbb N\,
    (u\preccurlyeq\eta_K(n)\land u\neq\eta_K(n))
  }[CompleteOrderedFieldNaturalImageNotUpperWitness]
    \proofstepwidestarbreakkeep[1]{\TotOrd{K,\preccurlyeq}}{\rA}
    \proofstepwidestarbreakkeep[2]{\eta_K\colon\mathbb N\to K}{\rA}
    \proofstepwidestarbreakkeep[3]{u\in K}{\rA}
    \proofstepwidestarbreakkeep[4]{%
      u\notin\UB_{K,\preccurlyeq}(\eta_K[\mathbb N])}{\rA}
    \proofstepwidestarbreakkeep[2]{\eta_K[\mathbb N]\subseteq K}{%
      \FormulaRefAuto{F\colon A\to B\vdash F[A]\subseteq B}{2}}
    \proofstepwidestarbreakkeep[6]{%
      \forall x\in\eta_K[\mathbb N]\,(x\preccurlyeq u)}{\rA}
    \proofstepwidestarbreakkeep[2,3,6]{%
      u\in\UB_{K,\preccurlyeq}(\eta_K[\mathbb N])}{%
      \FormulaRefAuto{UpperBoundSetIntro}[thm]{5,3,6}}
    \proofstepwidestarbreakkeep[2,3,4,6]{\bot}{\rBI{4,7}}
    \proofstepwidestarbreakkeep[2,3,4]{%
      \neg\forall x\in\eta_K[\mathbb N]\,(x\preccurlyeq u)}{%
      \rCI{6,8}}
    \proofstepwidestarbreakkeep[2,3,4]{%
      \exists x\in\eta_K[\mathbb N]\,\neg(x\preccurlyeq u)}{%
      \FormulaRefAuto{\neg\forall x(P(x)\rightarrow Q(x))
        \eqvdash \exists x(P(x)\land \neg Q(x))}[thm]{9}}
    \proofstepwidestarbreakkeep[11]{%
      x\in\eta_K[\mathbb N]\land\neg(x\preccurlyeq u)}{\rA}
    \proofstepwidestarbreakkeep[2,11]{x\in K}{%
      \FormulaRefAuto{A\subseteq B\dsep x\in A\vdash x\in B}[thm]{%
        5,\rAEa{11}}}
    \proofstepwidestarbreakkeep[1,3,11]{%
      x\preccurlyeq u\lor u\preccurlyeq x}{%
      \FormulaRefAuto{Totale Ordnung}[def]{1,3,12}}
    \proofstepwidestarbreakkeep[14]{x\preccurlyeq u}{\rA}
    \proofstepwidestarbreakkeep[11,14]{\bot}{\rBI{\rAEb{11},14}}
    \proofstepwidestarbreakkeep[11,14]{u\preccurlyeq x}{%
      \FormulaRefAuto{\bot\vdash P}[thm]{15}}
    \proofstepwidestarbreakkeep[17]{u\preccurlyeq x}{\rA}
    \proofstepwidestarbreakkeep[1,3,11]{u\preccurlyeq x}{%
      \rOE{13,14,16,17,17}}
    \proofstepwidestarbreakkeep[19]{u=x}{\rA}
    \proofstepwidestarbreakkeep[1,3]{u\preccurlyeq u}{%
      \FormulaRefAuto{Totale Ordnung}[def]{1,3}}
    \proofstepwidestarbreakkeep[1,3,19]{x\preccurlyeq u}{\rIE{19,20}}
    \proofstepwidestarbreakkeep[1,3,11,19]{\bot}{\rBI{\rAEb{11},21}}
    \proofstepwidestarbreakkeep[1,3,11]{u\neq x}{\rCI{19,22}}
    \proofstepwidestarbreakkeep[2,11]{%
      \exists n\in\mathbb N\,x=\eta_K(n)}{%
      \FormulaRefAuto{F\colon A\to B\dsep y\in F[A]
        \vdash \exists x\in A\,y=F(x)}{2,\rAEa{11}}}
    \proofstepwidestarbreakkeep[25]{n\in\mathbb N\land x=\eta_K(n)}{\rA}
    \proofstepwidestarbreakkeep[1,3,11,25]{%
      u\preccurlyeq\eta_K(n)\land u\neq\eta_K(n)}{%
      \rIE{\rAEb{25},\rAI{18,23}}}
    \proofstepwidestarbreakkeep[1,2,3,11,25]{%
      \exists n\in\mathbb N\,
      (u\preccurlyeq\eta_K(n)\land u\neq\eta_K(n))}{%
      \rEI{\rAI{\rAEa{25},26}}}
    \proofstepwidestarbreakkeep[1,2,3,11]{%
      \exists n\in\mathbb N\,
      (u\preccurlyeq\eta_K(n)\land u\neq\eta_K(n))}{%
      \rEE{24,25,27}}
    \proofstepwidestarbreakkeep[1,2,3,4]{%
      \exists n\in\mathbb N\,
      (u\preccurlyeq\eta_K(n)\land u\neq\eta_K(n))}{%
      \rEE{10,11,28}}
  \closeproofpartwideindR

  \proofpartwideindR[Das natürliche Bild ist unbeschränkt]{%
    \begin{aligned}[t]
      &\CompleteOrderedField{K,\boxplus,\boxtimes,\boxminus,
        \operatorname{inv},0_K,1_K,\preccurlyeq}\dsep
        \eta_K\colon\mathbb N\to K\\[-2pt]
      &\eta_K(0)=0_K\dsep
        \forall n\in\mathbb N\,
        (\eta_K(n+1)=\eta_K(n)\boxplus1_K)\\[-2pt]
      &\quad\vdash\UB_{K,\preccurlyeq}(\eta_K[\mathbb N])=\varnothing
    \end{aligned}
  }{%
    \CompleteOrderedField{K,\boxplus,\boxtimes,\boxminus,
      \operatorname{inv},0_K,1_K,\preccurlyeq}\dsep
    \eta_K\colon\mathbb N\to K\dsep\eta_K(0)=0_K\dsep
    \forall n\in\mathbb N\,
    (\eta_K(n+1)=\eta_K(n)\boxplus1_K)
    \vdash\UB_{K,\preccurlyeq}(\eta_K[\mathbb N])=\varnothing
  }[CompleteOrderedFieldNaturalImageUnbounded]
    \proofstepwidestar[1]{\CompleteOrderedField{K,\boxplus,\boxtimes,
      \boxminus,\operatorname{inv},0_K,1_K,\preccurlyeq}}{\rA}
    \proofstepwidestar[2]{\eta_K\colon\mathbb N\to K}{\rA}
    \proofstepwidestar[3]{\eta_K(0)=0_K}{\rA}
    \proofstepwidestar[4]{\forall n\in\mathbb N\,
      (\eta_K(n+1)=\eta_K(n)\boxplus1_K)}{\rA}
    \proofstepwidestar[5]{\UB_{K,\preccurlyeq}(\eta_K[\mathbb N])\neq\varnothing}{\rA}
    \proofstepwidestar[]{0 \in \mathbb N}{%
      \FormulaRefAuto{PeanoZero}[ax]}
    \proofstepwidestarbreakkeep[]{\mathbb N\in\PowNE{\mathbb N}}{%
      \FormulaRefAuto{NonemptyPowersetMembership}[thm]{%
        6}}
    \proofstepwidestarbreakkeep[2]{\eta_K[\mathbb N]\in\PowNE K}{%
      \FormulaRefAuto{NonemptyDirectImage}[thm]{2,7}}
    \proofstepwidestarbreakkeep[2]{%
      \eta_K[\mathbb N]\subseteq K\land
      \eta_K[\mathbb N]\neq\varnothing}{%
      \FormulaRefAuto{NonemptyPowersetMembership}[thm]{8}}
    \proofstepwidestar[1,2,5]{\varnothing \neq { \eta _ K [ \mathbb N ] }}{%
      \FormulaRefAuto{a\neq b\vdash b\neq a}[thm]{\rAEb{9}}}
    \proofstepwidestarbreakkeep[1,2,5]{%
      \exists s\in\UB_{K,\preccurlyeq}(\eta_K[\mathbb N])\,
      \forall u\in\UB_{K,\preccurlyeq}(\eta_K[\mathbb N])\,(s\preccurlyeq u)}{%
      \FormulaRefAuto{CompleteOrderedFieldAxioms}[def]{%
        1,\rAI{%
          10,
          \rAEa{9}},5}}
    \proofstepwidestar[12]{%
      s\in\UB_{K,\preccurlyeq}(\eta_K[\mathbb N])\land
      \forall u\in\UB_{K,\preccurlyeq}(\eta_K[\mathbb N])
        (s\preccurlyeq u)}{\rA}
    \proofstepwidestarbreakkeep[1,2,12]{%
      s\in K\land
      \forall x\in\eta_K[\mathbb N]\,(x\preccurlyeq s)}{%
      \FormulaRefAuto{UpperBoundSetElim}[thm]{\rAEa{9},\rAEa{12}}}
    \proofstepwidestarbreakkeep[1]{\boxminus1_K\in K}{%
      \FormulaRefAuto{CompleteOrderedFieldAxioms}[def]{1}}
    \proofstepwidestarbreakkeep[1,12]{s\boxplus(\boxminus1_K)\in K}{%
      \FormulaRefAuto{CompleteOrderedFieldAxioms}[def]{1,\rAEa{13},14}}
    \proofstepwidestarbreakkeep[1]{0_K\preccurlyeq1_K\land0_K\neq1_K}{%
      \FormulaRefAuto{CompleteOrderedFieldAxioms}[def]{1}}
    \proofstepwidestar[1,12]{0_K,1_K\in K}{%
      \FormulaRefAuto{CompleteOrderedFieldAxioms}[def]{1}}
    \proofstepwidestar[1,12]{\begin{aligned}[t]&0_K\boxplus(s\boxplus(\boxminus1_K))=s\boxplus(\boxminus1_K)\land{}\\[-2pt]&1_K\boxplus(s\boxplus(\boxminus1_K))=s\end{aligned}}{%
      \FormulaRefAuto{CompleteOrderedFieldAxioms}[def]{1,\rAEa{13},14}}
    \proofstepwidestar[1,12]{\begin{aligned}[t]&0_K\boxplus(s\boxplus(\boxminus1_K))\preccurlyeq1_K\boxplus(s\boxplus(\boxminus1_K))\land{}\\[-2pt]&0_K\boxplus(s\boxplus(\boxminus1_K))\neq1_K\boxplus(s\boxplus(\boxminus1_K))\end{aligned}}{%
      \FormulaRefAuto{CompleteOrderedFieldStrictTranslation}[thm]{%
          1,17,15,
          \rAEa{16},\rAEb{16}}}
    \proofstepwidestarbreakkeep[1,12]{%
      \begin{aligned}[t]
      &s\boxplus(\boxminus1_K)\preccurlyeq s\land{}\\[-2pt]
      &s\boxplus(\boxminus1_K)\neq s
      \end{aligned}}{%
      \rIE{%
        19,
        18}}
    \proofstepwidestar[21]{%
      s\boxplus(\boxminus1_K)\in
      \UB_{K,\preccurlyeq}(\eta_K[\mathbb N])}{\rA}
    \proofstepwidestarbreakkeep[12,21]{s\preccurlyeq s\boxplus(\boxminus1_K)}{%
      \rRE{\rUE{\rAEb{12}},21}}
    \proofstepwidestar[1,12,21]{\TotOrd{K,\preccurlyeq}}{%
      \FormulaRefAuto{CompleteOrderedFieldAxioms}[def]{1}}
    \proofstepwidestarbreakkeep[1,12,21]{s=s\boxplus(\boxminus1_K)}{%
      \FormulaRefAuto{Totale Ordnung}[def]{%
        23,
        \rAEa{13},15,22,\rAEa{20}}}
    \proofstepwidestar[1,12,21]{{ s \boxplus ( \boxminus 1 _ K ) } = s}{%
      \FormulaRefAuto{a=b\vdash b=a}[thm]{24}}
    \proofstepwidestarbreakkeep[1,12,21]{\bot}{%
      \rBI{\rAEb{20},
        25}}
    \proofstepwidestarbreakkeep[1,12]{%
      s\boxplus(\boxminus1_K)\notin
      \UB_{K,\preccurlyeq}(\eta_K[\mathbb N])}{\rCI{21,26}}
    \proofstepwidestar[1,2,12]{\TotOrd{K,\preccurlyeq}}{%
      \FormulaRefAuto{CompleteOrderedFieldAxioms}[def]{1}}
    \proofstepwidestarbreakkeep[1,2,12]{%
      \exists n\in\mathbb N\,
      (s\boxplus(\boxminus1_K)\preccurlyeq\eta_K(n)\land
       s\boxplus(\boxminus1_K)\neq\eta_K(n))}{%
      \FormulaRefAuto{CompleteOrderedFieldNaturalImageNotUpperWitness}[thm]{%
        28,2,15,27}}
    \proofstepwidestar[30]{n\in\mathbb N\land
      s\boxplus(\boxminus1_K)\preccurlyeq\eta_K(n)\land
      s\boxplus(\boxminus1_K)\neq\eta_K(n)}{\rA}
    \proofstepwidestarbreakkeep[2,30]{\eta_K(n)\in K}{%
      \FormulaRefAuto{F\colon A\to B\dsep x\in A\vdash F(x)\in B}{%
        2,\rAEa{30}}}
    \proofstepwidestarbreakkeep[30]{(n+1)\in\mathbb N}{%
      \FormulaRefAuto{PeanoAddOneClosure}[thm]{\rAEa{30}}}
    \proofstepwidestar[2,30]{\mathbb N \subseteq \mathbb N}{%
      \FormulaRefAuto{A\subseteq A}[thm]}
    \proofstepwidestarbreakkeep[2,30]{\eta_K(n+1)\in\eta_K[\mathbb N]}{%
      \FormulaRefAuto{C\subseteq A\dsep F\colon A\to B\dsep
        x\in C\vdash F(x)\in F[C]}[thm]{%
        33,2,32}}
    \proofstepwidestar[1,2,30]{1_K\in K}{%
      \FormulaRefAuto{CompleteOrderedFieldAxioms}[def]{1}}
    \proofstepwidestarbreakkeep[1,2,30]{%
      \begin{aligned}[t]
      &(s\boxplus(\boxminus1_K))\boxplus1_K
        \preccurlyeq\eta_K(n)\boxplus1_K\land{}\\[-2pt]
      &(s\boxplus(\boxminus1_K))\boxplus1_K
        \neq\eta_K(n)\boxplus1_K
      \end{aligned}}{%
      \FormulaRefAuto{CompleteOrderedFieldStrictTranslation}[thm]{%
        1,15,31,35,
        \rAEa{\rAEb{30}},\rAEb{\rAEb{30}}}}
    \proofstepwidestar[1,2,4,12,30]{(s\boxplus(\boxminus1_K))\boxplus1_K=s}{%
      \FormulaRefAuto{CompleteOrderedFieldAxioms}[def]{1,\rAEa{13},14}}
    \proofstepwidestarbreakkeep[1,2,4,12,30]{%
      s\preccurlyeq\eta_K(n+1)\land
      s\neq\eta_K(n+1)}{%
      \rIE{36,
        37,
        \rRE{\rUE{4},\rAEa{30}}}}
    \proofstepwidestarbreakkeep[1,2,12,30]{\eta_K(n+1)\preccurlyeq s}{%
      \rRE{\rUE{\rAEb{13}},34}}
    \proofstepwidestarbreakkeep[1,2,12,30]{s=\eta_K(n+1)}{%
      \FormulaRefAuto{Totale Ordnung}[def]{1,\rAEa{13},34,\rAEa{38},39}}
    \proofstepwidestarbreakkeep[1,2,4,12,30]{\bot}{\rBI{\rAEb{38},40}}
    \proofstepwidestarbreakkeep[1,2,4,12]{\bot}{\rEE{29,30,41}}
    \proofstepwidestarbreakkeep[1,2,3,4,5,12]{\bot}{\rEE{11,12,42}}
    \proofstepwidestarbreakkeep[1,2,3,4]{%
      \UB_{K,\preccurlyeq}(\eta_K[\mathbb N])=\varnothing}{\rCI{5,43}}
  \closeproofpartwideindR

  \proofpartwide{Archimedisches Prädikat}
    \proofstepwidestar[1]{\CompleteOrderedField{K,\boxplus,\boxtimes,
      \boxminus,\operatorname{inv},0_K,1_K,\preccurlyeq}}{\rA}
    \proofstepwidestar[2]{\eta_K\colon\mathbb N\to K}{\rA}
    \proofstepwidestar[3]{\eta_K(0)=0_K}{\rA}
    \proofstepwidestar[4]{\forall n\in\mathbb N\,
      (\eta_K(n+1)=\eta_K(n)\boxplus1_K)}{\rA}
    \proofstepwidestarbreakkeep[1,2,3,4]{%
      \UB_{K,\preccurlyeq}(\eta_K[\mathbb N])=\varnothing}{%
      \FormulaRefAuto{CompleteOrderedFieldNaturalImageUnbounded}[thm]{1,2,3,4}}
    \proofstepwidestar[6]{x\in K}{\rA}
    \proofstepwidestar[1,2,3,4,6]{x\notin\varnothing}{%
      \FormulaRefAuto{x\not\in\varnothing}[thm]}
    \proofstepwidestarbreakkeep[1,2,3,4,6]{%
      x\notin\UB_{K,\preccurlyeq}(\eta_K[\mathbb N])}{%
      \rIE{5,7}}
    \proofstepwidestar[1,2,3,4,6]{\TotOrd{K,\preccurlyeq}}{%
      \FormulaRefAuto{CompleteOrderedFieldAxioms}[def]{1}}
    \proofstepwidestarbreakkeep[1,2,3,4,6]{%
      \exists n\in\mathbb N\,
      (x\preccurlyeq\eta_K(n)\land x\neq\eta_K(n))}{%
      \FormulaRefAuto{CompleteOrderedFieldNaturalImageNotUpperWitness}[thm]{%
        9,2,6,8}}
    \proofstepwidestarbreakkeep[1,2,3,4]{%
      \forall x\in K\,\exists n\in\mathbb N\,
      (x\preccurlyeq\eta_K(n)\land x\neq\eta_K(n))}{%
      \rUI{\rRI{6,10}}}
    \proofstepwidestarbreakkeep[1,2,3,4]{%
      \operatorname{Arch}(K,\preccurlyeq,\eta_K)}{%
      \FormulaRefAuto{ArchimedeanOrderedEmbedding}[def]{11}}
  \closeproofpartwide
\end{tabproofsplitwide}

\FormulaDefDeltaK[Gesetze der ganzzahligen Primkörpereinbettung]{%
  \begin{aligned}[t]
  \IntegerPrimeEmbeddingLaws{\theta_{\mathbb Z}}\coloneqq{}&
    \theta_{\mathbb Z}(\IntZero)=0_K\land
    \theta_{\mathbb Z}(\IntOne)=1_K\land{}\\[-2pt]
  &\forall z,w\in\mathbb Z\,\bigl(
    \theta_{\mathbb Z}(z+w)=
      \theta_{\mathbb Z}(z)\boxplus\theta_{\mathbb Z}(w)\land{}\\[-2pt]
  &\qquad\theta_{\mathbb Z}(zw)=
      \theta_{\mathbb Z}(z)\boxtimes\theta_{\mathbb Z}(w)\land
    \theta_{\mathbb Z}(-z)=\boxminus\theta_{\mathbb Z}(z)\land{}\\[-2pt]
  &\qquad\bigl(z<w\leftrightarrow
    (\theta_{\mathbb Z}(z)\preccurlyeq\theta_{\mathbb Z}(w)\land
     \theta_{\mathbb Z}(z)\neq\theta_{\mathbb Z}(w))\bigr)\bigr).
  \end{aligned}%
}{CompleteOrderedFieldIntegerPrimeEmbeddingLaws}{%
  \DeltaRow{Mengen}{K\dsep z\dsep w}%
  \DeltaRow{Funktionen}{\theta_{\mathbb Z}}%
  \DeltaRow{Relationsoperatoren}{<\dsep\preccurlyeq}%
  \DeltaRow{Funktionsoperatoren}{\boxplus\dsep\boxtimes\dsep\boxminus}%
  \DeltaRow{Prädikat}{\mathsf{ZEmb}_{K}}%
}

\FormulaDefDeltaK[Gesetze der rationalen Primkörpereinbettung]{%
  \begin{aligned}[t]
  \RationalPrimeEmbeddingLaws{\theta_{\mathbb Q}}\coloneqq{}&
    \theta_{\mathbb Q}(\RatZero)=0_K\land
    \theta_{\mathbb Q}(\RatOne)=1_K\land{}\\[-2pt]
  &\forall p,q\in\mathbb Q\,\bigl(
    \theta_{\mathbb Q}(p+q)=
      \theta_{\mathbb Q}(p)\boxplus\theta_{\mathbb Q}(q)\land{}\\[-2pt]
  &\qquad\theta_{\mathbb Q}(pq)=
      \theta_{\mathbb Q}(p)\boxtimes\theta_{\mathbb Q}(q)\land{}\\[-2pt]
  &\qquad\bigl(p<q\leftrightarrow
    (\theta_{\mathbb Q}(p)\preccurlyeq\theta_{\mathbb Q}(q)\land
     \theta_{\mathbb Q}(p)\neq\theta_{\mathbb Q}(q))\bigr)\bigr)\land{}\\[-2pt]
  &\forall p\in\mathbb Q\,\bigl(
    \theta_{\mathbb Q}(-p)=\boxminus\theta_{\mathbb Q}(p)\land{}\\[-2pt]
  &\qquad(p\neq\RatZero\rightarrow
    \theta_{\mathbb Q}(p^{-1})=
      \operatorname{inv}(\theta_{\mathbb Q}(p)))\bigr).
  \end{aligned}%
}{CompleteOrderedFieldRationalPrimeEmbeddingLaws}{%
  \DeltaRow{Mengen}{K\dsep p\dsep q}%
  \DeltaRow{Funktionen}{\theta_{\mathbb Q}}%
  \DeltaRow{Relationsoperatoren}{<\dsep\preccurlyeq}%
  \DeltaRow{Funktionsoperatoren}{\boxplus\dsep\boxtimes\dsep\boxminus}%
  \DeltaRow{Funktionsoperatoren}{\operatorname{inv}}%
  \DeltaRow{Prädikat}{\mathsf{QEmb}_{K}}%
}

\FormulaThmDeltaK[Eindeutigkeit des vollständig angeordneten reellen Körpers]{%
  \begin{aligned}[t]
  &\CompleteOrderedField{K,\boxplus,\boxtimes,
    \boxminus,\operatorname{inv},0_K,1_K,\preccurlyeq}\\[-2pt]
  &\quad\vdash\exists!\Phi\colon\mathbb R\bij K\,
    \left(
      \begin{aligned}[t]
      &\forall x,y\in\mathbb R\,
        \bigl(\Phi(x\RealAdd y)=\Phi(x)\boxplus\Phi(y)\bigr)\land{}\\[-2pt]
      &\forall x,y\in\mathbb R\,
        \bigl(\Phi(x\RealMul y)=\Phi(x)\boxtimes\Phi(y)\bigr)\land{}\\[-2pt]
      &\forall x,y\in\mathbb R\,
        \bigl(x\RealLe y\leftrightarrow\Phi(x)\preccurlyeq\Phi(y)\bigr)\land{}\\[-2pt]
      &\Phi(\RealZero)=0_K\land\Phi(\RealOne)=1_K
      \end{aligned}
    \right).
  \end{aligned}%
}{CompleteOrderedFieldUniqueness}{%
  \DeltaRow{Mengen}{K\dsep x\dsep y}%
  \DeltaRow{Funktionen}{\Phi}%
  \DeltaRow{Relationsoperatoren}{\RealLe\dsep\preccurlyeq}%
  \DeltaRow{Funktionsoperatoren}{\RealAdd\dsep\RealMul
    \dsep\boxplus\dsep\boxtimes}%
}
\begin{tabproofsplitwide}
  \proofpartwideindR[Arithmetik der nat\"urlichen Einbettung]{%
    \begin{aligned}[t]
      &\CompleteOrderedField{K,\boxplus,\boxtimes,\boxminus,
        \operatorname{inv},0_K,1_K,\preccurlyeq}\dsep{}\\[-2pt]
      &\eta_K\colon\mathbb N\to K\dsep
        \eta_K(0)=0_K\dsep{}\\[-2pt]
      &\forall n\in\mathbb N\,
        \bigl(\eta_K(n+1)=\eta_K(n)\boxplus1_K\bigr)\\[-2pt]
      &\quad\vdash\forall m,n\in\mathbb N\,\bigl(
        \eta_K(m+n)=\eta_K(m)\boxplus\eta_K(n)\land{}\\[-2pt]
      &\qquad\eta_K(mn)=\eta_K(m)\boxtimes\eta_K(n)\land{}\\[-2pt]
      &\qquad(m<n\leftrightarrow
        (\eta_K(m)\preccurlyeq\eta_K(n)\land
         \eta_K(m)\neq\eta_K(n)))\bigr)
    \end{aligned}
  }{%
    \CompleteOrderedField{K,\boxplus,\boxtimes,\boxminus,
      \operatorname{inv},0_K,1_K,\preccurlyeq}\dsep
    \eta_K\colon\mathbb N\to K\dsep
    \eta_K(0)=0_K\dsep
    \forall n\in\mathbb N\,
      (\eta_K(n+1)=\eta_K(n)\boxplus1_K)
    \vdash\forall m,n\in\mathbb N\,
    (\eta_K(m+n)=\eta_K(m)\boxplus\eta_K(n)\land
     \eta_K(mn)=\eta_K(m)\boxtimes\eta_K(n)\land
     (m<n\leftrightarrow
      (\eta_K(m)\preccurlyeq\eta_K(n)\land
       \eta_K(m)\neq\eta_K(n))))
  }[CompleteOrderedFieldNaturalEmbeddingArithmetic]
    \proofstepwidestarbreakkeep[1]{%
      \CompleteOrderedField{K,\boxplus,\boxtimes,\boxminus,
        \operatorname{inv},0_K,1_K,\preccurlyeq}}{\rA}
    \proofstepwidestarbreakkeep[2]{\eta_K\colon\mathbb N\to K}{\rA}
    \proofstepwidestarbreakkeep[3]{\eta_K(0)=0_K}{\rA}
    \proofstepwidestarbreakkeep[4]{%
      \forall n\in\mathbb N\,
      \eta_K(n+1)=\eta_K(n)\boxplus1_K}{\rA}

    \proofstepwidestarbreakkeep[1,2,3,4]{\eta_K(1)=1_K}{%
      \FormulaRefAuto{PeanoAddLeftZero}[thm],4,3,
      \FormulaRefAuto{CompleteOrderedFieldAxioms}[def]{1}}
    \proofstepwidestarbreakkeep[1,2,3]{%
      \forall m\in\mathbb N\,
      \eta_K(m+0)=\eta_K(m)\boxplus\eta_K(0)}{%
      \FormulaRefAuto{PeanoAddRightZero}[thm],3,
      \FormulaRefAuto{CompleteOrderedFieldAxioms}[def]{1}}
    \proofstepwidestarbreakkeep[1,2,4]{\begin{aligned}[t]
      &\forall n\in\mathbb N\,\Bigl(\\[-2pt]
      &\quad\forall m\in\mathbb N\,(\eta_K(m+n)=\eta_K(m)\boxplus\eta_K(n))\\[-2pt]
      &\quad\rightarrow\forall m\in\mathbb N\,(\eta_K(m+(n+1))\\[-2pt]
      &\qquad=\eta_K(m)\boxplus\eta_K(n+1))\Bigr)
      \end{aligned}}{%
      \FormulaRefAuto{PeanoAddAddOneRight}[thm],4,
      \FormulaRefAuto{CompleteOrderedFieldAxioms}[def]{1}}
    \proofstepwidestarbreakkeep[1,2,3,4]{%
      \forall m,n\in\mathbb N\,
      \eta_K(m+n)=\eta_K(m)\boxplus\eta_K(n)}{%
      \FormulaRefAuto{PeanoInductionAddOne}[thm]{6,7}}
    \proofstepwidestarbreakkeep[1,2,3]{%
      \forall m\in\mathbb N\,
      \eta_K(m\cdot0)=\eta_K(m)\boxtimes\eta_K(0)}{%
      \FormulaRefAuto{PeanoMulRightZero}[thm],3,
      \FormulaRefAuto{CompleteOrderedFieldAxioms}[def]{1}}
    \proofstepwidestarbreakkeep[1,2,3,4]{\begin{aligned}[t]
      &\forall n\in\mathbb N\,\Bigl(\\[-2pt]
      &\quad\forall m\in\mathbb N\,(\eta_K(mn)=\eta_K(m)\boxtimes\eta_K(n))\\[-2pt]
      &\quad\rightarrow\forall m\in\mathbb N\,(\eta_K(m\cdot(n+1))\\[-2pt]
      &\qquad=\eta_K(m)\boxtimes\eta_K(n+1))\Bigr)
      \end{aligned}}{%
      \FormulaRefAuto{PeanoMulAddOneRight}[thm],4,8,5,
      \FormulaRefAuto{CompleteOrderedFieldAxioms}[def]{1}}
    \proofstepwidestarbreakkeep[1,2,3,4]{%
      \forall m,n\in\mathbb N\,
      \eta_K(mn)=\eta_K(m)\boxtimes\eta_K(n)}{%
      \FormulaRefAuto{PeanoInductionAddOne}[thm]{9,10}}
    \proofstepwidestarbreakkeep[1,2,4]{%
      \forall n\in\mathbb N\,
      (\eta_K(n)\preccurlyeq\eta_K(n+1)\land
       \eta_K(n)\neq\eta_K(n+1))}{%
      \FormulaRefAuto{CompleteOrderedFieldStrictTranslation}[thm]{1},4,
      \FormulaRefAuto{CompleteOrderedFieldAxioms}[def]{1}}
    \proofstepwidestarbreakkeep[]{%
      \forall m\in\mathbb N\,
      \bigl(m<0\rightarrow
       (\eta_K(m)\preccurlyeq\eta_K(0)\land
        \eta_K(m)\neq\eta_K(0))\bigr)}{%
      \FormulaRefAuto{PeanoNoLtZero}[thm],
      \FormulaRefAuto{\bot\vdash P}[thm]}
    \proofstepwidestarbreakkeep[1,2,4]{\begin{aligned}[t]
      &\forall n\in\mathbb N\,\Bigl(\\[-2pt]
      &\quad\forall m\in\mathbb N\,\bigl(m<n\rightarrow(\\[-2pt]
      &\qquad\eta_K(m)\preccurlyeq\eta_K(n)\land\eta_K(m)\neq\eta_K(n))\bigr)\\[-2pt]
      &\quad\rightarrow\forall m\in\mathbb N\,\bigl(m<n+1\rightarrow(\\[-2pt]
      &\qquad\eta_K(m)\preccurlyeq\eta_K(n+1)\land{}\\[-2pt]
      &\qquad\eta_K(m)\neq\eta_K(n+1))\bigr)\Bigr)
      \end{aligned}}{%
      \FormulaRefAuto{PeanoStrictAddOneSplit}[thm],12,
      \FormulaRefAuto{Totale Ordnung}[def]{1}}
    \proofstepwidestarbreakkeep[1,2,4]{%
      \forall m,n\in\mathbb N\,
      \bigl(m<n\rightarrow
       (\eta_K(m)\preccurlyeq\eta_K(n)\land
        \eta_K(m)\neq\eta_K(n))\bigr)}{%
      \FormulaRefAuto{PeanoInductionAddOne}[thm]{13,14}}
    \proofstepwidestarbreakkeep[1,2,4]{%
      \forall m,n\in\mathbb N\,
      \bigl((\eta_K(m)\preccurlyeq\eta_K(n)\land
        \eta_K(m)\neq\eta_K(n))\rightarrow m<n\bigr)}{%
      \FormulaRefAuto{PeanoOrderTrichotomy}[thm]{15},
      \FormulaRefAuto{Totale Ordnung}[def]{1}}
    \proofstepwidestarbreakkeep[1,2,3,4]{\begin{aligned}[t]
      &\forall m,n\in\mathbb N\,\bigl(\\[-2pt]
      &\quad\eta_K(m+n)=\eta_K(m)\boxplus\eta_K(n)\land{}\\[-2pt]
      &\quad\eta_K(mn)=\eta_K(m)\boxtimes\eta_K(n)\land{}\\[-2pt]
      &\quad(m<n\leftrightarrow(\\[-2pt]
      &\qquad\eta_K(m)\preccurlyeq\eta_K(n)\land\eta_K(m)\neq\eta_K(n)))\bigr)
      \end{aligned}}{\rAI{8,\rAI{11,\rLRI{15,16}}}}
  \closeproofpartwideindR

  \proofpartwideindR[Fortsetzung auf die ganzen Zahlen]{%
    \begin{aligned}[t]
      &\CompleteOrderedField{K,\boxplus,\boxtimes,\boxminus,
        \operatorname{inv},0_K,1_K,\preccurlyeq}\dsep{}\\[-2pt]
      &\quad\vdash\exists!\eta_{\mathbb Z}\colon\mathbb Z\inj K\,\bigl(
        \eta_{\mathbb Z}(\IntZero)=0_K\land
        \eta_{\mathbb Z}(\IntOne)=1_K\land{}\\[-2pt]
      &\qquad\forall z,w\in\mathbb Z\,(
        \eta_{\mathbb Z}(z+w)=
          \eta_{\mathbb Z}(z)\boxplus\eta_{\mathbb Z}(w)\land{}\\[-2pt]
      &\qquad\eta_{\mathbb Z}(zw)=
          \eta_{\mathbb Z}(z)\boxtimes\eta_{\mathbb Z}(w)\land
        \eta_{\mathbb Z}(-z)=\boxminus\eta_{\mathbb Z}(z)\land{}\\[-2pt]
      &\qquad(z<w\leftrightarrow
        (\eta_{\mathbb Z}(z)\preccurlyeq\eta_{\mathbb Z}(w)\land
         \eta_{\mathbb Z}(z)\neq\eta_{\mathbb Z}(w))))\bigr)
    \end{aligned}
  }{%
    \CompleteOrderedField{K,\boxplus,\boxtimes,\boxminus,
      \operatorname{inv},0_K,1_K,\preccurlyeq}
    \vdash\exists!\eta_{\mathbb Z}\colon\mathbb Z\inj K\,
    (\eta_{\mathbb Z}(\IntZero)=0_K\land
     \eta_{\mathbb Z}(\IntOne)=1_K\land
     \forall z,w\in\mathbb Z\,(
     \eta_{\mathbb Z}(z+w)=
       \eta_{\mathbb Z}(z)\boxplus\eta_{\mathbb Z}(w)\land
     \eta_{\mathbb Z}(zw)=
       \eta_{\mathbb Z}(z)\boxtimes\eta_{\mathbb Z}(w)\land
     \eta_{\mathbb Z}(-z)=\boxminus\eta_{\mathbb Z}(z)\land
     (z<w\leftrightarrow
      (\eta_{\mathbb Z}(z)\preccurlyeq\eta_{\mathbb Z}(w)\land
       \eta_{\mathbb Z}(z)\neq\eta_{\mathbb Z}(w)))))
  }[CompleteOrderedFieldIntegerPrimeEmbedding]
    \proofstepwidestarbreakkeep[1]{%
      \CompleteOrderedField{K,\boxplus,\boxtimes,\boxminus,
        \operatorname{inv},0_K,1_K,\preccurlyeq}}{\rA}
    \proofstepwidestarbreakkeep[1]{\begin{aligned}[t]
      &\eta_K\colon\mathbb N\inj K\land{}\\[-2pt]
      &\eta_K(0)=0_K\land\eta_K(1)=1_K\land{}\\[-2pt]
      &\forall m,n\in\mathbb N\,\bigl(\\[-2pt]
      &\quad\eta_K(m+n)=\eta_K(m)\boxplus\eta_K(n)\land{}\\[-2pt]
      &\quad\eta_K(mn)=\eta_K(m)\boxtimes\eta_K(n)\land{}\\[-2pt]
      &\quad(m<n\leftrightarrow(\\[-2pt]
      &\qquad\eta_K(m)\preccurlyeq\eta_K(n)\land\eta_K(m)\neq\eta_K(n)))\bigr)
      \end{aligned}}{%
      \FormulaRefAuto{CompleteOrderedFieldNaturalEmbedding}[thm]{1},
      \FormulaRefAuto{CompleteOrderedFieldNaturalEmbeddingArithmetic}[thm]{1}}
    \proofstepwidestarbreakkeep[1,2]{%
      \eta_{\mathbb Z}(\IntClass{a}{b})\coloneqq
      \eta_K(a)\boxplus\boxminus\eta_K(b)}{\rII}
    \proofstepwidestarbreakkeep[1,2]{%
      \begin{aligned}[t]
        &\IntClass{a}{b}=\IntClass{c}{d}\vdash{}\\[-2pt]
        &\eta_K(a)\boxplus\boxminus\eta_K(b)
        =\eta_K(c)\boxplus\boxminus\eta_K(d)
      \end{aligned}}{%
      \FormulaRefAuto{IntegerClassEquality}[thm],
      \FormulaRefAuto{CompleteOrderedFieldNaturalEmbeddingArithmetic}[thm]{1,2},
      \FormulaRefAuto{CompleteOrderedFieldAxioms}[def]{1}}
    \proofstepwidestarbreakkeep[1,2]{%
      \eta_{\mathbb Z}\colon\mathbb Z\to K}{%
      \FormulaRefAuto{IntegerPairRepresentative}[thm]{3,4}}
    \proofstepwidestarbreakkeep[1,2]{%
      \forall z,w\in\mathbb Z\,
      \eta_{\mathbb Z}(z+w)=
      \eta_{\mathbb Z}(z)\boxplus\eta_{\mathbb Z}(w)}{%
      \FormulaRefAuto{IntegerAdditionDef}[def],2,3,
      \FormulaRefAuto{IntegerPairRepresentative}[thm]}
    \proofstepwidestarbreakkeep[1,2]{%
      \forall z,w\in\mathbb Z\,
      \eta_{\mathbb Z}(zw)=
      \eta_{\mathbb Z}(z)\boxtimes\eta_{\mathbb Z}(w)}{%
      \FormulaRefAuto{IntegerMultiplicationDef}[def],2,3,
      \FormulaRefAuto{IntegerPairRepresentative}[thm],
      \FormulaRefAuto{CompleteOrderedFieldAxioms}[def]{1}}
    \proofstepwidestarbreakkeep[1,2]{%
      \eta_{\mathbb Z}(\IntZero)=0_K\land
      \eta_{\mathbb Z}(\IntOne)=1_K\land
      \forall z\in\mathbb Z\,
      \eta_{\mathbb Z}(-z)=\boxminus\eta_{\mathbb Z}(z)}{%
      \FormulaRefAuto{IntegerZeroClass}[thm],
      \FormulaRefAuto{IntegerOneDef}[def],
      \FormulaRefAuto{IntegerNegationDef}[def],2,3}
    \proofstepwidestarbreakkeep[1,2]{%
      \forall z,w\in\mathbb Z\,\bigl(
      z<w\leftrightarrow
      (\eta_{\mathbb Z}(z)\preccurlyeq\eta_{\mathbb Z}(w)\land
       \eta_{\mathbb Z}(z)\neq\eta_{\mathbb Z}(w))\bigr)}{%
      \FormulaRefAuto{IntegerPairRepresentative}[thm],
      \FormulaRefAuto{IntegerOrderClassChar}[thm],
      \FormulaRefAuto{IntegerStrictOrderDef}[def],2,3}
    \proofstepwidestarbreakkeep[1,5,9]{%
      \eta_{\mathbb Z}\colon\mathbb Z\inj K}{%
      \FormulaRefAuto{IntegerOrderTotal}[thm]{9},
      \FormulaRefAuto{Injektive Funktion}[def]{5}}
    \proofstepwidestarbreakkeep[\shortstack{1,2,5,6\\7,8,9}]{%
      \theta_{\mathbb Z}\colon\mathbb Z\to K\dsep
      \IntegerPrimeEmbeddingLaws{\theta_{\mathbb Z}}\vdash
      \forall z\in\mathbb Z\,
      \theta_{\mathbb Z}(z)=\eta_{\mathbb Z}(z)}{\text{\parbox[t]{\linewidth}{\raggedright \FormulaRefAuto{CompleteOrderedFieldIntegerPrimeEmbeddingLaws}[def] , \FormulaRefAuto{IntegerPairRepresentative}[thm] , \FormulaRefAuto{CompleteOrderedFieldNaturalEmbedding}[thm]{1}, \FormulaRefAuto{IntegerAdditionDef}[def]{3}}}}
    \proofstepwidestarbreakkeep[1]{%
      \exists!\eta_{\mathbb Z}\colon\mathbb Z\inj K\,
      \IntegerPrimeEmbeddingLaws{\eta_{\mathbb Z}}}{%
      \FormulaRefAuto{CompleteOrderedFieldIntegerPrimeEmbeddingLaws}[def],
      \UEI{5,6,7,8,9,10,11}}
  \closeproofpartwideindR

  \proofpartwideindR[Rationaler Primkörper in \(K\)]{%
    \begin{aligned}[t]
      &\CompleteOrderedField{K,\boxplus,\boxtimes,\boxminus,
        \operatorname{inv},0_K,1_K,\preccurlyeq}\\[-2pt]
      &\quad\vdash\exists!\eta_{\mathbb Q}\colon\mathbb Q\inj K\,
        \RationalPrimeEmbeddingLaws{\eta_{\mathbb Q}}
    \end{aligned}
  }{%
    \CompleteOrderedField{K,\boxplus,\boxtimes,\boxminus,
      \operatorname{inv},0_K,1_K,\preccurlyeq}
    \vdash\exists!\eta_{\mathbb Q}\colon\mathbb Q\inj K\,
    \RationalPrimeEmbeddingLaws{\eta_{\mathbb Q}}
  }[CompleteOrderedFieldRationalPrimeEmbedding]
    \proofstepwidestarbreakkeep[1]{\CompleteOrderedField{K,\boxplus,\boxtimes,
      \boxminus,\operatorname{inv},0_K,1_K,\preccurlyeq}}{\rA}
    \proofstepwidestarbreakkeep[1]{%
      \exists!\eta_{\mathbb Z}\colon\mathbb Z\inj K\,
      \IntegerPrimeEmbeddingLaws{\eta_{\mathbb Z}}}{%
      \FormulaRefAuto{CompleteOrderedFieldIntegerPrimeEmbedding}[thm]{1}}
    \proofstepwidestarbreakkeep[2]{%
      \eta_{\mathbb Z}\colon\mathbb Z\inj K\land
      \IntegerPrimeEmbeddingLaws{\eta_{\mathbb Z}}}{\rA}
    \proofstepwidestarbreakkeep[1,3]{%
      a\in\mathbb Z\dsep b\in\PositiveIntegers\vdash
      \eta_{\mathbb Z}(b)\neq0_K}{%
      \FormulaRefAuto{PositiveIntegersDef}[def],
      \FormulaRefAuto{CompleteOrderedFieldIntegerPrimeEmbedding}[thm]{1,3}}
    \proofstepwidestarbreakkeep[1,3]{%
      \eta_{\mathbb Q}(\RatClass{a}{b})\coloneqq
      \eta_{\mathbb Z}(a)\boxtimes
      \operatorname{inv}(\eta_{\mathbb Z}(b))}{\rII}
    \proofstepwidestarbreakkeep[1,3,4,5]{%
      \begin{aligned}[t]
      &\RatClass{a}{b}=\RatClass{c}{d}\vdash{}\\[-2pt]
      &\eta_{\mathbb Z}(a)\boxtimes
        \operatorname{inv}(\eta_{\mathbb Z}(b))
       =\eta_{\mathbb Z}(c)\boxtimes
        \operatorname{inv}(\eta_{\mathbb Z}(d))
      \end{aligned}}{%
      \FormulaRefAuto{RationalClassEquality}[thm],
      \FormulaRefAuto{RationalPairRepresentative}[thm],3,4,5,
      \FormulaRefAuto{CompleteOrderedFieldAxioms}[def]{1}}
    \proofstepwidestarbreakkeep[1,3,5,6]{%
      \eta_{\mathbb Q}\colon\mathbb Q\to K\quad\text{und}\quad
      \eta_{\mathbb Q}(\RatZero)=0_K\dsep
      \eta_{\mathbb Q}(\RatOne)=1_K}{\text{\parbox[t]{\linewidth}{\raggedright \FormulaRefAuto{RationalPairRepresentative}[thm]{5,6}, \FormulaRefAuto{RationalZeroOneClasses}[thm] , \FormulaRefAuto{CompleteOrderedFieldIntegerPrimeEmbedding}[thm]{1,3}, \FormulaRefAuto{CompleteOrderedFieldAxioms}[def]{1}}}}
    \proofstepwidestarbreakkeep[1,3,5,6,7]{%
      \begin{aligned}[t]
      &\forall p,q\in\mathbb Q\,\bigl(
        \eta_{\mathbb Q}(p+q)
          =\eta_{\mathbb Q}(p)\boxplus\eta_{\mathbb Q}(q)\land{}\\[-2pt]
      &\qquad\eta_{\mathbb Q}(pq)
          =\eta_{\mathbb Q}(p)\boxtimes\eta_{\mathbb Q}(q)\bigr),\\[-2pt]
      &\forall p\in\mathbb Q\,\bigl(
        \eta_{\mathbb Q}(-p)=\boxminus\eta_{\mathbb Q}(p)\land{}\\[-2pt]
      &\qquad(p\neq\RatZero\rightarrow
        \eta_{\mathbb Q}(p^{-1})=
        \operatorname{inv}(\eta_{\mathbb Q}(p)))\bigr)
      \end{aligned}}{\text{\parbox[t]{\linewidth}{\raggedright \FormulaRefAuto{RationalAdditionClassValue}[thm] , \FormulaRefAuto{RationalMultiplicationClassValue}[thm] , \FormulaRefAuto{RationalAddInverse}[thm] , \FormulaRefAuto{RationalMulInverse}[thm]{5,6}, \FormulaRefAuto{CompleteOrderedFieldIntegerPrimeEmbedding}[thm]{1,3}, \FormulaRefAuto{CompleteOrderedFieldAxioms}[def]{1}}}}
    \proofstepwidestarbreakkeep[1,3,5,6,7]{%
      \forall p,q\in\mathbb Q\,\bigl(p<q\leftrightarrow
      (\eta_{\mathbb Q}(p)\preccurlyeq\eta_{\mathbb Q}(q)
      \land\eta_{\mathbb Q}(p)\neq\eta_{\mathbb Q}(q))\bigr)}{\text{\parbox[t]{\linewidth}{\raggedright \FormulaRefAuto{RationalStrictOrderDef}[def] , \FormulaRefAuto{RationalOrderClassChar}[thm] \ensuremath{,\allowbreak 3,\allowbreak 5,\allowbreak 6,\allowbreak } \FormulaRefAuto{CompleteOrderedFieldIntegerPrimeEmbedding}[thm]{1,3}, \FormulaRefAuto{CompleteOrderedFieldAxioms}[def]{1}}}}
    \proofstepwidestarbreakkeep[1,7,8,9]{%
      \eta_{\mathbb Q}\colon\mathbb Q\inj K\land
      \RationalPrimeEmbeddingLaws{\eta_{\mathbb Q}}}{\text{\parbox[t]{\linewidth}{\raggedright \FormulaRefAuto{CompleteOrderedFieldRationalPrimeEmbeddingLaws}[def] , \FormulaRefAuto{Injektive Funktion}[def]{7,9}, punktweise Konjunktion der vollständig quantifizierten
          Regeln aus \ensuremath{7,\allowbreak 8,\allowbreak 9}}}}
    \proofstepwidestarbreakkeep[1,3,10]{%
      \theta\colon\mathbb Q\to K\dsep
      \RationalPrimeEmbeddingLaws{\theta}
      \vdash\forall q\in\mathbb Q\,
      \theta(q)=\eta_{\mathbb Q}(q)}{%
      \FormulaRefAuto{CompleteOrderedFieldRationalPrimeEmbeddingLaws}[def],
      \FormulaRefAuto{RationalPairRepresentative}[thm],
      \FormulaRefAuto{CompleteOrderedFieldIntegerPrimeEmbedding}[thm]{1,3},
      \FormulaRefAuto{CompleteOrderedFieldAxioms}[def]{1}}
    \proofstepwidestarbreakkeep[1]{%
      \exists!\eta_{\mathbb Q}\colon\mathbb Q\inj K\,
      \RationalPrimeEmbeddingLaws{\eta_{\mathbb Q}}}{%
      \UEI{2,3,10,11}}
  \closeproofpartwideindR

\end{tabproofsplitwide}

Für den im Folgenden jeweils festen vollständig angeordneten Körper schreiben
wir \(\eta_{\mathbb Z}\) und \(\eta_{\mathbb Q}\) für die durch
\FormulaRefAuto{CompleteOrderedFieldIntegerPrimeEmbedding} beziehungsweise
\FormulaRefAuto{CompleteOrderedFieldRationalPrimeEmbedding} eindeutig
bestimmten Zeugen. Zusammen mit der bereits festgelegten natürlichen
Einbettung \(\eta_K\) gelten damit stets deren Typ-, Rekursions-, Konstanten-,
Operations-, Inversen- und Ordnungsregeln; die folgenden Teilbeweise führen
die jeweils benötigten Regeln durch eine Formelreferenz ausdrücklich ein.

\begin{tabproofsplitwide}
\setcounter{proofpartnr}{3}% Fortsetzung des vorangehenden Teilbeweisblocks

  \proofpartwideindR[Dichtheit des rationalen Primk\"orpers]{%
    \begin{aligned}[t]
      &\CompleteOrderedField{K,\boxplus,\boxtimes,\boxminus,
        \operatorname{inv},0_K,1_K,\preccurlyeq}\dsep x,y\in K\\[-2pt]
      &\quad\dsep x\preccurlyeq y\dsep x\neq y\vdash
        \exists q\in\mathbb Q\,\bigl({}\\[-2pt]
      &\qquad x\preccurlyeq\eta_{\mathbb Q}(q)\land
        x\neq\eta_{\mathbb Q}(q)\land
        \eta_{\mathbb Q}(q)\preccurlyeq y\land
        \eta_{\mathbb Q}(q)\neq y\bigr)
    \end{aligned}
  }{%
    \CompleteOrderedField{K,\boxplus,\boxtimes,\boxminus,
      \operatorname{inv},0_K,1_K,\preccurlyeq}\dsep x,y\in K
    \dsep x\preccurlyeq y\dsep x\neq y\vdash
    \exists q\in\mathbb Q\,
    (x\preccurlyeq\eta_{\mathbb Q}(q)\land
     x\neq\eta_{\mathbb Q}(q)\land
     \eta_{\mathbb Q}(q)\preccurlyeq y\land
     \eta_{\mathbb Q}(q)\neq y)
  }[CompleteOrderedFieldRationalDense]
    \proofstepwidestarbreakkeep[1]{\CompleteOrderedField{K,\boxplus,\boxtimes,
      \boxminus,\operatorname{inv},0_K,1_K,\preccurlyeq}}{\rA}
    \proofstepwidestarbreakkeep[2]{x,y\in K}{\rA}
    \proofstepwidestarbreakkeep[3]{x\preccurlyeq y}{\rA}
    \proofstepwidestarbreakkeep[4]{x\neq y}{\rA}
    \proofstepwidestarbreakkeep[1]{%
      \begin{aligned}[t]
      &(\eta_K\colon\mathbb N\to K\land\eta_K(0)=0_K\land{}\\[-2pt]
      &\quad\forall n\in\mathbb N\,
        (\eta_K(n+1)=\eta_K(n)\boxplus1_K))\land{}\\[-2pt]
      &(\eta_{\mathbb Z}\colon\mathbb Z\inj K\land
        \IntegerPrimeEmbeddingLaws{\eta_{\mathbb Z}})\land{}\\[-2pt]
      &(\eta_{\mathbb Q}\colon\mathbb Q\inj K\land
        \RationalPrimeEmbeddingLaws{\eta_{\mathbb Q}})
      \end{aligned}}{%
      \FormulaRefAuto{CompleteOrderedFieldNaturalEmbedding}[thm]{1},
      \FormulaRefAuto{CompleteOrderedFieldIntegerPrimeEmbedding}[thm]{1},
      \FormulaRefAuto{CompleteOrderedFieldRationalPrimeEmbedding}[thm]{1}}
    \proofstepwidestarbreakkeep[1,2,3,4]{%
      0_K\preccurlyeq\delta\coloneqq y\boxplus(\boxminus x)
      \land0_K\neq\delta}{%
      \FormulaRefAuto{CompleteOrderedFieldStrictTranslation}[thm]{1,2,3,4}}
    \proofstepwidestarbreakkeep[1,6]{%
      \operatorname{inv}(\delta)\in K\land
      0_K\preccurlyeq\operatorname{inv}(\delta)
      \land0_K\neq\operatorname{inv}(\delta)}{%
      \FormulaRefAuto{CompleteOrderedFieldAxioms}[def]{1,6}}
    \proofstepwidestarbreakkeep[1,7]{%
      \exists n\in\mathbb N\,
      \bigl(\operatorname{inv}(\delta)\preccurlyeq\eta_K(n)
      \land\operatorname{inv}(\delta)\neq\eta_K(n)\bigr)}{%
      \FormulaRefAuto{CompleteOrderedFieldArchimedean}[thm]{1,5,7}}
    \proofstepwidestarbreakkeep[8]{%
      n\in\mathbb N\land
      \operatorname{inv}(\delta)\preccurlyeq\eta_K(n)
      \land\operatorname{inv}(\delta)\neq\eta_K(n)}{\rA}
    \proofstepwidestarbreakkeep[1,6\text{--}8]{%
      \begin{aligned}[t]
        &n\neq0\land
        0_K\preccurlyeq\operatorname{inv}(\eta_K(n))
        \land\operatorname{inv}(\eta_K(n))\neq0_K\land{}\\[-2pt]
        &\operatorname{inv}(\eta_K(n))\preccurlyeq\delta
        \land\operatorname{inv}(\eta_K(n))\neq\delta
      \end{aligned}}{%
      \FormulaRefAuto{CompleteOrderedFieldAxioms}[def]{1},
      \FormulaRefAuto{CompleteOrderedFieldNaturalEmbeddingArithmetic}[thm]{1,5}}
    \proofstepwidestarbreakkeep[1,2,9]{%
      \exists N\in\mathbb N\,
      \bigl(\boxminus\eta_K(N)\preccurlyeq
        x\boxtimes\eta_K(n)\land
        x\boxtimes\eta_K(n)\preccurlyeq\eta_K(N)\bigr)}{\text{\parbox[t]{\linewidth}{\raggedright Wende \FormulaRefAuto{CompleteOrderedFieldArchimedean}[thm] mit den Daten aus \ensuremath{5} auf \ensuremath{z=\allowbreak x\boxtimes\allowbreak \eta_K(n)} und auf \ensuremath{\boxminus z} an. Dies liefert \ensuremath{N_+,\allowbreak N_-\in\mathbb N} mit \ensuremath{z<\allowbreak \eta_K(N_+)} und \ensuremath{-z<\allowbreak \eta_K(N_-).} Setze \ensuremath{N=\allowbreak N_++N_-.} Die Positivität der natürlichen Bilder und \ensuremath{\eta_K(N)=\allowbreak \eta_K(N_+)\boxplus\allowbreak \eta_K(N_-)} liefern beide Schranken. \FormulaRefAuto{CompleteOrderedFieldNaturalEmbeddingArithmetic}[thm]{1,5}, \FormulaRefAuto{CompleteOrderedFieldAxioms}[def]{1}}}}
    \proofstepwidestarbreakkeep[11]{%
      N\in\mathbb N\land
      \boxminus\eta_K(N)\preccurlyeq x\boxtimes\eta_K(n)
      \land x\boxtimes\eta_K(n)\preccurlyeq\eta_K(N)}{\rA}
    \proofstepwidestarbreakkeep[1,2,9,12]{%
      \begin{aligned}[t]
      &M\coloneqq\{k\in\mathbb N\mid{}\\[-2pt]
       &\quad x\boxtimes\eta_K(n)
        \preccurlyeq\eta_K(k)\boxplus\boxminus\eta_K(N)\land{}\\[-2pt]
       &\quad x\boxtimes\eta_K(n)\neq
        \eta_K(k)\boxplus\boxminus\eta_K(N)\},\\[-2pt]
      &(N+N+1)\in M\quad\text{und damit}\quad M\neq\varnothing
      \end{aligned}}{\text{\parbox[t]{\linewidth}{\raggedright \ensuremath{x\boxtimes\allowbreak \eta_K(n)\preccurlyeq\allowbreak \eta_K(N)
          <\allowbreak \eta_K(N)\boxplus1_K} \ensuremath{=\allowbreak \eta_K(N+N+1)\boxplus\allowbreak \boxminus\eta_K(N),\allowbreak } \FormulaRefAuto{PeanoAddOneClosure}[thm] , \FormulaRefAuto{CompleteOrderedFieldNaturalEmbeddingArithmetic}[thm]{1,5}, \FormulaRefAuto{CompleteOrderedFieldAxioms}[def]{1,12}}}}
    \proofstepwidestarbreakkeep[1,2,9,12]{%
      \exists k_0\in M\,\forall k\in M\,(k_0\leq k)}{%
      \FormulaRefAuto{NaturalWellOrderingPrinciple}[thm]{13}}
    \proofstepwidestarbreakkeep[14]{%
      k_0\in M\land\forall k\in M\,(k_0\leq k)}{\rA}
    \proofstepwidestarbreakkeep[\shortstack{1,2,9\\12,15}]{k_0\neq0}{\text{\parbox[t]{\linewidth}{\raggedright Für \ensuremath{k_0=\allowbreak 0} wäre nach \ensuremath{12} und der Definition von \ensuremath{M} zugleich \ensuremath{\boxminus\eta_K(N)\preccurlyeq\allowbreak  x\boxtimes\allowbreak \eta_K(n)} sowie \ensuremath{x\boxtimes\allowbreak \eta_K(n)\preccurlyeq\allowbreak \boxminus\eta_K(N)} mit \ensuremath{x\boxtimes\allowbreak \eta_K(n)\neq\allowbreak \boxminus\eta_K(N),\allowbreak } also einen Widerspruch.}}}
    \proofstepwidestarbreakkeep[\shortstack{1,2,9\\12,15}]{%
      \exists \ell\in\mathbb N\,(k_0=(\ell+1))}{%
      \FormulaRefAuto{PeanoNonzeroIsAddOne}[thm]{\rAEa{15},16}}
    \proofstepwidestarbreakkeep[17]{%
      \ell\in\mathbb N\land k_0=(\ell+1)}{\rA}
    \proofstepwidestar[\shortstack{1,2,9\\12,15,18}]{\ell < ( \ell + 1 )}{%
      \FormulaRefAuto{PeanoLtAddOne}[thm]{\rAEa{18}}}
    \proofstepwidestarbreakkeep[\shortstack{1,2,9\\12,15,18}]{\ell<k_0}{%
      \rIE{\rAEb{18},
        19}}
    \proofstepwidestarbreakkeep[\shortstack{1,2,9\\12,15,18}]{\ell\notin M}{\text{\parbox[t]{\linewidth}{\raggedright \ensuremath{\ell\in M} lieferte aus der Minimalität von \ensuremath{k_0} die Ungleichung \ensuremath{k_0\leq\allowbreak \ell,\allowbreak } im Widerspruch zu \ensuremath{\ell<\allowbreak k_0.}}}}
    \proofstepwidestarbreakkeep[\shortstack{1--2,9,12\\15,18}]{%
      \eta_K(\ell)\boxplus\boxminus\eta_K(N)
      \preccurlyeq x\boxtimes\eta_K(n)}{\text{\parbox[t]{\linewidth}{\raggedright \ensuremath{\ell\in\mathbb N\land\allowbreak \ell\notin M} und die Definition von \ensuremath{M,\allowbreak } \FormulaRefAuto{Totale Ordnung}[def] , \FormulaRefAuto{CompleteOrderedFieldAxioms}[def]{1}}}}
    \proofstepwidestarbreakkeep[\shortstack{1--2,9,12\\15,18}]{%
      \begin{aligned}[t]
        &m\coloneqq\IntEmbed(k_0)-\IntEmbed(N)\in\mathbb Z\land{}\\[-2pt]
        &\eta_{\mathbb Z}(m-\IntOne)
        \preccurlyeq x\boxtimes\eta_K(n)
        \land x\boxtimes\eta_K(n)\preccurlyeq\eta_{\mathbb Z}(m)\land{}\\[-2pt]
        &x\boxtimes\eta_K(n)\neq\eta_{\mathbb Z}(m)
      \end{aligned}}{\text{\parbox[t]{\linewidth}{\raggedright \ensuremath{\rAEa{15},\allowbreak \rAEb{15},\allowbreak \rAEb{18},\allowbreak 20,\allowbreak 22,\allowbreak } \FormulaRefAuto{CompleteOrderedFieldIntegerPrimeEmbedding}[thm]{1,5}, \FormulaRefAuto{CompleteOrderedFieldNaturalEmbeddingArithmetic}[thm]{1,5}}}}
    \proofstepwidestarbreakkeep[\shortstack{1--4,9,12\\15,18}]{%
      \begin{aligned}[t]
      &x\preccurlyeq
        \eta_{\mathbb Q}(\RatClass{m}{\IntEmbed(n)})
        \land x\neq
        \eta_{\mathbb Q}(\RatClass{m}{\IntEmbed(n)})\land{}\\[-2pt]
      &\eta_{\mathbb Q}(\RatClass{m}{\IntEmbed(n)})
        \preccurlyeq x\boxplus\operatorname{inv}(\eta_K(n))\land{}\\[-2pt]
      &x\boxplus\operatorname{inv}(\eta_K(n))\preccurlyeq y
        \land x\boxplus\operatorname{inv}(\eta_K(n))\neq y
      \end{aligned}}{%
      \FormulaRefAuto{CompleteOrderedFieldAxioms}[def]{1},
      \FormulaRefAuto{CompleteOrderedFieldRationalPrimeEmbedding}[thm]{1},
      6,10,23}
    \proofstepwidestarbreakkeep[\shortstack{1--4,9,12\\15}]{%
      \begin{aligned}[t]
      &x\preccurlyeq
        \eta_{\mathbb Q}(\RatClass{m}{\IntEmbed(n)})
        \land x\neq
        \eta_{\mathbb Q}(\RatClass{m}{\IntEmbed(n)})\land{}\\[-2pt]
      &\eta_{\mathbb Q}(\RatClass{m}{\IntEmbed(n)})
        \preccurlyeq y\land
        \eta_{\mathbb Q}(\RatClass{m}{\IntEmbed(n)})\neq y
      \end{aligned}}{%
      \rEE{17,18,24}}
    \proofstepwidestarbreakkeep[1\text{--}4]{%
      \exists q\in\mathbb Q\,
      (x\preccurlyeq\eta_{\mathbb Q}(q)\land
       x\neq\eta_{\mathbb Q}(q)\land
       \eta_{\mathbb Q}(q)\preccurlyeq y\land
       \eta_{\mathbb Q}(q)\neq y)}{%
      \rEE{8,9,\rEE{11,12,\rEE{14,15,
        \rEI{\RatClass{m}{\IntEmbed(n)},25}}}}}
  \closeproofpartwideindR

  \proofpartwideindR[Definition durch Suprema]{%
    \begin{aligned}[t]
      &\CompleteOrderedField{K,\boxplus,\boxtimes,\boxminus,
        \operatorname{inv},0_K,1_K,\preccurlyeq}\dsep A\in\mathbb R\\[-2pt]
      &\quad\vdash
      \eta_{\mathbb Q}[A]\neq\varnothing\land
      \UB_{K,\preccurlyeq}(\eta_{\mathbb Q}[A])\neq\varnothing
    \end{aligned}
  }{%
    \CompleteOrderedField{K,\boxplus,\boxtimes,\boxminus,
      \operatorname{inv},0_K,1_K,\preccurlyeq}\dsep A\in\mathbb R
    \vdash
    \eta_{\mathbb Q}[A]\neq\varnothing\land
    \UB_{K,\preccurlyeq}(\eta_{\mathbb Q}[A])\neq\varnothing
  }[CompleteOrderedFieldCutImageBounded]
    \proofstepwidestar[1]{\CompleteOrderedField{K,\boxplus,\boxtimes,
      \boxminus,\operatorname{inv},0_K,1_K,\preccurlyeq}}{\rA}
    \proofstepwidestar[2]{A\in\mathbb R}{\rA}
    \proofstepwidestar[1]{\eta_{\mathbb Q}\colon\mathbb Q\inj K}{%
      \FormulaRefAuto{CompleteOrderedFieldRationalPrimeEmbedding}[thm]{1}}
    \proofstepwidestar[2]{\RealCut A}{%
      \FormulaRefAuto{RealCutMembership}[thm]{2}}
    \proofstepwidestar[2]{\exists a\in A\land\exists r\in\mathbb Q\setminus A}{%
      \FormulaRefAuto{DedekindCutCharacterization}[thm]{%
        4}}
    \proofstepwidestar[1,2]{\eta_{\mathbb Q}[A]\neq\varnothing}{%
      \FormulaRefAuto{NonemptyDirectImage}[thm]{3,5}}
    \proofstepwidestar[5]{a\in A\land r\in\mathbb Q\setminus A}{\rA}
    \proofstepwidestarbreakkeep[1,2,7]{%
      \forall q\in A\,(q<r)}{%
      \FormulaRefAuto{RealCutMemberBelowExterior}[thm]{2,7}}
    \proofstepwidestarbreakkeep[1,2,7]{%
      \forall q\in A\,
      \eta_{\mathbb Q}(q)\preccurlyeq\eta_{\mathbb Q}(r)}{%
      \FormulaRefAuto{CompleteOrderedFieldRationalPrimeEmbedding}[thm]{1,8}}
    \proofstepwidestar[1,2,7]{%
      \eta_{\mathbb Q}(r)\in
      \UB_{K,\preccurlyeq}(\eta_{\mathbb Q}[A])}{%
      \FormulaRefAuto{UpperBoundSetIntro}[thm]{3,9}}
    \proofstepwidestar[1,2,7]{%
      \UB_{K,\preccurlyeq}(\eta_{\mathbb Q}[A])\neq\varnothing}{%
      \FormulaRefAuto{x\in A\vdash A\neq\varnothing}{10}}
    \proofstepwidestar[1,2]{%
      \eta_{\mathbb Q}[A]\neq\varnothing\land
      \UB_{K,\preccurlyeq}(\eta_{\mathbb Q}[A])\neq\varnothing}{%
      \rEE{5,7,\rAI{6,11}}}
  \closeproofpartwideindR

  \proofpartwideindR[Strikte Approximation eines Supremums]{%
    \begin{aligned}[t]
      &\CompleteOrderedField{K,\boxplus,\boxtimes,\boxminus,
        \operatorname{inv},0_K,1_K,\preccurlyeq}\dsep
        \varnothing\neq S\subseteq K\dsep
        \UB_{K,\preccurlyeq}(S)\neq\varnothing\\[-2pt]
      &\quad\dsep \alpha=\sup_{K,\preccurlyeq}(S)\dsep u\in K\dsep
        u\preccurlyeq\alpha\dsep u\neq\alpha\\[-2pt]
      &\quad\vdash\exists s\in S\,
        (u\preccurlyeq s\land u\neq s)
    \end{aligned}
  }{%
    \CompleteOrderedField{K,\boxplus,\boxtimes,\boxminus,
      \operatorname{inv},0_K,1_K,\preccurlyeq}\dsep
    \varnothing\neq S\subseteq K\dsep
    \UB_{K,\preccurlyeq}(S)\neq\varnothing\dsep
    \alpha=\sup_{K,\preccurlyeq}(S)\dsep u\in K\dsep
    u\preccurlyeq\alpha\dsep u\neq\alpha
    \vdash\exists s\in S\,(u\preccurlyeq s\land u\neq s)
  }[CompleteOrderedFieldSupremumStrictWitness]
    \proofstepwidestarbreakkeep[1]{%
      \CompleteOrderedField{K,\boxplus,\boxtimes,\boxminus,
        \operatorname{inv},0_K,1_K,\preccurlyeq}}{\rA}
    \proofstepwidestarbreakkeep[2]{\varnothing\neq S\subseteq K}{\rA}
    \proofstepwidestarbreakkeep[3]{%
      \UB_{K,\preccurlyeq}(S)\neq\varnothing}{\rA}
    \proofstepwidestarbreakkeep[4]{%
      \alpha=\sup_{K,\preccurlyeq}(S)}{\rA}
    \proofstepwidestarbreakkeep[5]{u\in K}{\rA}
    \proofstepwidestarbreakkeep[6]{u\preccurlyeq\alpha}{\rA}
    \proofstepwidestarbreakkeep[7]{u\neq\alpha}{\rA}
    \proofstepwidestarbreakkeep[1,2,3,4]{%
      \alpha\in\UB_{K,\preccurlyeq}(S)\land
      \forall v\in\UB_{K,\preccurlyeq}(S)\,
        (\alpha\preccurlyeq v)}{%
      \FormulaRefAuto{SupremumLeastUpperBound}[thm]{1,2,3,4}}
    \proofstepwidestarbreakkeep[9]{%
      u\in\UB_{K,\preccurlyeq}(S)}{\rA}
    \proofstepwidestarbreakkeep[1,2,3,4,9]{\alpha\preccurlyeq u}{%
      \rUE{\rAEb{8},9}}
    \proofstepwidestarbreakkeep[1,5,6,9]{u=\alpha}{%
      \FormulaRefAuto{Totale Ordnung}[def]{1,5,6,10}}
    \proofstepwidestarbreakkeep[1,2,3,4,5,6,7,9]{\bot}{%
      \rBI{7,11}}
    \proofstepwidestarbreakkeep[1,2,3,4,5,6,7]{%
      u\notin\UB_{K,\preccurlyeq}(S)}{\rCI{9,12}}
    \proofstepwidestarbreakkeep[1,2,5,13]{%
      \exists s\in S\,\neg(s\preccurlyeq u)}{%
      \FormulaRefAuto{UpperBoundSetCharacterization}[thm]{1,2,5,13}}
    \proofstepwidestarbreakkeep[14]{%
      s\in S\land\neg(s\preccurlyeq u)}{\rA}
    \proofstepwidestarbreakkeep[1,2,5,15]{%
      u\preccurlyeq s\land u\neq s}{%
      \FormulaRefAuto{Totale Ordnung}[def]{1,2,5,15}}
    \proofstepwidestarbreakkeep[1,2,3,4,5,6,7]{%
      \exists s\in S\,(u\preccurlyeq s\land u\neq s)}{%
      \rEE{14,15,\rEI{s,16}}}
  \closeproofpartwideindR

  \proofpartwide{Ordnungserhaltung und Injektivität}
\proofstepwidestarbreakkeep[1]{%
      \CompleteOrderedField{K,\boxplus,\boxtimes,\boxminus,
        \operatorname{inv},0_K,1_K,\preccurlyeq}\dsep
      \Phi(A)\coloneqq
        \sup_{K,\preccurlyeq}(\eta_{\mathbb Q}[A])}{\rA}
    \proofstepwidestar[]{A\in\mathbb R\vdash\eta_{\mathbb Q}[A]\neq\varnothing\land\UB_{K,\preccurlyeq}(\eta_{\mathbb Q}[A])\neq\varnothing}{%
      \FormulaRefAuto{CompleteOrderedFieldCutImageBounded}[thm]{1}}
    \proofstepwidestarbreakkeep[]{\Phi\colon\mathbb R\to K}{%
      \FormulaRefAuto{CompleteOrderedFieldAxioms}[def]{%
        2}}
    \proofstepwidestarbreakkeep[4]{A,B\in\mathbb R\land A\subsetneq B}{\rA}
    \proofstepwidestarbreakkeep[4]{%
      \exists q\in B\setminus A\,\exists r\in B\,(q<r)}{%
      \FormulaRefAuto{A\subseteq B\dsep A\neq B\vdash
        \exists x\in B\setminus A}[thm]{4},
      \FormulaRefAuto{DedekindCutCharacterization}[thm]{4}}
    \proofstepwidestarbreakkeep[5]{q\in B\setminus A\land r\in B\land q<r}{\rA}
    \proofstepwidestarbreakkeep[4,6]{%
      \forall a\in A\,(a<q)}{%
      \FormulaRefAuto{RealCutMemberBelowExterior}[thm]{4,6}}
    \proofstepwidestarbreakkeep[1,4,6]{%
      \Phi(A)\preccurlyeq\eta_{\mathbb Q}(q)}{%
      \FormulaRefAuto{UpperBoundSetIntro}[thm]{7},
      \FormulaRefAuto{SupremumLeastUpperBound}[thm]{1}}
    \proofstepwidestarbreakkeep[1,4,6]{%
      \eta_{\mathbb Q}(q)\preccurlyeq\eta_{\mathbb Q}(r)
      \land\eta_{\mathbb Q}(q)\neq\eta_{\mathbb Q}(r)}{%
      \FormulaRefAuto{CompleteOrderedFieldRationalPrimeEmbedding}[thm]{1,6}}
    \proofstepwidestarbreakkeep[1,4,6]{%
      \eta_{\mathbb Q}(r)\preccurlyeq\Phi(B)}{%
      \FormulaRefAuto{SupremumLeastUpperBound}[thm]{1,6}}
    \proofstepwidestar[1,4]{\Phi(A)\preccurlyeq\Phi(B)\land\Phi(A)\neq\Phi(B)}{%
      \FormulaRefAuto{Totale Ordnung}[def]{8,9,10}}
    \proofstepwidestarbreakkeep[1,4]{%
      A\subsetneq B\vdash
      \Phi(A)\preccurlyeq\Phi(B)\land\Phi(A)\neq\Phi(B)}{%
      \rEE{5,6,11}}
    \proofstepwidestarbreakkeep[1]{%
      \Phi\text{ ist ordnungserhaltend und injektiv}}{%
      \FormulaRefAuto{RealCutTotalOrder}[thm]{3,12}}

\closeproofpartwide

\proofpartwide{Surjektivität}
  \setcounter{proofstepnr}{13}% Fortlaufende Schritte dieses Beweises.
\proofstepwidestarbreakkeep[14]{y\in K}{\rA}
    \proofstepwidestar[1,14]{\begin{aligned}[t]&\eta_K\colon\mathbb N\to K\land\eta_K(0)=0_K\land{}\\[-2pt]&\forall n\in\mathbb N\,(\eta_K(n+1)=\eta_K(n)\boxplus1_K)\end{aligned}}{%
      \FormulaRefAuto{CompleteOrderedFieldNaturalEmbedding}[thm]{1}}
    \proofstepwidestarbreakkeep[1,14]{%
      \exists p,r\in\mathbb Q\,\bigl(
      \eta_{\mathbb Q}(p)\preccurlyeq y\land
      \eta_{\mathbb Q}(p)\neq y\land
      y\preccurlyeq\eta_{\mathbb Q}(r)\land
       y\neq\eta_{\mathbb Q}(r)\bigr)}{\text{\parbox[t]{\linewidth}{\raggedright Das archimedische Lemma, auf \ensuremath{y} und \ensuremath{\boxminus y} angewandt, liefert \ensuremath{m,\allowbreak n\in\mathbb N} mit \ensuremath{\boxminus\eta_K(m)\preccurlyeq\allowbreak  y\land\allowbreak
          \boxminus\eta_K(m)\neq\allowbreak  y\land\allowbreak
          y\preccurlyeq\allowbreak \eta_K(n)\land\allowbreak  y\neq\allowbreak \eta_K(n).} Setze \ensuremath{p=\allowbreak -\RatEmbed(\IntEmbed(m))} und \ensuremath{r=\allowbreak \RatEmbed(\IntEmbed(n)).} \FormulaRefAuto{CompleteOrderedFieldNaturalEmbedding}[thm]{1}, \FormulaRefAuto{CompleteOrderedFieldArchimedean}[thm]{
            1,15,14}, \FormulaRefAuto{CompleteOrderedFieldRationalPrimeEmbedding}[thm]{1}, \FormulaRefAuto{CompleteOrderedFieldIntegerPrimeEmbedding}[thm]{1}}}}
    \proofstepwidestarbreakkeep[16]{%
      p,r\in\mathbb Q\land
      \eta_{\mathbb Q}(p)\preccurlyeq y\land
      \eta_{\mathbb Q}(p)\neq y\land
      y\preccurlyeq\eta_{\mathbb Q}(r)\land
      y\neq\eta_{\mathbb Q}(r)}{\rA}
    \proofstepwidestarbreakkeep[14]{%
      A_y\coloneqq\{q\in\mathbb Q\mid
      \eta_{\mathbb Q}(q)\preccurlyeq y\land
      \eta_{\mathbb Q}(q)\neq y\}}{\rII}
    \proofstepwidestarbreakkeep[14,17]{p\in A_y\land r\notin A_y}{%
      \FormulaRefAuto{CompleteOrderedFieldRationalPrimeEmbedding}[thm]{1,17,18}}
    \proofstepwidestarbreakkeep[14,18]{%
      q\in A_y\dsep s\in\mathbb Q\dsep s<q\vdash s\in A_y}{%
      \FormulaRefAuto{CompleteOrderedFieldRationalPrimeEmbedding}[thm]{1,18},
      \FormulaRefAuto{Totale Ordnung}[def]{14}}
    \proofstepwidestarbreakkeep[1,14,18]{%
      q\in A_y\vdash\exists s\in A_y\,(q<s)}{%
      \FormulaRefAuto{CompleteOrderedFieldRationalDense}[thm]{%
        1,14,18},
      \FormulaRefAuto{CompleteOrderedFieldRationalPrimeEmbedding}[thm]{1}}
    \proofstepwidestarbreakkeep[1,14,17]{\RealCut{A_y}}{%
      \FormulaRefAuto{DedekindCutCharacterization}[thm]{18,19,20,21}}
    \proofstepwidestarbreakkeep[14,18]{%
      y\in\UB_{K,\preccurlyeq}(\eta_{\mathbb Q}[A_y])}{%
      \FormulaRefAuto{UpperBoundSetIntro}[thm]{18}}
    \proofstepwidestarbreakkeep[1,14,18]{%
      u\in K\dsep u\preccurlyeq y\dsep u\neq y
      \vdash u\notin\UB_{K,\preccurlyeq}(\eta_{\mathbb Q}[A_y])}{%
      \FormulaRefAuto{CompleteOrderedFieldRationalDense}[thm]{1,14,18}}
    \proofstepwidestarbreakkeep[1,14,17]{%
      \Phi(A_y)=y}{%
      \FormulaRefAuto{SupremumCandidateEqualsSupremum}[thm]{1,22,23,24}}
    \proofstepwidestarbreakkeep[1]{\Phi\colon\mathbb R\bij K}{%
      \FormulaRefAuto{Bijektive Funktion}[def]{3,13,14,25}}

\closeproofpartwide

\proofpartwide{Ordnungsreflexion}
  \setcounter{proofstepnr}{26}% Fortlaufende Schritte dieses Beweises.
\proofstepwidestarbreakkeep[1]{%
      \forall A,B\in\mathbb R\,
      (A\RealLe B\leftrightarrow\Phi(A)\preccurlyeq\Phi(B))}{%
      \FormulaRefAuto{RealCutOrderCharacterization}[thm]{12,13}}

\closeproofpartwide

\proofpartwideindR[Ordnung und Bijektivität der Supremumsabbildung]{%
    \begin{aligned}[t]
      &\CompleteOrderedField{K,\boxplus,\boxtimes,\boxminus,
        \operatorname{inv},0_K,1_K,\preccurlyeq}\dsep{}\\[-2pt]
      &\Phi(A)\coloneqq
        \sup_{K,\preccurlyeq}(\eta_{\mathbb Q}[A])\\[-2pt]
      &\quad\vdash\Phi\colon\mathbb R\bij K\land
      \forall A,B\in\mathbb R\,
      (A\RealLe B\leftrightarrow\Phi(A)\preccurlyeq\Phi(B))
    \end{aligned}
  }{%
    \CompleteOrderedField{K,\boxplus,\boxtimes,\boxminus,
      \operatorname{inv},0_K,1_K,\preccurlyeq}\dsep
    \Phi(A)\coloneqq\sup_{K,\preccurlyeq}(\eta_{\mathbb Q}[A])
    \vdash\Phi\colon\mathbb R\bij K\land
    \forall A,B\in\mathbb R\,
    (A\RealLe B\leftrightarrow\Phi(A)\preccurlyeq\Phi(B))
  }[CompleteOrderedFieldCutSupremumBijection]
  \setcounter{proofstepnr}{27}% Fortlaufende Schritte dieses Beweises.
\proofstepwidestarbreakkeep[1]{\Phi\colon\mathbb R\bij K\land
      \forall A,B\in\mathbb R\,
      (A\RealLe B\leftrightarrow\Phi(A)\preccurlyeq\Phi(B))}{\rAI{26,27}}

\closeproofpartwideindR


  \proofpartwideindR[Schnittrekonstruktion aus dem Supremum]{%
    \begin{aligned}[t]
      &\CompleteOrderedField{K,\boxplus,\boxtimes,\boxminus,
        \operatorname{inv},0_K,1_K,\preccurlyeq}\dsep{}\\[-2pt]
      &\Phi(A)\coloneqq
        \sup_{K,\preccurlyeq}(\eta_{\mathbb Q}[A])\\[-2pt]
      &\quad\vdash\forall A\in\mathbb R\,\forall q\in\mathbb Q\,
        \Bigl(q\in A\leftrightarrow{}\\[-2pt]
      &\qquad
        (\eta_{\mathbb Q}(q)\preccurlyeq\Phi(A)\land
         \eta_{\mathbb Q}(q)\neq\Phi(A))\Bigr)
    \end{aligned}
  }{%
    \CompleteOrderedField{K,\boxplus,\boxtimes,\boxminus,
      \operatorname{inv},0_K,1_K,\preccurlyeq}\dsep
    \Phi(A)\coloneqq
      \sup_{K,\preccurlyeq}(\eta_{\mathbb Q}[A])
    \vdash\forall A\in\mathbb R\,\forall q\in\mathbb Q\,
    (q\in A\leftrightarrow
     (\eta_{\mathbb Q}(q)\preccurlyeq\Phi(A)\land
      \eta_{\mathbb Q}(q)\neq\Phi(A)))
  }[CompleteOrderedFieldCutRecovery]
    \proofstepwidestarbreakkeep[1]{%
      \CompleteOrderedField{K,\boxplus,\boxtimes,\boxminus,
        \operatorname{inv},0_K,1_K,\preccurlyeq}\dsep
      \Phi(A)\coloneqq
        \sup_{K,\preccurlyeq}(\eta_{\mathbb Q}[A])}{\rA}
    \proofstepwidestarbreakkeep[2]{A\in\mathbb R\land q\in\mathbb Q}{\rA}
    \proofstepwidestarbreakkeep[3]{q\in A}{\rA}
    \proofstepwidestar[2,3]{\RealCut A}{%
      \FormulaRefAuto{RealCutMembership}[thm]{\rAEa{2}}}
    \proofstepwidestarbreakkeep[2,3]{%
      \exists r\in A\,(q<r)}{%
      \FormulaRefAuto{DedekindCutCharacterization}[thm]{%
        4,3}}
    \proofstepwidestarbreakkeep[5]{r\in A\land q<r}{\rA}
    \proofstepwidestarbreakkeep[2,6]{%
      \eta_{\mathbb Q}(q)\preccurlyeq\eta_{\mathbb Q}(r)\land
      \eta_{\mathbb Q}(q)\neq\eta_{\mathbb Q}(r)}{%
      \FormulaRefAuto{CompleteOrderedFieldRationalPrimeEmbedding}[thm]{%
        1,\rAEb{2},\rAEb{6}}}
    \proofstepwidestar[1,2,6]{\eta_{\mathbb Q}[A]\neq\varnothing\land\UB_{K,\preccurlyeq}(\eta_{\mathbb Q}[A])\neq\varnothing}{%
      \FormulaRefAuto{CompleteOrderedFieldCutImageBounded}[thm]{\rAEa{2}}}
    \proofstepwidestarbreakkeep[1,2,6]{%
      \eta_{\mathbb Q}(r)\preccurlyeq\Phi(A)}{%
      \FormulaRefAuto{SupremumLeastUpperBound}[thm]{%
        1,8,
        \rAEa{6}}}
    \proofstepwidestarbreakkeep[1,2,6]{%
      \eta_{\mathbb Q}(q)\preccurlyeq\Phi(A)\land
      \eta_{\mathbb Q}(q)\neq\Phi(A)}{%
      \FormulaRefAuto{Totale Ordnung}[def]{7,9}}
    \proofstepwidestarbreakkeep[1,2]{%
      q\in A\rightarrow
      (\eta_{\mathbb Q}(q)\preccurlyeq\Phi(A)\land
       \eta_{\mathbb Q}(q)\neq\Phi(A))}{%
      \rEE{5,6,\rRI{3,10}}}
    \proofstepwidestarbreakkeep[12]{%
      \eta_{\mathbb Q}(q)\preccurlyeq\Phi(A)\land
      \eta_{\mathbb Q}(q)\neq\Phi(A)}{\rA}
    \proofstepwidestarbreakkeep[13]{q\notin A}{\rA}
    \proofstepwidestarbreakkeep[2,13]{%
      \forall a\in A\,(a<q)}{%
      \FormulaRefAuto{RealCutMemberBelowExterior}[thm]{\rAEa{2},13}}
    \proofstepwidestar[1,2,13]{\eta_{\mathbb Q}\colon\mathbb Q\inj K\land\RationalPrimeEmbeddingLaws{\eta_{\mathbb Q}}}{%
      \FormulaRefAuto{CompleteOrderedFieldRationalPrimeEmbedding}[thm]{1}}
    \proofstepwidestarbreakkeep[1,2,13]{%
      \eta_{\mathbb Q}(q)\in
      \UB_{K,\preccurlyeq}(\eta_{\mathbb Q}[A])}{%
      \FormulaRefAuto{UpperBoundSetIntro}[thm]{%
        15,
        14}}
    \proofstepwidestarbreakkeep[1,2,13]{%
      \Phi(A)\preccurlyeq\eta_{\mathbb Q}(q)}{%
      \FormulaRefAuto{SupremumLeastUpperBound}[thm]{1,16}}
    \proofstepwidestarbreakkeep[1,2,13]{%
      \Phi(A)=\eta_{\mathbb Q}(q)}{%
      \FormulaRefAuto{Totale Ordnung}[def]{12,17}}
    \proofstepwidestar[1,2,13]{{ \eta _ { \mathbb Q } ( q ) } = { \Phi ( A ) }}{%
      \FormulaRefAuto{a=b\vdash b=a}[thm]{18}}
    \proofstepwidestarbreakkeep[1,2,13]{\bot}{%
      \rBI{\rAEb{12},19}}
    \proofstepwidestarbreakkeep[1,2]{q\in A}{\rCI{13,20}}
    \proofstepwidestarbreakkeep[1,2]{%
      (\eta_{\mathbb Q}(q)\preccurlyeq\Phi(A)\land
       \eta_{\mathbb Q}(q)\neq\Phi(A))\rightarrow q\in A}{%
      \rRI{12,21}}
    \proofstepwidestarbreakkeep[1,2]{%
      q\in A\leftrightarrow
      (\eta_{\mathbb Q}(q)\preccurlyeq\Phi(A)\land
       \eta_{\mathbb Q}(q)\neq\Phi(A))}{\rBI{11,22}}
    \proofstepwidestarbreakkeep[1]{%
      \forall A\in\mathbb R\,\forall q\in\mathbb Q\,
      \bigl(q\in A\leftrightarrow
      (\eta_{\mathbb Q}(q)\preccurlyeq\Phi(A)\land
       \eta_{\mathbb Q}(q)\neq\Phi(A))\bigr)}{%
      \rUI{\rRI{2,23}}}
  \closeproofpartwideindR

  \proofpartwideindR[Hauptschnitte werden auf rationale Elemente abgebildet]{%
    \begin{aligned}[t]
      &\CompleteOrderedField{K,\boxplus,\boxtimes,\boxminus,
        \operatorname{inv},0_K,1_K,\preccurlyeq}\dsep{}\\[-2pt]
      &\Phi(A)\coloneqq
        \sup_{K,\preccurlyeq}(\eta_{\mathbb Q}[A])\\[-2pt]
      &\quad\vdash\forall q\in\mathbb Q\,
        \Phi(\RealEmb(q))=\eta_{\mathbb Q}(q)
    \end{aligned}
  }{%
    \CompleteOrderedField{K,\boxplus,\boxtimes,\boxminus,
      \operatorname{inv},0_K,1_K,\preccurlyeq}\dsep
    \Phi(A)\coloneqq
      \sup_{K,\preccurlyeq}(\eta_{\mathbb Q}[A])
    \vdash\forall q\in\mathbb Q\,
      \Phi(\RealEmb(q))=\eta_{\mathbb Q}(q)
  }[CompleteOrderedFieldPrincipalCutImage]
    \proofstepwidestarbreakkeep[1]{%
      \CompleteOrderedField{K,\boxplus,\boxtimes,\boxminus,
        \operatorname{inv},0_K,1_K,\preccurlyeq}\dsep
      \Phi(A)\coloneqq
        \sup_{K,\preccurlyeq}(\eta_{\mathbb Q}[A])}{\rA}
    \proofstepwidestarbreakkeep[2]{q\in\mathbb Q}{\rA}
    \proofstepwidestarbreakkeep[1,2]{\RealEmb(q)=\RealHat q\in\mathbb R}{%
      \FormulaRefAuto{RationalToRealEmbedding}[def]{2},
      \FormulaRefAuto{RealPrincipalCutIsCut}[thm]{2}}
    \proofstepwidestarbreakkeep[4]{%
      \Phi(\RealEmb(q))\neq\eta_{\mathbb Q}(q)}{\rA}
    \proofstepwidestarbreakkeep[1,2,4]{%
      \begin{aligned}[t]
      &(\Phi(\RealEmb(q))\preccurlyeq\eta_{\mathbb Q}(q)\land
        \Phi(\RealEmb(q))\neq\eta_{\mathbb Q}(q))\lor{}\\[-2pt]
      &(\eta_{\mathbb Q}(q)\preccurlyeq\Phi(\RealEmb(q))\land
        \eta_{\mathbb Q}(q)\neq\Phi(\RealEmb(q)))
      \end{aligned}}{%
      \FormulaRefAuto{Totale Ordnung}[def]{4}}
    \proofstepwidestarbreakkeep[6]{%
      \Phi(\RealEmb(q))\preccurlyeq\eta_{\mathbb Q}(q)\land
      \Phi(\RealEmb(q))\neq\eta_{\mathbb Q}(q)}{\rA}
    \proofstepwidestarbreakkeep[1,2,6]{%
      \exists r\in\mathbb Q\,\bigl(
      \Phi(\RealEmb(q))\preccurlyeq\eta_{\mathbb Q}(r)\land
      \Phi(\RealEmb(q))\neq\eta_{\mathbb Q}(r)\land
      \eta_{\mathbb Q}(r)\preccurlyeq\eta_{\mathbb Q}(q)\land
      \eta_{\mathbb Q}(r)\neq\eta_{\mathbb Q}(q)\bigr)}{%
      \FormulaRefAuto{CompleteOrderedFieldRationalDense}[thm]{1,6}}
    \proofstepwidestarbreakkeep[7]{%
      \begin{aligned}[t]
      &r\in\mathbb Q\land
       \Phi(\RealEmb(q))\preccurlyeq\eta_{\mathbb Q}(r)\land
       \Phi(\RealEmb(q))\neq\eta_{\mathbb Q}(r)\land{}\\[-2pt]
      &\eta_{\mathbb Q}(r)\preccurlyeq\eta_{\mathbb Q}(q)\land
       \eta_{\mathbb Q}(r)\neq\eta_{\mathbb Q}(q)
      \end{aligned}}{\rA}
    \proofstepwidestarbreakkeep[1,2,8]{r<q\land r\in\RealEmb(q)}{%
      \FormulaRefAuto{CompleteOrderedFieldRationalPrimeEmbedding}[thm]{1,8},
      \FormulaRefAuto{RationalToRealEmbedding}[def]{2}}
    \proofstepwidestarbreakkeep[1,2,8]{%
      \eta_{\mathbb Q}(r)\preccurlyeq\Phi(\RealEmb(q))\land
      \eta_{\mathbb Q}(r)\neq\Phi(\RealEmb(q))}{%
      \FormulaRefAuto{CompleteOrderedFieldCutRecovery}[thm]{1,3,9}}
    \proofstepwidestarbreakkeep[1,2,8]{\bot}{%
      \FormulaRefAuto{Totale Ordnung}[def]{8,10}}
    \proofstepwidestarbreakkeep[12]{%
      \eta_{\mathbb Q}(q)\preccurlyeq\Phi(\RealEmb(q))\land
      \eta_{\mathbb Q}(q)\neq\Phi(\RealEmb(q))}{\rA}
    \proofstepwidestarbreakkeep[1,2,12]{%
      \exists r\in\mathbb Q\,\bigl(
      \eta_{\mathbb Q}(q)\preccurlyeq\eta_{\mathbb Q}(r)\land
      \eta_{\mathbb Q}(q)\neq\eta_{\mathbb Q}(r)\land
      \eta_{\mathbb Q}(r)\preccurlyeq\Phi(\RealEmb(q))\land
      \eta_{\mathbb Q}(r)\neq\Phi(\RealEmb(q))\bigr)}{%
      \FormulaRefAuto{CompleteOrderedFieldRationalDense}[thm]{1,12}}
    \proofstepwidestarbreakkeep[13]{%
      \begin{aligned}[t]
      &r\in\mathbb Q\land
       \eta_{\mathbb Q}(q)\preccurlyeq\eta_{\mathbb Q}(r)\land
       \eta_{\mathbb Q}(q)\neq\eta_{\mathbb Q}(r)\land{}\\[-2pt]
      &\eta_{\mathbb Q}(r)\preccurlyeq\Phi(\RealEmb(q))\land
       \eta_{\mathbb Q}(r)\neq\Phi(\RealEmb(q))
      \end{aligned}}{\rA}
    \proofstepwidestarbreakkeep[1,2,14]{q<r\land r<q}{%
      \FormulaRefAuto{CompleteOrderedFieldRationalPrimeEmbedding}[thm]{1,14},
      \FormulaRefAuto{CompleteOrderedFieldCutRecovery}[thm]{1,3,14},
      \FormulaRefAuto{RationalToRealEmbedding}[def]{2}}
    \proofstepwidestarbreakkeep[1,2,14]{\bot}{%
      \FormulaRefAuto{RationalOrderTotal}[thm]{15}}
    \proofstepwidestarbreakkeep[1,2,4]{\bot}{%
      \rOE{5,6,\rEE{7,8,11},12,\rEE{13,14,16}}}
    \proofstepwidestarbreakkeep[1,2]{%
      \Phi(\RealEmb(q))=\eta_{\mathbb Q}(q)}{\rCI{4,17}}
    \proofstepwidestarbreakkeep[1]{%
      \forall q\in\mathbb Q\,
      \Phi(\RealEmb(q))=\eta_{\mathbb Q}(q)}{%
      \rUI{\rRI{2,18}}}
  \closeproofpartwideindR

  \proofpartwideindR[Erhaltung der Addition]{%
    \begin{aligned}[t]
      &\CompleteOrderedField{K,\boxplus,\boxtimes,\boxminus,
        \operatorname{inv},0_K,1_K,\preccurlyeq}\dsep
        \Phi(A)\coloneqq
        \sup_{K,\preccurlyeq}(\eta_{\mathbb Q}[A])\\[-2pt]
      &\quad\vdash\forall A,B\in\mathbb R\,
        \Phi(A\RealAdd B)=\Phi(A)\boxplus\Phi(B)
    \end{aligned}
  }{%
    \CompleteOrderedField{K,\boxplus,\boxtimes,\boxminus,
      \operatorname{inv},0_K,1_K,\preccurlyeq}\dsep
    \Phi(A)\coloneqq
      \sup_{K,\preccurlyeq}(\eta_{\mathbb Q}[A])
    \vdash\forall A,B\in\mathbb R\,
      \Phi(A\RealAdd B)=\Phi(A)\boxplus\Phi(B)
  }[CompleteOrderedFieldCutSupremumAddition]
    \proofstepwidestarbreakkeep[1]{%
      \CompleteOrderedField{K,\boxplus,\boxtimes,\boxminus,
        \operatorname{inv},0_K,1_K,\preccurlyeq}}{\rA}
    \proofstepwidestarbreakkeep[2]{%
      \Phi(A)\coloneqq
        \sup_{K,\preccurlyeq}(\eta_{\mathbb Q}[A])}{\rA}
    \proofstepwidestarbreakkeep[1,2]{%
      \Phi\colon\mathbb R\bij K\land
      \forall X,Y\in\mathbb R\,
      (X\RealLe Y\leftrightarrow\Phi(X)\preccurlyeq\Phi(Y))}{%
      \FormulaRefAuto{CompleteOrderedFieldCutSupremumBijection}[thm]{1,2}}
    \proofstepwidestarbreakkeep[1,2]{%
      \forall C\in\mathbb R\,\forall q\in\mathbb Q\,
      \bigl(q\in C\leftrightarrow
      (\eta_{\mathbb Q}(q)\preccurlyeq\Phi(C)\land
       \eta_{\mathbb Q}(q)\neq\Phi(C))\bigr)}{%
      \FormulaRefAuto{CompleteOrderedFieldCutRecovery}[thm]{1,2}}
    \proofstepwidestarbreakkeep[5]{A,B\in\mathbb R}{\rA}
    \proofstepwidestarbreakkeep[6]{q\in\mathbb Q}{\rA}
    \proofstepwidestarbreakkeep[7]{q\in A\RealAdd B}{\rA}
    \proofstepwidestarbreakkeep[5,6,7]{%
      \exists a\in A\,\exists b\in B\,(q<a+b)}{%
      \FormulaRefAuto{RealCutAddition}[def]{5,6,7}}
    \proofstepwidestarbreakkeep[8]{%
      a\in A\land b\in B\land q<a+b}{\rA}
    \proofstepwidestarbreakkeep[1,5,6,9]{%
      \eta_{\mathbb Q}(q)\preccurlyeq
        \eta_{\mathbb Q}(a)\boxplus\eta_{\mathbb Q}(b)
      \land
      \eta_{\mathbb Q}(q)\neq
        \eta_{\mathbb Q}(a)\boxplus\eta_{\mathbb Q}(b)}{%
      \FormulaRefAuto{CompleteOrderedFieldRationalPrimeEmbedding}[thm]{%
        1,6,9}}
    \proofstepwidestar[1,2,5,9]{\begin{aligned}[t]&\eta_{\mathbb Q}[A]\neq\varnothing\land\UB_{K,\preccurlyeq}(\eta_{\mathbb Q}[A])\neq\varnothing\land{}\\[-2pt]&\eta_{\mathbb Q}[B]\neq\varnothing\land\UB_{K,\preccurlyeq}(\eta_{\mathbb Q}[B])\neq\varnothing\end{aligned}}{%
      \FormulaRefAuto{CompleteOrderedFieldCutImageBounded}[thm]{1,5}}
    \proofstepwidestarbreakkeep[1,2,5,9]{%
      \eta_{\mathbb Q}(a)\preccurlyeq\Phi(A)\land
      \eta_{\mathbb Q}(b)\preccurlyeq\Phi(B)}{%
      \FormulaRefAuto{SupremumLeastUpperBound}[thm]{%
        2,11,9}}
    \proofstepwidestarbreakkeep[1,2,5,6,9]{%
      \eta_{\mathbb Q}(q)\preccurlyeq
        \Phi(A)\boxplus\Phi(B)\land
      \eta_{\mathbb Q}(q)\neq
        \Phi(A)\boxplus\Phi(B)}{%
      \FormulaRefAuto{CompleteOrderedFieldAxioms}[def]{1,10,12}}
    \proofstepwidestarbreakkeep[1,2,5,6]{%
      q\in A\RealAdd B\rightarrow
      (\eta_{\mathbb Q}(q)\preccurlyeq\Phi(A)\boxplus\Phi(B)\land
       \eta_{\mathbb Q}(q)\neq\Phi(A)\boxplus\Phi(B))}{%
      \rEE{8,9,\rRI{7,13}}}
    \proofstepwidestarbreakkeep[15]{%
      \eta_{\mathbb Q}(q)\preccurlyeq\Phi(A)\boxplus\Phi(B)\land
      \eta_{\mathbb Q}(q)\neq\Phi(A)\boxplus\Phi(B)}{\rA}
    \proofstepwidestarbreakkeep[1,2,5,6,15]{%
      \begin{aligned}[t]
      &u\coloneqq\eta_{\mathbb Q}(q)\boxplus\boxminus\Phi(B)
        \in K,\\[-2pt]
      &u\preccurlyeq\Phi(A)\land u\neq\Phi(A)
      \end{aligned}}{%
      \FormulaRefAuto{CompleteOrderedFieldStrictTranslation}[thm]{1,15},
      \FormulaRefAuto{CompleteOrderedFieldAxioms}[def]{1}}
    \proofstepwidestar[1,2,5,6,15]{\eta_{\mathbb Q}[A]\neq\varnothing\land\UB_{K,\preccurlyeq}(\eta_{\mathbb Q}[A])\neq\varnothing}{%
      \FormulaRefAuto{CompleteOrderedFieldCutImageBounded}[thm]{1,5}}
    \proofstepwidestarbreakkeep[1,2,5,6,15]{%
      \exists a\in A\,
      (u\preccurlyeq\eta_{\mathbb Q}(a)\land
       u\neq\eta_{\mathbb Q}(a))}{%
      \FormulaRefAuto{CompleteOrderedFieldSupremumStrictWitness}[thm]{%
        1,2,17,16}}
    \proofstepwidestarbreakkeep[18]{%
      a\in A\land u\preccurlyeq\eta_{\mathbb Q}(a)
      \land u\neq\eta_{\mathbb Q}(a)}{\rA}
    \proofstepwidestarbreakkeep[1,2,5,6,15,19]{%
      \eta_{\mathbb Q}(q)\preccurlyeq
        \eta_{\mathbb Q}(a)\boxplus\Phi(B)
      \land
      \eta_{\mathbb Q}(q)\neq
        \eta_{\mathbb Q}(a)\boxplus\Phi(B)}{%
      \FormulaRefAuto{CompleteOrderedFieldStrictTranslation}[thm]{1,16,19}}
    \proofstepwidestarbreakkeep[1,2,5,6,15,19]{%
      \begin{aligned}[t]
      &v\coloneqq\eta_{\mathbb Q}(q)
        \boxplus\boxminus\eta_{\mathbb Q}(a)\in K,\\[-2pt]
      &v\preccurlyeq\Phi(B)\land v\neq\Phi(B)
      \end{aligned}}{%
      \FormulaRefAuto{CompleteOrderedFieldStrictTranslation}[thm]{1,20},
      \FormulaRefAuto{CompleteOrderedFieldAxioms}[def]{1}}
    \proofstepwidestar[1,2,5,6,15,19]{\eta_{\mathbb Q}[B]\neq\varnothing\land\UB_{K,\preccurlyeq}(\eta_{\mathbb Q}[B])\neq\varnothing}{%
      \FormulaRefAuto{CompleteOrderedFieldCutImageBounded}[thm]{1,5}}
    \proofstepwidestarbreakkeep[1,2,5,6,15,19]{%
      \exists b\in B\,
      (v\preccurlyeq\eta_{\mathbb Q}(b)\land
       v\neq\eta_{\mathbb Q}(b))}{%
      \FormulaRefAuto{CompleteOrderedFieldSupremumStrictWitness}[thm]{%
        1,2,22,21}}
    \proofstepwidestarbreakkeep[23]{%
      b\in B\land v\preccurlyeq\eta_{\mathbb Q}(b)
      \land v\neq\eta_{\mathbb Q}(b)}{\rA}
    \proofstepwidestarbreakkeep[1,2,5,6,15,19,24]{%
      q<a+b\land q\in A\RealAdd B}{%
      \FormulaRefAuto{CompleteOrderedFieldStrictTranslation}[thm]{1,21,24},
      \FormulaRefAuto{CompleteOrderedFieldRationalPrimeEmbedding}[thm]{1},
      \FormulaRefAuto{RealCutAddition}[def]{19,24}}
    \proofstepwidestarbreakkeep[1,2,5,6]{%
      (\eta_{\mathbb Q}(q)\preccurlyeq\Phi(A)\boxplus\Phi(B)\land
       \eta_{\mathbb Q}(q)\neq\Phi(A)\boxplus\Phi(B))
      \rightarrow q\in A\RealAdd B}{%
      \rEE{18,19,\rEE{23,24,\rRI{15,25}}}}
    \proofstepwidestarbreakkeep[1,2,5,6]{%
      q\in A\RealAdd B\leftrightarrow
      (\eta_{\mathbb Q}(q)\preccurlyeq\Phi(A)\boxplus\Phi(B)\land
       \eta_{\mathbb Q}(q)\neq\Phi(A)\boxplus\Phi(B))}{%
      \rBI{14,26}}
    \proofstepwidestarbreakkeep[1,2,5]{%
      \forall q\in\mathbb Q\,\bigl(q\in A\RealAdd B\leftrightarrow
      (\eta_{\mathbb Q}(q)\preccurlyeq\Phi(A)\boxplus\Phi(B)\land
       \eta_{\mathbb Q}(q)\neq\Phi(A)\boxplus\Phi(B))\bigr)}{%
      \rUI{\rRI{6,27}}}
    \proofstepwidestarbreakkeep[1,2,3,5]{%
      \exists C\in\mathbb R\,
      \Phi(C)=\Phi(A)\boxplus\Phi(B)}{%
      \FormulaRefAuto{Surjektive Funktion}[def]{3},
      \FormulaRefAuto{CompleteOrderedFieldAxioms}[def]{1}}
    \proofstepwidestarbreakkeep[29]{%
      C\in\mathbb R\land\Phi(C)=\Phi(A)\boxplus\Phi(B)}{\rA}
    \proofstepwidestarbreakkeep[1,2,4,30]{%
      \forall q\in\mathbb Q\,\bigl(q\in C\leftrightarrow
      (\eta_{\mathbb Q}(q)\preccurlyeq\Phi(A)\boxplus\Phi(B)\land
       \eta_{\mathbb Q}(q)\neq\Phi(A)\boxplus\Phi(B))\bigr)}{%
      \rIE{\rUE{4,\rAEa{30}},\rAEb{30}}}
    \proofstepwidestarbreakkeep[1,2,5,30]{C=A\RealAdd B}{%
      \FormulaRefAuto{\forall x\,(x\in A\leftrightarrow x\in B)
        \vdash A=B}[ax]{28,31}}
    \proofstepwidestarbreakkeep[1,2,5,30]{%
      \Phi(A\RealAdd B)=\Phi(A)\boxplus\Phi(B)}{%
      \rIE{32,\rAEb{30}}}
    \proofstepwidestarbreakkeep[1,2,5]{%
      \Phi(A\RealAdd B)=\Phi(A)\boxplus\Phi(B)}{%
      \rEE{29,30,33}}
    \proofstepwidestarbreakkeep[1,2]{%
      \forall A,B\in\mathbb R\,
      \Phi(A\RealAdd B)=\Phi(A)\boxplus\Phi(B)}{%
      \rUI{\rRI{5,34}}}
  \closeproofpartwideindR

  \proofpartwideindR[Erhaltung des positiven Produkts]{%
    \begin{aligned}[t]
      &\CompleteOrderedField{K,\boxplus,\boxtimes,\boxminus,
        \operatorname{inv},0_K,1_K,\preccurlyeq}\dsep
        \Phi(A)\coloneqq
        \sup_{K,\preccurlyeq}(\eta_{\mathbb Q}[A])\\[-2pt]
      &\quad\vdash\forall A,B\in\mathbb R\,
        \bigl((\RealZero\RealLe A\land\RealZero\RealLe B)\rightarrow{}\\[-2pt]
      &\qquad\Phi(A\RealPosMul B)=\Phi(A)\boxtimes\Phi(B)\bigr)
    \end{aligned}
  }{%
    \CompleteOrderedField{K,\boxplus,\boxtimes,\boxminus,
      \operatorname{inv},0_K,1_K,\preccurlyeq}\dsep
    \Phi(A)\coloneqq
      \sup_{K,\preccurlyeq}(\eta_{\mathbb Q}[A])
    \vdash\forall A,B\in\mathbb R\,
    ((\RealZero\RealLe A\land\RealZero\RealLe B)\rightarrow
     \Phi(A\RealPosMul B)=\Phi(A)\boxtimes\Phi(B))
  }[CompleteOrderedFieldCutSupremumPositiveProduct]
    \proofstepwidestarbreakkeep[1]{%
      \CompleteOrderedField{K,\boxplus,\boxtimes,\boxminus,
        \operatorname{inv},0_K,1_K,\preccurlyeq}}{\rA}
    \proofstepwidestarbreakkeep[2]{%
      \Phi(A)\coloneqq
        \sup_{K,\preccurlyeq}(\eta_{\mathbb Q}[A])}{\rA}
    \proofstepwidestarbreakkeep[1,2]{%
      \Phi\colon\mathbb R\bij K\quad\text{ordnungserhaltend}}{%
      \FormulaRefAuto{CompleteOrderedFieldCutSupremumBijection}[thm]{1,2}}
    \proofstepwidestarbreakkeep[1,2]{%
      \Phi(\RealZero)=0_K\quad\text{und}\quad
      \forall C\in\mathbb R\,\forall q\in\mathbb Q\,
      \bigl(q\in C\leftrightarrow
      (\eta_{\mathbb Q}(q)\preccurlyeq\Phi(C)\land
       \eta_{\mathbb Q}(q)\neq\Phi(C))\bigr)}{%
      \FormulaRefAuto{CompleteOrderedFieldPrincipalCutImage}[thm]{1,2},
      \FormulaRefAuto{RealCutZero}[def],
      \FormulaRefAuto{CompleteOrderedFieldCutRecovery}[thm]{1,2}}
    \proofstepwidestarbreakkeep[5]{A,B\in\mathbb R}{\rA}
    \proofstepwidestarbreakkeep[6]{%
      \RealZero\RealLe A\land\RealZero\RealLe B}{\rA}
    \proofstepwidestarbreakkeep[1\text{--}6]{%
      \begin{aligned}[t]
      &\alpha\coloneqq\Phi(A)\dsep\beta\coloneqq\Phi(B),\\[-2pt]
      &0_K\preccurlyeq\alpha\land0_K\preccurlyeq\beta
      \end{aligned}}{\text{\parbox[t]{\linewidth}{\raggedright Aus der Ordnungsäquivalenz in \ensuremath{3} und \ensuremath{6} folgen \ensuremath{\Phi(\RealZero)\preccurlyeq\allowbreak \Phi(A)} und \ensuremath{\Phi(\RealZero)\preccurlyeq\allowbreak \Phi(B);} mit \ensuremath{4} und den Definitionen von \ensuremath{\alpha,\allowbreak \beta} werden daraus \ensuremath{0_K\preccurlyeq\allowbreak \alpha} und \ensuremath{0_K\preccurlyeq\allowbreak \beta.}}}}
    \proofstepwidestarbreakkeep[8]{q\in\mathbb Q}{\rA}
    \proofstepwidestarbreakkeep[9]{q\in A\RealPosMul B}{\rA}
    \proofstepwidestarbreakkeep[5,6,8,9]{%
      q<\RatZero\lor
      \exists a\in A\,\exists b\in B\,
      (\RatZero<a\land\RatZero<b\land q<ab)}{%
      \FormulaRefAuto{RealCutPositiveProduct}[def]{5,6,8,9},
      \FormulaRefAuto{RealCutZero}[def],
      \FormulaRefAuto{RealPrincipalCut}[def]}
    \proofstepwidestarbreakkeep[11]{q<\RatZero}{\rA}
    \proofstepwidestarbreakkeep[1,7,8,11]{%
      \eta_{\mathbb Q}(q)\preccurlyeq\alpha\boxtimes\beta\land
      \eta_{\mathbb Q}(q)\neq\alpha\boxtimes\beta}{%
      \FormulaRefAuto{CompleteOrderedFieldRationalPrimeEmbedding}[thm]{1,8,11},
      \FormulaRefAuto{CompleteOrderedFieldAxioms}[def]{1,7}}
    \proofstepwidestarbreakkeep[11]{%
      \exists a\in A\,\exists b\in B\,
      (\RatZero<a\land\RatZero<b\land q<ab)}{\rA}
    \proofstepwidestarbreakkeep[13]{%
      a\in A\land b\in B\land\RatZero<a\land\RatZero<b\land q<ab}{\rA}
    \proofstepwidestarbreakkeep[1,7,8,14]{%
      \begin{aligned}[t]
      &\eta_{\mathbb Q}(q)\preccurlyeq
        \eta_{\mathbb Q}(a)\boxtimes\eta_{\mathbb Q}(b)
        \land
        \eta_{\mathbb Q}(q)\neq
        \eta_{\mathbb Q}(a)\boxtimes\eta_{\mathbb Q}(b),\\[-2pt]
      &0_K\preccurlyeq\eta_{\mathbb Q}(a)\land
        0_K\preccurlyeq\eta_{\mathbb Q}(b)
      \end{aligned}}{%
      \FormulaRefAuto{CompleteOrderedFieldRationalPrimeEmbedding}[thm]{1,14}}
    \proofstepwidestar[\shortstack{1,2,5\\7,14}]{\begin{aligned}[t]&\eta_{\mathbb Q}[A]\neq\varnothing\land\UB_{K,\preccurlyeq}(\eta_{\mathbb Q}[A])\neq\varnothing\land{}\\[-2pt]&\eta_{\mathbb Q}[B]\neq\varnothing\land\UB_{K,\preccurlyeq}(\eta_{\mathbb Q}[B])\neq\varnothing\end{aligned}}{%
      \FormulaRefAuto{CompleteOrderedFieldCutImageBounded}[thm]{1,5}}
    \proofstepwidestarbreakkeep[\shortstack{1,2,5\\7,14}]{%
      \eta_{\mathbb Q}(a)\preccurlyeq\alpha\land
      \eta_{\mathbb Q}(b)\preccurlyeq\beta}{%
      \FormulaRefAuto{SupremumLeastUpperBound}[thm]{%
        2,16,14}}
    \proofstepwidestarbreakkeep[1,7,14]{%
      \eta_{\mathbb Q}(a)\boxtimes\eta_{\mathbb Q}(b)
      \preccurlyeq\alpha\boxtimes\beta}{%
      \FormulaRefAuto{CompleteOrderedFieldAxioms}[def]{1,7,15,17}}
    \proofstepwidestarbreakkeep[\shortstack{1--2,5\\7--8,14}]{%
      \eta_{\mathbb Q}(q)\preccurlyeq\alpha\boxtimes\beta\land
      \eta_{\mathbb Q}(q)\neq\alpha\boxtimes\beta}{%
      \FormulaRefAuto{Totale Ordnung}[def]{15,18}}
    \proofstepwidestarbreakkeep[1--2,5--8]{%
      q\in A\RealPosMul B\rightarrow
      (\eta_{\mathbb Q}(q)\preccurlyeq\alpha\boxtimes\beta\land
       \eta_{\mathbb Q}(q)\neq\alpha\boxtimes\beta)}{%
      \rRI{9,\rOE{10,11,12,13,\rEE{13,14,19}}}}
    \proofstepwidestarbreakkeep[21]{%
      \eta_{\mathbb Q}(q)\preccurlyeq\alpha\boxtimes\beta\land
      \eta_{\mathbb Q}(q)\neq\alpha\boxtimes\beta}{\rA}
    \proofstepwidestarbreakkeep[8]{%
      q<\RatZero\lor\RatZero\leq q}{%
      \FormulaRefAuto{RationalOrderTotal}[thm]{8}}
    \proofstepwidestarbreakkeep[23]{q<\RatZero}{\rA}
    \proofstepwidestarbreakkeep[5,6,8,23]{q\in A\RealPosMul B}{%
      \FormulaRefAuto{RealCutPositiveProduct}[def]{5,6,8},
      \FormulaRefAuto{RealCutZero}[def],
      \FormulaRefAuto{RealPrincipalCut}[def]{23}}
    \proofstepwidestarbreakkeep[25]{\RatZero\leq q}{\rA}
    \proofstepwidestarbreakkeep[\shortstack{1,7,8\\21,25}]{%
      \alpha\neq0_K\land\beta\neq0_K}{\text{\parbox[t]{\linewidth}{\raggedright Aus \ensuremath{\alpha=\allowbreak 0_K} oder \ensuremath{\beta=\allowbreak 0_K} folgte \ensuremath{\alpha\boxtimes\allowbreak \beta=\allowbreak 0_K,\allowbreak } im Widerspruch zu \ensuremath{0_K\preccurlyeq\allowbreak \eta_{\mathbb Q}(q)} und \ensuremath{21.}}}}
    \proofstepwidestarbreakkeep[\shortstack{1,7,8\\21,25}]{%
      \begin{aligned}[t]
      &u\coloneqq\eta_{\mathbb Q}(q)\boxtimes
        \operatorname{inv}(\beta)\in K,\\[-2pt]
      &0_K\preccurlyeq u\land u\preccurlyeq\alpha\land u\neq\alpha
      \end{aligned}}{%
      \FormulaRefAuto{CompleteOrderedFieldAxioms}[def]{1,7,21,25,26}}
    \proofstepwidestar[\shortstack{1--2,5,7--8\\21,25}]{\eta_{\mathbb Q}[A]\neq\varnothing\land\UB_{K,\preccurlyeq}(\eta_{\mathbb Q}[A])\neq\varnothing}{%
      \FormulaRefAuto{CompleteOrderedFieldCutImageBounded}[thm]{1,5}}
    \proofstepwidestarbreakkeep[\shortstack{1--2,5,7--8\\21,25}]{%
      \exists a\in A\,
      (u\preccurlyeq\eta_{\mathbb Q}(a)\land
       u\neq\eta_{\mathbb Q}(a))}{%
      \FormulaRefAuto{CompleteOrderedFieldSupremumStrictWitness}[thm]{%
        1,2,28,27}}
    \proofstepwidestarbreakkeep[29]{%
      a\in A\land u\preccurlyeq\eta_{\mathbb Q}(a)
      \land u\neq\eta_{\mathbb Q}(a)}{\rA}
    \proofstepwidestarbreakkeep[\shortstack{1,7--8,21\\25,30}]{%
      \RatZero<a\land\eta_{\mathbb Q}(a)\neq0_K}{%
      \FormulaRefAuto{CompleteOrderedFieldRationalPrimeEmbedding}[thm]{%
        1,8,25,27,30}}
    \proofstepwidestarbreakkeep[\shortstack{1,7--8,21\\25,30}]{%
      \begin{aligned}[t]
      &v\coloneqq\eta_{\mathbb Q}(q)\boxtimes
        \operatorname{inv}(\eta_{\mathbb Q}(a))\in K,\\[-2pt]
      &0_K\preccurlyeq v\land v\preccurlyeq\beta\land v\neq\beta
      \end{aligned}}{%
      \FormulaRefAuto{CompleteOrderedFieldAxioms}[def]{1,21,27,30,31}}
    \proofstepwidestar[\shortstack{1--2,5,7--8\\21,25,30}]{\eta_{\mathbb Q}[B]\neq\varnothing\land\UB_{K,\preccurlyeq}(\eta_{\mathbb Q}[B])\neq\varnothing}{%
      \FormulaRefAuto{CompleteOrderedFieldCutImageBounded}[thm]{1,5}}
    \proofstepwidestarbreakkeep[\shortstack{1--2,5,7--8\\21,25,30}]{%
      \exists b\in B\,
      (v\preccurlyeq\eta_{\mathbb Q}(b)\land
       v\neq\eta_{\mathbb Q}(b))}{%
      \FormulaRefAuto{CompleteOrderedFieldSupremumStrictWitness}[thm]{%
        1,2,33,32}}
    \proofstepwidestarbreakkeep[34]{%
      b\in B\land v\preccurlyeq\eta_{\mathbb Q}(b)
      \land v\neq\eta_{\mathbb Q}(b)}{\rA}
    \proofstepwidestarbreakkeep[\shortstack{1,7--8,21\\25,30,35}]{%
      \RatZero<b\land q<ab}{%
      \FormulaRefAuto{CompleteOrderedFieldAxioms}[def]{1,21,31,32,35},
      \FormulaRefAuto{CompleteOrderedFieldRationalPrimeEmbedding}[thm]{1}}
    \proofstepwidestarbreakkeep[\shortstack{5--6,8,21\\25,30,35}]{%
      q\in A\RealPosMul B}{%
      \FormulaRefAuto{RealCutPositiveProduct}[def]{5,6,8,30,35,36}}
    \proofstepwidestarbreakkeep[\shortstack{1--2,5--8\\21,25}]{%
      q\in A\RealPosMul B}{%
      \rEE{29,30,\rEE{34,35,37}}}
    \proofstepwidestarbreakkeep[1--2,5--8]{%
      (\eta_{\mathbb Q}(q)\preccurlyeq\alpha\boxtimes\beta\land
       \eta_{\mathbb Q}(q)\neq\alpha\boxtimes\beta)
      \rightarrow q\in A\RealPosMul B}{%
      \rRI{21,\rOE{22,23,24,25,38}}}
    \proofstepwidestarbreakkeep[1--2,5--8]{%
      q\in A\RealPosMul B\leftrightarrow
      (\eta_{\mathbb Q}(q)\preccurlyeq\alpha\boxtimes\beta\land
       \eta_{\mathbb Q}(q)\neq\alpha\boxtimes\beta)}{%
      \rBI{20,39}}
    \proofstepwidestarbreakkeep[1--2,5--7]{%
      \forall q\in\mathbb Q\,\bigl(q\in A\RealPosMul B\leftrightarrow
      (\eta_{\mathbb Q}(q)\preccurlyeq\alpha\boxtimes\beta\land
       \eta_{\mathbb Q}(q)\neq\alpha\boxtimes\beta)\bigr)}{%
      \rUI{\rRI{8,40}}}
    \proofstepwidestarbreakkeep[1--3,5,7]{%
      \exists C\in\mathbb R\,\Phi(C)=\alpha\boxtimes\beta}{%
      \FormulaRefAuto{Surjektive Funktion}[def]{\rAEa{3}},
      \FormulaRefAuto{CompleteOrderedFieldAxioms}[def]{1,7}}
    \proofstepwidestarbreakkeep[42]{%
      C\in\mathbb R\land\Phi(C)=\alpha\boxtimes\beta}{\rA}
    \proofstepwidestarbreakkeep[1,2,4,43]{%
      \forall q\in\mathbb Q\,\bigl(q\in C\leftrightarrow
      (\eta_{\mathbb Q}(q)\preccurlyeq\alpha\boxtimes\beta\land
       \eta_{\mathbb Q}(q)\neq\alpha\boxtimes\beta)\bigr)}{%
      \rIE{\rUE{4,\rAEa{43}},\rAEb{43}}}
    \proofstepwidestarbreakkeep[\shortstack{1--2,5--7\\43}]{%
      C=A\RealPosMul B}{%
      \FormulaRefAuto{\forall x\,(x\in A\leftrightarrow x\in B)
        \vdash A=B}[ax]{41,44}}
    \proofstepwidestarbreakkeep[\shortstack{1--2,5--7\\43}]{%
      \Phi(A\RealPosMul B)=\Phi(A)\boxtimes\Phi(B)}{%
      \rIE{7,45,\rAEb{43}}}
    \proofstepwidestarbreakkeep[1\text{--}2,5\text{--}7]{%
      \Phi(A\RealPosMul B)=\Phi(A)\boxtimes\Phi(B)}{%
      \rEE{42,43,46}}
    \proofstepwidestarbreakkeep[1,2,5]{%
      (\RealZero\RealLe A\land\RealZero\RealLe B)\rightarrow
      \Phi(A\RealPosMul B)=\Phi(A)\boxtimes\Phi(B)}{%
      \rRI{6,47}}
    \proofstepwidestarbreakkeep[1,2]{%
      \forall A,B\in\mathbb R\,
      ((\RealZero\RealLe A\land\RealZero\RealLe B)\rightarrow
       \Phi(A\RealPosMul B)=\Phi(A)\boxtimes\Phi(B))}{%
      \rUI{\rRI{5,48}}}
  \closeproofpartwideindR

  \proofpartwide{Erhaltung der Addition}
\proofstepwidestarbreakkeep[1]{%
      \CompleteOrderedField{K,\boxplus,\boxtimes,\boxminus,
        \operatorname{inv},0_K,1_K,\preccurlyeq}}{\rA}
    \proofstepwidestarbreakkeep[2]{\Phi(A)\coloneqq
      \sup_{K,\preccurlyeq}(\eta_{\mathbb Q}[A])}{\rA}
    \proofstepwidestarbreakkeep[1,2]{%
      \forall A,B\in\mathbb R\,
      \Phi(A\RealAdd B)=\Phi(A)\boxplus\Phi(B)}{%
      \FormulaRefAuto{CompleteOrderedFieldCutSupremumAddition}[thm]{1,2}}

\closeproofpartwide

\proofpartwide{Erhaltung von Null und Eins}
  \setcounter{proofstepnr}{3}% Fortlaufende Schritte dieses Beweises.
\proofstepwidestarbreakkeep[1,2]{%
      \forall q\in\mathbb Q\,
      \Phi(\RealEmb(q))=\eta_{\mathbb Q}(q)}{%
      \FormulaRefAuto{CompleteOrderedFieldPrincipalCutImage}[thm]{1,2}}
    \proofstepwidestarbreakkeep[1,2]{%
      \Phi(\RealZero)=0_K\land\Phi(\RealOne)=1_K}{%
      \FormulaRefAuto{RealCutZero}[def],
      \FormulaRefAuto{RealCutOne}[def],
      \FormulaRefAuto{RationalZeroOneClasses}[thm],
      \FormulaRefAuto{CompleteOrderedFieldRationalPrimeEmbedding}[thm]{1},4}

\closeproofpartwide

\proofpartwide{Erhaltung der Negation}
  \setcounter{proofstepnr}{5}% Fortlaufende Schritte dieses Beweises.
\proofstepwidestarbreakkeep[6]{A\in\mathbb R}{\rA}
    \proofstepwidestar[1,2,3,5,6]{A\RealAdd\RealNeg A=\RealZero}{%
      \FormulaRefAuto{RealCutAdditionInverse}[thm]{6}}
    \proofstepwidestarbreakkeep[1,2,3,5,6]{%
      \Phi(A)\boxplus\Phi(\RealNeg A)=0_K}{%
      \rIE{\rUE{3,A,\RealNeg A},
        7,5}}
    \proofstepwidestarbreakkeep[1,2,6]{%
      \Phi(\RealNeg A)=\boxminus\Phi(A)}{%
      \FormulaRefAuto{CompleteOrderedFieldAxioms}[def]{1,8}}
    \proofstepwidestarbreakkeep[1,2]{%
      \forall A\in\mathbb R\,
      \Phi(\RealNeg A)=\boxminus\Phi(A)}{%
      \rUI{\rRI{6,9}}}

\closeproofpartwide

\proofpartwide{Erhaltung der Multiplikation}
  \setcounter{proofstepnr}{10}% Fortlaufende Schritte dieses Beweises.
\proofstepwidestarbreakkeep[1,2]{%
      \forall A,B\in\mathbb R\,
      ((\RealZero\RealLe A\land\RealZero\RealLe B)\rightarrow
       \Phi(A\RealPosMul B)=\Phi(A)\boxtimes\Phi(B))}{%
      \FormulaRefAuto{CompleteOrderedFieldCutSupremumPositiveProduct}[thm]{1,2}}
    \proofstepwidestarbreakkeep[12]{A,B\in\mathbb R}{\rA}
    \proofstepwidestarbreakkeep[1,12]{%
      \RealZero\RealLe A\lor A\RealLt\RealZero}{%
      \FormulaRefAuto{RealCutTotalOrder}[thm]{12}}
    \proofstepwidestarbreakkeep[1,12]{%
      \RealZero\RealLe B\lor B\RealLt\RealZero}{%
      \FormulaRefAuto{RealCutTotalOrder}[thm]{12}}
    \proofstepwidestarbreakkeep[15]{\RealZero\RealLe A}{\rA}
    \proofstepwidestarbreakkeep[16]{\RealZero\RealLe B}{\rA}
    \proofstepwidestarbreakkeep[1,2,12,15,16]{%
      \Phi(A\RealMul B)=\Phi(A)\boxtimes\Phi(B)}{%
      \FormulaRefAuto{RealCutMultiplication}[def]{15,16},
      \rUE{11,12,\rAI{15,16}}}
    \proofstepwidestarbreakkeep[18]{B\RealLt\RealZero}{\rA}
    \proofstepwidestarbreakkeep[1,2,12,15,18]{%
      \Phi(A\RealMul B)=\Phi(A)\boxtimes\Phi(B)}{\text{\parbox[t]{\linewidth}{\raggedright \FormulaRefAuto{RealCutMultiplication}[def]{15,18}, \FormulaRefAuto{RealCutNegationOrder}[thm]{18}, \ensuremath{9,\allowbreak 11,\allowbreak } \FormulaRefAuto{CompleteOrderedFieldAxioms}[def]{1}}}}
    \proofstepwidestarbreakkeep[1,2,12,15]{%
      \Phi(A\RealMul B)=\Phi(A)\boxtimes\Phi(B)}{%
      \rOE{14,16,17,18,19}}
    \proofstepwidestarbreakkeep[21]{A\RealLt\RealZero}{\rA}
    \proofstepwidestarbreakkeep[22]{\RealZero\RealLe B}{\rA}
    \proofstepwidestarbreakkeep[1,2,12,21,22]{%
      \Phi(A\RealMul B)=\Phi(A)\boxtimes\Phi(B)}{\text{\parbox[t]{\linewidth}{\raggedright \FormulaRefAuto{RealCutMultiplication}[def]{21,22}, \FormulaRefAuto{RealCutNegationOrder}[thm]{21}, \ensuremath{9,\allowbreak 11,\allowbreak } \FormulaRefAuto{CompleteOrderedFieldAxioms}[def]{1}}}}
    \proofstepwidestarbreakkeep[24]{B\RealLt\RealZero}{\rA}
    \proofstepwidestarbreakkeep[1,2,12,21,24]{%
      \Phi(A\RealMul B)=\Phi(A)\boxtimes\Phi(B)}{\text{\parbox[t]{\linewidth}{\raggedright \FormulaRefAuto{RealCutMultiplication}[def]{21,24}, \FormulaRefAuto{RealCutNegationOrder}[thm]{21,24}, \ensuremath{9,\allowbreak 11,\allowbreak } \FormulaRefAuto{CompleteOrderedFieldAxioms}[def]{1}}}}
    \proofstepwidestarbreakkeep[1,2,12,21]{%
      \Phi(A\RealMul B)=\Phi(A)\boxtimes\Phi(B)}{%
      \rOE{14,22,23,24,25}}
    \proofstepwidestarbreakkeep[1,2,12]{%
      \Phi(A\RealMul B)=\Phi(A)\boxtimes\Phi(B)}{%
      \rOE{13,15,20,21,26}}
    \proofstepwidestarbreakkeep[1,2]{%
      \forall A,B\in\mathbb R\,
      \Phi(A\RealMul B)=\Phi(A)\boxtimes\Phi(B)}{%
      \rUI{\rRI{12,27}}}

\closeproofpartwide

\proofpartwideindR[Suprema und Schnittoperationen]{%
    \begin{aligned}[t]
      &\CompleteOrderedField{K,\boxplus,\boxtimes,\boxminus,
        \operatorname{inv},0_K,1_K,\preccurlyeq}\dsep
        \Phi(A)\coloneqq
        \sup_{K,\preccurlyeq}(\eta_{\mathbb Q}[A])\vdash{}\\[-2pt]
      &\forall A,B\in\mathbb R\,\bigl(
        \Phi(A\RealAdd B)=\Phi(A)\boxplus\Phi(B)\land{}\\[-2pt]
      &\qquad\Phi(A\RealMul B)=\Phi(A)\boxtimes\Phi(B)\bigr)
        \land\Phi(\RealZero)=0_K\land\Phi(\RealOne)=1_K
    \end{aligned}
  }{%
    \CompleteOrderedField{K,\boxplus,\boxtimes,\boxminus,
      \operatorname{inv},0_K,1_K,\preccurlyeq}\dsep
    \Phi(A)\coloneqq
      \sup_{K,\preccurlyeq}(\eta_{\mathbb Q}[A])\vdash
    \forall A,B\in\mathbb R\,
    (\Phi(A\RealAdd B)=\Phi(A)\boxplus\Phi(B)\land
     \Phi(A\RealMul B)=\Phi(A)\boxtimes\Phi(B))
    \land\Phi(\RealZero)=0_K\land\Phi(\RealOne)=1_K
  }[CompleteOrderedFieldCutSupremumArithmetic]
  \setcounter{proofstepnr}{28}% Fortlaufende Schritte dieses Beweises.
\proofstepwidestarbreakkeep[1,2]{\begin{aligned}[t]
      &\forall A,B\in\mathbb R\,\bigl(\\[-2pt]
      &\quad\Phi(A\RealAdd B)=\Phi(A)\boxplus\Phi(B)\land{}\\[-2pt]
      &\quad\Phi(A\RealMul B)=\Phi(A)\boxtimes\Phi(B)\bigr)\\[-2pt]
      &\quad\land\Phi(\RealZero)=0_K\land\Phi(\RealOne)=1_K
      \end{aligned}}{%
      \rAI{3,\rAI{28,5}}}

\closeproofpartwideindR


  \proofpartwideindR[Supremumsdichte der Hauptschnitte]{%
    \begin{aligned}[t]
      &A\in\mathbb R\vdash
        \varnothing\neq\RealEmb[A]\subseteq\mathbb R\land
        \UB_{\mathbb R,\RealLe}(\RealEmb[A])\neq\varnothing\land{}\\[-2pt]
      &\hspace{31mm}
        A=\sup_{\mathbb R,\RealLe}(\RealEmb[A])
    \end{aligned}
  }{%
    A\in\mathbb R\vdash
    \varnothing\neq\RealEmb[A]\subseteq\mathbb R\land
    \UB_{\mathbb R,\RealLe}(\RealEmb[A])\neq\varnothing\land
    A=\sup_{\mathbb R,\RealLe}(\RealEmb[A])
  }[RealCutPrincipalFamilySupremum]
    \proofstepwidestarbreakkeep[1]{A\in\mathbb R}{\rA}
    \proofstepwidestar[1]{\RealCut A}{%
      \FormulaRefAuto{RealCutMembership}[thm]{1}}
    \proofstepwidestarbreakkeep[1]{\exists q\in A}{%
      \FormulaRefAuto{DedekindCutCharacterization}[thm]{%
        2}}
    \proofstepwidestarbreakkeep[1]{%
      \varnothing\neq\RealEmb[A]\subseteq\mathbb R}{%
      \FormulaRefAuto{NonemptyDirectImage}[thm]{3},
      \FormulaRefAuto{RationalToRealEmbeddingFunction}[thm]}
    \proofstepwidestar[1]{\forall q\in\mathbb Q\,(\RealHat q=\{p\in\mathbb Q\mid p<q\})}{%
      \FormulaRefAuto{RealPrincipalCut}[def]}
    \proofstepwidestar[1]{\forall p\in A\,\forall q\in\mathbb Q\,(q<p\rightarrow q\in A)}{%
      \FormulaRefAuto{DedekindCutCharacterization}[thm]{1}}
    \proofstepwidestarbreakkeep[1]{%
      \forall q\in A\,(\RealEmb(q)\RealLe A)}{%
      \FormulaRefAuto{RealCutOrderCharacterization}[thm]{%
        5,
        6}}
    \proofstepwidestarbreakkeep[1]{%
      A\in\UB_{\mathbb R,\RealLe}(\RealEmb[A])\land
      \UB_{\mathbb R,\RealLe}(\RealEmb[A])\neq\varnothing}{%
      \FormulaRefAuto{UpperBoundSetIntro}[thm]{1,4,7},
      \FormulaRefAuto{x\in A\vdash A\neq\varnothing}[thm]}
    \proofstepwidestarbreakkeep[9]{x\in\bigcup\RealEmb[A]}{\rA}
    \proofstepwidestarbreakkeep[9]{%
      \exists q\in A\,(x\in\RealEmb(q))}{%
      \FormulaRefAuto{x\in\bigcup A\eqvdash\exists B\in A\,(x\in B)}{9}}
    \proofstepwidestarbreakkeep[11]{q\in A\land x\in\RealEmb(q)}{\rA}
    \proofstepwidestarbreakkeep[1,9,11]{x\in A}{%
      \FormulaRefAuto{RealPrincipalCut}[def]{11},
      \FormulaRefAuto{DedekindCutCharacterization}[thm]{1,11}}
    \proofstepwidestarbreakkeep[1]{\bigcup\RealEmb[A]\subseteq A}{%
      \FormulaRefAuto{SubsetIntroductionRule}[thm]{[9]\vdots \rEE{10,11,12}}}
    \proofstepwidestarbreakkeep[14]{x\in A}{\rA}
    \proofstepwidestarbreakkeep[1,14]{%
      \exists q\in A\,(x<q)}{%
      \FormulaRefAuto{DedekindCutCharacterization}[thm]{1,14}}
    \proofstepwidestarbreakkeep[16]{q\in A\land x<q}{\rA}
    \proofstepwidestarbreakkeep[1,14,16]{%
      x\in\RealEmb(q)\land\RealEmb(q)\in\RealEmb[A]}{%
      \FormulaRefAuto{RationalToRealEmbedding}[def]{16},
      \FormulaRefAuto{F\colon A\to B\dsep
        \exists x\in A\,y=F(x)\vdash y\in F[A]}[thm]{16}}
    \proofstepwidestarbreakkeep[1,14,16]{x\in\bigcup\RealEmb[A]}{%
      \FormulaRefAuto{x\in\bigcup A\eqvdash\exists B\in A\,(x\in B)}{17}}
    \proofstepwidestarbreakkeep[1]{A\subseteq\bigcup\RealEmb[A]}{%
      \FormulaRefAuto{SubsetIntroductionRule}[thm]{[14]\vdots \rEE{15,16,18}}}
    \proofstepwidestarbreakkeep[1]{\bigcup\RealEmb[A]=A}{%
      \FormulaRefAuto{\forall x\,(x\in A\leftrightarrow x\in B)
        \vdash A=B}[ax]{13,19}}
    \proofstepwidestarbreakkeep[1]{%
      \sup_{\mathbb R,\RealLe}(\RealEmb[A])=\bigcup\RealEmb[A]}{%
      \FormulaRefAuto{RealCutSupremumUnion}[thm]{4,1,7}}
    \proofstepwidestarbreakkeep[1]{%
      \varnothing\neq\RealEmb[A]\subseteq\mathbb R\land
      \UB_{\mathbb R,\RealLe}(\RealEmb[A])\neq\varnothing\land
      A=\sup_{\mathbb R,\RealLe}(\RealEmb[A])}{%
      \rAI{4,\rAI{\rAEb{8},\rIE{20,21}}}}
  \closeproofpartwideindR

  \proofpartwideindR[Ordnungsisomorphismen erhalten vorhandene Suprema]{%
    \begin{aligned}[t]
      &\Psi\colon\mathbb R\bij K\dsep
        \forall X,Y\in\mathbb R\,
        (X\RealLe Y\leftrightarrow\Psi(X)\preccurlyeq\Psi(Y))\\[-2pt]
      &\quad\dsep\varnothing\neq S\subseteq\mathbb R\dsep
        \UB_{\mathbb R,\RealLe}(S)\neq\varnothing\\[-2pt]
      &\quad\vdash
        \Psi(\sup_{\mathbb R,\RealLe}S)
        =\sup_{K,\preccurlyeq}(\Psi[S])
    \end{aligned}
  }{%
    \Psi\colon\mathbb R\bij K\dsep
    \forall X,Y\in\mathbb R\,
    (X\RealLe Y\leftrightarrow\Psi(X)\preccurlyeq\Psi(Y))\dsep
    \varnothing\neq S\subseteq\mathbb R\dsep
    \UB_{\mathbb R,\RealLe}(S)\neq\varnothing
    \vdash
    \Psi(\sup_{\mathbb R,\RealLe}S)
    =\sup_{K,\preccurlyeq}(\Psi[S])
  }[CompleteOrderedFieldOrderIsomorphismPreservesSupremum]
    \proofstepwidestarbreakkeep[1]{\Psi\colon\mathbb R\bij K}{\rA}
    \proofstepwidestarbreakkeep[2]{%
      \forall X,Y\in\mathbb R\,
      (X\RealLe Y\leftrightarrow\Psi(X)\preccurlyeq\Psi(Y))}{\rA}
    \proofstepwidestarbreakkeep[1,2]{\PartOrd{K,\preccurlyeq}}{\text{\parbox[t]{\linewidth}{\raggedright Surjektivität aus \ensuremath{1} schreibt jedes \ensuremath{x,\allowbreak y,\allowbreak z\in K} als \ensuremath{\Psi(X),\allowbreak \Psi(Y),\allowbreak \Psi(Z).} Mit der Äquivalenz \ensuremath{2} übertragen sich Reflexivität,
        Antisymmetrie und Transitivität von \ensuremath{\RealLe\allowbreak .} \FormulaRefAuto{RealCutInclusionPartialOrder}[thm] , \FormulaRefAuto{Bijektive Funktion}[def]{1}}}}
    \proofstepwidestarbreakkeep[4]{\varnothing\neq S\subseteq\mathbb R}{\rA}
    \proofstepwidestarbreakkeep[5]{%
      \UB_{\mathbb R,\RealLe}(S)\neq\varnothing}{\rA}
    \proofstepwidestarbreakkeep[4,5]{%
      \begin{aligned}[t]
      &s\coloneqq\sup_{\mathbb R,\RealLe}S\in
        \UB_{\mathbb R,\RealLe}(S),\\[-2pt]
      &\forall U\in\UB_{\mathbb R,\RealLe}(S)\,(s\RealLe U)
      \end{aligned}}{%
      \FormulaRefAuto{SupremumLeastUpperBound}[thm]{4,5}}
    \proofstepwidestarbreakkeep[1,2,4,5,6]{%
      \Psi[S]\subseteq K\land
      \Psi(s)\in\UB_{K,\preccurlyeq}(\Psi[S])}{%
      \FormulaRefAuto{C\subseteq A\dsep F\colon A\to B
        \vdash F[C]\subseteq B}[thm]{4,1},
      \FormulaRefAuto{UpperBoundSetIntro}[thm]{1,2,4,6}}
    \proofstepwidestarbreakkeep[8]{%
      v\in\UB_{K,\preccurlyeq}(\Psi[S])}{\rA}
    \proofstepwidestarbreakkeep[1,7,8]{%
      \exists V\in\mathbb R\,(\Psi(V)=v)}{%
      \FormulaRefAuto{Surjektive Funktion}[def]{1},
      \FormulaRefAuto{UpperBoundSetElim}[thm]{\rAEa{7},8}}
    \proofstepwidestarbreakkeep[9]{V\in\mathbb R\land\Psi(V)=v}{\rA}
    \proofstepwidestarbreakkeep[1,2,4,7,8,10]{%
      V\in\UB_{\mathbb R,\RealLe}(S)}{%
      \FormulaRefAuto{UpperBoundSetElim}[thm]{\rAEa{7},8},
      \FormulaRefAuto{UpperBoundSetIntro}[thm]{2,4,10}}
    \proofstepwidestarbreakkeep[1,2,4,5,6,8,10]{s\RealLe V}{%
      \rUE{\rAEb{6},11}}
    \proofstepwidestarbreakkeep[1,2,4,5,6,8,10]{%
      \Psi(s)\preccurlyeq v}{%
      \rIE{\rUE{2,s,V},12,\rAEb{10}}}
    \proofstepwidestarbreakkeep[1,2,4,5,6,8]{%
      \Psi(s)\preccurlyeq v}{%
      \rEE{9,10,13}}
    \proofstepwidestarbreakkeep[1,2,4,5,6]{%
      \forall v\in\UB_{K,\preccurlyeq}(\Psi[S])\,
      (\Psi(s)\preccurlyeq v)}{%
      \rUI{\rRI{8,14}}}
    \proofstepwidestarbreakkeep[1,2,3,4,5,6]{%
      \Psi(s)=\sup_{K,\preccurlyeq}(\Psi[S])}{%
      \FormulaRefAuto{SupremumCandidateEqualsSupremum}[thm]{%
        3,\rAEa{7},\rAEb{7},15}}
    \proofstepwidestarbreakkeep[1,2,4,5]{%
      \Psi(\sup_{\mathbb R,\RealLe}S)
      =\sup_{K,\preccurlyeq}(\Psi[S])}{\rIE{6,16}}
  \closeproofpartwideindR

  \proofpartwide{Existenz und Operationserhaltung}
\proofstepwidestar[1]{\CompleteOrderedField{K,\boxplus,\boxtimes,
      \boxminus,\operatorname{inv},0_K,1_K,\preccurlyeq}}{\rA}
    \proofstepwidestar[1]{\Phi(A)=
      \sup_{K,\preccurlyeq}(\eta_{\mathbb Q}[A])}{\rII}
    \proofstepwidestarbreakkeep[1]{%
      \begin{aligned}[t]
        &\Phi\colon\mathbb R\bij K\land{}\\[-2pt]
        &\forall A,B\in\mathbb R\,\\[-2pt]
        &(A\RealLe B\leftrightarrow\Phi(A)\preccurlyeq\Phi(B))
      \end{aligned}}{%
      \FormulaRefAuto{CompleteOrderedFieldCutSupremumBijection}[thm]{1,2}}
    \proofstepwidestarbreakkeep[1]{\begin{aligned}[t]
      &\forall A,B\in\mathbb R\,\bigl(\\[-2pt]
      &\quad\Phi(A\RealAdd B)=\Phi(A)\boxplus\Phi(B)\land{}\\[-2pt]
      &\quad\Phi(A\RealMul B)=\Phi(A)\boxtimes\Phi(B)\bigr)\\[-2pt]
      &\quad\land\Phi(\RealZero)=0_K\land\Phi(\RealOne)=1_K
      \end{aligned}}{%
      \FormulaRefAuto{CompleteOrderedFieldCutSupremumArithmetic}[thm]{1,2}}
    \proofstepwidestarbreakkeep[1]{%
      \begin{aligned}[t]
        &\forall A,B\in\mathbb R\,\\[-2pt]
        &\quad \Phi(A\RealAdd B)=\Phi(A)\boxplus\Phi(B)
      \end{aligned}}{%
      \rAEa{4}}
    \proofstepwidestarbreakkeep[1]{%
      \begin{aligned}[t]
        &\forall A,B\in\mathbb R\,\\[-2pt]
        &\quad \Phi(A\RealMul B)=\Phi(A)\boxtimes\Phi(B)
      \end{aligned}}{%
      \rAEa{\rAEb{4}}}
    \proofstepwidestar[1]{\Phi(\RealZero)=0_K\land\Phi(\RealOne)=1_K}{%
      \rAI{\rAEa{\rAEb{4}},\rAEb{\rAEb{4}}}}

\closeproofpartwide

\proofpartwide{Eindeutigkeit}
  \setcounter{proofstepnr}{7}% Fortlaufende Schritte dieses Beweises.
\proofstepwidestarbreakkeep[8]{\begin{aligned}[t]
      &\Psi\colon\mathbb R\bij K\dsep{}\\[-2pt]
      &\forall X,Y\in\mathbb R\,\bigl(\\[-2pt]
      &\quad\Psi(X\RealAdd Y)=\Psi(X)\boxplus\Psi(Y)\land{}\\[-2pt]
      &\quad\Psi(X\RealMul Y)=\Psi(X)\boxtimes\Psi(Y)\bigr)\dsep{}\\[-2pt]
      &\forall X,Y\in\mathbb R\,(X\RealLe Y\leftrightarrow\Psi(X)\preccurlyeq\Psi(Y))\dsep{}\\[-2pt]
      &\Psi(\RealZero)=0_K\dsep\Psi(\RealOne)=1_K
      \end{aligned}}{\rA}
    \proofstepwidestarbreakkeep[1,8]{%
      (\Psi\circ\RealEmb)\colon\mathbb Q\inj K\land
      \RationalPrimeEmbeddingLaws{\Psi\circ\RealEmb}}{\text{\parbox[t]{\linewidth}{\raggedright \FormulaRefAuto{RationalToRealEmbeddingFunction}[thm] , \FormulaRefAuto{RationalToRealEmbeddingInjective}[thm] , \FormulaRefAuto{RationalToRealEmbeddingArithmetic}[thm] , \FormulaRefAuto{RationalToRealEmbeddingOrder}[thm] , \FormulaRefAuto{CompleteOrderedFieldRationalPrimeEmbeddingLaws}[def] , die Additions- und Nullregeln aus \ensuremath{8} geben \ensuremath{\Psi(\RealNeg X)=\allowbreak \boxminus\Psi(X),\allowbreak } die Produkt- und Einsregeln aus \ensuremath{8} geben für \ensuremath{X\neq0} ebenso \ensuremath{\Psi(\RealInv X)=\allowbreak \operatorname{inv}(\Psi(X));} \FormulaRefAuto{RealCutAdditiveLaws}[thm] , \FormulaRefAuto{RealCutMultiplicativeLaws}[thm] , \FormulaRefAuto{CompleteOrderedFieldAxioms}[def]{1,8}}}}
    \proofstepwidestarbreakkeep[1,8]{%
      \forall q\in\mathbb Q\,
      \Psi(\RealEmb(q))=\eta_{\mathbb Q}(q)}{%
      \FormulaRefAuto{CompleteOrderedFieldRationalPrimeEmbedding}[thm]{1,9}}
    \proofstepwidestarbreakkeep[1]{%
      \begin{aligned}[t]
      \forall A\in\mathbb R\,\bigl(&
        \varnothing\neq\RealEmb[A]\subseteq\mathbb R\land{}\\[-2pt]
      &\UB_{\mathbb R,\RealLe}(\RealEmb[A])\neq\varnothing\land{}\\[-2pt]
      &A=\sup_{\mathbb R,\RealLe}(\RealEmb[A])\bigr)
      \end{aligned}}{%
      \FormulaRefAuto{RealCutPrincipalFamilySupremum}[thm]}
    \proofstepwidestarbreakkeep[1,8]{%
      \begin{aligned}[t]
      &\varnothing\neq S\subseteq\mathbb R\dsep
        \UB_{\mathbb R,\RealLe}(S)\neq\varnothing\vdash{}\\[-2pt]
      &\Psi(\sup_{\mathbb R,\RealLe}S)
       =\sup_{K,\preccurlyeq}(\Psi[S])
      \end{aligned}}{%
      \FormulaRefAuto{CompleteOrderedFieldOrderIsomorphismPreservesSupremum}[thm]{8}}
    \proofstepwidestarbreakkeep[1,8]{%
      \forall A\in\mathbb R\,
      \Psi[\RealEmb[A]]=\eta_{\mathbb Q}[A]}{%
      \FormulaRefAuto{F\colon A\to B\vdash
        y\in F[A]\leftrightarrow\exists x\in A\,y=F(x)}[thm]{10},
      \FormulaRefAuto{\forall x\,(x\in A\leftrightarrow x\in B)
        \vdash A=B}[ax]}
    \proofstepwidestarbreakkeep[1,2,8]{%
      \forall A\in\mathbb R\,\Psi(A)=\Phi(A)}{%
      \rIE{10,11,12,13,2}}

\closeproofpartwide

\proofpartwide{Eindeutige Existenz}
  \setcounter{proofstepnr}{14}% Fortlaufende Schritte dieses Beweises.
\proofstepwidestarbreakkeep[1]{%
      \begin{aligned}[t]
        &\exists!\Phi\colon\mathbb R\bij K\,\bigl(\\[-2pt]
        &\quad \forall x,y\in\mathbb R\,\\[-2pt]
        &\qquad \Phi(x\RealAdd y)=\Phi(x)\boxplus\Phi(y)\land{}\\[-2pt]
        &\quad \forall x,y\in\mathbb R\,\\[-2pt]
        &\qquad \Phi(x\RealMul y)=\Phi(x)\boxtimes\Phi(y)\land{}\\[-2pt]
        &\quad \forall x,y\in\mathbb R\,\\[-2pt]
        &\qquad (x\RealLe y\leftrightarrow
          \Phi(x)\preccurlyeq\Phi(y))\land{}\\[-2pt]
        &\quad \Phi(\RealZero)=0_K\land\Phi(\RealOne)=1_K\bigr)
      \end{aligned}}{\UEI{3,5,6,7,\rRI{8,14}}}

\closeproofpartwide
\end{tabproofsplitwide}

Die Eindeutigkeit ist als Eindeutigkeit des ordnungs- und
körperstrukturerhaltenden Isomorphismus zu verstehen. Damit kann in späteren
Bänden entweder mit dem konkreten Schnittmodell oder ausschließlich mit den
registrierten Axiomen eines vollständig angeordneten Körpers gearbeitet
werden. Für Grenzwertaussagen stehen insbesondere die reelle Ordnung, der
Abstand und positive Epsilon-Umgebungen bereits zur Verfügung.

\end{document}
